% Options for packages loaded elsewhere
% Options for packages loaded elsewhere
\PassOptionsToPackage{unicode}{hyperref}
\PassOptionsToPackage{hyphens}{url}
\PassOptionsToPackage{dvipsnames,svgnames,x11names}{xcolor}
%
\documentclass[
  letterpaper,
]{scrbook}
\usepackage{xcolor}
\usepackage{amsmath,amssymb}
\setcounter{secnumdepth}{5}
\usepackage{iftex}
\ifPDFTeX
  \usepackage[T1]{fontenc}
  \usepackage[utf8]{inputenc}
  \usepackage{textcomp} % provide euro and other symbols
\else % if luatex or xetex
  \usepackage{unicode-math} % this also loads fontspec
  \defaultfontfeatures{Scale=MatchLowercase}
  \defaultfontfeatures[\rmfamily]{Ligatures=TeX,Scale=1}
\fi
\usepackage{lmodern}
\ifPDFTeX\else
  % xetex/luatex font selection
\fi
% Use upquote if available, for straight quotes in verbatim environments
\IfFileExists{upquote.sty}{\usepackage{upquote}}{}
\IfFileExists{microtype.sty}{% use microtype if available
  \usepackage[]{microtype}
  \UseMicrotypeSet[protrusion]{basicmath} % disable protrusion for tt fonts
}{}
\makeatletter
\@ifundefined{KOMAClassName}{% if non-KOMA class
  \IfFileExists{parskip.sty}{%
    \usepackage{parskip}
  }{% else
    \setlength{\parindent}{0pt}
    \setlength{\parskip}{6pt plus 2pt minus 1pt}}
}{% if KOMA class
  \KOMAoptions{parskip=half}}
\makeatother
% Make \paragraph and \subparagraph free-standing
\makeatletter
\ifx\paragraph\undefined\else
  \let\oldparagraph\paragraph
  \renewcommand{\paragraph}{
    \@ifstar
      \xxxParagraphStar
      \xxxParagraphNoStar
  }
  \newcommand{\xxxParagraphStar}[1]{\oldparagraph*{#1}\mbox{}}
  \newcommand{\xxxParagraphNoStar}[1]{\oldparagraph{#1}\mbox{}}
\fi
\ifx\subparagraph\undefined\else
  \let\oldsubparagraph\subparagraph
  \renewcommand{\subparagraph}{
    \@ifstar
      \xxxSubParagraphStar
      \xxxSubParagraphNoStar
  }
  \newcommand{\xxxSubParagraphStar}[1]{\oldsubparagraph*{#1}\mbox{}}
  \newcommand{\xxxSubParagraphNoStar}[1]{\oldsubparagraph{#1}\mbox{}}
\fi
\makeatother

\usepackage{color}
\usepackage{fancyvrb}
\newcommand{\VerbBar}{|}
\newcommand{\VERB}{\Verb[commandchars=\\\{\}]}
\DefineVerbatimEnvironment{Highlighting}{Verbatim}{commandchars=\\\{\}}
% Add ',fontsize=\small' for more characters per line
\usepackage{framed}
\definecolor{shadecolor}{RGB}{241,243,245}
\newenvironment{Shaded}{\begin{snugshade}}{\end{snugshade}}
\newcommand{\AlertTok}[1]{\textcolor[rgb]{0.68,0.00,0.00}{#1}}
\newcommand{\AnnotationTok}[1]{\textcolor[rgb]{0.37,0.37,0.37}{#1}}
\newcommand{\AttributeTok}[1]{\textcolor[rgb]{0.40,0.45,0.13}{#1}}
\newcommand{\BaseNTok}[1]{\textcolor[rgb]{0.68,0.00,0.00}{#1}}
\newcommand{\BuiltInTok}[1]{\textcolor[rgb]{0.00,0.23,0.31}{#1}}
\newcommand{\CharTok}[1]{\textcolor[rgb]{0.13,0.47,0.30}{#1}}
\newcommand{\CommentTok}[1]{\textcolor[rgb]{0.37,0.37,0.37}{#1}}
\newcommand{\CommentVarTok}[1]{\textcolor[rgb]{0.37,0.37,0.37}{\textit{#1}}}
\newcommand{\ConstantTok}[1]{\textcolor[rgb]{0.56,0.35,0.01}{#1}}
\newcommand{\ControlFlowTok}[1]{\textcolor[rgb]{0.00,0.23,0.31}{\textbf{#1}}}
\newcommand{\DataTypeTok}[1]{\textcolor[rgb]{0.68,0.00,0.00}{#1}}
\newcommand{\DecValTok}[1]{\textcolor[rgb]{0.68,0.00,0.00}{#1}}
\newcommand{\DocumentationTok}[1]{\textcolor[rgb]{0.37,0.37,0.37}{\textit{#1}}}
\newcommand{\ErrorTok}[1]{\textcolor[rgb]{0.68,0.00,0.00}{#1}}
\newcommand{\ExtensionTok}[1]{\textcolor[rgb]{0.00,0.23,0.31}{#1}}
\newcommand{\FloatTok}[1]{\textcolor[rgb]{0.68,0.00,0.00}{#1}}
\newcommand{\FunctionTok}[1]{\textcolor[rgb]{0.28,0.35,0.67}{#1}}
\newcommand{\ImportTok}[1]{\textcolor[rgb]{0.00,0.46,0.62}{#1}}
\newcommand{\InformationTok}[1]{\textcolor[rgb]{0.37,0.37,0.37}{#1}}
\newcommand{\KeywordTok}[1]{\textcolor[rgb]{0.00,0.23,0.31}{\textbf{#1}}}
\newcommand{\NormalTok}[1]{\textcolor[rgb]{0.00,0.23,0.31}{#1}}
\newcommand{\OperatorTok}[1]{\textcolor[rgb]{0.37,0.37,0.37}{#1}}
\newcommand{\OtherTok}[1]{\textcolor[rgb]{0.00,0.23,0.31}{#1}}
\newcommand{\PreprocessorTok}[1]{\textcolor[rgb]{0.68,0.00,0.00}{#1}}
\newcommand{\RegionMarkerTok}[1]{\textcolor[rgb]{0.00,0.23,0.31}{#1}}
\newcommand{\SpecialCharTok}[1]{\textcolor[rgb]{0.37,0.37,0.37}{#1}}
\newcommand{\SpecialStringTok}[1]{\textcolor[rgb]{0.13,0.47,0.30}{#1}}
\newcommand{\StringTok}[1]{\textcolor[rgb]{0.13,0.47,0.30}{#1}}
\newcommand{\VariableTok}[1]{\textcolor[rgb]{0.07,0.07,0.07}{#1}}
\newcommand{\VerbatimStringTok}[1]{\textcolor[rgb]{0.13,0.47,0.30}{#1}}
\newcommand{\WarningTok}[1]{\textcolor[rgb]{0.37,0.37,0.37}{\textit{#1}}}

\usepackage{longtable,booktabs,array}
\usepackage{calc} % for calculating minipage widths
% Correct order of tables after \paragraph or \subparagraph
\usepackage{etoolbox}
\makeatletter
\patchcmd\longtable{\par}{\if@noskipsec\mbox{}\fi\par}{}{}
\makeatother
% Allow footnotes in longtable head/foot
\IfFileExists{footnotehyper.sty}{\usepackage{footnotehyper}}{\usepackage{footnote}}
\makesavenoteenv{longtable}
\usepackage{graphicx}
\makeatletter
\newsavebox\pandoc@box
\newcommand*\pandocbounded[1]{% scales image to fit in text height/width
  \sbox\pandoc@box{#1}%
  \Gscale@div\@tempa{\textheight}{\dimexpr\ht\pandoc@box+\dp\pandoc@box\relax}%
  \Gscale@div\@tempb{\linewidth}{\wd\pandoc@box}%
  \ifdim\@tempb\p@<\@tempa\p@\let\@tempa\@tempb\fi% select the smaller of both
  \ifdim\@tempa\p@<\p@\scalebox{\@tempa}{\usebox\pandoc@box}%
  \else\usebox{\pandoc@box}%
  \fi%
}
% Set default figure placement to htbp
\def\fps@figure{htbp}
\makeatother





\setlength{\emergencystretch}{3em} % prevent overfull lines

\providecommand{\tightlist}{%
  \setlength{\itemsep}{0pt}\setlength{\parskip}{0pt}}



 


\makeatletter
\@ifpackageloaded{bookmark}{}{\usepackage{bookmark}}
\makeatother
\makeatletter
\@ifpackageloaded{caption}{}{\usepackage{caption}}
\AtBeginDocument{%
\ifdefined\contentsname
  \renewcommand*\contentsname{Table of contents}
\else
  \newcommand\contentsname{Table of contents}
\fi
\ifdefined\listfigurename
  \renewcommand*\listfigurename{List of Figures}
\else
  \newcommand\listfigurename{List of Figures}
\fi
\ifdefined\listtablename
  \renewcommand*\listtablename{List of Tables}
\else
  \newcommand\listtablename{List of Tables}
\fi
\ifdefined\figurename
  \renewcommand*\figurename{Figure}
\else
  \newcommand\figurename{Figure}
\fi
\ifdefined\tablename
  \renewcommand*\tablename{Table}
\else
  \newcommand\tablename{Table}
\fi
}
\@ifpackageloaded{float}{}{\usepackage{float}}
\floatstyle{ruled}
\@ifundefined{c@chapter}{\newfloat{codelisting}{h}{lop}}{\newfloat{codelisting}{h}{lop}[chapter]}
\floatname{codelisting}{Listing}
\newcommand*\listoflistings{\listof{codelisting}{List of Listings}}
\makeatother
\makeatletter
\makeatother
\makeatletter
\@ifpackageloaded{caption}{}{\usepackage{caption}}
\@ifpackageloaded{subcaption}{}{\usepackage{subcaption}}
\makeatother
\usepackage{bookmark}
\IfFileExists{xurl.sty}{\usepackage{xurl}}{} % add URL line breaks if available
\urlstyle{same}
\hypersetup{
  pdftitle={Additory Documentation},
  pdfauthor={Krishnamoorthy Sankaran},
  colorlinks=true,
  linkcolor={Maroon},
  filecolor={Maroon},
  citecolor={Blue},
  urlcolor={Blue},
  pdfcreator={LaTeX via pandoc}}


\title{Additory Documentation}
\usepackage{etoolbox}
\makeatletter
\providecommand{\subtitle}[1]{% add subtitle to \maketitle
  \apptocmd{\@title}{\par {\large #1 \par}}{}{}
}
\makeatother
\subtitle{Elegant data operations for DataFrames}
\author{Krishnamoorthy Sankaran}
\date{2026-03-12}
\begin{document}
\frontmatter
\maketitle

\renewcommand*\contentsname{Table of contents}
{
\hypersetup{linkcolor=}
\setcounter{tocdepth}{2}
\tableofcontents
}

\mainmatter
\bookmarksetup{startatroot}

\chapter*{Welcome}\label{welcome}
\addcontentsline{toc}{chapter}{Welcome}

\markboth{Welcome}{Welcome}

Welcome to the \textbf{Additory Documentation} - your comprehensive
guide to elegant data operations for DataFrames.

\section*{About Additory}\label{about-additory}
\addcontentsline{toc}{section}{About Additory}

\markright{About Additory}

\textbf{Additory} is a Rust-powered Python library that provides
intuitive data transformations, lookups, and synthetic data generation
for both Polars and Pandas DataFrames.

\subsection*{Key Features}\label{key-features}
\addcontentsline{toc}{subsection}{Key Features}

\begin{itemize}
\tightlist
\item
  🔗 \textbf{Intuitive Lookups} - Add columns from external sources with
  simple syntax
\item
  ⚡ \textbf{Powerful Transforms} - Calculate, filter, sort, aggregate
  with mode-based operations
\item
  🎲 \textbf{Synthetic Data} - Generate realistic test data or augment
  existing datasets
\item
  📊 \textbf{Lineage Tracking} - Track data transformations and view
  operation history
\item
  🔍 \textbf{Data Scanning} - Analyze data quality and inspect
  DataFrames
\item
  🚀 \textbf{Rust Performance} - Built with Rust for blazing-fast
  operations
\item
  🐼 \textbf{Polars \& Pandas} - Works seamlessly with both DataFrame
  libraries
\end{itemize}

\section*{Installation}\label{installation}
\addcontentsline{toc}{section}{Installation}

\markright{Installation}

\begin{Shaded}
\begin{Highlighting}[]
\ExtensionTok{pip}\NormalTok{ install additory}
\end{Highlighting}
\end{Shaded}

\textbf{Requirements:} - Python 3.9+ - Polars (required) - Pandas
(optional)

\section*{Quick Start}\label{quick-start}
\addcontentsline{toc}{section}{Quick Start}

\markright{Quick Start}

\begin{Shaded}
\begin{Highlighting}[]
\ImportTok{import}\NormalTok{ additory }\ImportTok{as}\NormalTok{ add}
\ImportTok{import}\NormalTok{ polars }\ImportTok{as}\NormalTok{ pl}

\CommentTok{\# Add data from external sources}
\NormalTok{orders }\OperatorTok{=}\NormalTok{ pl.DataFrame(\{}\StringTok{\textquotesingle{}id\textquotesingle{}}\NormalTok{: [}\DecValTok{1}\NormalTok{, }\DecValTok{2}\NormalTok{], }\StringTok{\textquotesingle{}customer\_id\textquotesingle{}}\NormalTok{: [}\DecValTok{101}\NormalTok{, }\DecValTok{102}\NormalTok{]\})}
\NormalTok{customers }\OperatorTok{=}\NormalTok{ pl.DataFrame(\{}\StringTok{\textquotesingle{}customer\_id\textquotesingle{}}\NormalTok{: [}\DecValTok{101}\NormalTok{, }\DecValTok{102}\NormalTok{], }\StringTok{\textquotesingle{}name\textquotesingle{}}\NormalTok{: [}\StringTok{\textquotesingle{}Alice\textquotesingle{}}\NormalTok{, }\StringTok{\textquotesingle{}Bob\textquotesingle{}}\NormalTok{]\})}
\NormalTok{result }\OperatorTok{=}\NormalTok{ add.to(orders, bring\_from}\OperatorTok{=}\NormalTok{customers, bring}\OperatorTok{=}\NormalTok{[}\StringTok{\textquotesingle{}name\textquotesingle{}}\NormalTok{], against}\OperatorTok{=}\StringTok{\textquotesingle{}customer\_id\textquotesingle{}}\NormalTok{)}

\CommentTok{\# Transform data}
\NormalTok{df }\OperatorTok{=}\NormalTok{ pl.DataFrame(\{}\StringTok{\textquotesingle{}x\textquotesingle{}}\NormalTok{: [}\DecValTok{1}\NormalTok{, }\DecValTok{2}\NormalTok{, }\DecValTok{3}\NormalTok{]\})}
\NormalTok{result }\OperatorTok{=}\NormalTok{ add.transform(}\StringTok{\textquotesingle{}@calc\textquotesingle{}}\NormalTok{, df, strategy}\OperatorTok{=}\NormalTok{\{}\StringTok{\textquotesingle{}x\_squared\textquotesingle{}}\NormalTok{: }\StringTok{\textquotesingle{}x ** 2\textquotesingle{}}\NormalTok{\})}

\CommentTok{\# Generate synthetic data}
\NormalTok{result }\OperatorTok{=}\NormalTok{ add.synthetic(}\StringTok{\textquotesingle{}@new\textquotesingle{}}\NormalTok{, n}\OperatorTok{=}\DecValTok{100}\NormalTok{, strategy}\OperatorTok{=}\NormalTok{\{}\StringTok{\textquotesingle{}age\textquotesingle{}}\NormalTok{: }\StringTok{\textquotesingle{}normal(40, 10)\textquotesingle{}}\NormalTok{\})}
\end{Highlighting}
\end{Shaded}

\section*{Documentation Structure}\label{documentation-structure}
\addcontentsline{toc}{section}{Documentation Structure}

\markright{Documentation Structure}

This book is organized into the following sections:

\subsection*{Core Functions}\label{core-functions}
\addcontentsline{toc}{subsection}{Core Functions}

\begin{enumerate}
\def\labelenumi{\arabic{enumi}.}
\tightlist
\item
  \textbf{add.to()} - Data lookups and joins

  \begin{itemize}
  \tightlist
  \item
    Basic lookups
  \item
    Multiple columns and keys
  \item
    Relationship patterns
  \item
    Aggregation strategies
  \item
    Real-world scenarios
  \end{itemize}
\item
  \textbf{add.transform()} - Data transformations

  \begin{itemize}
  \tightlist
  \item
    Calculations
  \item
    Filtering and sorting
  \item
    Aggregation
  \item
    Advanced modes
  \item
    Real-world workflows
  \end{itemize}
\item
  \textbf{add.synthetic()} - Synthetic data generation

  \begin{itemize}
  \tightlist
  \item
    Basic generation
  \item
    Statistical distributions
  \item
    Data augmentation
  \item
    Real-world use cases
  \end{itemize}
\item
  \textbf{add.scan()} - Data analysis and inspection

  \begin{itemize}
  \tightlist
  \item
    Basic analysis
  \item
    Lineage tracking
  \item
    Real-world analysis
  \end{itemize}
\end{enumerate}

\subsection*{Advanced Topics}\label{advanced-topics}
\addcontentsline{toc}{subsection}{Advanced Topics}

\begin{enumerate}
\def\labelenumi{\arabic{enumi}.}
\setcounter{enumi}{4}
\tightlist
\item
  \textbf{Lineage Tracking} - Understanding data provenance

  \begin{itemize}
  \tightlist
  \item
    Basics of lineage
  \item
    Advanced lineage workflows
  \end{itemize}
\item
  \textbf{Guides \& Troubleshooting} - Common issues and solutions
\end{enumerate}

\section*{Version}\label{version}
\addcontentsline{toc}{section}{Version}

\markright{Version}

Current version: \textbf{0.1.3a10} (Alpha Release)

\section*{Links}\label{links}
\addcontentsline{toc}{section}{Links}

\markright{Links}

\begin{itemize}
\tightlist
\item
  \textbf{PyPI}: \url{https://pypi.org/project/additory/}
\item
  \textbf{GitHub}: \url{https://github.com/sekarkrishna/additory}
\item
  \textbf{Issues}: \url{https://github.com/sekarkrishna/additory/issues}
\end{itemize}

\section*{License}\label{license}
\addcontentsline{toc}{section}{License}

\markright{License}

MIT License - see
\href{https://github.com/sekarkrishna/additory/blob/main/LICENSE}{LICENSE}
for details.

\section*{Author}\label{author}
\addcontentsline{toc}{section}{Author}

\markright{Author}

\textbf{Krishnamoorthy Sankaran}\\
Email: krishnamoorthy.sankaran@sekrad.org

\begin{center}\rule{0.5\linewidth}{0.5pt}\end{center}

\textbf{Made with ❤️ for the data science community}

\part{Getting Started}

\chapter{Introduction}\label{introduction}

\section{What is Additory?}\label{what-is-additory}

Additory is a Rust-powered Python library designed to make data
operations elegant, intuitive, and fast. It provides a unified API for
common data tasks that work seamlessly with both Polars and Pandas
DataFrames.

\section{Philosophy}\label{philosophy}

Additory follows three core principles:

\begin{enumerate}
\def\labelenumi{\arabic{enumi}.}
\tightlist
\item
  \textbf{Simplicity} - Natural language function names and intuitive
  parameters
\item
  \textbf{Performance} - Rust-powered core for blazing-fast operations
\item
  \textbf{Flexibility} - Works with both Polars and Pandas, supports
  multiple modes
\end{enumerate}

\section{Core Functions}\label{core-functions-1}

Additory provides four main functions:

\subsection{add.to()}\label{add.to}

Add columns from external DataFrames with intelligent lookups and
aggregation.

\begin{Shaded}
\begin{Highlighting}[]
\ImportTok{import}\NormalTok{ additory }\ImportTok{as}\NormalTok{ add}
\ImportTok{import}\NormalTok{ polars }\ImportTok{as}\NormalTok{ pl}

\NormalTok{orders }\OperatorTok{=}\NormalTok{ pl.DataFrame(\{}\StringTok{\textquotesingle{}id\textquotesingle{}}\NormalTok{: [}\DecValTok{1}\NormalTok{, }\DecValTok{2}\NormalTok{], }\StringTok{\textquotesingle{}customer\_id\textquotesingle{}}\NormalTok{: [}\DecValTok{101}\NormalTok{, }\DecValTok{102}\NormalTok{]\})}
\NormalTok{customers }\OperatorTok{=}\NormalTok{ pl.DataFrame(\{}\StringTok{\textquotesingle{}customer\_id\textquotesingle{}}\NormalTok{: [}\DecValTok{101}\NormalTok{, }\DecValTok{102}\NormalTok{], }\StringTok{\textquotesingle{}name\textquotesingle{}}\NormalTok{: [}\StringTok{\textquotesingle{}Alice\textquotesingle{}}\NormalTok{, }\StringTok{\textquotesingle{}Bob\textquotesingle{}}\NormalTok{]\})}

\NormalTok{result }\OperatorTok{=}\NormalTok{ add.to(orders, bring\_from}\OperatorTok{=}\NormalTok{customers, bring}\OperatorTok{=}\NormalTok{[}\StringTok{\textquotesingle{}name\textquotesingle{}}\NormalTok{], against}\OperatorTok{=}\StringTok{\textquotesingle{}customer\_id\textquotesingle{}}\NormalTok{)}
\end{Highlighting}
\end{Shaded}

\subsection{add.transform()}\label{add.transform}

Transform data with mode-based operations: calculate, filter, sort,
aggregate, and more.

\begin{Shaded}
\begin{Highlighting}[]
\NormalTok{df }\OperatorTok{=}\NormalTok{ pl.DataFrame(\{}\StringTok{\textquotesingle{}price\textquotesingle{}}\NormalTok{: [}\DecValTok{100}\NormalTok{, }\DecValTok{200}\NormalTok{, }\DecValTok{300}\NormalTok{], }\StringTok{\textquotesingle{}quantity\textquotesingle{}}\NormalTok{: [}\DecValTok{2}\NormalTok{, }\DecValTok{3}\NormalTok{, }\DecValTok{1}\NormalTok{]\})}

\CommentTok{\# Calculate new columns}
\NormalTok{result }\OperatorTok{=}\NormalTok{ add.transform(}\StringTok{\textquotesingle{}@calc\textquotesingle{}}\NormalTok{, df, strategy}\OperatorTok{=}\NormalTok{\{}\StringTok{\textquotesingle{}total\textquotesingle{}}\NormalTok{: }\StringTok{\textquotesingle{}price * quantity\textquotesingle{}}\NormalTok{\})}

\CommentTok{\# Filter data}
\NormalTok{result }\OperatorTok{=}\NormalTok{ add.transform(}\StringTok{\textquotesingle{}@filter\textquotesingle{}}\NormalTok{, df, where}\OperatorTok{=}\StringTok{\textquotesingle{}price \textgreater{} 150\textquotesingle{}}\NormalTok{)}

\CommentTok{\# Aggregate data}
\NormalTok{result }\OperatorTok{=}\NormalTok{ add.transform(}\StringTok{\textquotesingle{}@aggregate\textquotesingle{}}\NormalTok{, df, by}\OperatorTok{=}\StringTok{\textquotesingle{}category\textquotesingle{}}\NormalTok{, strategy}\OperatorTok{=}\NormalTok{\{}\StringTok{\textquotesingle{}price\textquotesingle{}}\NormalTok{: }\StringTok{\textquotesingle{}sum\textquotesingle{}}\NormalTok{\})}
\end{Highlighting}
\end{Shaded}

\subsection{add.synthetic()}\label{add.synthetic}

Generate synthetic data for testing, ML training, or data augmentation.

\begin{Shaded}
\begin{Highlighting}[]
\CommentTok{\# Create synthetic DataFrame}
\NormalTok{result }\OperatorTok{=}\NormalTok{ add.synthetic(}\StringTok{\textquotesingle{}@new\textquotesingle{}}\NormalTok{, n}\OperatorTok{=}\DecValTok{1000}\NormalTok{, strategy}\OperatorTok{=}\NormalTok{\{}
    \StringTok{\textquotesingle{}age\textquotesingle{}}\NormalTok{: }\StringTok{\textquotesingle{}normal(40, 10)\textquotesingle{}}\NormalTok{,}
    \StringTok{\textquotesingle{}salary\textquotesingle{}}\NormalTok{: }\StringTok{\textquotesingle{}normal(75000, 15000)\textquotesingle{}}\NormalTok{,}
    \StringTok{\textquotesingle{}score\textquotesingle{}}\NormalTok{: }\StringTok{\textquotesingle{}uniform(0, 100)\textquotesingle{}}
\NormalTok{\})}

\CommentTok{\# Augment existing data}
\NormalTok{result }\OperatorTok{=}\NormalTok{ add.synthetic(}\StringTok{\textquotesingle{}@augment\textquotesingle{}}\NormalTok{, df, n}\OperatorTok{=}\DecValTok{100}\NormalTok{)}
\end{Highlighting}
\end{Shaded}

\subsection{add.scan()}\label{add.scan}

Analyze data quality and track lineage.

\begin{Shaded}
\begin{Highlighting}[]
\CommentTok{\# Analyze data}
\NormalTok{result }\OperatorTok{=}\NormalTok{ add.scan(}\StringTok{\textquotesingle{}@analyze\textquotesingle{}}\NormalTok{, df)}

\CommentTok{\# View lineage (requires lineage=True in operations)}
\NormalTok{result }\OperatorTok{=}\NormalTok{ add.scan(}\StringTok{\textquotesingle{}@lineage\textquotesingle{}}\NormalTok{, df)}
\end{Highlighting}
\end{Shaded}

\section{Key Features}\label{key-features-1}

\subsection{Lineage Tracking}\label{lineage-tracking}

Track data transformations across operations to understand data
provenance.

\begin{Shaded}
\begin{Highlighting}[]
\CommentTok{\# Enable lineage tracking}
\NormalTok{result }\OperatorTok{=}\NormalTok{ add.to(customers, bring\_from}\OperatorTok{=}\NormalTok{orders, bring}\OperatorTok{=}\NormalTok{[}\StringTok{\textquotesingle{}amount\textquotesingle{}}\NormalTok{], }
\NormalTok{                against}\OperatorTok{=}\StringTok{\textquotesingle{}customer\_id\textquotesingle{}}\NormalTok{, lineage}\OperatorTok{=}\VariableTok{True}\NormalTok{)}

\NormalTok{result }\OperatorTok{=}\NormalTok{ add.transform(}\StringTok{\textquotesingle{}@calc\textquotesingle{}}\NormalTok{, result, strategy}\OperatorTok{=}\NormalTok{\{}\StringTok{\textquotesingle{}total\textquotesingle{}}\NormalTok{: }\StringTok{\textquotesingle{}amount * 1.1\textquotesingle{}}\NormalTok{\}, }
\NormalTok{                       lineage}\OperatorTok{=}\VariableTok{True}\NormalTok{)}

\CommentTok{\# View lineage report}
\NormalTok{lineage\_report }\OperatorTok{=}\NormalTok{ add.scan(}\StringTok{\textquotesingle{}@lineage\textquotesingle{}}\NormalTok{, result)}
\end{Highlighting}
\end{Shaded}

\subsection{Strategy Parameter}\label{strategy-parameter}

Fine-grained control over operations with the \texttt{strategy}
parameter.

\begin{Shaded}
\begin{Highlighting}[]
\CommentTok{\# Aggregation strategies in add.to()}
\NormalTok{strategy}\OperatorTok{=}\NormalTok{\{}\StringTok{\textquotesingle{}amount\textquotesingle{}}\NormalTok{: }\StringTok{\textquotesingle{}sum\textquotesingle{}}\NormalTok{, }\StringTok{\textquotesingle{}date\textquotesingle{}}\NormalTok{: }\StringTok{\textquotesingle{}last\textquotesingle{}}\NormalTok{\}}

\CommentTok{\# Calculation strategies in add.transform()}
\NormalTok{strategy}\OperatorTok{=}\NormalTok{\{}\StringTok{\textquotesingle{}total\textquotesingle{}}\NormalTok{: }\StringTok{\textquotesingle{}price * quantity\textquotesingle{}}\NormalTok{, }\StringTok{\textquotesingle{}discount\textquotesingle{}}\NormalTok{: }\StringTok{\textquotesingle{}total * 0.1\textquotesingle{}}\NormalTok{\}}

\CommentTok{\# Generation strategies in add.synthetic()}
\NormalTok{strategy}\OperatorTok{=}\NormalTok{\{}\StringTok{\textquotesingle{}id\textquotesingle{}}\NormalTok{: }\StringTok{\textquotesingle{}increment\textquotesingle{}}\NormalTok{, }\StringTok{\textquotesingle{}age\textquotesingle{}}\NormalTok{: }\StringTok{\textquotesingle{}normal(40, 10)\textquotesingle{}}\NormalTok{\}}
\end{Highlighting}
\end{Shaded}

\subsection{Mode-Based Operations}\label{mode-based-operations}

Use \texttt{@mode} syntax for different operation types:

\begin{itemize}
\tightlist
\item
  \texttt{@calc} - Calculate new columns
\item
  \texttt{@filter} - Filter rows
\item
  \texttt{@sort} - Sort data
\item
  \texttt{@aggregate} - Group and aggregate
\item
  \texttt{@round} - Round numbers
\item
  \texttt{@transpose} - Transpose DataFrame
\item
  \texttt{@extract} - Extract patterns
\item
  \texttt{@onehotencode} - One-hot encoding
\item
  \texttt{@deduce} - Fill missing values
\item
  \texttt{@new} - Create synthetic data
\item
  \texttt{@augment} - Augment existing data
\item
  \texttt{@analyze} - Analyze data quality
\item
  \texttt{@lineage} - View lineage
\end{itemize}

\section{Performance}\label{performance}

Additory is built with Rust (\textasciitilde95\% of core operations) for
optimal performance:

\begin{itemize}
\tightlist
\item
  3-5x faster transformations vs pure Python
\item
  5-10x faster data joining operations
\item
  10-20x faster synthetic data generation
\item
  Efficient memory usage with Arrow IPC serialization
\end{itemize}

\section{Next Steps}\label{next-steps}

\begin{itemize}
\tightlist
\item
  \textbf{Installation} - Get started with installing additory
\item
  \textbf{Examples} - Explore comprehensive examples for each function
\item
  \textbf{API Reference} - Detailed documentation of all functions and
  parameters
\end{itemize}

\chapter{Installation}\label{installation-1}

\section{Quick Install}\label{quick-install}

Install additory from PyPI using pip:

\begin{Shaded}
\begin{Highlighting}[]
\ExtensionTok{pip}\NormalTok{ install additory}
\end{Highlighting}
\end{Shaded}

\section{Requirements}\label{requirements}

\subsection{Python Version}\label{python-version}

\begin{itemize}
\tightlist
\item
  Python 3.9 or higher
\end{itemize}

\subsection{Required Dependencies}\label{required-dependencies}

\begin{itemize}
\tightlist
\item
  \textbf{polars} \textgreater= 0.19.0 (required for internal
  operations)
\item
  \textbf{pyarrow} \textgreater= 10.0.0 (required for DataFrame
  conversions)
\end{itemize}

\subsection{Optional Dependencies}\label{optional-dependencies}

\begin{itemize}
\tightlist
\item
  \textbf{pandas} \textgreater= 1.5.0 (optional, for pandas DataFrame
  support)
\end{itemize}

\section{Installation Options}\label{installation-options}

\subsection{Standard Installation}\label{standard-installation}

For Polars-only usage:

\begin{Shaded}
\begin{Highlighting}[]
\ExtensionTok{pip}\NormalTok{ install additory}
\end{Highlighting}
\end{Shaded}

\subsection{With Pandas Support}\label{with-pandas-support}

To use additory with pandas DataFrames:

\begin{Shaded}
\begin{Highlighting}[]
\ExtensionTok{pip}\NormalTok{ install additory}\PreprocessorTok{[}\SpecialStringTok{pandas}\PreprocessorTok{]}
\end{Highlighting}
\end{Shaded}

\subsection{With All Optional
Dependencies}\label{with-all-optional-dependencies}

\begin{Shaded}
\begin{Highlighting}[]
\ExtensionTok{pip}\NormalTok{ install additory}\PreprocessorTok{[}\SpecialStringTok{all}\PreprocessorTok{]}
\end{Highlighting}
\end{Shaded}

\subsection{Development Installation}\label{development-installation}

For development with testing tools:

\begin{Shaded}
\begin{Highlighting}[]
\ExtensionTok{pip}\NormalTok{ install additory}\PreprocessorTok{[}\SpecialStringTok{dev}\PreprocessorTok{]}
\end{Highlighting}
\end{Shaded}

\section{Building from Source}\label{building-from-source}

If you want to build additory from source:

\subsection{Prerequisites}\label{prerequisites}

\begin{enumerate}
\def\labelenumi{\arabic{enumi}.}
\tightlist
\item
  \textbf{Rust} - Install from \url{https://rustup.rs/}
\end{enumerate}

\begin{Shaded}
\begin{Highlighting}[]
\ExtensionTok{curl} \AttributeTok{{-}{-}proto} \StringTok{\textquotesingle{}=https\textquotesingle{}} \AttributeTok{{-}{-}tlsv1.2} \AttributeTok{{-}sSf}\NormalTok{ https://sh.rustup.rs }\KeywordTok{|} \FunctionTok{sh}
\end{Highlighting}
\end{Shaded}

\begin{enumerate}
\def\labelenumi{\arabic{enumi}.}
\setcounter{enumi}{1}
\tightlist
\item
  \textbf{Maturin} - Python/Rust build tool
\end{enumerate}

\begin{Shaded}
\begin{Highlighting}[]
\ExtensionTok{pip}\NormalTok{ install maturin}
\end{Highlighting}
\end{Shaded}

\subsection{Build Steps}\label{build-steps}

\begin{Shaded}
\begin{Highlighting}[]
\CommentTok{\# Clone the repository}
\FunctionTok{git}\NormalTok{ clone https://github.com/sekarkrishna/additory.git}
\BuiltInTok{cd}\NormalTok{ additory}

\CommentTok{\# Build the package}
\BuiltInTok{export} \VariableTok{PYO3\_USE\_ABI3\_FORWARD\_COMPATIBILITY}\OperatorTok{=}\NormalTok{1}
\ExtensionTok{maturin}\NormalTok{ build }\AttributeTok{{-}{-}release}

\CommentTok{\# Install locally}
\ExtensionTok{pip}\NormalTok{ install target/wheels/}\PreprocessorTok{*}\NormalTok{.whl}
\end{Highlighting}
\end{Shaded}

\section{Verification}\label{verification}

Verify your installation:

\begin{Shaded}
\begin{Highlighting}[]
\ImportTok{import}\NormalTok{ additory }\ImportTok{as}\NormalTok{ add}
\BuiltInTok{print}\NormalTok{(add.\_\_version\_\_)}
\end{Highlighting}
\end{Shaded}

Expected output:

\begin{verbatim}
0.1.3a10
\end{verbatim}

\section{Quick Test}\label{quick-test}

Run a quick test to ensure everything works:

\begin{Shaded}
\begin{Highlighting}[]
\ImportTok{import}\NormalTok{ additory }\ImportTok{as}\NormalTok{ add}
\ImportTok{import}\NormalTok{ polars }\ImportTok{as}\NormalTok{ pl}

\CommentTok{\# Test add.to()}
\NormalTok{df1 }\OperatorTok{=}\NormalTok{ pl.DataFrame(\{}\StringTok{\textquotesingle{}id\textquotesingle{}}\NormalTok{: [}\DecValTok{1}\NormalTok{, }\DecValTok{2}\NormalTok{], }\StringTok{\textquotesingle{}value\textquotesingle{}}\NormalTok{: [}\DecValTok{10}\NormalTok{, }\DecValTok{20}\NormalTok{]\})}
\NormalTok{df2 }\OperatorTok{=}\NormalTok{ pl.DataFrame(\{}\StringTok{\textquotesingle{}id\textquotesingle{}}\NormalTok{: [}\DecValTok{1}\NormalTok{, }\DecValTok{2}\NormalTok{], }\StringTok{\textquotesingle{}name\textquotesingle{}}\NormalTok{: [}\StringTok{\textquotesingle{}A\textquotesingle{}}\NormalTok{, }\StringTok{\textquotesingle{}B\textquotesingle{}}\NormalTok{]\})}
\NormalTok{result }\OperatorTok{=}\NormalTok{ add.to(df1, bring\_from}\OperatorTok{=}\NormalTok{df2, bring}\OperatorTok{=}\NormalTok{[}\StringTok{\textquotesingle{}name\textquotesingle{}}\NormalTok{], against}\OperatorTok{=}\StringTok{\textquotesingle{}id\textquotesingle{}}\NormalTok{)}
\BuiltInTok{print}\NormalTok{(}\StringTok{"add.to() works:"}\NormalTok{, result.shape)}

\CommentTok{\# Test add.transform()}
\NormalTok{result }\OperatorTok{=}\NormalTok{ add.transform(}\StringTok{\textquotesingle{}@calc\textquotesingle{}}\NormalTok{, df1, strategy}\OperatorTok{=}\NormalTok{\{}\StringTok{\textquotesingle{}doubled\textquotesingle{}}\NormalTok{: }\StringTok{\textquotesingle{}value * 2\textquotesingle{}}\NormalTok{\})}
\BuiltInTok{print}\NormalTok{(}\StringTok{"add.transform() works:"}\NormalTok{, result.shape)}

\CommentTok{\# Test add.synthetic()}
\NormalTok{result }\OperatorTok{=}\NormalTok{ add.synthetic(}\StringTok{\textquotesingle{}@new\textquotesingle{}}\NormalTok{, n}\OperatorTok{=}\DecValTok{10}\NormalTok{, strategy}\OperatorTok{=}\NormalTok{\{}\StringTok{\textquotesingle{}x\textquotesingle{}}\NormalTok{: }\StringTok{\textquotesingle{}normal(0, 1)\textquotesingle{}}\NormalTok{\})}
\BuiltInTok{print}\NormalTok{(}\StringTok{"add.synthetic() works:"}\NormalTok{, result.shape)}

\CommentTok{\# Test add.scan()}
\NormalTok{result }\OperatorTok{=}\NormalTok{ add.scan(}\StringTok{\textquotesingle{}@analyze\textquotesingle{}}\NormalTok{, df1)}
\BuiltInTok{print}\NormalTok{(}\StringTok{"add.scan() works"}\NormalTok{)}
\end{Highlighting}
\end{Shaded}

\section{Troubleshooting}\label{troubleshooting}

\subsection{Python 3.13 Support}\label{python-3.13-support}

If you're using Python 3.13, you may need to set an environment
variable:

\begin{Shaded}
\begin{Highlighting}[]
\BuiltInTok{export} \VariableTok{PYO3\_USE\_ABI3\_FORWARD\_COMPATIBILITY}\OperatorTok{=}\NormalTok{1}
\ExtensionTok{pip}\NormalTok{ install additory}
\end{Highlighting}
\end{Shaded}

\subsection{Import Errors}\label{import-errors}

If you encounter import errors:

\begin{enumerate}
\def\labelenumi{\arabic{enumi}.}
\tightlist
\item
  Ensure polars and pyarrow are installed:
\end{enumerate}

\begin{Shaded}
\begin{Highlighting}[]
\ExtensionTok{pip}\NormalTok{ install polars pyarrow}
\end{Highlighting}
\end{Shaded}

\begin{enumerate}
\def\labelenumi{\arabic{enumi}.}
\setcounter{enumi}{1}
\tightlist
\item
  For pandas support:
\end{enumerate}

\begin{Shaded}
\begin{Highlighting}[]
\ExtensionTok{pip}\NormalTok{ install pandas}
\end{Highlighting}
\end{Shaded}

\subsection{Build Errors}\label{build-errors}

If building from source fails:

\begin{enumerate}
\def\labelenumi{\arabic{enumi}.}
\tightlist
\item
  Ensure Rust is properly installed:
\end{enumerate}

\begin{Shaded}
\begin{Highlighting}[]
\ExtensionTok{rustc} \AttributeTok{{-}{-}version}
\ExtensionTok{cargo} \AttributeTok{{-}{-}version}
\end{Highlighting}
\end{Shaded}

\begin{enumerate}
\def\labelenumi{\arabic{enumi}.}
\setcounter{enumi}{1}
\tightlist
\item
  Update maturin:
\end{enumerate}

\begin{Shaded}
\begin{Highlighting}[]
\ExtensionTok{pip}\NormalTok{ install }\AttributeTok{{-}{-}upgrade}\NormalTok{ maturin}
\end{Highlighting}
\end{Shaded}

\begin{enumerate}
\def\labelenumi{\arabic{enumi}.}
\setcounter{enumi}{2}
\tightlist
\item
  Check for platform-specific issues in the
  \href{https://github.com/sekarkrishna/additory/issues}{GitHub Issues}
\end{enumerate}

\section{Platform Support}\label{platform-support}

Additory is tested on:

\begin{itemize}
\tightlist
\item
  \textbf{Linux} (x86\_64, aarch64)
\item
  \textbf{macOS} (x86\_64, Apple Silicon)
\item
  \textbf{Windows} (x86\_64)
\end{itemize}

\section{Next Steps}\label{next-steps-1}

Now that you have additory installed, explore the examples:

\begin{itemize}
\tightlist
\item
  \textbf{add.to()} - Learn about data lookups and joins
\item
  \textbf{add.transform()} - Discover data transformation modes
\item
  \textbf{add.synthetic()} - Generate synthetic data
\item
  \textbf{add.scan()} - Analyze and inspect DataFrames
\end{itemize}

\part{add.to() - Data Lookups \& Joins}

\chapter{add.to() - Basic Lookup}\label{add.to---basic-lookup}

\section{Overview}\label{overview}

Learn the fundamentals of \texttt{add.to()} by bringing data from one
DataFrame to another using simple lookups. This page covers the most
basic use case: adding a single column from a reference DataFrame to a
target DataFrame based on a matching key.

\textbf{What you'll learn:} - How to bring a single column from one
DataFrame to another - How to specify the key column for matching - How
to read the full function signature - Basic lookup behavior and output

\begin{center}\rule{0.5\linewidth}{0.5pt}\end{center}

\section{Example 1: Add Customer Name to
Orders}\label{example-1-add-customer-name-to-orders}

\textbf{Business Context:} You have an orders table with customer IDs,
and you need to add customer names from a customers table for better
readability in reports.

\textbf{Code:}

\begin{Shaded}
\begin{Highlighting}[]
\ImportTok{import}\NormalTok{ additory }\ImportTok{as}\NormalTok{ add}
\ImportTok{import}\NormalTok{ pandas }\ImportTok{as}\NormalTok{ pd}

\CommentTok{\# Create sample data}
\NormalTok{customers }\OperatorTok{=}\NormalTok{ pd.DataFrame(\{}
    \StringTok{\textquotesingle{}customer\_id\textquotesingle{}}\NormalTok{: [}\DecValTok{1}\NormalTok{, }\DecValTok{2}\NormalTok{, }\DecValTok{3}\NormalTok{],}
    \StringTok{\textquotesingle{}name\textquotesingle{}}\NormalTok{: [}\StringTok{\textquotesingle{}Alice\textquotesingle{}}\NormalTok{, }\StringTok{\textquotesingle{}Bob\textquotesingle{}}\NormalTok{, }\StringTok{\textquotesingle{}Charlie\textquotesingle{}}\NormalTok{]}
\NormalTok{\})}

\NormalTok{orders }\OperatorTok{=}\NormalTok{ pd.DataFrame(\{}
    \StringTok{\textquotesingle{}order\_id\textquotesingle{}}\NormalTok{: [}\DecValTok{101}\NormalTok{, }\DecValTok{102}\NormalTok{, }\DecValTok{103}\NormalTok{],}
    \StringTok{\textquotesingle{}customer\_id\textquotesingle{}}\NormalTok{: [}\DecValTok{1}\NormalTok{, }\DecValTok{2}\NormalTok{, }\DecValTok{3}\NormalTok{],}
    \StringTok{\textquotesingle{}amount\textquotesingle{}}\NormalTok{: [}\DecValTok{100}\NormalTok{, }\DecValTok{200}\NormalTok{, }\DecValTok{150}\NormalTok{]}
\NormalTok{\})}

\CommentTok{\# Bring customer name to orders}
\NormalTok{result }\OperatorTok{=}\NormalTok{ add.to(}
\NormalTok{    orders,                      }\CommentTok{\# Target DataFrame}
\NormalTok{    bring\_from}\OperatorTok{=}\NormalTok{customers,        }\CommentTok{\# Source DataFrame}
\NormalTok{    bring}\OperatorTok{=}\StringTok{\textquotesingle{}name\textquotesingle{}}\NormalTok{,                }\CommentTok{\# Column to bring}
\NormalTok{    against}\OperatorTok{=}\StringTok{\textquotesingle{}customer\_id\textquotesingle{}}        \CommentTok{\# Key to match on}
\NormalTok{)}

\CommentTok{\# Positional parameters (also works without naming certain parameters):}
\CommentTok{\# result = add.to(orders, customers, \textquotesingle{}name\textquotesingle{}, \textquotesingle{}customer\_id\textquotesingle{})}

\BuiltInTok{print}\NormalTok{(result)}
\end{Highlighting}
\end{Shaded}

\textbf{Output:}

\begin{verbatim}
   order_id  customer_id  amount     name
0       101            1     100    Alice
1       102            2     200      Bob
2       103            3     150  Charlie
\end{verbatim}

\textbf{Explanation:} - \texttt{orders} is the target DataFrame where we
want to add data - \texttt{bring\_from=customers} specifies where to get
the data from - \texttt{bring=\textquotesingle{}name\textquotesingle{}}
tells additory to bring the `name' column -
\texttt{against=\textquotesingle{}customer\_id\textquotesingle{}}
specifies the key column to match on - The result has all original
columns from \texttt{orders} plus the new `name' column

\textbf{Note:} This also works with polars DataFrames. Simply replace
\texttt{pd.DataFrame} with \texttt{pl.DataFrame} and import
\texttt{polars\ as\ pl}.

\begin{center}\rule{0.5\linewidth}{0.5pt}\end{center}

\section{Example 2: Add Product Price to Cart
Items}\label{example-2-add-product-price-to-cart-items}

\textbf{Business Context:} Your shopping cart has product IDs, and you
need to add the current price from the products catalog.

\textbf{Code:}

\begin{Shaded}
\begin{Highlighting}[]
\ImportTok{import}\NormalTok{ additory }\ImportTok{as}\NormalTok{ add}
\ImportTok{import}\NormalTok{ pandas }\ImportTok{as}\NormalTok{ pd}

\CommentTok{\# Products catalog}
\NormalTok{products }\OperatorTok{=}\NormalTok{ pd.DataFrame(\{}
    \StringTok{\textquotesingle{}product\_id\textquotesingle{}}\NormalTok{: [}\DecValTok{101}\NormalTok{, }\DecValTok{102}\NormalTok{, }\DecValTok{103}\NormalTok{, }\DecValTok{104}\NormalTok{],}
    \StringTok{\textquotesingle{}price\textquotesingle{}}\NormalTok{: [}\FloatTok{29.99}\NormalTok{, }\FloatTok{49.99}\NormalTok{, }\FloatTok{19.99}\NormalTok{, }\FloatTok{99.99}\NormalTok{]}
\NormalTok{\})}

\CommentTok{\# Shopping cart}
\NormalTok{cart }\OperatorTok{=}\NormalTok{ pd.DataFrame(\{}
    \StringTok{\textquotesingle{}cart\_item\_id\textquotesingle{}}\NormalTok{: [}\DecValTok{1}\NormalTok{, }\DecValTok{2}\NormalTok{, }\DecValTok{3}\NormalTok{],}
    \StringTok{\textquotesingle{}product\_id\textquotesingle{}}\NormalTok{: [}\DecValTok{102}\NormalTok{, }\DecValTok{101}\NormalTok{, }\DecValTok{104}\NormalTok{],}
    \StringTok{\textquotesingle{}quantity\textquotesingle{}}\NormalTok{: [}\DecValTok{2}\NormalTok{, }\DecValTok{1}\NormalTok{, }\DecValTok{1}\NormalTok{]}
\NormalTok{\})}

\CommentTok{\# Bring price to cart}
\NormalTok{result }\OperatorTok{=}\NormalTok{ add.to(}
\NormalTok{    cart,}
\NormalTok{    bring\_from}\OperatorTok{=}\NormalTok{products,}
\NormalTok{    bring}\OperatorTok{=}\StringTok{\textquotesingle{}price\textquotesingle{}}\NormalTok{,}
\NormalTok{    against}\OperatorTok{=}\StringTok{\textquotesingle{}product\_id\textquotesingle{}}
\NormalTok{)}

\CommentTok{\# Positional parameters (also works without naming certain parameters):}
\CommentTok{\# result = add.to(cart, products, \textquotesingle{}price\textquotesingle{}, \textquotesingle{}product\_id\textquotesingle{})}

\BuiltInTok{print}\NormalTok{(result)}
\end{Highlighting}
\end{Shaded}

\textbf{Output:}

\begin{verbatim}
   cart_item_id  product_id  quantity  price
0             1         102         2  49.99
1             2         101         1  29.99
2             3         104         1  99.99
\end{verbatim}

\textbf{Explanation:} - Each cart item now has its corresponding price -
The matching is done automatically based on \texttt{product\_id} - The
order of rows in the target DataFrame is preserved - Unmatched products
in the catalog (103) are ignored

\begin{center}\rule{0.5\linewidth}{0.5pt}\end{center}

\section{Example 3: Add Department to
Employees}\label{example-3-add-department-to-employees}

\textbf{Business Context:} Your employee records have department codes,
and you need to add the full department name from a departments lookup
table.

\textbf{Code:}

\begin{Shaded}
\begin{Highlighting}[]
\ImportTok{import}\NormalTok{ additory }\ImportTok{as}\NormalTok{ add}
\ImportTok{import}\NormalTok{ pandas }\ImportTok{as}\NormalTok{ pd}

\CommentTok{\# Departments lookup}
\NormalTok{departments }\OperatorTok{=}\NormalTok{ pd.DataFrame(\{}
    \StringTok{\textquotesingle{}dept\_code\textquotesingle{}}\NormalTok{: [}\StringTok{\textquotesingle{}ENG\textquotesingle{}}\NormalTok{, }\StringTok{\textquotesingle{}SAL\textquotesingle{}}\NormalTok{, }\StringTok{\textquotesingle{}MKT\textquotesingle{}}\NormalTok{, }\StringTok{\textquotesingle{}HR\textquotesingle{}}\NormalTok{],}
    \StringTok{\textquotesingle{}dept\_name\textquotesingle{}}\NormalTok{: [}\StringTok{\textquotesingle{}Engineering\textquotesingle{}}\NormalTok{, }\StringTok{\textquotesingle{}Sales\textquotesingle{}}\NormalTok{, }\StringTok{\textquotesingle{}Marketing\textquotesingle{}}\NormalTok{, }\StringTok{\textquotesingle{}Human Resources\textquotesingle{}}\NormalTok{]}
\NormalTok{\})}

\CommentTok{\# Employees (each in different department)}
\NormalTok{employees }\OperatorTok{=}\NormalTok{ pd.DataFrame(\{}
    \StringTok{\textquotesingle{}emp\_id\textquotesingle{}}\NormalTok{: [}\DecValTok{1001}\NormalTok{, }\DecValTok{1002}\NormalTok{, }\DecValTok{1003}\NormalTok{, }\DecValTok{1004}\NormalTok{],}
    \StringTok{\textquotesingle{}name\textquotesingle{}}\NormalTok{: [}\StringTok{\textquotesingle{}Alice\textquotesingle{}}\NormalTok{, }\StringTok{\textquotesingle{}Bob\textquotesingle{}}\NormalTok{, }\StringTok{\textquotesingle{}Charlie\textquotesingle{}}\NormalTok{, }\StringTok{\textquotesingle{}David\textquotesingle{}}\NormalTok{],}
    \StringTok{\textquotesingle{}dept\_code\textquotesingle{}}\NormalTok{: [}\StringTok{\textquotesingle{}ENG\textquotesingle{}}\NormalTok{, }\StringTok{\textquotesingle{}SAL\textquotesingle{}}\NormalTok{, }\StringTok{\textquotesingle{}MKT\textquotesingle{}}\NormalTok{, }\StringTok{\textquotesingle{}HR\textquotesingle{}}\NormalTok{]}
\NormalTok{\})}

\CommentTok{\# Bring department name to employees}
\NormalTok{result }\OperatorTok{=}\NormalTok{ add.to(}
\NormalTok{    employees,}
\NormalTok{    bring\_from}\OperatorTok{=}\NormalTok{departments,}
\NormalTok{    bring}\OperatorTok{=}\StringTok{\textquotesingle{}dept\_name\textquotesingle{}}\NormalTok{,}
\NormalTok{    against}\OperatorTok{=}\StringTok{\textquotesingle{}dept\_code\textquotesingle{}}
\NormalTok{)}

\CommentTok{\# Positional parameters (also works without naming certain parameters):}
\CommentTok{\# result = add.to(employees, departments, \textquotesingle{}dept\_name\textquotesingle{}, \textquotesingle{}dept\_code\textquotesingle{})}

\BuiltInTok{print}\NormalTok{(result)}
\end{Highlighting}
\end{Shaded}

\textbf{Output:}

\begin{verbatim}
   emp_id     name dept_code       dept_name
0    1001    Alice       ENG     Engineering
1    1002      Bob       SAL           Sales
2    1003  Charlie       MKT       Marketing
3    1004    David        HR  Human Resources
\end{verbatim}

\textbf{Explanation:} - Each employee has a unique department in this
example - The lookup table acts as a reference - each unique dept\_code
maps to one dept\_name - This is a classic ``lookup'' or ``dimension
table'' pattern in data warehousing - When multiple employees share the
same department, you'll need aggregation strategies (covered in
\href{examples-to-04-strategy.html}{Aggregation Strategies})

\begin{center}\rule{0.5\linewidth}{0.5pt}\end{center}

\section{Understanding the Signature}\label{understanding-the-signature}

The full \texttt{add.to()} signature for basic usage:

\begin{Shaded}
\begin{Highlighting}[]
\NormalTok{result }\OperatorTok{=}\NormalTok{ add.to(}
\NormalTok{    target\_df,                   }\CommentTok{\# Required: DataFrame to add columns to}
\NormalTok{    bring\_from}\OperatorTok{=}\NormalTok{source\_df,        }\CommentTok{\# Required: DataFrame to bring columns from}
\NormalTok{    bring}\OperatorTok{=}\StringTok{\textquotesingle{}column\_name\textquotesingle{}}\NormalTok{,         }\CommentTok{\# Required: Column to bring (string)}
\NormalTok{    against}\OperatorTok{=}\StringTok{\textquotesingle{}key\_column\textquotesingle{}}\NormalTok{,        }\CommentTok{\# Required: Key column to match on (string)}
\NormalTok{    position}\OperatorTok{=}\VariableTok{None}\NormalTok{,               }\CommentTok{\# Optional: Where to place new column}
\NormalTok{    strategy}\OperatorTok{=}\VariableTok{None}\NormalTok{,               }\CommentTok{\# Optional: Aggregation strategy (covered later)}
\NormalTok{    join\_type}\OperatorTok{=}\StringTok{\textquotesingle{}lookup\textquotesingle{}}\NormalTok{,          }\CommentTok{\# Optional: Type of join (default: \textquotesingle{}lookup\textquotesingle{})}
\NormalTok{    logging}\OperatorTok{=}\VariableTok{False}\NormalTok{,               }\CommentTok{\# Optional: Enable detailed logs}
\NormalTok{    lineage}\OperatorTok{=}\VariableTok{False}\NormalTok{,               }\CommentTok{\# Optional: Track data lineage}
\NormalTok{    as\_type}\OperatorTok{=}\VariableTok{None}                 \CommentTok{\# Optional: Force output type}
\NormalTok{)}
\end{Highlighting}
\end{Shaded}

\textbf{For basic lookups, you only need the first 4 parameters.}

\begin{center}\rule{0.5\linewidth}{0.5pt}\end{center}

\section{Key Takeaways}\label{key-takeaways}

\begin{itemize}
\tightlist
\item
  \texttt{add.to()} brings columns from one DataFrame to another based
  on matching keys
\item
  The first parameter is always the target DataFrame (where you want to
  add data)
\item
  Use \texttt{bring\_from=} to specify the source DataFrame
\item
  Use \texttt{bring=} to specify which column to bring
\item
  Use \texttt{against=} to specify the key column for matching
\item
  The result preserves the order and structure of the target DataFrame
\item
  Unmatched rows in the source are ignored (lookup behavior)
\end{itemize}

\begin{center}\rule{0.5\linewidth}{0.5pt}\end{center}

\section{Common Questions}\label{common-questions}

\textbf{Q: What if the key column has different names in each
DataFrame?}\\
A: The key column must have the same name in both DataFrames for basic
usage. We'll cover handling different key names in advanced examples.

\textbf{Q: What happens if a key in the target doesn't exist in the
source?}\\
A: The brought column will have \texttt{NaN} (null) for that row.

\textbf{Q: Can I bring multiple columns at once?}\\
A: Yes! We'll cover that in the next page:
\href{examples-to-02-multiple.html}{Multiple Columns \& Keys}.

\textbf{Q: What if multiple rows in the source match the same key?}\\
A: By default, you'll get an error. We'll cover aggregation strategies
in \href{examples-to-04-strategy.html}{Aggregation Strategies}.

\begin{center}\rule{0.5\linewidth}{0.5pt}\end{center}

\section{Next Steps}\label{next-steps-2}

\begin{itemize}
\tightlist
\item
  \textbf{\href{examples-to-02-multiple.html}{Multiple Columns \& Keys}}
  - Learn to bring multiple columns and use multiple keys
\item
  \textbf{\href{reference-to-basic.html}{API Reference}} - Complete
  \texttt{add.to()} documentation
\end{itemize}

\begin{center}\rule{0.5\linewidth}{0.5pt}\end{center}

\textbf{Version:} 0.1.3\\
\textbf{Last Updated:} March 9, 2026

\chapter{add.to() - Multiple Columns \&
Keys}\label{add.to---multiple-columns-keys}

\section{Overview}\label{overview-1}

Learn how to bring multiple columns at once and use multiple keys for
matching. This page builds on basic lookups by showing you how to work
with more complex data relationships.

\textbf{What you'll learn:} - How to bring multiple columns in a single
operation - How to use multiple keys for matching (composite keys) - How
to combine multiple columns and multiple keys

\textbf{Prerequisites:} - \href{examples-to-01-basic.html}{Basic Lookup}
- Understanding of single column lookups

\begin{center}\rule{0.5\linewidth}{0.5pt}\end{center}

\section{Example 1: Bring Multiple
Columns}\label{example-1-bring-multiple-columns}

\textbf{Business Context:} You have orders with customer IDs, and you
need to add both the customer name and email address for shipping
labels.

\textbf{Code:}

\begin{Shaded}
\begin{Highlighting}[]
\ImportTok{import}\NormalTok{ additory }\ImportTok{as}\NormalTok{ add}
\ImportTok{import}\NormalTok{ pandas }\ImportTok{as}\NormalTok{ pd}

\CommentTok{\# Customers with multiple attributes}
\NormalTok{customers }\OperatorTok{=}\NormalTok{ pd.DataFrame(\{}
    \StringTok{\textquotesingle{}customer\_id\textquotesingle{}}\NormalTok{: [}\DecValTok{1}\NormalTok{, }\DecValTok{2}\NormalTok{, }\DecValTok{3}\NormalTok{],}
    \StringTok{\textquotesingle{}name\textquotesingle{}}\NormalTok{: [}\StringTok{\textquotesingle{}Alice\textquotesingle{}}\NormalTok{, }\StringTok{\textquotesingle{}Bob\textquotesingle{}}\NormalTok{, }\StringTok{\textquotesingle{}Charlie\textquotesingle{}}\NormalTok{],}
    \StringTok{\textquotesingle{}email\textquotesingle{}}\NormalTok{: [}\StringTok{\textquotesingle{}alice@email.com\textquotesingle{}}\NormalTok{, }\StringTok{\textquotesingle{}bob@email.com\textquotesingle{}}\NormalTok{, }\StringTok{\textquotesingle{}charlie@email.com\textquotesingle{}}\NormalTok{],}
    \StringTok{\textquotesingle{}city\textquotesingle{}}\NormalTok{: [}\StringTok{\textquotesingle{}NYC\textquotesingle{}}\NormalTok{, }\StringTok{\textquotesingle{}LA\textquotesingle{}}\NormalTok{, }\StringTok{\textquotesingle{}Chicago\textquotesingle{}}\NormalTok{]}
\NormalTok{\})}

\CommentTok{\# Orders}
\NormalTok{orders }\OperatorTok{=}\NormalTok{ pd.DataFrame(\{}
    \StringTok{\textquotesingle{}order\_id\textquotesingle{}}\NormalTok{: [}\DecValTok{101}\NormalTok{, }\DecValTok{102}\NormalTok{, }\DecValTok{103}\NormalTok{],}
    \StringTok{\textquotesingle{}customer\_id\textquotesingle{}}\NormalTok{: [}\DecValTok{1}\NormalTok{, }\DecValTok{2}\NormalTok{, }\DecValTok{3}\NormalTok{],}
    \StringTok{\textquotesingle{}amount\textquotesingle{}}\NormalTok{: [}\DecValTok{100}\NormalTok{, }\DecValTok{200}\NormalTok{, }\DecValTok{150}\NormalTok{]}
\NormalTok{\})}

\CommentTok{\# Bring multiple columns using a list}
\NormalTok{result }\OperatorTok{=}\NormalTok{ add.to(}
\NormalTok{    orders,}
\NormalTok{    bring\_from}\OperatorTok{=}\NormalTok{customers,}
\NormalTok{    bring}\OperatorTok{=}\NormalTok{[}\StringTok{\textquotesingle{}name\textquotesingle{}}\NormalTok{, }\StringTok{\textquotesingle{}email\textquotesingle{}}\NormalTok{],      }\CommentTok{\# List of columns to bring}
\NormalTok{    against}\OperatorTok{=}\StringTok{\textquotesingle{}customer\_id\textquotesingle{}}
\NormalTok{)}

\CommentTok{\# Positional parameters (also works without naming certain parameters):}
\CommentTok{\# result = add.to(orders, customers, [\textquotesingle{}name\textquotesingle{}, \textquotesingle{}email\textquotesingle{}], \textquotesingle{}customer\_id\textquotesingle{})}

\BuiltInTok{print}\NormalTok{(result)}
\end{Highlighting}
\end{Shaded}

\textbf{Output:}

\begin{verbatim}
   order_id  customer_id  amount     name              email
0       101            1     100    Alice  alice@email.com
1       102            2     200      Bob    bob@email.com
2       103            3     150  Charlie charlie@email.com
\end{verbatim}

\textbf{Explanation:} - Use a list
\texttt{{[}\textquotesingle{}name\textquotesingle{},\ \textquotesingle{}email\textquotesingle{}{]}}
to bring multiple columns - Both columns are added in a single operation
- The `city' column is not brought because it's not in the list - This
is more efficient than calling \texttt{add.to()} twice

\textbf{Note:} This also works with polars DataFrames.

\begin{center}\rule{0.5\linewidth}{0.5pt}\end{center}

\section{Example 2: Multiple Keys (Composite
Keys)}\label{example-2-multiple-keys-composite-keys}

\textbf{Business Context:} Your sales data is uniquely identified by
both region AND product. You need to match targets using both keys
together.

\textbf{Code:}

\begin{Shaded}
\begin{Highlighting}[]
\ImportTok{import}\NormalTok{ additory }\ImportTok{as}\NormalTok{ add}
\ImportTok{import}\NormalTok{ pandas }\ImportTok{as}\NormalTok{ pd}

\CommentTok{\# Sales data (identified by region + product)}
\NormalTok{sales }\OperatorTok{=}\NormalTok{ pd.DataFrame(\{}
    \StringTok{\textquotesingle{}region\textquotesingle{}}\NormalTok{: [}\StringTok{\textquotesingle{}East\textquotesingle{}}\NormalTok{, }\StringTok{\textquotesingle{}West\textquotesingle{}}\NormalTok{, }\StringTok{\textquotesingle{}East\textquotesingle{}}\NormalTok{, }\StringTok{\textquotesingle{}West\textquotesingle{}}\NormalTok{],}
    \StringTok{\textquotesingle{}product\textquotesingle{}}\NormalTok{: [}\StringTok{\textquotesingle{}A\textquotesingle{}}\NormalTok{, }\StringTok{\textquotesingle{}A\textquotesingle{}}\NormalTok{, }\StringTok{\textquotesingle{}B\textquotesingle{}}\NormalTok{, }\StringTok{\textquotesingle{}B\textquotesingle{}}\NormalTok{],}
    \StringTok{\textquotesingle{}revenue\textquotesingle{}}\NormalTok{: [}\DecValTok{1000}\NormalTok{, }\DecValTok{1500}\NormalTok{, }\DecValTok{800}\NormalTok{, }\DecValTok{1200}\NormalTok{]}
\NormalTok{\})}

\CommentTok{\# Targets (also identified by region + product)}
\NormalTok{targets }\OperatorTok{=}\NormalTok{ pd.DataFrame(\{}
    \StringTok{\textquotesingle{}region\textquotesingle{}}\NormalTok{: [}\StringTok{\textquotesingle{}East\textquotesingle{}}\NormalTok{, }\StringTok{\textquotesingle{}West\textquotesingle{}}\NormalTok{, }\StringTok{\textquotesingle{}East\textquotesingle{}}\NormalTok{, }\StringTok{\textquotesingle{}West\textquotesingle{}}\NormalTok{],}
    \StringTok{\textquotesingle{}product\textquotesingle{}}\NormalTok{: [}\StringTok{\textquotesingle{}A\textquotesingle{}}\NormalTok{, }\StringTok{\textquotesingle{}A\textquotesingle{}}\NormalTok{, }\StringTok{\textquotesingle{}B\textquotesingle{}}\NormalTok{, }\StringTok{\textquotesingle{}B\textquotesingle{}}\NormalTok{],}
    \StringTok{\textquotesingle{}target\textquotesingle{}}\NormalTok{: [}\DecValTok{900}\NormalTok{, }\DecValTok{1400}\NormalTok{, }\DecValTok{750}\NormalTok{, }\DecValTok{1100}\NormalTok{]}
\NormalTok{\})}

\CommentTok{\# Match using both region AND product}
\NormalTok{result }\OperatorTok{=}\NormalTok{ add.to(}
\NormalTok{    sales,}
\NormalTok{    bring\_from}\OperatorTok{=}\NormalTok{targets,}
\NormalTok{    bring}\OperatorTok{=}\StringTok{\textquotesingle{}target\textquotesingle{}}\NormalTok{,}
\NormalTok{    against}\OperatorTok{=}\NormalTok{[}\StringTok{\textquotesingle{}region\textquotesingle{}}\NormalTok{, }\StringTok{\textquotesingle{}product\textquotesingle{}}\NormalTok{]    }\CommentTok{\# List of key columns}
\NormalTok{)}

\CommentTok{\# Positional parameters (also works without naming certain parameters):}
\CommentTok{\# result = add.to(sales, targets, \textquotesingle{}target\textquotesingle{}, [\textquotesingle{}region\textquotesingle{}, \textquotesingle{}product\textquotesingle{}])}

\BuiltInTok{print}\NormalTok{(result)}
\end{Highlighting}
\end{Shaded}

\textbf{Output:}

\begin{verbatim}
  region product  revenue  target
0   East       A     1000     900
1   West       A     1500    1400
2   East       B      800     750
3   West       B     1200    1100
\end{verbatim}

\textbf{Explanation:} - Use a list
\texttt{{[}\textquotesingle{}region\textquotesingle{},\ \textquotesingle{}product\textquotesingle{}{]}}
to match on multiple keys - Both keys must match for a successful lookup
- This is called a ``composite key'' in database terminology - East-A
matches East-A, but East-A does NOT match West-A

\textbf{Note:} This also works with polars DataFrames.

\begin{center}\rule{0.5\linewidth}{0.5pt}\end{center}

\section{Example 3: Multiple Columns AND Multiple
Keys}\label{example-3-multiple-columns-and-multiple-keys}

\textbf{Business Context:} Your transaction data is identified by date
and store\_id. You need to add manager name and region information using
both keys.

\textbf{Code:}

\begin{Shaded}
\begin{Highlighting}[]
\ImportTok{import}\NormalTok{ additory }\ImportTok{as}\NormalTok{ add}
\ImportTok{import}\NormalTok{ pandas }\ImportTok{as}\NormalTok{ pd}

\CommentTok{\# Transactions (identified by date + store\_id)}
\NormalTok{transactions }\OperatorTok{=}\NormalTok{ pd.DataFrame(\{}
    \StringTok{\textquotesingle{}date\textquotesingle{}}\NormalTok{: [}\StringTok{\textquotesingle{}2024{-}01{-}01\textquotesingle{}}\NormalTok{, }\StringTok{\textquotesingle{}2024{-}01{-}01\textquotesingle{}}\NormalTok{, }\StringTok{\textquotesingle{}2024{-}01{-}02\textquotesingle{}}\NormalTok{],}
    \StringTok{\textquotesingle{}store\_id\textquotesingle{}}\NormalTok{: [}\DecValTok{1}\NormalTok{, }\DecValTok{2}\NormalTok{, }\DecValTok{1}\NormalTok{],}
    \StringTok{\textquotesingle{}sales\textquotesingle{}}\NormalTok{: [}\DecValTok{500}\NormalTok{, }\DecValTok{600}\NormalTok{, }\DecValTok{550}\NormalTok{]}
\NormalTok{\})}

\CommentTok{\# Store information (also by date + store\_id)}
\NormalTok{store\_info }\OperatorTok{=}\NormalTok{ pd.DataFrame(\{}
    \StringTok{\textquotesingle{}date\textquotesingle{}}\NormalTok{: [}\StringTok{\textquotesingle{}2024{-}01{-}01\textquotesingle{}}\NormalTok{, }\StringTok{\textquotesingle{}2024{-}01{-}01\textquotesingle{}}\NormalTok{, }\StringTok{\textquotesingle{}2024{-}01{-}02\textquotesingle{}}\NormalTok{],}
    \StringTok{\textquotesingle{}store\_id\textquotesingle{}}\NormalTok{: [}\DecValTok{1}\NormalTok{, }\DecValTok{2}\NormalTok{, }\DecValTok{1}\NormalTok{],}
    \StringTok{\textquotesingle{}manager\textquotesingle{}}\NormalTok{: [}\StringTok{\textquotesingle{}Alice\textquotesingle{}}\NormalTok{, }\StringTok{\textquotesingle{}Bob\textquotesingle{}}\NormalTok{, }\StringTok{\textquotesingle{}Alice\textquotesingle{}}\NormalTok{],}
    \StringTok{\textquotesingle{}region\textquotesingle{}}\NormalTok{: [}\StringTok{\textquotesingle{}North\textquotesingle{}}\NormalTok{, }\StringTok{\textquotesingle{}South\textquotesingle{}}\NormalTok{, }\StringTok{\textquotesingle{}North\textquotesingle{}}\NormalTok{]}
\NormalTok{\})}

\CommentTok{\# Bring multiple columns using multiple keys}
\NormalTok{result }\OperatorTok{=}\NormalTok{ add.to(}
\NormalTok{    transactions,}
\NormalTok{    bring\_from}\OperatorTok{=}\NormalTok{store\_info,}
\NormalTok{    bring}\OperatorTok{=}\NormalTok{[}\StringTok{\textquotesingle{}manager\textquotesingle{}}\NormalTok{, }\StringTok{\textquotesingle{}region\textquotesingle{}}\NormalTok{],        }\CommentTok{\# Multiple columns}
\NormalTok{    against}\OperatorTok{=}\NormalTok{[}\StringTok{\textquotesingle{}date\textquotesingle{}}\NormalTok{, }\StringTok{\textquotesingle{}store\_id\textquotesingle{}}\NormalTok{]        }\CommentTok{\# Multiple keys}
\NormalTok{)}

\CommentTok{\# Positional parameters (also works without naming certain parameters):}
\CommentTok{\# result = add.to(transactions, store\_info, [\textquotesingle{}manager\textquotesingle{}, \textquotesingle{}region\textquotesingle{}], [\textquotesingle{}date\textquotesingle{}, \textquotesingle{}store\_id\textquotesingle{}])}

\BuiltInTok{print}\NormalTok{(result)}
\end{Highlighting}
\end{Shaded}

\textbf{Output:}

\begin{verbatim}
         date  store_id  sales manager region
0  2024-01-01         1    500   Alice  North
1  2024-01-01         2    600     Bob  South
2  2024-01-02         1    550   Alice  North
\end{verbatim}

\textbf{Explanation:} - Combines both techniques: multiple columns AND
multiple keys - Each row is matched on BOTH date AND store\_id - Both
manager and region are brought together - Store 1 on different dates can
have different managers (as shown)

\textbf{Note:} This also works with polars DataFrames.

\begin{center}\rule{0.5\linewidth}{0.5pt}\end{center}

\section{Understanding Lists in
add.to()}\label{understanding-lists-in-add.to}

\subsection{Single vs Multiple Values}\label{single-vs-multiple-values}

\begin{Shaded}
\begin{Highlighting}[]
\CommentTok{\# Single column (string)}
\NormalTok{bring}\OperatorTok{=}\StringTok{\textquotesingle{}name\textquotesingle{}}              \CommentTok{\# Brings one column}

\CommentTok{\# Multiple columns (list)}
\NormalTok{bring}\OperatorTok{=}\NormalTok{[}\StringTok{\textquotesingle{}name\textquotesingle{}}\NormalTok{, }\StringTok{\textquotesingle{}email\textquotesingle{}}\NormalTok{]   }\CommentTok{\# Brings two columns}

\CommentTok{\# Single key (string)}
\NormalTok{against}\OperatorTok{=}\StringTok{\textquotesingle{}customer\_id\textquotesingle{}}     \CommentTok{\# Matches on one column}

\CommentTok{\# Multiple keys (list)}
\NormalTok{against}\OperatorTok{=}\NormalTok{[}\StringTok{\textquotesingle{}region\textquotesingle{}}\NormalTok{, }\StringTok{\textquotesingle{}product\textquotesingle{}}\NormalTok{]  }\CommentTok{\# Matches on two columns}
\end{Highlighting}
\end{Shaded}

\subsection{When to Use Multiple Keys}\label{when-to-use-multiple-keys}

Use multiple keys when: - Your data has composite primary keys - A
single column isn't unique enough - You need to match on multiple
dimensions (e.g., time + location)

\textbf{Example scenarios:} - Sales by region + product - Transactions
by date + store - Inventory by warehouse + SKU - Prices by date +
product

\begin{center}\rule{0.5\linewidth}{0.5pt}\end{center}

\section{Key Takeaways}\label{key-takeaways-1}

\begin{itemize}
\tightlist
\item
  Use lists
  \texttt{{[}\textquotesingle{}col1\textquotesingle{},\ \textquotesingle{}col2\textquotesingle{}{]}}
  to bring multiple columns at once
\item
  Use lists
  \texttt{{[}\textquotesingle{}key1\textquotesingle{},\ \textquotesingle{}key2\textquotesingle{}{]}}
  to match on multiple keys (composite keys)
\item
  You can combine both: multiple columns AND multiple keys
\item
  All keys must match for a successful lookup
\item
  This is more efficient than multiple separate \texttt{add.to()} calls
\end{itemize}

\begin{center}\rule{0.5\linewidth}{0.5pt}\end{center}

\section{Common Questions}\label{common-questions-1}

\textbf{Q: Can I bring all columns from the reference DataFrame?}\\
A: Not directly. You must specify which columns to bring. This is
intentional to keep operations explicit and avoid accidentally bringing
unwanted columns.

\textbf{Q: What if the key columns have different names in each
DataFrame?}\\
A: For basic usage, key columns must have the same names. We'll cover
handling different key names in advanced examples.

\textbf{Q: Is there a limit to how many columns or keys I can use?}\\
A: No practical limit, but keep it reasonable for performance and
readability.

\textbf{Q: What happens if some rows don't match on all keys?}\\
A: Those rows will have \texttt{NaN} (null) values in the brought
columns, just like single-key lookups.

\begin{center}\rule{0.5\linewidth}{0.5pt}\end{center}

\section{Next Steps}\label{next-steps-3}

\begin{itemize}
\tightlist
\item
  \textbf{\href{examples-to-03-patterns.html}{One-to-Many \& Many-to-One
  Patterns}} - Learn advanced data relationship patterns
\item
  \textbf{\href{examples-to-01-basic.html}{Basic Lookup}} - Review
  single column lookups
\item
  \textbf{\href{reference-to-basic.html}{API Reference}} - Complete
  \texttt{add.to()} documentation
\end{itemize}

\begin{center}\rule{0.5\linewidth}{0.5pt}\end{center}

\textbf{Version:} 0.1.3\\
\textbf{Last Updated:} March 9, 2026

\chapter{add.to() - One-to-Many \& Many-to-One
Patterns}\label{add.to---one-to-many-many-to-one-patterns}

\section{Overview}\label{overview-2}

Learn how to work with lists of DataFrames to handle complex data
relationships. This page introduces powerful patterns for working with
multiple datasets simultaneously.

\textbf{What you'll learn:} - One-to-Many: Single target with multiple
reference DataFrames - Many-to-One: Multiple targets with a single
reference DataFrame - Many-to-Many: Multiple targets with multiple
references - When to use each pattern

\textbf{Prerequisites:} - \href{examples-to-01-basic.html}{Basic Lookup}
- Understanding of single DataFrame operations -
\href{examples-to-02-multiple.html}{Multiple Columns \& Keys} - Working
with multiple columns

\begin{center}\rule{0.5\linewidth}{0.5pt}\end{center}

\section{Example 1: One-to-Many
Pattern}\label{example-1-one-to-many-pattern}

\textbf{Business Context:} You have a customer list and need to
aggregate order amounts from multiple months (January and February).
Each month's orders are in separate DataFrames.

\textbf{Code:}

\begin{Shaded}
\begin{Highlighting}[]
\ImportTok{import}\NormalTok{ additory }\ImportTok{as}\NormalTok{ add}
\ImportTok{import}\NormalTok{ pandas }\ImportTok{as}\NormalTok{ pd}

\CommentTok{\# Single target: Customer list}
\NormalTok{customers }\OperatorTok{=}\NormalTok{ pd.DataFrame(\{}
    \StringTok{\textquotesingle{}customer\_id\textquotesingle{}}\NormalTok{: [}\DecValTok{1}\NormalTok{, }\DecValTok{2}\NormalTok{, }\DecValTok{3}\NormalTok{],}
    \StringTok{\textquotesingle{}name\textquotesingle{}}\NormalTok{: [}\StringTok{\textquotesingle{}Alice\textquotesingle{}}\NormalTok{, }\StringTok{\textquotesingle{}Bob\textquotesingle{}}\NormalTok{, }\StringTok{\textquotesingle{}Charlie\textquotesingle{}}\NormalTok{]}
\NormalTok{\})}

\CommentTok{\# Multiple references: Orders from different months}
\NormalTok{orders\_jan }\OperatorTok{=}\NormalTok{ pd.DataFrame(\{}
    \StringTok{\textquotesingle{}customer\_id\textquotesingle{}}\NormalTok{: [}\DecValTok{1}\NormalTok{, }\DecValTok{1}\NormalTok{, }\DecValTok{2}\NormalTok{],}
    \StringTok{\textquotesingle{}amount\textquotesingle{}}\NormalTok{: [}\DecValTok{100}\NormalTok{, }\DecValTok{150}\NormalTok{, }\DecValTok{200}\NormalTok{]}
\NormalTok{\})}

\NormalTok{orders\_feb }\OperatorTok{=}\NormalTok{ pd.DataFrame(\{}
    \StringTok{\textquotesingle{}customer\_id\textquotesingle{}}\NormalTok{: [}\DecValTok{1}\NormalTok{, }\DecValTok{3}\NormalTok{, }\DecValTok{3}\NormalTok{],}
    \StringTok{\textquotesingle{}amount\textquotesingle{}}\NormalTok{: [}\DecValTok{175}\NormalTok{, }\DecValTok{125}\NormalTok{, }\DecValTok{225}\NormalTok{]}
\NormalTok{\})}

\CommentTok{\# One{-}to{-}many: single target, list of references}
\NormalTok{result }\OperatorTok{=}\NormalTok{ add.to(}
\NormalTok{    customers,}
\NormalTok{    bring\_from}\OperatorTok{=}\NormalTok{[orders\_jan, orders\_feb],    }\CommentTok{\# List of reference DataFrames}
\NormalTok{    bring}\OperatorTok{=}\StringTok{\textquotesingle{}amount\textquotesingle{}}\NormalTok{,}
\NormalTok{    against}\OperatorTok{=}\StringTok{\textquotesingle{}customer\_id\textquotesingle{}}\NormalTok{,}
\NormalTok{    strategy}\OperatorTok{=}\NormalTok{\{}\StringTok{\textquotesingle{}amount\textquotesingle{}}\NormalTok{: }\StringTok{\textquotesingle{}sum\textquotesingle{}}\NormalTok{\}              }\CommentTok{\# Aggregate multiple matches}
\NormalTok{)}

\CommentTok{\# Positional parameters (also works without naming certain parameters):}
\CommentTok{\# result = add.to(customers, [orders\_jan, orders\_feb], \textquotesingle{}amount\textquotesingle{}, \textquotesingle{}customer\_id\textquotesingle{}, strategy=\{\textquotesingle{}amount\textquotesingle{}: \textquotesingle{}sum\textquotesingle{}\})}

\BuiltInTok{print}\NormalTok{(result)}
\end{Highlighting}
\end{Shaded}

\textbf{Output:}

\begin{verbatim}
   customer_id     name  amount
0            1    Alice     425
1            2      Bob     200
2            3  Charlie     350
\end{verbatim}

\textbf{Explanation:} - Pass a list
\texttt{{[}orders\_jan,\ orders\_feb{]}} to \texttt{bring\_from} -
additory combines data from both DataFrames - Alice (ID 1): 100 + 150
(Jan) + 175 (Feb) = 425 - Bob (ID 2): 200 (Jan only) = 200 - Charlie (ID
3): 125 + 225 (Feb only) = 350 - The \texttt{strategy} parameter handles
multiple matches (we'll cover this in detail next page)

\textbf{Note:} This also works with polars DataFrames.

\begin{center}\rule{0.5\linewidth}{0.5pt}\end{center}

\section{Example 2: Many-to-One
Pattern}\label{example-2-many-to-one-pattern}

\textbf{Business Context:} You have customer lists split by type
(internal and external), and you need to add order amounts from a single
orders DataFrame to both lists.

\textbf{Code:}

\begin{Shaded}
\begin{Highlighting}[]
\ImportTok{import}\NormalTok{ additory }\ImportTok{as}\NormalTok{ add}
\ImportTok{import}\NormalTok{ pandas }\ImportTok{as}\NormalTok{ pd}

\CommentTok{\# Multiple targets: Different customer segments}
\NormalTok{customers\_internal }\OperatorTok{=}\NormalTok{ pd.DataFrame(\{}
    \StringTok{\textquotesingle{}customer\_id\textquotesingle{}}\NormalTok{: [}\DecValTok{1}\NormalTok{, }\DecValTok{2}\NormalTok{],}
    \StringTok{\textquotesingle{}name\textquotesingle{}}\NormalTok{: [}\StringTok{\textquotesingle{}Alice\textquotesingle{}}\NormalTok{, }\StringTok{\textquotesingle{}Bob\textquotesingle{}}\NormalTok{]}
\NormalTok{\})}

\NormalTok{customers\_external }\OperatorTok{=}\NormalTok{ pd.DataFrame(\{}
    \StringTok{\textquotesingle{}customer\_id\textquotesingle{}}\NormalTok{: [}\DecValTok{3}\NormalTok{, }\DecValTok{4}\NormalTok{],}
    \StringTok{\textquotesingle{}name\textquotesingle{}}\NormalTok{: [}\StringTok{\textquotesingle{}Charlie\textquotesingle{}}\NormalTok{, }\StringTok{\textquotesingle{}David\textquotesingle{}}\NormalTok{]}
\NormalTok{\})}

\CommentTok{\# Single reference: All orders}
\NormalTok{orders }\OperatorTok{=}\NormalTok{ pd.DataFrame(\{}
    \StringTok{\textquotesingle{}customer\_id\textquotesingle{}}\NormalTok{: [}\DecValTok{1}\NormalTok{, }\DecValTok{2}\NormalTok{, }\DecValTok{3}\NormalTok{, }\DecValTok{4}\NormalTok{],}
    \StringTok{\textquotesingle{}amount\textquotesingle{}}\NormalTok{: [}\DecValTok{100}\NormalTok{, }\DecValTok{200}\NormalTok{, }\DecValTok{150}\NormalTok{, }\DecValTok{175}\NormalTok{]}
\NormalTok{\})}

\CommentTok{\# Many{-}to{-}one: list of targets, single reference}
\NormalTok{results }\OperatorTok{=}\NormalTok{ add.to(}
\NormalTok{    [customers\_internal, customers\_external],    }\CommentTok{\# List of target DataFrames}
\NormalTok{    bring\_from}\OperatorTok{=}\NormalTok{orders,}
\NormalTok{    bring}\OperatorTok{=}\StringTok{\textquotesingle{}amount\textquotesingle{}}\NormalTok{,}
\NormalTok{    against}\OperatorTok{=}\StringTok{\textquotesingle{}customer\_id\textquotesingle{}}
\NormalTok{)}

\CommentTok{\# Positional parameters (also works without naming certain parameters):}
\CommentTok{\# results = add.to([customers\_internal, customers\_external], orders, \textquotesingle{}amount\textquotesingle{}, \textquotesingle{}customer\_id\textquotesingle{})}

\CommentTok{\# Results is a list of DataFrames}
\BuiltInTok{print}\NormalTok{(}\StringTok{"Internal customers:"}\NormalTok{)}
\BuiltInTok{print}\NormalTok{(results[}\DecValTok{0}\NormalTok{])}
\BuiltInTok{print}\NormalTok{(}\StringTok{"}\CharTok{\textbackslash{}n}\StringTok{External customers:"}\NormalTok{)}
\BuiltInTok{print}\NormalTok{(results[}\DecValTok{1}\NormalTok{])}
\end{Highlighting}
\end{Shaded}

\textbf{Output:}

\begin{verbatim}
Internal customers:
   customer_id   name  amount
0            1  Alice     100
1            2    Bob     200

External customers:
   customer_id     name  amount
0            3  Charlie     150
1            4    David     175
\end{verbatim}

\textbf{Explanation:} - Pass a list
\texttt{{[}customers\_internal,\ customers\_external{]}} as the first
argument - additory returns a list of DataFrames (same order as input) -
Each target DataFrame gets matched against the same reference -
\texttt{results{[}0{]}} contains internal customers with amounts -
\texttt{results{[}1{]}} contains external customers with amounts - This
is useful for processing segmented data consistently

\textbf{Note:} This also works with polars DataFrames.

\begin{center}\rule{0.5\linewidth}{0.5pt}\end{center}

\section{Example 3: Many-to-Many
Pattern}\label{example-3-many-to-many-pattern}

\textbf{Business Context:} You have customer lists split by quarter (Q1
and Q2), and order data also split by month (January and February). You
need to match each customer list with all relevant orders.

\textbf{Code:}

\begin{Shaded}
\begin{Highlighting}[]
\ImportTok{import}\NormalTok{ additory }\ImportTok{as}\NormalTok{ add}
\ImportTok{import}\NormalTok{ pandas }\ImportTok{as}\NormalTok{ pd}

\CommentTok{\# Multiple targets: Customers by quarter}
\NormalTok{customers\_q1 }\OperatorTok{=}\NormalTok{ pd.DataFrame(\{}
    \StringTok{\textquotesingle{}customer\_id\textquotesingle{}}\NormalTok{: [}\DecValTok{1}\NormalTok{, }\DecValTok{2}\NormalTok{],}
    \StringTok{\textquotesingle{}name\textquotesingle{}}\NormalTok{: [}\StringTok{\textquotesingle{}Alice\textquotesingle{}}\NormalTok{, }\StringTok{\textquotesingle{}Bob\textquotesingle{}}\NormalTok{]}
\NormalTok{\})}

\NormalTok{customers\_q2 }\OperatorTok{=}\NormalTok{ pd.DataFrame(\{}
    \StringTok{\textquotesingle{}customer\_id\textquotesingle{}}\NormalTok{: [}\DecValTok{3}\NormalTok{, }\DecValTok{4}\NormalTok{],}
    \StringTok{\textquotesingle{}name\textquotesingle{}}\NormalTok{: [}\StringTok{\textquotesingle{}Charlie\textquotesingle{}}\NormalTok{, }\StringTok{\textquotesingle{}David\textquotesingle{}}\NormalTok{]}
\NormalTok{\})}

\CommentTok{\# Multiple references: Orders by month}
\NormalTok{orders\_jan }\OperatorTok{=}\NormalTok{ pd.DataFrame(\{}
    \StringTok{\textquotesingle{}customer\_id\textquotesingle{}}\NormalTok{: [}\DecValTok{1}\NormalTok{, }\DecValTok{2}\NormalTok{],}
    \StringTok{\textquotesingle{}amount\textquotesingle{}}\NormalTok{: [}\DecValTok{100}\NormalTok{, }\DecValTok{200}\NormalTok{]}
\NormalTok{\})}

\NormalTok{orders\_feb }\OperatorTok{=}\NormalTok{ pd.DataFrame(\{}
    \StringTok{\textquotesingle{}customer\_id\textquotesingle{}}\NormalTok{: [}\DecValTok{3}\NormalTok{, }\DecValTok{4}\NormalTok{],}
    \StringTok{\textquotesingle{}amount\textquotesingle{}}\NormalTok{: [}\DecValTok{150}\NormalTok{, }\DecValTok{175}\NormalTok{]}
\NormalTok{\})}

\CommentTok{\# Many{-}to{-}many: list of targets, list of references}
\NormalTok{results }\OperatorTok{=}\NormalTok{ add.to(}
\NormalTok{    [customers\_q1, customers\_q2],           }\CommentTok{\# List of targets}
\NormalTok{    bring\_from}\OperatorTok{=}\NormalTok{[orders\_jan, orders\_feb],    }\CommentTok{\# List of references}
\NormalTok{    bring}\OperatorTok{=}\StringTok{\textquotesingle{}amount\textquotesingle{}}\NormalTok{,}
\NormalTok{    against}\OperatorTok{=}\StringTok{\textquotesingle{}customer\_id\textquotesingle{}}\NormalTok{,}
\NormalTok{    strategy}\OperatorTok{=}\NormalTok{\{}\StringTok{\textquotesingle{}amount\textquotesingle{}}\NormalTok{: }\StringTok{\textquotesingle{}sum\textquotesingle{}}\NormalTok{\}}
\NormalTok{)}

\CommentTok{\# Positional parameters (also works without naming certain parameters):}
\CommentTok{\# results = add.to([customers\_q1, customers\_q2], [orders\_jan, orders\_feb], \textquotesingle{}amount\textquotesingle{}, \textquotesingle{}customer\_id\textquotesingle{}, strategy=\{\textquotesingle{}amount\textquotesingle{}: \textquotesingle{}sum\textquotesingle{}\})}

\BuiltInTok{print}\NormalTok{(}\StringTok{"Q1 customers:"}\NormalTok{)}
\BuiltInTok{print}\NormalTok{(results[}\DecValTok{0}\NormalTok{])}
\BuiltInTok{print}\NormalTok{(}\StringTok{"}\CharTok{\textbackslash{}n}\StringTok{Q2 customers:"}\NormalTok{)}
\BuiltInTok{print}\NormalTok{(results[}\DecValTok{1}\NormalTok{])}
\end{Highlighting}
\end{Shaded}

\textbf{Output:}

\begin{verbatim}
Q1 customers:
   customer_id  name  amount
0            1 Alice     100
1            2   Bob     200

Q2 customers:
   customer_id     name  amount
0            3  Charlie     150
1            4    David     175
\end{verbatim}

\textbf{Explanation:} - Pass lists for both targets and references -
Each target DataFrame is matched against ALL reference DataFrames -
additory returns a list of DataFrames (one per target) - Q1 customers
get amounts from both Jan and Feb orders (where they exist) - Q2
customers get amounts from both Jan and Feb orders (where they exist) -
This pattern is powerful for time-series or segmented data analysis

\textbf{Note:} This also works with polars DataFrames.

\begin{center}\rule{0.5\linewidth}{0.5pt}\end{center}

\section{Pattern Comparison}\label{pattern-comparison}

\subsection{When to Use Each Pattern}\label{when-to-use-each-pattern}

\begin{longtable}[]{@{}
  >{\raggedright\arraybackslash}p{(\linewidth - 8\tabcolsep) * \real{0.1915}}
  >{\raggedright\arraybackslash}p{(\linewidth - 8\tabcolsep) * \real{0.1702}}
  >{\raggedright\arraybackslash}p{(\linewidth - 8\tabcolsep) * \real{0.2340}}
  >{\raggedright\arraybackslash}p{(\linewidth - 8\tabcolsep) * \real{0.1915}}
  >{\raggedright\arraybackslash}p{(\linewidth - 8\tabcolsep) * \real{0.2128}}@{}}
\toprule\noalign{}
\begin{minipage}[b]{\linewidth}\raggedright
Pattern
\end{minipage} & \begin{minipage}[b]{\linewidth}\raggedright
Target
\end{minipage} & \begin{minipage}[b]{\linewidth}\raggedright
Reference
\end{minipage} & \begin{minipage}[b]{\linewidth}\raggedright
Returns
\end{minipage} & \begin{minipage}[b]{\linewidth}\raggedright
Use Case
\end{minipage} \\
\midrule\noalign{}
\endhead
\bottomrule\noalign{}
\endlastfoot
\textbf{Basic} & Single DF & Single DF & Single DF & Simple lookup \\
\textbf{One-to-Many} & Single DF & List of DFs & Single DF & Aggregate
from multiple sources \\
\textbf{Many-to-One} & List of DFs & Single DF & List of DFs & Apply
same reference to multiple targets \\
\textbf{Many-to-Many} & List of DFs & List of DFs & List of DFs &
Complex multi-source matching \\
\end{longtable}

\subsection{Visual Guide}\label{visual-guide}

\begin{verbatim}
Basic (1:1)
Target ──> Reference
  ↓
Result

One-to-Many (1:N)
Target ──> [Ref1, Ref2, Ref3]
  ↓
Result

Many-to-One (N:1)
[Target1, Target2] ──> Reference
  ↓
[Result1, Result2]

Many-to-Many (N:N)
[Target1, Target2] ──> [Ref1, Ref2]
  ↓
[Result1, Result2]
\end{verbatim}

\begin{center}\rule{0.5\linewidth}{0.5pt}\end{center}

\section{Working with Result Lists}\label{working-with-result-lists}

\subsection{Accessing Individual
Results}\label{accessing-individual-results}

\begin{Shaded}
\begin{Highlighting}[]
\CommentTok{\# Many{-}to{-}one or many{-}to{-}many returns a list}
\NormalTok{results }\OperatorTok{=}\NormalTok{ add.to(}
\NormalTok{    [df1, df2, df3],}
\NormalTok{    bring\_from}\OperatorTok{=}\NormalTok{ref\_df,}
\NormalTok{    bring}\OperatorTok{=}\StringTok{\textquotesingle{}amount\textquotesingle{}}\NormalTok{,}
\NormalTok{    against}\OperatorTok{=}\StringTok{\textquotesingle{}id\textquotesingle{}}
\NormalTok{)}

\CommentTok{\# Access by index}
\NormalTok{first\_result }\OperatorTok{=}\NormalTok{ results[}\DecValTok{0}\NormalTok{]}
\NormalTok{second\_result }\OperatorTok{=}\NormalTok{ results[}\DecValTok{1}\NormalTok{]}

\CommentTok{\# Iterate over results}
\ControlFlowTok{for}\NormalTok{ i, result }\KeywordTok{in} \BuiltInTok{enumerate}\NormalTok{(results):}
    \BuiltInTok{print}\NormalTok{(}\SpecialStringTok{f"Result }\SpecialCharTok{\{}\NormalTok{i}\SpecialCharTok{\}}\SpecialStringTok{: }\SpecialCharTok{\{}\BuiltInTok{len}\NormalTok{(result)}\SpecialCharTok{\}}\SpecialStringTok{ rows"}\NormalTok{)}
\end{Highlighting}
\end{Shaded}

\subsection{Combining Results}\label{combining-results}

\begin{Shaded}
\begin{Highlighting}[]
\CommentTok{\# If you want to combine all results into one DataFrame}
\ImportTok{import}\NormalTok{ pandas }\ImportTok{as}\NormalTok{ pd}

\NormalTok{combined }\OperatorTok{=}\NormalTok{ pd.concat(results, ignore\_index}\OperatorTok{=}\VariableTok{True}\NormalTok{)}
\end{Highlighting}
\end{Shaded}

\begin{center}\rule{0.5\linewidth}{0.5pt}\end{center}

\section{Key Takeaways}\label{key-takeaways-2}

\begin{itemize}
\tightlist
\item
  Use lists \texttt{{[}df1,\ df2{]}} for targets or references to enable
  advanced patterns
\item
  One-to-Many: Single target + multiple references → Single result
\item
  Many-to-One: Multiple targets + single reference → List of results
\item
  Many-to-Many: Multiple targets + multiple references → List of results
\item
  When you pass a list of targets, you always get a list of results back
\item
  These patterns are essential for working with segmented or time-series
  data
\end{itemize}

\begin{center}\rule{0.5\linewidth}{0.5pt}\end{center}

\section{Common Questions}\label{common-questions-2}

\textbf{Q: Can I mix pandas and polars DataFrames in the lists?}\\
A: No, all DataFrames in a single operation must be the same type (all
pandas or all polars).

\textbf{Q: What if the reference DataFrames have overlapping data?}\\
A: Use the \texttt{strategy} parameter to control how duplicates are
handled (sum, mean, first, etc.). We cover this in detail on the next
page.

\textbf{Q: Is there a limit to how many DataFrames I can include in a
list?}\\
A: No hard limit, but performance may degrade with very large lists.
Consider combining DataFrames first if you have dozens of them.

\textbf{Q: Do the DataFrames in a list need to have the same columns?}\\
A: They need to have the key columns (\texttt{against}) and the columns
you're bringing (\texttt{bring}). Other columns can differ.

\begin{center}\rule{0.5\linewidth}{0.5pt}\end{center}

\section{Next Steps}\label{next-steps-4}

\begin{itemize}
\tightlist
\item
  \textbf{\href{examples-to-04-strategy.html}{Aggregation Strategies}} -
  Learn how to handle multiple matches with strategy parameter
\item
  \textbf{\href{examples-to-02-multiple.html}{Multiple Columns \& Keys}}
  - Review multiple column operations
\item
  \textbf{\href{reference-to-basic.html}{API Reference}} - Complete
  \texttt{add.to()} documentation
\end{itemize}

\begin{center}\rule{0.5\linewidth}{0.5pt}\end{center}

\textbf{Version:} 0.1.3\\
\textbf{Last Updated:} March 9, 2026

\chapter{add.to() - Aggregation
Strategies}\label{add.to---aggregation-strategies}

\section{Overview}\label{overview-3}

Learn how to control what happens when multiple rows match during a
lookup. The \texttt{strategy} parameter lets you aggregate, select, or
combine matching values in powerful ways.

\textbf{What you'll learn:} - How to use aggregation strategies (sum,
mean, count, etc.) - How to select specific matches (first, last) - How
to control column positioning with the \texttt{position} parameter - How
to use different join types with \texttt{join\_type}

\textbf{Prerequisites:} -
\href{examples-to-03-patterns.html}{One-to-Many \& Many-to-One Patterns}
- Understanding of multiple matches -
\href{examples-to-02-multiple.html}{Multiple Columns \& Keys} - Working
with multiple columns

\begin{center}\rule{0.5\linewidth}{0.5pt}\end{center}

\section{Example 1: Aggregation Strategies (Sum, Mean,
Count)}\label{example-1-aggregation-strategies-sum-mean-count}

\textbf{Business Context:} You have a customer list and multiple orders
per customer. You need to calculate total revenue (sum), average order
value (mean), and number of orders (count) for each customer.

\textbf{Code:}

\begin{Shaded}
\begin{Highlighting}[]
\ImportTok{import}\NormalTok{ additory }\ImportTok{as}\NormalTok{ add}
\ImportTok{import}\NormalTok{ pandas }\ImportTok{as}\NormalTok{ pd}

\CommentTok{\# Customers}
\NormalTok{customers }\OperatorTok{=}\NormalTok{ pd.DataFrame(\{}
    \StringTok{\textquotesingle{}customer\_id\textquotesingle{}}\NormalTok{: [}\DecValTok{1}\NormalTok{, }\DecValTok{2}\NormalTok{, }\DecValTok{3}\NormalTok{],}
    \StringTok{\textquotesingle{}name\textquotesingle{}}\NormalTok{: [}\StringTok{\textquotesingle{}Alice\textquotesingle{}}\NormalTok{, }\StringTok{\textquotesingle{}Bob\textquotesingle{}}\NormalTok{, }\StringTok{\textquotesingle{}Charlie\textquotesingle{}}\NormalTok{]}
\NormalTok{\})}

\CommentTok{\# Orders (multiple per customer)}
\NormalTok{orders }\OperatorTok{=}\NormalTok{ pd.DataFrame(\{}
    \StringTok{\textquotesingle{}customer\_id\textquotesingle{}}\NormalTok{: [}\DecValTok{1}\NormalTok{, }\DecValTok{1}\NormalTok{, }\DecValTok{1}\NormalTok{, }\DecValTok{2}\NormalTok{, }\DecValTok{2}\NormalTok{, }\DecValTok{3}\NormalTok{],}
    \StringTok{\textquotesingle{}amount\textquotesingle{}}\NormalTok{: [}\DecValTok{100}\NormalTok{, }\DecValTok{150}\NormalTok{, }\DecValTok{200}\NormalTok{, }\DecValTok{300}\NormalTok{, }\DecValTok{250}\NormalTok{, }\DecValTok{175}\NormalTok{]}
\NormalTok{\})}

\CommentTok{\# Sum: Total revenue per customer}
\NormalTok{result\_sum }\OperatorTok{=}\NormalTok{ add.to(}
\NormalTok{    customers,}
\NormalTok{    bring\_from}\OperatorTok{=}\NormalTok{orders,}
\NormalTok{    bring}\OperatorTok{=}\StringTok{\textquotesingle{}amount\textquotesingle{}}\NormalTok{,}
\NormalTok{    against}\OperatorTok{=}\StringTok{\textquotesingle{}customer\_id\textquotesingle{}}\NormalTok{,}
\NormalTok{    strategy}\OperatorTok{=}\NormalTok{\{}\StringTok{\textquotesingle{}amount\textquotesingle{}}\NormalTok{: }\StringTok{\textquotesingle{}sum\textquotesingle{}}\NormalTok{\}}
\NormalTok{)}

\CommentTok{\# Positional parameters (also works without naming certain parameters):}
\CommentTok{\# result\_sum = add.to(customers, orders, \textquotesingle{}amount\textquotesingle{}, \textquotesingle{}customer\_id\textquotesingle{}, strategy=\{\textquotesingle{}amount\textquotesingle{}: \textquotesingle{}sum\textquotesingle{}\})}

\BuiltInTok{print}\NormalTok{(}\StringTok{"Total Revenue:"}\NormalTok{)}
\BuiltInTok{print}\NormalTok{(result\_sum)}

\CommentTok{\# Mean: Average order value per customer}
\NormalTok{result\_mean }\OperatorTok{=}\NormalTok{ add.to(}
\NormalTok{    customers,}
\NormalTok{    bring\_from}\OperatorTok{=}\NormalTok{orders,}
\NormalTok{    bring}\OperatorTok{=}\StringTok{\textquotesingle{}amount\textquotesingle{}}\NormalTok{,}
\NormalTok{    against}\OperatorTok{=}\StringTok{\textquotesingle{}customer\_id\textquotesingle{}}\NormalTok{,}
\NormalTok{    strategy}\OperatorTok{=}\NormalTok{\{}\StringTok{\textquotesingle{}amount\textquotesingle{}}\NormalTok{: }\StringTok{\textquotesingle{}mean\textquotesingle{}}\NormalTok{\}}
\NormalTok{)}

\BuiltInTok{print}\NormalTok{(}\StringTok{"}\CharTok{\textbackslash{}n}\StringTok{Average Order Value:"}\NormalTok{)}
\BuiltInTok{print}\NormalTok{(result\_mean)}

\CommentTok{\# Count: Number of orders per customer}
\NormalTok{result\_count }\OperatorTok{=}\NormalTok{ add.to(}
\NormalTok{    customers,}
\NormalTok{    bring\_from}\OperatorTok{=}\NormalTok{orders,}
\NormalTok{    bring}\OperatorTok{=}\StringTok{\textquotesingle{}amount\textquotesingle{}}\NormalTok{,}
\NormalTok{    against}\OperatorTok{=}\StringTok{\textquotesingle{}customer\_id\textquotesingle{}}\NormalTok{,}
\NormalTok{    strategy}\OperatorTok{=}\NormalTok{\{}\StringTok{\textquotesingle{}amount\textquotesingle{}}\NormalTok{: }\StringTok{\textquotesingle{}count\textquotesingle{}}\NormalTok{\}}
\NormalTok{)}

\BuiltInTok{print}\NormalTok{(}\StringTok{"}\CharTok{\textbackslash{}n}\StringTok{Order Count:"}\NormalTok{)}
\BuiltInTok{print}\NormalTok{(result\_count)}
\end{Highlighting}
\end{Shaded}

\textbf{Output:}

\begin{verbatim}
Total Revenue:
   customer_id     name  amount
0            1    Alice     450
1            2      Bob     550
2            3  Charlie     175

Average Order Value:
   customer_id     name  amount
0            1    Alice     150.0
1            2      Bob     275.0
2            3  Charlie     175.0

Order Count:
   customer_id     name  amount
0            1    Alice       3
1            2      Bob       2
2            3  Charlie       1
\end{verbatim}

\textbf{Explanation:} -
\texttt{strategy=\{\textquotesingle{}amount\textquotesingle{}:\ \textquotesingle{}sum\textquotesingle{}\}}
adds up all matching values - Alice has 3 orders: 100 + 150 + 200 = 450
-
\texttt{strategy=\{\textquotesingle{}amount\textquotesingle{}:\ \textquotesingle{}mean\textquotesingle{}\}}
calculates the average - Alice's average: (100 + 150 + 200) / 3 = 150 -
\texttt{strategy=\{\textquotesingle{}amount\textquotesingle{}:\ \textquotesingle{}count\textquotesingle{}\}}
counts the number of matches - Alice has 3 matching orders - The
strategy is specified as a dictionary:
\texttt{\{column\_name:\ strategy\_type\}}

\textbf{Available Strategies:} -
\texttt{\textquotesingle{}sum\textquotesingle{}} - Add all values -
\texttt{\textquotesingle{}mean\textquotesingle{}} - Calculate average -
\texttt{\textquotesingle{}count\textquotesingle{}} - Count number of
matches - \texttt{\textquotesingle{}min\textquotesingle{}} - Take
minimum value - \texttt{\textquotesingle{}max\textquotesingle{}} - Take
maximum value - \texttt{\textquotesingle{}first\textquotesingle{}} -
Take first match (see Example 2) -
\texttt{\textquotesingle{}last\textquotesingle{}} - Take last match (see
Example 2)

\textbf{Note:} This also works with polars DataFrames.

\begin{center}\rule{0.5\linewidth}{0.5pt}\end{center}

\section{Example 2: First and Last
Strategies}\label{example-2-first-and-last-strategies}

\textbf{Business Context:} You have a product catalog and a price
history table. You need to show both the original price (first) and the
current price (last) for each product.

\textbf{Code:}

\begin{Shaded}
\begin{Highlighting}[]
\ImportTok{import}\NormalTok{ additory }\ImportTok{as}\NormalTok{ add}
\ImportTok{import}\NormalTok{ pandas }\ImportTok{as}\NormalTok{ pd}

\CommentTok{\# Products}
\NormalTok{products }\OperatorTok{=}\NormalTok{ pd.DataFrame(\{}
    \StringTok{\textquotesingle{}product\_id\textquotesingle{}}\NormalTok{: [}\DecValTok{101}\NormalTok{, }\DecValTok{102}\NormalTok{, }\DecValTok{103}\NormalTok{],}
    \StringTok{\textquotesingle{}name\textquotesingle{}}\NormalTok{: [}\StringTok{\textquotesingle{}Widget\textquotesingle{}}\NormalTok{, }\StringTok{\textquotesingle{}Gadget\textquotesingle{}}\NormalTok{, }\StringTok{\textquotesingle{}Doohickey\textquotesingle{}}\NormalTok{]}
\NormalTok{\})}

\CommentTok{\# Price history (multiple prices per product over time)}
\NormalTok{price\_history }\OperatorTok{=}\NormalTok{ pd.DataFrame(\{}
    \StringTok{\textquotesingle{}product\_id\textquotesingle{}}\NormalTok{: [}\DecValTok{101}\NormalTok{, }\DecValTok{101}\NormalTok{, }\DecValTok{101}\NormalTok{, }\DecValTok{102}\NormalTok{, }\DecValTok{102}\NormalTok{, }\DecValTok{103}\NormalTok{],}
    \StringTok{\textquotesingle{}date\textquotesingle{}}\NormalTok{: [}\StringTok{\textquotesingle{}2024{-}01{-}01\textquotesingle{}}\NormalTok{, }\StringTok{\textquotesingle{}2024{-}02{-}01\textquotesingle{}}\NormalTok{, }\StringTok{\textquotesingle{}2024{-}03{-}01\textquotesingle{}}\NormalTok{, }
             \StringTok{\textquotesingle{}2024{-}01{-}01\textquotesingle{}}\NormalTok{, }\StringTok{\textquotesingle{}2024{-}02{-}01\textquotesingle{}}\NormalTok{, }\StringTok{\textquotesingle{}2024{-}01{-}01\textquotesingle{}}\NormalTok{],}
    \StringTok{\textquotesingle{}price\textquotesingle{}}\NormalTok{: [}\FloatTok{29.99}\NormalTok{, }\FloatTok{34.99}\NormalTok{, }\FloatTok{39.99}\NormalTok{, }\FloatTok{49.99}\NormalTok{, }\FloatTok{54.99}\NormalTok{, }\FloatTok{19.99}\NormalTok{]}
\NormalTok{\})}

\CommentTok{\# Get first price (original/oldest)}
\NormalTok{result\_first }\OperatorTok{=}\NormalTok{ add.to(}
\NormalTok{    products,}
\NormalTok{    bring\_from}\OperatorTok{=}\NormalTok{price\_history,}
\NormalTok{    bring}\OperatorTok{=}\StringTok{\textquotesingle{}price\textquotesingle{}}\NormalTok{,}
\NormalTok{    against}\OperatorTok{=}\StringTok{\textquotesingle{}product\_id\textquotesingle{}}\NormalTok{,}
\NormalTok{    strategy}\OperatorTok{=}\NormalTok{\{}\StringTok{\textquotesingle{}price\textquotesingle{}}\NormalTok{: }\StringTok{\textquotesingle{}first\textquotesingle{}}\NormalTok{\}}
\NormalTok{)}

\CommentTok{\# Positional parameters (also works without naming certain parameters):}
\CommentTok{\# result\_first = add.to(products, price\_history, \textquotesingle{}price\textquotesingle{}, \textquotesingle{}product\_id\textquotesingle{}, strategy=\{\textquotesingle{}price\textquotesingle{}: \textquotesingle{}first\textquotesingle{}\})}

\BuiltInTok{print}\NormalTok{(}\StringTok{"Original Prices:"}\NormalTok{)}
\BuiltInTok{print}\NormalTok{(result\_first)}

\CommentTok{\# Get last price (current/newest)}
\NormalTok{result\_last }\OperatorTok{=}\NormalTok{ add.to(}
\NormalTok{    products,}
\NormalTok{    bring\_from}\OperatorTok{=}\NormalTok{price\_history,}
\NormalTok{    bring}\OperatorTok{=}\StringTok{\textquotesingle{}price\textquotesingle{}}\NormalTok{,}
\NormalTok{    against}\OperatorTok{=}\StringTok{\textquotesingle{}product\_id\textquotesingle{}}\NormalTok{,}
\NormalTok{    strategy}\OperatorTok{=}\NormalTok{\{}\StringTok{\textquotesingle{}price\textquotesingle{}}\NormalTok{: }\StringTok{\textquotesingle{}last\textquotesingle{}}\NormalTok{\}}
\NormalTok{)}

\BuiltInTok{print}\NormalTok{(}\StringTok{"}\CharTok{\textbackslash{}n}\StringTok{Current Prices:"}\NormalTok{)}
\BuiltInTok{print}\NormalTok{(result\_last)}
\end{Highlighting}
\end{Shaded}

\textbf{Output:}

\begin{verbatim}
Original Prices:
   product_id       name  price
0         101     Widget  29.99
1         102     Gadget  49.99
2         103  Doohickey  19.99

Current Prices:
   product_id       name  price
0         101     Widget  39.99
1         102     Gadget  54.99
2         103  Doohickey  19.99
\end{verbatim}

\textbf{Explanation:} -
\texttt{strategy=\{\textquotesingle{}price\textquotesingle{}:\ \textquotesingle{}first\textquotesingle{}\}}
takes the first matching row - Widget (101) has 3 prices: 29.99 (first),
34.99, 39.99 - First strategy returns 29.99 -
\texttt{strategy=\{\textquotesingle{}price\textquotesingle{}:\ \textquotesingle{}last\textquotesingle{}\}}
takes the last matching row - Last strategy returns 39.99 - Order
matters! The order in the reference DataFrame determines first/last -
Doohickey (103) has only one price, so first and last are the same

\textbf{Note:} This also works with polars DataFrames.

\begin{center}\rule{0.5\linewidth}{0.5pt}\end{center}

\section{Example 3: Position and Join Type
Parameters}\label{example-3-position-and-join-type-parameters}

\textbf{Business Context:} You have an employee list and a skills table.
You want to control where the skill column appears and handle employees
without skills differently.

\textbf{Code:}

\begin{Shaded}
\begin{Highlighting}[]
\ImportTok{import}\NormalTok{ additory }\ImportTok{as}\NormalTok{ add}
\ImportTok{import}\NormalTok{ pandas }\ImportTok{as}\NormalTok{ pd}

\CommentTok{\# Employees}
\NormalTok{employees }\OperatorTok{=}\NormalTok{ pd.DataFrame(\{}
    \StringTok{\textquotesingle{}emp\_id\textquotesingle{}}\NormalTok{: [}\DecValTok{1}\NormalTok{, }\DecValTok{2}\NormalTok{, }\DecValTok{3}\NormalTok{],}
    \StringTok{\textquotesingle{}name\textquotesingle{}}\NormalTok{: [}\StringTok{\textquotesingle{}Alice\textquotesingle{}}\NormalTok{, }\StringTok{\textquotesingle{}Bob\textquotesingle{}}\NormalTok{, }\StringTok{\textquotesingle{}Charlie\textquotesingle{}}\NormalTok{]}
\NormalTok{\})}

\CommentTok{\# Skills (not all employees have skills listed)}
\NormalTok{skills }\OperatorTok{=}\NormalTok{ pd.DataFrame(\{}
    \StringTok{\textquotesingle{}emp\_id\textquotesingle{}}\NormalTok{: [}\DecValTok{1}\NormalTok{, }\DecValTok{2}\NormalTok{, }\DecValTok{4}\NormalTok{],  }\CommentTok{\# Note: emp\_id 4 doesn\textquotesingle{}t exist in employees}
    \StringTok{\textquotesingle{}skill\textquotesingle{}}\NormalTok{: [}\StringTok{\textquotesingle{}Python\textquotesingle{}}\NormalTok{, }\StringTok{\textquotesingle{}Java\textquotesingle{}}\NormalTok{, }\StringTok{\textquotesingle{}JavaScript\textquotesingle{}}\NormalTok{]}
\NormalTok{\})}

\CommentTok{\# Left join (default) {-} keeps all employees}
\NormalTok{result\_left }\OperatorTok{=}\NormalTok{ add.to(}
\NormalTok{    employees,}
\NormalTok{    bring\_from}\OperatorTok{=}\NormalTok{skills,}
\NormalTok{    bring}\OperatorTok{=}\StringTok{\textquotesingle{}skill\textquotesingle{}}\NormalTok{,}
\NormalTok{    against}\OperatorTok{=}\StringTok{\textquotesingle{}emp\_id\textquotesingle{}}\NormalTok{,}
\NormalTok{    join\_type}\OperatorTok{=}\StringTok{\textquotesingle{}left\textquotesingle{}}
\NormalTok{)}

\CommentTok{\# Positional parameters (also works without naming certain parameters):}
\CommentTok{\# result\_left = add.to(employees, skills, \textquotesingle{}skill\textquotesingle{}, \textquotesingle{}emp\_id\textquotesingle{}, join\_type=\textquotesingle{}left\textquotesingle{})}

\BuiltInTok{print}\NormalTok{(}\StringTok{"Left Join (all employees):"}\NormalTok{)}
\BuiltInTok{print}\NormalTok{(result\_left)}

\CommentTok{\# Inner join {-} only matching employees}
\NormalTok{result\_inner }\OperatorTok{=}\NormalTok{ add.to(}
\NormalTok{    employees,}
\NormalTok{    bring\_from}\OperatorTok{=}\NormalTok{skills,}
\NormalTok{    bring}\OperatorTok{=}\StringTok{\textquotesingle{}skill\textquotesingle{}}\NormalTok{,}
\NormalTok{    against}\OperatorTok{=}\StringTok{\textquotesingle{}emp\_id\textquotesingle{}}\NormalTok{,}
\NormalTok{    join\_type}\OperatorTok{=}\StringTok{\textquotesingle{}inner\textquotesingle{}}
\NormalTok{)}

\BuiltInTok{print}\NormalTok{(}\StringTok{"}\CharTok{\textbackslash{}n}\StringTok{Inner Join (only employees with skills):"}\NormalTok{)}
\BuiltInTok{print}\NormalTok{(result\_inner)}

\CommentTok{\# Position parameter {-} control column placement}
\NormalTok{result\_position }\OperatorTok{=}\NormalTok{ add.to(}
\NormalTok{    employees,}
\NormalTok{    bring\_from}\OperatorTok{=}\NormalTok{skills,}
\NormalTok{    bring}\OperatorTok{=}\StringTok{\textquotesingle{}skill\textquotesingle{}}\NormalTok{,}
\NormalTok{    against}\OperatorTok{=}\StringTok{\textquotesingle{}emp\_id\textquotesingle{}}\NormalTok{,}
\NormalTok{    position}\OperatorTok{=}\DecValTok{1}  \CommentTok{\# Insert after first column (0{-}indexed)}
\NormalTok{)}

\BuiltInTok{print}\NormalTok{(}\StringTok{"}\CharTok{\textbackslash{}n}\StringTok{With Position Control:"}\NormalTok{)}
\BuiltInTok{print}\NormalTok{(result\_position)}
\end{Highlighting}
\end{Shaded}

\textbf{Output:}

\begin{verbatim}
Left Join (all employees):
   emp_id     name    skill
0       1    Alice   Python
1       2      Bob     Java
2       3  Charlie      NaN

Inner Join (only employees with skills):
   emp_id   name   skill
0       1  Alice  Python
1       2    Bob    Java

With Position Control:
   emp_id    skill     name
0       1   Python    Alice
1       2     Java      Bob
2       3      NaN  Charlie
\end{verbatim}

\textbf{Explanation:} -
\texttt{join\_type=\textquotesingle{}left\textquotesingle{}} (default)
keeps all rows from the target DataFrame - Charlie (emp\_id 3) has no
skill, so gets NaN -
\texttt{join\_type=\textquotesingle{}inner\textquotesingle{}} only keeps
rows that match in both DataFrames - Charlie is excluded because he has
no skill - \texttt{position=1} inserts the new column at index 1 (after
emp\_id) - Without position, new columns are added at the end - Position
is 0-indexed: 0 = first column, 1 = second column, etc.

\textbf{Available Join Types:} -
\texttt{\textquotesingle{}left\textquotesingle{}} (default) - Keep all
target rows, add NaN for non-matches -
\texttt{\textquotesingle{}inner\textquotesingle{}} - Only keep rows that
match in both DataFrames -
\texttt{\textquotesingle{}right\textquotesingle{}} - Keep all reference
rows (rarely used) - \texttt{\textquotesingle{}outer\textquotesingle{}}
- Keep all rows from both DataFrames (rarely used)

\textbf{Note:} This also works with polars DataFrames.

\begin{center}\rule{0.5\linewidth}{0.5pt}\end{center}

\section{Strategy Dictionary Format}\label{strategy-dictionary-format}

\subsection{Single Column Strategy}\label{single-column-strategy}

\begin{Shaded}
\begin{Highlighting}[]
\NormalTok{strategy}\OperatorTok{=}\NormalTok{\{}\StringTok{\textquotesingle{}amount\textquotesingle{}}\NormalTok{: }\StringTok{\textquotesingle{}sum\textquotesingle{}}\NormalTok{\}}
\end{Highlighting}
\end{Shaded}

\subsection{Multiple Column
Strategies}\label{multiple-column-strategies}

\begin{Shaded}
\begin{Highlighting}[]
\NormalTok{strategy}\OperatorTok{=}\NormalTok{\{}
    \StringTok{\textquotesingle{}amount\textquotesingle{}}\NormalTok{: }\StringTok{\textquotesingle{}sum\textquotesingle{}}\NormalTok{,}
    \StringTok{\textquotesingle{}quantity\textquotesingle{}}\NormalTok{: }\StringTok{\textquotesingle{}count\textquotesingle{}}\NormalTok{,}
    \StringTok{\textquotesingle{}price\textquotesingle{}}\NormalTok{: }\StringTok{\textquotesingle{}mean\textquotesingle{}}
\NormalTok{\}}
\end{Highlighting}
\end{Shaded}

\subsection{When Strategy is Required}\label{when-strategy-is-required}

You MUST provide a strategy when: - Using one-to-many pattern (single
target, multiple references) - Using many-to-many pattern (multiple
targets, multiple references) - Any situation where multiple rows might
match

You DON'T need a strategy when: - Each key has exactly one match - Using
basic one-to-one lookups

\begin{center}\rule{0.5\linewidth}{0.5pt}\end{center}

\section{Combining Parameters}\label{combining-parameters}

You can combine multiple parameters for powerful control:

\begin{Shaded}
\begin{Highlighting}[]
\NormalTok{result }\OperatorTok{=}\NormalTok{ add.to(}
\NormalTok{    employees,}
\NormalTok{    bring\_from}\OperatorTok{=}\NormalTok{orders,}
\NormalTok{    bring}\OperatorTok{=}\NormalTok{[}\StringTok{\textquotesingle{}amount\textquotesingle{}}\NormalTok{, }\StringTok{\textquotesingle{}quantity\textquotesingle{}}\NormalTok{],}
\NormalTok{    against}\OperatorTok{=}\StringTok{\textquotesingle{}emp\_id\textquotesingle{}}\NormalTok{,}
\NormalTok{    strategy}\OperatorTok{=}\NormalTok{\{}
        \StringTok{\textquotesingle{}amount\textquotesingle{}}\NormalTok{: }\StringTok{\textquotesingle{}sum\textquotesingle{}}\NormalTok{,}
        \StringTok{\textquotesingle{}quantity\textquotesingle{}}\NormalTok{: }\StringTok{\textquotesingle{}count\textquotesingle{}}
\NormalTok{    \},}
\NormalTok{    join\_type}\OperatorTok{=}\StringTok{\textquotesingle{}left\textquotesingle{}}\NormalTok{,}
\NormalTok{    position}\OperatorTok{=}\DecValTok{2}
\NormalTok{)}
\end{Highlighting}
\end{Shaded}

This example: - Brings two columns - Sums the amount column - Counts the
quantity column - Keeps all employees (left join) - Inserts new columns
at position 2

\begin{center}\rule{0.5\linewidth}{0.5pt}\end{center}

\section{Key Takeaways}\label{key-takeaways-3}

\begin{itemize}
\tightlist
\item
  Use \texttt{strategy} parameter to control how multiple matches are
  handled
\item
  Common strategies: sum, mean, count, first, last, min, max
\item
  Strategy is a dictionary: \texttt{\{column\_name:\ strategy\_type\}}
\item
  Use \texttt{join\_type} to control which rows are kept
\item
  Use \texttt{position} to control where new columns are inserted
\item
  Combine parameters for precise control over the operation
\end{itemize}

\begin{center}\rule{0.5\linewidth}{0.5pt}\end{center}

\section{Common Questions}\label{common-questions-3}

\textbf{Q: What happens if I don't specify a strategy when there are
multiple matches?}\\
A: additory will raise an error asking you to specify a strategy. This
prevents unexpected behavior.

\textbf{Q: Can I use different strategies for different columns?}\\
A: Yes! Use a dictionary with multiple entries:
\texttt{strategy=\{\textquotesingle{}amount\textquotesingle{}:\ \textquotesingle{}sum\textquotesingle{},\ \textquotesingle{}count\textquotesingle{}:\ \textquotesingle{}count\textquotesingle{}\}}

\textbf{Q: What's the difference between `first' and `min'?}\\
A: `first' takes the first row in order, `min' takes the row with the
minimum value. They can give different results.

\textbf{Q: Can I create custom aggregation strategies?}\\
A: Not directly, but you can use \texttt{add.transform()} after
\texttt{add.to()} to apply custom calculations.

\begin{center}\rule{0.5\linewidth}{0.5pt}\end{center}

\section{Next Steps}\label{next-steps-5}

\begin{itemize}
\tightlist
\item
  \textbf{\href{examples-to-05-real-world.html}{Real-World Scenarios}} -
  See complete business workflows using these techniques
\item
  \textbf{\href{examples-to-03-patterns.html}{One-to-Many \& Many-to-One
  Patterns}} - Review when strategies are needed
\item
  \textbf{\href{reference-to-basic.html}{API Reference}} - Complete
  \texttt{add.to()} documentation
\end{itemize}

\begin{center}\rule{0.5\linewidth}{0.5pt}\end{center}

\textbf{Version:} 0.1.3\\
\textbf{Last Updated:} March 9, 2026

\chapter{add.to() - Real-World
Scenarios}\label{add.to---real-world-scenarios}

\section{Overview}\label{overview-4}

See how to apply \texttt{add.to()} in complete business workflows. These
examples combine multiple techniques you've learned to solve real-world
data problems.

\textbf{What you'll learn:} - Multi-step data enrichment workflows -
Combining lookups from multiple reference tables - Sales analytics with
aggregation - Best practices for production use

\textbf{Prerequisites:} - \href{examples-to-01-basic.html}{Basic Lookup}
- Single column operations -
\href{examples-to-02-multiple.html}{Multiple Columns \& Keys} - Working
with multiple columns - \href{examples-to-04-strategy.html}{Aggregation
Strategies} - Using strategy parameter

\begin{center}\rule{0.5\linewidth}{0.5pt}\end{center}

\section{Example 1: E-commerce Order
Enrichment}\label{example-1-e-commerce-order-enrichment}

\textbf{Business Context:} You have a raw orders table with just IDs.
You need to create a complete order report with customer names, emails,
customer tiers, product names, and prices for your fulfillment team.

\textbf{Code:}

\begin{Shaded}
\begin{Highlighting}[]
\ImportTok{import}\NormalTok{ additory }\ImportTok{as}\NormalTok{ add}
\ImportTok{import}\NormalTok{ pandas }\ImportTok{as}\NormalTok{ pd}

\CommentTok{\# Raw orders (just IDs and quantities)}
\NormalTok{orders }\OperatorTok{=}\NormalTok{ pd.DataFrame(\{}
    \StringTok{\textquotesingle{}order\_id\textquotesingle{}}\NormalTok{: [}\DecValTok{1001}\NormalTok{, }\DecValTok{1002}\NormalTok{, }\DecValTok{1003}\NormalTok{],}
    \StringTok{\textquotesingle{}customer\_id\textquotesingle{}}\NormalTok{: [}\DecValTok{101}\NormalTok{, }\DecValTok{102}\NormalTok{, }\DecValTok{103}\NormalTok{],}
    \StringTok{\textquotesingle{}product\_id\textquotesingle{}}\NormalTok{: [}\DecValTok{501}\NormalTok{, }\DecValTok{502}\NormalTok{, }\DecValTok{503}\NormalTok{],}
    \StringTok{\textquotesingle{}quantity\textquotesingle{}}\NormalTok{: [}\DecValTok{2}\NormalTok{, }\DecValTok{1}\NormalTok{, }\DecValTok{3}\NormalTok{]}
\NormalTok{\})}

\CommentTok{\# Customer master data}
\NormalTok{customers }\OperatorTok{=}\NormalTok{ pd.DataFrame(\{}
    \StringTok{\textquotesingle{}customer\_id\textquotesingle{}}\NormalTok{: [}\DecValTok{101}\NormalTok{, }\DecValTok{102}\NormalTok{, }\DecValTok{103}\NormalTok{],}
    \StringTok{\textquotesingle{}name\textquotesingle{}}\NormalTok{: [}\StringTok{\textquotesingle{}Alice\textquotesingle{}}\NormalTok{, }\StringTok{\textquotesingle{}Bob\textquotesingle{}}\NormalTok{, }\StringTok{\textquotesingle{}Charlie\textquotesingle{}}\NormalTok{],}
    \StringTok{\textquotesingle{}email\textquotesingle{}}\NormalTok{: [}\StringTok{\textquotesingle{}alice@email.com\textquotesingle{}}\NormalTok{, }\StringTok{\textquotesingle{}bob@email.com\textquotesingle{}}\NormalTok{, }\StringTok{\textquotesingle{}charlie@email.com\textquotesingle{}}\NormalTok{],}
    \StringTok{\textquotesingle{}tier\textquotesingle{}}\NormalTok{: [}\StringTok{\textquotesingle{}Gold\textquotesingle{}}\NormalTok{, }\StringTok{\textquotesingle{}Silver\textquotesingle{}}\NormalTok{, }\StringTok{\textquotesingle{}Bronze\textquotesingle{}}\NormalTok{]}
\NormalTok{\})}

\CommentTok{\# Product catalog}
\NormalTok{products }\OperatorTok{=}\NormalTok{ pd.DataFrame(\{}
    \StringTok{\textquotesingle{}product\_id\textquotesingle{}}\NormalTok{: [}\DecValTok{501}\NormalTok{, }\DecValTok{502}\NormalTok{, }\DecValTok{503}\NormalTok{],}
    \StringTok{\textquotesingle{}product\_name\textquotesingle{}}\NormalTok{: [}\StringTok{\textquotesingle{}Widget\textquotesingle{}}\NormalTok{, }\StringTok{\textquotesingle{}Gadget\textquotesingle{}}\NormalTok{, }\StringTok{\textquotesingle{}Doohickey\textquotesingle{}}\NormalTok{],}
    \StringTok{\textquotesingle{}price\textquotesingle{}}\NormalTok{: [}\FloatTok{29.99}\NormalTok{, }\FloatTok{49.99}\NormalTok{, }\FloatTok{19.99}\NormalTok{]}
\NormalTok{\})}

\CommentTok{\# Step 1: Add customer information}
\NormalTok{enriched\_orders }\OperatorTok{=}\NormalTok{ add.to(}
\NormalTok{    orders,}
\NormalTok{    bring\_from}\OperatorTok{=}\NormalTok{customers,}
\NormalTok{    bring}\OperatorTok{=}\NormalTok{[}\StringTok{\textquotesingle{}name\textquotesingle{}}\NormalTok{, }\StringTok{\textquotesingle{}email\textquotesingle{}}\NormalTok{, }\StringTok{\textquotesingle{}tier\textquotesingle{}}\NormalTok{],}
\NormalTok{    against}\OperatorTok{=}\StringTok{\textquotesingle{}customer\_id\textquotesingle{}}
\NormalTok{)}

\CommentTok{\# Positional parameters (also works without naming certain parameters):}
\CommentTok{\# enriched\_orders = add.to(orders, customers, [\textquotesingle{}name\textquotesingle{}, \textquotesingle{}email\textquotesingle{}, \textquotesingle{}tier\textquotesingle{}], \textquotesingle{}customer\_id\textquotesingle{})}

\BuiltInTok{print}\NormalTok{(}\StringTok{"After adding customer info:"}\NormalTok{)}
\BuiltInTok{print}\NormalTok{(enriched\_orders)}

\CommentTok{\# Step 2: Add product information}
\NormalTok{enriched\_orders }\OperatorTok{=}\NormalTok{ add.to(}
\NormalTok{    enriched\_orders,}
\NormalTok{    bring\_from}\OperatorTok{=}\NormalTok{products,}
\NormalTok{    bring}\OperatorTok{=}\NormalTok{[}\StringTok{\textquotesingle{}product\_name\textquotesingle{}}\NormalTok{, }\StringTok{\textquotesingle{}price\textquotesingle{}}\NormalTok{],}
\NormalTok{    against}\OperatorTok{=}\StringTok{\textquotesingle{}product\_id\textquotesingle{}}
\NormalTok{)}

\BuiltInTok{print}\NormalTok{(}\StringTok{"}\CharTok{\textbackslash{}n}\StringTok{Fully enriched orders:"}\NormalTok{)}
\BuiltInTok{print}\NormalTok{(enriched\_orders)}

\CommentTok{\# Step 3: Calculate total (using pandas)}
\NormalTok{enriched\_orders[}\StringTok{\textquotesingle{}total\textquotesingle{}}\NormalTok{] }\OperatorTok{=}\NormalTok{ enriched\_orders[}\StringTok{\textquotesingle{}quantity\textquotesingle{}}\NormalTok{] }\OperatorTok{*}\NormalTok{ enriched\_orders[}\StringTok{\textquotesingle{}price\textquotesingle{}}\NormalTok{]}

\BuiltInTok{print}\NormalTok{(}\StringTok{"}\CharTok{\textbackslash{}n}\StringTok{Final order report:"}\NormalTok{)}
\BuiltInTok{print}\NormalTok{(enriched\_orders[[}\StringTok{\textquotesingle{}order\_id\textquotesingle{}}\NormalTok{, }\StringTok{\textquotesingle{}name\textquotesingle{}}\NormalTok{, }\StringTok{\textquotesingle{}email\textquotesingle{}}\NormalTok{, }\StringTok{\textquotesingle{}tier\textquotesingle{}}\NormalTok{, }
                        \StringTok{\textquotesingle{}product\_name\textquotesingle{}}\NormalTok{, }\StringTok{\textquotesingle{}quantity\textquotesingle{}}\NormalTok{, }\StringTok{\textquotesingle{}price\textquotesingle{}}\NormalTok{, }\StringTok{\textquotesingle{}total\textquotesingle{}}\NormalTok{]])}
\end{Highlighting}
\end{Shaded}

\textbf{Output:}

\begin{verbatim}
After adding customer info:
   order_id  customer_id  product_id  quantity     name              email     tier
0      1001          101         501         2    Alice  alice@email.com     Gold
1      1002          102         502         1      Bob    bob@email.com   Silver
2      1003          103         503         3  Charlie charlie@email.com   Bronze

Fully enriched orders:
   order_id  customer_id  product_id  quantity     name              email     tier product_name  price
0      1001          101         501         2    Alice  alice@email.com     Gold       Widget  29.99
1      1002          102         502         1      Bob    bob@email.com   Silver       Gadget  49.99
2      1003          103         503         3  Charlie charlie@email.com   Bronze    Doohickey  19.99

Final order report:
   order_id     name              email     tier product_name  quantity  price  total
0      1001    Alice  alice@email.com     Gold       Widget         2  29.99  59.98
1      1002      Bob    bob@email.com   Silver       Gadget         1  49.99  49.99
2      1003  Charlie charlie@email.com   Bronze    Doohickey         3  19.99  59.97
\end{verbatim}

\textbf{Explanation:} - Start with minimal order data (just IDs) - First
enrichment: Add customer details (name, email, tier) - Second
enrichment: Add product details (name, price) - Final step: Calculate
totals using pandas - Each \texttt{add.to()} builds on the previous
result - This pattern is common in ETL (Extract, Transform, Load)
workflows

\textbf{Key Takeaways:} - Chain multiple \texttt{add.to()} calls for
complex enrichment - Each step adds more context to your data - Keep
reference tables separate for maintainability - Combine with
pandas/polars operations for calculations

\textbf{Note:} This also works with polars DataFrames.

\begin{center}\rule{0.5\linewidth}{0.5pt}\end{center}

\section{Example 2: Sales Analytics with
Aggregation}\label{example-2-sales-analytics-with-aggregation}

\textbf{Business Context:} You have detailed transaction data and need
to create a store performance report showing total revenue by store,
along with store metadata (name and region).

\textbf{Code:}

\begin{Shaded}
\begin{Highlighting}[]
\ImportTok{import}\NormalTok{ additory }\ImportTok{as}\NormalTok{ add}
\ImportTok{import}\NormalTok{ pandas }\ImportTok{as}\NormalTok{ pd}

\CommentTok{\# Detailed transactions}
\NormalTok{transactions }\OperatorTok{=}\NormalTok{ pd.DataFrame(\{}
    \StringTok{\textquotesingle{}transaction\_id\textquotesingle{}}\NormalTok{: [}\DecValTok{1}\NormalTok{, }\DecValTok{2}\NormalTok{, }\DecValTok{3}\NormalTok{, }\DecValTok{4}\NormalTok{, }\DecValTok{5}\NormalTok{, }\DecValTok{6}\NormalTok{],}
    \StringTok{\textquotesingle{}store\_id\textquotesingle{}}\NormalTok{: [}\DecValTok{1}\NormalTok{, }\DecValTok{1}\NormalTok{, }\DecValTok{2}\NormalTok{, }\DecValTok{2}\NormalTok{, }\DecValTok{3}\NormalTok{, }\DecValTok{3}\NormalTok{],}
    \StringTok{\textquotesingle{}amount\textquotesingle{}}\NormalTok{: [}\DecValTok{100}\NormalTok{, }\DecValTok{150}\NormalTok{, }\DecValTok{200}\NormalTok{, }\DecValTok{250}\NormalTok{, }\DecValTok{175}\NormalTok{, }\DecValTok{225}\NormalTok{],}
    \StringTok{\textquotesingle{}date\textquotesingle{}}\NormalTok{: [}\StringTok{\textquotesingle{}2024{-}01{-}15\textquotesingle{}}\NormalTok{, }\StringTok{\textquotesingle{}2024{-}01{-}16\textquotesingle{}}\NormalTok{, }\StringTok{\textquotesingle{}2024{-}01{-}15\textquotesingle{}}\NormalTok{, }
             \StringTok{\textquotesingle{}2024{-}01{-}16\textquotesingle{}}\NormalTok{, }\StringTok{\textquotesingle{}2024{-}01{-}15\textquotesingle{}}\NormalTok{, }\StringTok{\textquotesingle{}2024{-}01{-}16\textquotesingle{}}\NormalTok{]}
\NormalTok{\})}

\CommentTok{\# Store master data}
\NormalTok{stores }\OperatorTok{=}\NormalTok{ pd.DataFrame(\{}
    \StringTok{\textquotesingle{}store\_id\textquotesingle{}}\NormalTok{: [}\DecValTok{1}\NormalTok{, }\DecValTok{2}\NormalTok{, }\DecValTok{3}\NormalTok{],}
    \StringTok{\textquotesingle{}store\_name\textquotesingle{}}\NormalTok{: [}\StringTok{\textquotesingle{}Downtown\textquotesingle{}}\NormalTok{, }\StringTok{\textquotesingle{}Mall\textquotesingle{}}\NormalTok{, }\StringTok{\textquotesingle{}Airport\textquotesingle{}}\NormalTok{],}
    \StringTok{\textquotesingle{}region\textquotesingle{}}\NormalTok{: [}\StringTok{\textquotesingle{}North\textquotesingle{}}\NormalTok{, }\StringTok{\textquotesingle{}South\textquotesingle{}}\NormalTok{, }\StringTok{\textquotesingle{}East\textquotesingle{}}\NormalTok{]}
\NormalTok{\})}

\CommentTok{\# Aggregate transactions to store level}
\NormalTok{store\_performance }\OperatorTok{=}\NormalTok{ add.to(}
\NormalTok{    stores,}
\NormalTok{    bring\_from}\OperatorTok{=}\NormalTok{transactions,}
\NormalTok{    bring}\OperatorTok{=}\StringTok{\textquotesingle{}amount\textquotesingle{}}\NormalTok{,}
\NormalTok{    against}\OperatorTok{=}\StringTok{\textquotesingle{}store\_id\textquotesingle{}}\NormalTok{,}
\NormalTok{    strategy}\OperatorTok{=}\NormalTok{\{}\StringTok{\textquotesingle{}amount\textquotesingle{}}\NormalTok{: }\StringTok{\textquotesingle{}sum\textquotesingle{}}\NormalTok{\}  }\CommentTok{\# Sum all transactions per store}
\NormalTok{)}

\CommentTok{\# Positional parameters (also works without naming certain parameters):}
\CommentTok{\# store\_performance = add.to(stores, transactions, \textquotesingle{}amount\textquotesingle{}, \textquotesingle{}store\_id\textquotesingle{}, strategy=\{\textquotesingle{}amount\textquotesingle{}: \textquotesingle{}sum\textquotesingle{}\})}

\BuiltInTok{print}\NormalTok{(}\StringTok{"Store Performance Report:"}\NormalTok{)}
\BuiltInTok{print}\NormalTok{(store\_performance)}

\CommentTok{\# Calculate percentage of total}
\NormalTok{total\_revenue }\OperatorTok{=}\NormalTok{ store\_performance[}\StringTok{\textquotesingle{}amount\textquotesingle{}}\NormalTok{].}\BuiltInTok{sum}\NormalTok{()}
\NormalTok{store\_performance[}\StringTok{\textquotesingle{}pct\_of\_total\textquotesingle{}}\NormalTok{] }\OperatorTok{=}\NormalTok{ (store\_performance[}\StringTok{\textquotesingle{}amount\textquotesingle{}}\NormalTok{] }\OperatorTok{/}\NormalTok{ total\_revenue }\OperatorTok{*} \DecValTok{100}\NormalTok{).}\BuiltInTok{round}\NormalTok{(}\DecValTok{2}\NormalTok{)}

\BuiltInTok{print}\NormalTok{(}\StringTok{"}\CharTok{\textbackslash{}n}\StringTok{With Percentage:"}\NormalTok{)}
\BuiltInTok{print}\NormalTok{(store\_performance)}
\end{Highlighting}
\end{Shaded}

\textbf{Output:}

\begin{verbatim}
Store Performance Report:
   store_id store_name region  amount
0         1   Downtown  North     250
1         2       Mall  South     450
2         3    Airport   East     400

With Percentage:
   store_id store_name region  amount  pct_of_total
0         1   Downtown  North     250         22.73
1         2       Mall  South     450         40.91
2         3    Airport   East     400         36.36
\end{verbatim}

\textbf{Explanation:} - Start with store master data (the target) -
Aggregate transactions using
\texttt{strategy=\{\textquotesingle{}amount\textquotesingle{}:\ \textquotesingle{}sum\textquotesingle{}\}}
- Downtown: 100 + 150 = 250 - Mall: 200 + 250 = 450 - Airport: 175 + 225
= 400 - Add percentage calculations using pandas - This pattern is
common for creating summary reports

\textbf{Key Takeaways:} - Use stores as target (not transactions) for
aggregation - Strategy parameter handles multiple transactions per store
- Result has one row per store with aggregated values - Combine with
pandas for additional calculations

\textbf{Note:} This also works with polars DataFrames.

\begin{center}\rule{0.5\linewidth}{0.5pt}\end{center}

\section{Best Practices for Production
Use}\label{best-practices-for-production-use}

\subsection{1. Validate Your Data First}\label{validate-your-data-first}

\begin{Shaded}
\begin{Highlighting}[]
\CommentTok{\# Check for missing keys before joining}
\NormalTok{missing\_customers }\OperatorTok{=}\NormalTok{ orders[}\OperatorTok{\textasciitilde{}}\NormalTok{orders[}\StringTok{\textquotesingle{}customer\_id\textquotesingle{}}\NormalTok{].isin(customers[}\StringTok{\textquotesingle{}customer\_id\textquotesingle{}}\NormalTok{])]}
\ControlFlowTok{if} \BuiltInTok{len}\NormalTok{(missing\_customers) }\OperatorTok{\textgreater{}} \DecValTok{0}\NormalTok{:}
    \BuiltInTok{print}\NormalTok{(}\SpecialStringTok{f"Warning: }\SpecialCharTok{\{}\BuiltInTok{len}\NormalTok{(missing\_customers)}\SpecialCharTok{\}}\SpecialStringTok{ orders have invalid customer IDs"}\NormalTok{)}
\end{Highlighting}
\end{Shaded}

\subsection{2. Use Explicit Column
Names}\label{use-explicit-column-names}

\begin{Shaded}
\begin{Highlighting}[]
\CommentTok{\# Good: Explicit and clear}
\NormalTok{result }\OperatorTok{=}\NormalTok{ add.to(}
\NormalTok{    orders,}
\NormalTok{    bring\_from}\OperatorTok{=}\NormalTok{customers,}
\NormalTok{    bring}\OperatorTok{=}\NormalTok{[}\StringTok{\textquotesingle{}name\textquotesingle{}}\NormalTok{, }\StringTok{\textquotesingle{}email\textquotesingle{}}\NormalTok{],}
\NormalTok{    against}\OperatorTok{=}\StringTok{\textquotesingle{}customer\_id\textquotesingle{}}
\NormalTok{)}

\CommentTok{\# Avoid: Relying on column order or implicit behavior}
\end{Highlighting}
\end{Shaded}

\subsection{3. Handle Missing Matches}\label{handle-missing-matches}

\begin{Shaded}
\begin{Highlighting}[]
\CommentTok{\# Use left join to keep all orders}
\NormalTok{result }\OperatorTok{=}\NormalTok{ add.to(}
\NormalTok{    orders,}
\NormalTok{    bring\_from}\OperatorTok{=}\NormalTok{customers,}
\NormalTok{    bring}\OperatorTok{=}\StringTok{\textquotesingle{}name\textquotesingle{}}\NormalTok{,}
\NormalTok{    against}\OperatorTok{=}\StringTok{\textquotesingle{}customer\_id\textquotesingle{}}\NormalTok{,}
\NormalTok{    join\_type}\OperatorTok{=}\StringTok{\textquotesingle{}left\textquotesingle{}}
\NormalTok{)}

\CommentTok{\# Check for missing matches}
\NormalTok{missing }\OperatorTok{=}\NormalTok{ result[result[}\StringTok{\textquotesingle{}name\textquotesingle{}}\NormalTok{].isna()]}
\ControlFlowTok{if} \BuiltInTok{len}\NormalTok{(missing) }\OperatorTok{\textgreater{}} \DecValTok{0}\NormalTok{:}
    \BuiltInTok{print}\NormalTok{(}\SpecialStringTok{f"Warning: }\SpecialCharTok{\{}\BuiltInTok{len}\NormalTok{(missing)}\SpecialCharTok{\}}\SpecialStringTok{ orders have no customer match"}\NormalTok{)}
\end{Highlighting}
\end{Shaded}

\subsection{4. Chain Operations
Carefully}\label{chain-operations-carefully}

\begin{Shaded}
\begin{Highlighting}[]
\CommentTok{\# Good: Clear intermediate steps}
\NormalTok{step1 }\OperatorTok{=}\NormalTok{ add.to(orders, customers, }\StringTok{\textquotesingle{}name\textquotesingle{}}\NormalTok{, }\StringTok{\textquotesingle{}customer\_id\textquotesingle{}}\NormalTok{)}
\NormalTok{step2 }\OperatorTok{=}\NormalTok{ add.to(step1, products, }\StringTok{\textquotesingle{}price\textquotesingle{}}\NormalTok{, }\StringTok{\textquotesingle{}product\_id\textquotesingle{}}\NormalTok{)}
\NormalTok{final }\OperatorTok{=}\NormalTok{ step2}

\CommentTok{\# Also good: Chain with clear variable names}
\NormalTok{result }\OperatorTok{=}\NormalTok{ add.to(orders, customers, }\StringTok{\textquotesingle{}name\textquotesingle{}}\NormalTok{, }\StringTok{\textquotesingle{}customer\_id\textquotesingle{}}\NormalTok{)}
\NormalTok{result }\OperatorTok{=}\NormalTok{ add.to(result, products, }\StringTok{\textquotesingle{}price\textquotesingle{}}\NormalTok{, }\StringTok{\textquotesingle{}product\_id\textquotesingle{}}\NormalTok{)}
\end{Highlighting}
\end{Shaded}

\subsection{5. Document Your Strategy
Choices}\label{document-your-strategy-choices}

\begin{Shaded}
\begin{Highlighting}[]
\CommentTok{\# Good: Comment explains why}
\NormalTok{result }\OperatorTok{=}\NormalTok{ add.to(}
\NormalTok{    customers,}
\NormalTok{    bring\_from}\OperatorTok{=}\NormalTok{orders,}
\NormalTok{    bring}\OperatorTok{=}\StringTok{\textquotesingle{}amount\textquotesingle{}}\NormalTok{,}
\NormalTok{    against}\OperatorTok{=}\StringTok{\textquotesingle{}customer\_id\textquotesingle{}}\NormalTok{,}
\NormalTok{    strategy}\OperatorTok{=}\NormalTok{\{}\StringTok{\textquotesingle{}amount\textquotesingle{}}\NormalTok{: }\StringTok{\textquotesingle{}sum\textquotesingle{}}\NormalTok{\}  }\CommentTok{\# Total revenue per customer}
\NormalTok{)}
\end{Highlighting}
\end{Shaded}

\begin{center}\rule{0.5\linewidth}{0.5pt}\end{center}

\section{Common Patterns}\label{common-patterns}

\subsection{Pattern 1: Dimension Table
Enrichment}\label{pattern-1-dimension-table-enrichment}

\begin{Shaded}
\begin{Highlighting}[]
\CommentTok{\# Add dimension attributes to fact table}
\NormalTok{fact\_enriched }\OperatorTok{=}\NormalTok{ add.to(fact\_table, dim\_table, [}\StringTok{\textquotesingle{}attr1\textquotesingle{}}\NormalTok{, }\StringTok{\textquotesingle{}attr2\textquotesingle{}}\NormalTok{], }\StringTok{\textquotesingle{}key\textquotesingle{}}\NormalTok{)}
\end{Highlighting}
\end{Shaded}

\subsection{Pattern 2: Aggregation to Summary
Level}\label{pattern-2-aggregation-to-summary-level}

\begin{Shaded}
\begin{Highlighting}[]
\CommentTok{\# Aggregate detail to summary}
\NormalTok{summary }\OperatorTok{=}\NormalTok{ add.to(master, detail, }\StringTok{\textquotesingle{}metric\textquotesingle{}}\NormalTok{, }\StringTok{\textquotesingle{}key\textquotesingle{}}\NormalTok{, strategy}\OperatorTok{=}\NormalTok{\{}\StringTok{\textquotesingle{}metric\textquotesingle{}}\NormalTok{: }\StringTok{\textquotesingle{}sum\textquotesingle{}}\NormalTok{\})}
\end{Highlighting}
\end{Shaded}

\subsection{Pattern 3: Multi-Step
Enrichment}\label{pattern-3-multi-step-enrichment}

\begin{Shaded}
\begin{Highlighting}[]
\CommentTok{\# Chain multiple enrichments}
\NormalTok{result }\OperatorTok{=}\NormalTok{ add.to(base, ref1, }\StringTok{\textquotesingle{}col1\textquotesingle{}}\NormalTok{, }\StringTok{\textquotesingle{}key1\textquotesingle{}}\NormalTok{)}
\NormalTok{result }\OperatorTok{=}\NormalTok{ add.to(result, ref2, }\StringTok{\textquotesingle{}col2\textquotesingle{}}\NormalTok{, }\StringTok{\textquotesingle{}key2\textquotesingle{}}\NormalTok{)}
\NormalTok{result }\OperatorTok{=}\NormalTok{ add.to(result, ref3, }\StringTok{\textquotesingle{}col3\textquotesingle{}}\NormalTok{, }\StringTok{\textquotesingle{}key3\textquotesingle{}}\NormalTok{)}
\end{Highlighting}
\end{Shaded}

\subsection{Pattern 4: Lookup with
Fallback}\label{pattern-4-lookup-with-fallback}

\begin{Shaded}
\begin{Highlighting}[]
\CommentTok{\# Use left join and fill missing values}
\NormalTok{result }\OperatorTok{=}\NormalTok{ add.to(target, ref, }\StringTok{\textquotesingle{}col\textquotesingle{}}\NormalTok{, }\StringTok{\textquotesingle{}key\textquotesingle{}}\NormalTok{, join\_type}\OperatorTok{=}\StringTok{\textquotesingle{}left\textquotesingle{}}\NormalTok{)}
\NormalTok{result[}\StringTok{\textquotesingle{}col\textquotesingle{}}\NormalTok{] }\OperatorTok{=}\NormalTok{ result[}\StringTok{\textquotesingle{}col\textquotesingle{}}\NormalTok{].fillna(}\StringTok{\textquotesingle{}Unknown\textquotesingle{}}\NormalTok{)}
\end{Highlighting}
\end{Shaded}

\begin{center}\rule{0.5\linewidth}{0.5pt}\end{center}

\section{Performance Tips}\label{performance-tips}

\begin{enumerate}
\def\labelenumi{\arabic{enumi}.}
\tightlist
\item
  \textbf{Filter before joining}: Reduce data size before
  \texttt{add.to()}
\item
  \textbf{Use appropriate join types}: \texttt{inner} is faster than
  \texttt{left} when you don't need all rows
\item
  \textbf{Aggregate early}: Summarize large tables before joining
\item
  \textbf{Index your keys}: Ensure key columns are indexed (for large
  datasets)
\item
  \textbf{Consider memory}: For very large datasets, process in chunks
\end{enumerate}

\begin{center}\rule{0.5\linewidth}{0.5pt}\end{center}

\section{Key Takeaways}\label{key-takeaways-4}

\begin{itemize}
\tightlist
\item
  Chain multiple \texttt{add.to()} calls for complex workflows
\item
  Start with master/dimension tables, add details as needed
\item
  Use aggregation strategies for summary reports
\item
  Validate data and handle missing matches explicitly
\item
  Document your strategy choices for maintainability
\item
  Combine with pandas/polars for calculations and transformations
\end{itemize}

\begin{center}\rule{0.5\linewidth}{0.5pt}\end{center}

\section{Common Questions}\label{common-questions-4}

\textbf{Q: Should I add all columns at once or in multiple steps?}\\
A: It depends. Multiple steps are clearer and easier to debug, but a
single step is more efficient. For production, prefer clarity unless
performance is critical.

\textbf{Q: How do I handle slowly changing dimensions (SCD)?}\\
A: Use date-based keys or add date filters before joining. We'll cover
this in advanced examples.

\textbf{Q: Can I use add.to() with very large datasets (millions of
rows)?}\\
A: Yes, but consider filtering and aggregating first. For extremely
large data, consider database joins instead.

\textbf{Q: What if my keys have different names in different tables?}\\
A: Rename columns first using pandas/polars, or use advanced
\texttt{add.to()} features (covered in advanced examples).

\begin{center}\rule{0.5\linewidth}{0.5pt}\end{center}

\section{Next Steps}\label{next-steps-6}

\begin{itemize}
\tightlist
\item
  \textbf{\href{examples-transform-01-calc.html}{add.transform()
  Examples}} - Learn data transformation operations
\item
  \textbf{\href{examples-to-04-strategy.html}{Aggregation Strategies}} -
  Review strategy parameter details
\item
  \textbf{\href{reference-to-basic.html}{API Reference}} - Complete
  \texttt{add.to()} documentation
\end{itemize}

\begin{center}\rule{0.5\linewidth}{0.5pt}\end{center}

\section{Congratulations!}\label{congratulations}

You've completed the \texttt{add.to()} tutorial series. You now know how
to: - ✅ Perform basic lookups - ✅ Work with multiple columns and keys
- ✅ Handle one-to-many and many-to-one patterns - ✅ Use aggregation
strategies - ✅ Build complete real-world workflows

Ready to learn more? Explore \texttt{add.transform()} for data
transformations, or dive into the API reference for advanced features.

\begin{center}\rule{0.5\linewidth}{0.5pt}\end{center}

\textbf{Version:} 0.1.3\\
\textbf{Last Updated:} March 9, 2026

\part{add.transform() - Data Transformations}

\chapter{add.transform() -
Calculations}\label{add.transform---calculations}

\section{Overview}\label{overview-5}

Learn how to create new columns using calculations with
\texttt{add.transform()} in \texttt{@calc} mode. This is one of the most
powerful features for deriving new insights from your data.

\textbf{What you'll learn:} - How to perform simple calculations - How
to create multiple calculated columns at once - How to use different
arithmetic operations - Expression syntax and best practices

\textbf{Prerequisites:} - Basic understanding of DataFrames (pandas or
polars) - Familiarity with arithmetic operations

\begin{center}\rule{0.5\linewidth}{0.5pt}\end{center}

\section{Example 1: Simple
Calculation}\label{example-1-simple-calculation}

\textbf{Business Context:} You have product inventory with prices and
quantities. You need to calculate the total value of each product line.

\textbf{Code:}

\begin{Shaded}
\begin{Highlighting}[]
\ImportTok{import}\NormalTok{ additory }\ImportTok{as}\NormalTok{ add}
\ImportTok{import}\NormalTok{ pandas }\ImportTok{as}\NormalTok{ pd}

\CommentTok{\# Product inventory}
\NormalTok{df }\OperatorTok{=}\NormalTok{ pd.DataFrame(\{}
    \StringTok{\textquotesingle{}product\textquotesingle{}}\NormalTok{: [}\StringTok{\textquotesingle{}Widget\textquotesingle{}}\NormalTok{, }\StringTok{\textquotesingle{}Gadget\textquotesingle{}}\NormalTok{, }\StringTok{\textquotesingle{}Doohickey\textquotesingle{}}\NormalTok{],}
    \StringTok{\textquotesingle{}price\textquotesingle{}}\NormalTok{: [}\FloatTok{29.99}\NormalTok{, }\FloatTok{49.99}\NormalTok{, }\FloatTok{19.99}\NormalTok{],}
    \StringTok{\textquotesingle{}quantity\textquotesingle{}}\NormalTok{: [}\DecValTok{10}\NormalTok{, }\DecValTok{5}\NormalTok{, }\DecValTok{15}\NormalTok{]}
\NormalTok{\})}

\CommentTok{\# Calculate total value}
\NormalTok{result }\OperatorTok{=}\NormalTok{ add.transform(}
    \StringTok{\textquotesingle{}@calc\textquotesingle{}}\NormalTok{,}
\NormalTok{    df,}
\NormalTok{    columns}\OperatorTok{=}\NormalTok{[}\StringTok{\textquotesingle{}price\textquotesingle{}}\NormalTok{, }\StringTok{\textquotesingle{}quantity\textquotesingle{}}\NormalTok{],}
\NormalTok{    expression}\OperatorTok{=}\StringTok{\textquotesingle{}price * quantity\textquotesingle{}}\NormalTok{,}
\NormalTok{    as\_}\OperatorTok{=}\StringTok{\textquotesingle{}total\_value\textquotesingle{}}
\NormalTok{)}

\CommentTok{\# Positional parameters (also works without naming certain parameters):}
\CommentTok{\# result = add.transform(\textquotesingle{}@calc\textquotesingle{}, df, [\textquotesingle{}price\textquotesingle{}, \textquotesingle{}quantity\textquotesingle{}], expression=\textquotesingle{}price * quantity\textquotesingle{}, as\_=\textquotesingle{}total\_value\textquotesingle{})}

\BuiltInTok{print}\NormalTok{(result)}
\end{Highlighting}
\end{Shaded}

\textbf{Output:}

\begin{verbatim}
      product  price  quantity  total_value
0      Widget  29.99        10       299.90
1      Gadget  49.99         5       249.95
2  Doohickey  19.99        15       299.85
\end{verbatim}

\textbf{Explanation:} -
\texttt{\textquotesingle{}@calc\textquotesingle{}} mode creates new
calculated columns - \texttt{columns} specifies which columns are used
in the calculation - \texttt{expression} defines the calculation formula
- \texttt{as\_} names the new column - The calculation is applied to
every row - Original columns are preserved

\textbf{Note:} This also works with polars DataFrames.

\begin{center}\rule{0.5\linewidth}{0.5pt}\end{center}

\section{Example 2: Multiple
Calculations}\label{example-2-multiple-calculations}

\textbf{Business Context:} You have product pricing data and need to
calculate both profit and total value for each product.

\textbf{Code:}

\begin{Shaded}
\begin{Highlighting}[]
\ImportTok{import}\NormalTok{ additory }\ImportTok{as}\NormalTok{ add}
\ImportTok{import}\NormalTok{ pandas }\ImportTok{as}\NormalTok{ pd}

\CommentTok{\# Product data}
\NormalTok{df }\OperatorTok{=}\NormalTok{ pd.DataFrame(\{}
    \StringTok{\textquotesingle{}product\textquotesingle{}}\NormalTok{: [}\StringTok{\textquotesingle{}Widget\textquotesingle{}}\NormalTok{, }\StringTok{\textquotesingle{}Gadget\textquotesingle{}}\NormalTok{, }\StringTok{\textquotesingle{}Doohickey\textquotesingle{}}\NormalTok{],}
    \StringTok{\textquotesingle{}price\textquotesingle{}}\NormalTok{: [}\DecValTok{100}\NormalTok{, }\DecValTok{200}\NormalTok{, }\DecValTok{150}\NormalTok{],}
    \StringTok{\textquotesingle{}cost\textquotesingle{}}\NormalTok{: [}\DecValTok{60}\NormalTok{, }\DecValTok{120}\NormalTok{, }\DecValTok{90}\NormalTok{],}
    \StringTok{\textquotesingle{}quantity\textquotesingle{}}\NormalTok{: [}\DecValTok{10}\NormalTok{, }\DecValTok{5}\NormalTok{, }\DecValTok{15}\NormalTok{]}
\NormalTok{\})}

\CommentTok{\# Calculate profit and total value}
\NormalTok{result }\OperatorTok{=}\NormalTok{ add.transform(}
    \StringTok{\textquotesingle{}@calc\textquotesingle{}}\NormalTok{,}
\NormalTok{    df,}
\NormalTok{    columns}\OperatorTok{=}\NormalTok{[}\StringTok{\textquotesingle{}price\textquotesingle{}}\NormalTok{, }\StringTok{\textquotesingle{}cost\textquotesingle{}}\NormalTok{, }\StringTok{\textquotesingle{}quantity\textquotesingle{}}\NormalTok{],}
\NormalTok{    expression}\OperatorTok{=}\NormalTok{[}\StringTok{\textquotesingle{}price {-} cost\textquotesingle{}}\NormalTok{, }\StringTok{\textquotesingle{}price * quantity\textquotesingle{}}\NormalTok{],}
\NormalTok{    as\_}\OperatorTok{=}\NormalTok{[}\StringTok{\textquotesingle{}profit\textquotesingle{}}\NormalTok{, }\StringTok{\textquotesingle{}total\_value\textquotesingle{}}\NormalTok{]}
\NormalTok{)}

\CommentTok{\# Positional parameters (also works without naming certain parameters):}
\CommentTok{\# result = add.transform(\textquotesingle{}@calc\textquotesingle{}, df, [\textquotesingle{}price\textquotesingle{}, \textquotesingle{}cost\textquotesingle{}, \textquotesingle{}quantity\textquotesingle{}], }
\CommentTok{\#                        expression=[\textquotesingle{}price {-} cost\textquotesingle{}, \textquotesingle{}price * quantity\textquotesingle{}], }
\CommentTok{\#                        as\_=[\textquotesingle{}profit\textquotesingle{}, \textquotesingle{}total\_value\textquotesingle{}])}

\BuiltInTok{print}\NormalTok{(result)}
\end{Highlighting}
\end{Shaded}

\textbf{Output:}

\begin{verbatim}
      product  price  cost  quantity  profit  total_value
0      Widget    100    60        10      40         1000
1      Gadget    200   120         5      80         1000
2  Doohickey    150    90        15      60         2250
\end{verbatim}

\textbf{Explanation:} - Use lists for \texttt{expression} and
\texttt{as\_} to create multiple columns - Each expression creates one
new column - The number of expressions must match the number of names in
\texttt{as\_} - All calculations happen in a single operation - More
efficient than calling \texttt{add.transform()} multiple times

\textbf{Note:} This also works with polars DataFrames.

\begin{center}\rule{0.5\linewidth}{0.5pt}\end{center}

\section{Example 3: Division
Operations}\label{example-3-division-operations}

\textbf{Business Context:} You need to calculate the price per unit for
products sold in bulk.

\textbf{Code:}

\begin{Shaded}
\begin{Highlighting}[]
\ImportTok{import}\NormalTok{ additory }\ImportTok{as}\NormalTok{ add}
\ImportTok{import}\NormalTok{ pandas }\ImportTok{as}\NormalTok{ pd}

\CommentTok{\# Bulk products}
\NormalTok{df }\OperatorTok{=}\NormalTok{ pd.DataFrame(\{}
    \StringTok{\textquotesingle{}product\textquotesingle{}}\NormalTok{: [}\StringTok{\textquotesingle{}Widget\textquotesingle{}}\NormalTok{, }\StringTok{\textquotesingle{}Gadget\textquotesingle{}}\NormalTok{, }\StringTok{\textquotesingle{}Doohickey\textquotesingle{}}\NormalTok{],}
    \StringTok{\textquotesingle{}price\textquotesingle{}}\NormalTok{: [}\DecValTok{100}\NormalTok{, }\DecValTok{200}\NormalTok{, }\DecValTok{150}\NormalTok{],}
    \StringTok{\textquotesingle{}quantity\textquotesingle{}}\NormalTok{: [}\DecValTok{3}\NormalTok{, }\DecValTok{7}\NormalTok{, }\DecValTok{5}\NormalTok{]}
\NormalTok{\})}

\CommentTok{\# Calculate price per unit}
\NormalTok{result }\OperatorTok{=}\NormalTok{ add.transform(}
    \StringTok{\textquotesingle{}@calc\textquotesingle{}}\NormalTok{,}
\NormalTok{    df,}
\NormalTok{    columns}\OperatorTok{=}\NormalTok{[}\StringTok{\textquotesingle{}price\textquotesingle{}}\NormalTok{, }\StringTok{\textquotesingle{}quantity\textquotesingle{}}\NormalTok{],}
\NormalTok{    expression}\OperatorTok{=}\StringTok{\textquotesingle{}price / quantity\textquotesingle{}}\NormalTok{,}
\NormalTok{    as\_}\OperatorTok{=}\StringTok{\textquotesingle{}price\_per\_unit\textquotesingle{}}
\NormalTok{)}

\CommentTok{\# Positional parameters (also works without naming certain parameters):}
\CommentTok{\# result = add.transform(\textquotesingle{}@calc\textquotesingle{}, df, [\textquotesingle{}price\textquotesingle{}, \textquotesingle{}quantity\textquotesingle{}], }
\CommentTok{\#                        expression=\textquotesingle{}price / quantity\textquotesingle{}, as\_=\textquotesingle{}price\_per\_unit\textquotesingle{})}

\BuiltInTok{print}\NormalTok{(result)}
\end{Highlighting}
\end{Shaded}

\textbf{Output:}

\begin{verbatim}
      product  price  quantity  price_per_unit
0      Widget    100         3       33.333333
1      Gadget    200         7       28.571429
2  Doohickey    150         5       30.000000
\end{verbatim}

\textbf{Explanation:} - Division operations work just like
multiplication - Results are floating-point numbers - Widget: 100 / 3 =
33.33\ldots{} - Gadget: 200 / 7 = 28.57\ldots{} - Doohickey: 150 / 5 =
30.00

\textbf{Note:} This also works with polars DataFrames.

\begin{center}\rule{0.5\linewidth}{0.5pt}\end{center}

\section{Supported Operations}\label{supported-operations}

\subsection{Arithmetic Operators}\label{arithmetic-operators}

\begin{longtable}[]{@{}lll@{}}
\toprule\noalign{}
Operator & Description & Example \\
\midrule\noalign{}
\endhead
\bottomrule\noalign{}
\endlastfoot
\texttt{+} & Addition & \texttt{price\ +\ tax} \\
\texttt{-} & Subtraction & \texttt{price\ -\ discount} \\
\texttt{*} & Multiplication & \texttt{price\ *\ quantity} \\
\texttt{/} & Division & \texttt{total\ /\ count} \\
\end{longtable}

\subsection{Expression Syntax}\label{expression-syntax}

\begin{Shaded}
\begin{Highlighting}[]
\CommentTok{\# Single column reference}
\NormalTok{expression}\OperatorTok{=}\StringTok{\textquotesingle{}price\textquotesingle{}}

\CommentTok{\# Binary operation}
\NormalTok{expression}\OperatorTok{=}\StringTok{\textquotesingle{}price * quantity\textquotesingle{}}

\CommentTok{\# Multiple operations (left to right)}
\NormalTok{expression}\OperatorTok{=}\StringTok{\textquotesingle{}price * quantity + tax\textquotesingle{}}

\CommentTok{\# Using multiple columns}
\NormalTok{expression}\OperatorTok{=}\StringTok{\textquotesingle{}a + b {-} c\textquotesingle{}}
\end{Highlighting}
\end{Shaded}

\subsection{Important Notes}\label{important-notes}

\begin{itemize}
\tightlist
\item
  Expressions are evaluated left to right
\item
  Column names must match exactly (case-sensitive)
\item
  All columns in the expression must be listed in \texttt{columns}
  parameter
\item
  Currently, parentheses for grouping are not supported
\item
  For complex calculations, break them into multiple steps
\end{itemize}

\begin{center}\rule{0.5\linewidth}{0.5pt}\end{center}

\section{Common Patterns}\label{common-patterns-1}

\subsection{Pattern 1: Calculate Total}\label{pattern-1-calculate-total}

\begin{Shaded}
\begin{Highlighting}[]
\NormalTok{result }\OperatorTok{=}\NormalTok{ add.transform(}\StringTok{\textquotesingle{}@calc\textquotesingle{}}\NormalTok{, df, [}\StringTok{\textquotesingle{}price\textquotesingle{}}\NormalTok{, }\StringTok{\textquotesingle{}qty\textquotesingle{}}\NormalTok{], }
\NormalTok{                       expression}\OperatorTok{=}\StringTok{\textquotesingle{}price * qty\textquotesingle{}}\NormalTok{, as\_}\OperatorTok{=}\StringTok{\textquotesingle{}total\textquotesingle{}}\NormalTok{)}
\end{Highlighting}
\end{Shaded}

\subsection{Pattern 2: Calculate
Difference}\label{pattern-2-calculate-difference}

\begin{Shaded}
\begin{Highlighting}[]
\NormalTok{result }\OperatorTok{=}\NormalTok{ add.transform(}\StringTok{\textquotesingle{}@calc\textquotesingle{}}\NormalTok{, df, [}\StringTok{\textquotesingle{}actual\textquotesingle{}}\NormalTok{, }\StringTok{\textquotesingle{}target\textquotesingle{}}\NormalTok{], }
\NormalTok{                       expression}\OperatorTok{=}\StringTok{\textquotesingle{}actual {-} target\textquotesingle{}}\NormalTok{, as\_}\OperatorTok{=}\StringTok{\textquotesingle{}variance\textquotesingle{}}\NormalTok{)}
\end{Highlighting}
\end{Shaded}

\subsection{Pattern 3: Multiple
Metrics}\label{pattern-3-multiple-metrics}

\begin{Shaded}
\begin{Highlighting}[]
\NormalTok{result }\OperatorTok{=}\NormalTok{ add.transform(}\StringTok{\textquotesingle{}@calc\textquotesingle{}}\NormalTok{, df, [}\StringTok{\textquotesingle{}revenue\textquotesingle{}}\NormalTok{, }\StringTok{\textquotesingle{}cost\textquotesingle{}}\NormalTok{], }
\NormalTok{                       expression}\OperatorTok{=}\NormalTok{[}\StringTok{\textquotesingle{}revenue {-} cost\textquotesingle{}}\NormalTok{, }\StringTok{\textquotesingle{}revenue / cost\textquotesingle{}}\NormalTok{],}
\NormalTok{                       as\_}\OperatorTok{=}\NormalTok{[}\StringTok{\textquotesingle{}profit\textquotesingle{}}\NormalTok{, }\StringTok{\textquotesingle{}margin\_ratio\textquotesingle{}}\NormalTok{])}
\end{Highlighting}
\end{Shaded}

\subsection{Pattern 4: Chain
Calculations}\label{pattern-4-chain-calculations}

\begin{Shaded}
\begin{Highlighting}[]
\CommentTok{\# Step 1: Calculate profit}
\NormalTok{df }\OperatorTok{=}\NormalTok{ add.transform(}\StringTok{\textquotesingle{}@calc\textquotesingle{}}\NormalTok{, df, [}\StringTok{\textquotesingle{}price\textquotesingle{}}\NormalTok{, }\StringTok{\textquotesingle{}cost\textquotesingle{}}\NormalTok{], }
\NormalTok{                   expression}\OperatorTok{=}\StringTok{\textquotesingle{}price {-} cost\textquotesingle{}}\NormalTok{, as\_}\OperatorTok{=}\StringTok{\textquotesingle{}profit\textquotesingle{}}\NormalTok{)}

\CommentTok{\# Step 2: Calculate profit margin using the new column}
\NormalTok{df }\OperatorTok{=}\NormalTok{ add.transform(}\StringTok{\textquotesingle{}@calc\textquotesingle{}}\NormalTok{, df, [}\StringTok{\textquotesingle{}profit\textquotesingle{}}\NormalTok{, }\StringTok{\textquotesingle{}price\textquotesingle{}}\NormalTok{], }
\NormalTok{                   expression}\OperatorTok{=}\StringTok{\textquotesingle{}profit / price\textquotesingle{}}\NormalTok{, as\_}\OperatorTok{=}\StringTok{\textquotesingle{}margin\textquotesingle{}}\NormalTok{)}
\end{Highlighting}
\end{Shaded}

\begin{center}\rule{0.5\linewidth}{0.5pt}\end{center}

\section{Best Practices}\label{best-practices}

\begin{enumerate}
\def\labelenumi{\arabic{enumi}.}
\item
  \textbf{List all columns used}: Always include all columns referenced
  in your expression in the \texttt{columns} parameter
\item
  \textbf{Use descriptive names}: Choose clear names for calculated
  columns

\begin{Shaded}
\begin{Highlighting}[]
\CommentTok{\# Good}
\NormalTok{as\_}\OperatorTok{=}\StringTok{\textquotesingle{}total\_revenue\textquotesingle{}}

\CommentTok{\# Avoid}
\NormalTok{as\_}\OperatorTok{=}\StringTok{\textquotesingle{}col1\textquotesingle{}}
\end{Highlighting}
\end{Shaded}
\item
  \textbf{Break complex calculations}: For readability, split complex
  calculations into steps

\begin{Shaded}
\begin{Highlighting}[]
\CommentTok{\# Instead of one complex expression}
\CommentTok{\# Do this:}
\NormalTok{df }\OperatorTok{=}\NormalTok{ add.transform(}\StringTok{\textquotesingle{}@calc\textquotesingle{}}\NormalTok{, df, [}\StringTok{\textquotesingle{}a\textquotesingle{}}\NormalTok{, }\StringTok{\textquotesingle{}b\textquotesingle{}}\NormalTok{], expression}\OperatorTok{=}\StringTok{\textquotesingle{}a + b\textquotesingle{}}\NormalTok{, as\_}\OperatorTok{=}\StringTok{\textquotesingle{}sum\textquotesingle{}}\NormalTok{)}
\NormalTok{df }\OperatorTok{=}\NormalTok{ add.transform(}\StringTok{\textquotesingle{}@calc\textquotesingle{}}\NormalTok{, df, [}\StringTok{\textquotesingle{}sum\textquotesingle{}}\NormalTok{, }\StringTok{\textquotesingle{}c\textquotesingle{}}\NormalTok{], expression}\OperatorTok{=}\StringTok{\textquotesingle{}sum * c\textquotesingle{}}\NormalTok{, as\_}\OperatorTok{=}\StringTok{\textquotesingle{}result\textquotesingle{}}\NormalTok{)}
\end{Highlighting}
\end{Shaded}
\item
  \textbf{Check for division by zero}: Validate your data before
  division operations

\begin{Shaded}
\begin{Highlighting}[]
\CommentTok{\# Check for zeros first}
\ControlFlowTok{if}\NormalTok{ (df[}\StringTok{\textquotesingle{}quantity\textquotesingle{}}\NormalTok{] }\OperatorTok{==} \DecValTok{0}\NormalTok{).}\BuiltInTok{any}\NormalTok{():}
    \BuiltInTok{print}\NormalTok{(}\StringTok{"Warning: Zero quantities found"}\NormalTok{)}
\end{Highlighting}
\end{Shaded}
\item
  \textbf{Match expression and as\_ lengths}: When using lists, ensure
  they have the same length

\begin{Shaded}
\begin{Highlighting}[]
\CommentTok{\# Correct}
\NormalTok{expression}\OperatorTok{=}\NormalTok{[}\StringTok{\textquotesingle{}a + b\textquotesingle{}}\NormalTok{, }\StringTok{\textquotesingle{}a * b\textquotesingle{}}\NormalTok{]}
\NormalTok{as\_}\OperatorTok{=}\NormalTok{[}\StringTok{\textquotesingle{}sum\textquotesingle{}}\NormalTok{, }\StringTok{\textquotesingle{}product\textquotesingle{}}\NormalTok{]}

\CommentTok{\# Wrong {-} will error}
\NormalTok{expression}\OperatorTok{=}\NormalTok{[}\StringTok{\textquotesingle{}a + b\textquotesingle{}}\NormalTok{, }\StringTok{\textquotesingle{}a * b\textquotesingle{}}\NormalTok{]}
\NormalTok{as\_}\OperatorTok{=}\NormalTok{[}\StringTok{\textquotesingle{}sum\textquotesingle{}}\NormalTok{]  }\CommentTok{\# Missing second name}
\end{Highlighting}
\end{Shaded}
\end{enumerate}

\begin{center}\rule{0.5\linewidth}{0.5pt}\end{center}

\section{Key Takeaways}\label{key-takeaways-5}

\begin{itemize}
\tightlist
\item
  Use \texttt{@calc} mode to create calculated columns
\item
  Specify source columns in \texttt{columns} parameter
\item
  Define calculations in \texttt{expression} parameter
\item
  Name new columns with \texttt{as\_} parameter
\item
  Use lists for multiple calculations
\item
  Supports basic arithmetic: +, -, *, /
\item
  Original columns are preserved
\item
  Works with both pandas and polars
\end{itemize}

\begin{center}\rule{0.5\linewidth}{0.5pt}\end{center}

\section{Common Questions}\label{common-questions-5}

\textbf{Q: Can I use parentheses in expressions?}\\
A: Currently, parentheses are not supported. Break complex calculations
into multiple steps instead.

\textbf{Q: Can I reference a calculated column in the same operation?}\\
A: No, you need to chain operations. Calculate the first column, then
use it in a second \texttt{add.transform()} call.

\textbf{Q: What happens if a column name doesn't exist?}\\
A: You'll get an error. Make sure all column names in your expression
exist in the DataFrame.

\textbf{Q: Can I use functions like sqrt() or abs()?}\\
A: Currently, only basic arithmetic operators are supported. Use
pandas/polars functions for advanced operations.

\textbf{Q: How do I handle null/NaN values?}\\
A: Calculations with NaN values will result in NaN. Filter or fill NaN
values before calculating if needed.

\begin{center}\rule{0.5\linewidth}{0.5pt}\end{center}

\section{Next Steps}\label{next-steps-7}

\begin{itemize}
\tightlist
\item
  \textbf{\href{examples-transform-02-filter-sort.html}{Filter \& Sort}}
  - Learn to filter and sort data
\item
  \textbf{\href{examples-to-01-basic.html}{add.to() Examples}} - Review
  data enrichment operations
\item
  \textbf{\href{reference-transform.html}{API Reference}} - Complete
  \texttt{add.transform()} documentation
\end{itemize}

\begin{center}\rule{0.5\linewidth}{0.5pt}\end{center}

\textbf{Version:} 0.1.3\\
\textbf{Last Updated:} March 9, 2026

\chapter{add.transform() - Filter \&
Sort}\label{add.transform---filter-sort}

\section{Overview}\label{overview-6}

Learn how to filter rows and sort data using \texttt{add.transform()}.
These are essential operations for data cleaning and analysis.

\textbf{What you'll learn:} - How to filter rows based on conditions -
How to sort data in ascending or descending order - How to combine
filtering and sorting - Best practices for data selection

\textbf{Prerequisites:} -
\href{examples-transform-01-calc.html}{Calculations} - Understanding of
basic transform operations - Basic understanding of comparison operators

\begin{center}\rule{0.5\linewidth}{0.5pt}\end{center}

\section{Example 1: Filter Rows}\label{example-1-filter-rows}

\textbf{Business Context:} You have a product inventory and need to see
only products that are currently in stock (stock \textgreater{} 0).

\textbf{Code:}

\begin{Shaded}
\begin{Highlighting}[]
\ImportTok{import}\NormalTok{ additory }\ImportTok{as}\NormalTok{ add}
\ImportTok{import}\NormalTok{ pandas }\ImportTok{as}\NormalTok{ pd}

\CommentTok{\# Product inventory}
\NormalTok{df }\OperatorTok{=}\NormalTok{ pd.DataFrame(\{}
    \StringTok{\textquotesingle{}product\textquotesingle{}}\NormalTok{: [}\StringTok{\textquotesingle{}Widget\textquotesingle{}}\NormalTok{, }\StringTok{\textquotesingle{}Gadget\textquotesingle{}}\NormalTok{, }\StringTok{\textquotesingle{}Doohickey\textquotesingle{}}\NormalTok{, }\StringTok{\textquotesingle{}Thingamajig\textquotesingle{}}\NormalTok{],}
    \StringTok{\textquotesingle{}price\textquotesingle{}}\NormalTok{: [}\FloatTok{29.99}\NormalTok{, }\FloatTok{49.99}\NormalTok{, }\FloatTok{19.99}\NormalTok{, }\FloatTok{39.99}\NormalTok{],}
    \StringTok{\textquotesingle{}stock\textquotesingle{}}\NormalTok{: [}\DecValTok{10}\NormalTok{, }\DecValTok{5}\NormalTok{, }\DecValTok{0}\NormalTok{, }\DecValTok{15}\NormalTok{]}
\NormalTok{\})}

\CommentTok{\# Filter products in stock}
\NormalTok{result }\OperatorTok{=}\NormalTok{ add.transform(}
    \StringTok{\textquotesingle{}@filter\textquotesingle{}}\NormalTok{,}
\NormalTok{    df,}
\NormalTok{    columns}\OperatorTok{=}\NormalTok{[}\StringTok{\textquotesingle{}product\textquotesingle{}}\NormalTok{, }\StringTok{\textquotesingle{}price\textquotesingle{}}\NormalTok{, }\StringTok{\textquotesingle{}stock\textquotesingle{}}\NormalTok{],}
\NormalTok{    where}\OperatorTok{=}\StringTok{\textquotesingle{}stock \textgreater{} 0\textquotesingle{}}
\NormalTok{)}

\CommentTok{\# Positional parameters (also works without naming certain parameters):}
\CommentTok{\# result = add.transform(\textquotesingle{}@filter\textquotesingle{}, df, [\textquotesingle{}product\textquotesingle{}, \textquotesingle{}price\textquotesingle{}, \textquotesingle{}stock\textquotesingle{}], where=\textquotesingle{}stock \textgreater{} 0\textquotesingle{})}

\BuiltInTok{print}\NormalTok{(result)}
\end{Highlighting}
\end{Shaded}

\textbf{Output:}

\begin{verbatim}
       product  price  stock
0       Widget  29.99     10
1       Gadget  49.99      5
2  Thingamajig  39.99     15
\end{verbatim}

\textbf{Explanation:} -
\texttt{\textquotesingle{}@filter\textquotesingle{}} mode removes rows
that don't match the condition - \texttt{columns} specifies which
columns to keep in the result - \texttt{where} defines the filter
condition - Doohickey (stock = 0) is filtered out - Only 3 products
remain (stock \textgreater{} 0)

\textbf{Note:} This also works with polars DataFrames.

\begin{center}\rule{0.5\linewidth}{0.5pt}\end{center}

\section{Example 2: Sort Rows}\label{example-2-sort-rows}

\textbf{Business Context:} You want to see your products ranked by sales
performance, with the best-selling products first.

\textbf{Code:}

\begin{Shaded}
\begin{Highlighting}[]
\ImportTok{import}\NormalTok{ additory }\ImportTok{as}\NormalTok{ add}
\ImportTok{import}\NormalTok{ pandas }\ImportTok{as}\NormalTok{ pd}

\CommentTok{\# Product sales data}
\NormalTok{df }\OperatorTok{=}\NormalTok{ pd.DataFrame(\{}
    \StringTok{\textquotesingle{}product\textquotesingle{}}\NormalTok{: [}\StringTok{\textquotesingle{}Widget\textquotesingle{}}\NormalTok{, }\StringTok{\textquotesingle{}Gadget\textquotesingle{}}\NormalTok{, }\StringTok{\textquotesingle{}Doohickey\textquotesingle{}}\NormalTok{],}
    \StringTok{\textquotesingle{}price\textquotesingle{}}\NormalTok{: [}\FloatTok{29.99}\NormalTok{, }\FloatTok{49.99}\NormalTok{, }\FloatTok{19.99}\NormalTok{],}
    \StringTok{\textquotesingle{}sales\textquotesingle{}}\NormalTok{: [}\DecValTok{100}\NormalTok{, }\DecValTok{150}\NormalTok{, }\DecValTok{75}\NormalTok{]}
\NormalTok{\})}

\CommentTok{\# Sort by sales descending (highest first)}
\NormalTok{result }\OperatorTok{=}\NormalTok{ add.transform(}
    \StringTok{\textquotesingle{}@sort\textquotesingle{}}\NormalTok{,}
\NormalTok{    df,}
\NormalTok{    columns}\OperatorTok{=}\NormalTok{[}\StringTok{\textquotesingle{}product\textquotesingle{}}\NormalTok{, }\StringTok{\textquotesingle{}price\textquotesingle{}}\NormalTok{, }\StringTok{\textquotesingle{}sales\textquotesingle{}}\NormalTok{],}
\NormalTok{    by}\OperatorTok{=}\StringTok{\textquotesingle{}sales\textquotesingle{}}\NormalTok{,}
\NormalTok{    order}\OperatorTok{=}\StringTok{\textquotesingle{}desc\textquotesingle{}}
\NormalTok{)}

\CommentTok{\# Positional parameters (also works without naming certain parameters):}
\CommentTok{\# result = add.transform(\textquotesingle{}@sort\textquotesingle{}, df, [\textquotesingle{}product\textquotesingle{}, \textquotesingle{}price\textquotesingle{}, \textquotesingle{}sales\textquotesingle{}], by=\textquotesingle{}sales\textquotesingle{}, order=\textquotesingle{}desc\textquotesingle{})}

\BuiltInTok{print}\NormalTok{(result)}
\end{Highlighting}
\end{Shaded}

\textbf{Output:}

\begin{verbatim}
      product  price  sales
0      Gadget  49.99    150
1      Widget  29.99    100
2  Doohickey  19.99     75
\end{verbatim}

\textbf{Explanation:} -
\texttt{\textquotesingle{}@sort\textquotesingle{}} mode reorders rows -
\texttt{columns} specifies which columns to include - \texttt{by}
specifies the column to sort by -
\texttt{order=\textquotesingle{}desc\textquotesingle{}} sorts from
highest to lowest - Gadget (150 sales) is now first - Use
\texttt{order=\textquotesingle{}asc\textquotesingle{}} for ascending
(lowest to highest)

\textbf{Note:} This also works with polars DataFrames.

\begin{center}\rule{0.5\linewidth}{0.5pt}\end{center}

\section{Example 3: Filter and Sort
Combined}\label{example-3-filter-and-sort-combined}

\textbf{Business Context:} You want to see in-stock products sorted by
price, from cheapest to most expensive.

\textbf{Code:}

\begin{Shaded}
\begin{Highlighting}[]
\ImportTok{import}\NormalTok{ additory }\ImportTok{as}\NormalTok{ add}
\ImportTok{import}\NormalTok{ pandas }\ImportTok{as}\NormalTok{ pd}

\CommentTok{\# Product inventory}
\NormalTok{df }\OperatorTok{=}\NormalTok{ pd.DataFrame(\{}
    \StringTok{\textquotesingle{}product\textquotesingle{}}\NormalTok{: [}\StringTok{\textquotesingle{}Widget\textquotesingle{}}\NormalTok{, }\StringTok{\textquotesingle{}Gadget\textquotesingle{}}\NormalTok{, }\StringTok{\textquotesingle{}Doohickey\textquotesingle{}}\NormalTok{, }\StringTok{\textquotesingle{}Thingamajig\textquotesingle{}}\NormalTok{],}
    \StringTok{\textquotesingle{}price\textquotesingle{}}\NormalTok{: [}\FloatTok{29.99}\NormalTok{, }\FloatTok{49.99}\NormalTok{, }\FloatTok{19.99}\NormalTok{, }\FloatTok{39.99}\NormalTok{],}
    \StringTok{\textquotesingle{}stock\textquotesingle{}}\NormalTok{: [}\DecValTok{10}\NormalTok{, }\DecValTok{5}\NormalTok{, }\DecValTok{0}\NormalTok{, }\DecValTok{15}\NormalTok{]}
\NormalTok{\})}

\CommentTok{\# Step 1: Filter in{-}stock products}
\NormalTok{result }\OperatorTok{=}\NormalTok{ add.transform(}
    \StringTok{\textquotesingle{}@filter\textquotesingle{}}\NormalTok{,}
\NormalTok{    df,}
\NormalTok{    columns}\OperatorTok{=}\NormalTok{[}\StringTok{\textquotesingle{}product\textquotesingle{}}\NormalTok{, }\StringTok{\textquotesingle{}price\textquotesingle{}}\NormalTok{, }\StringTok{\textquotesingle{}stock\textquotesingle{}}\NormalTok{],}
\NormalTok{    where}\OperatorTok{=}\StringTok{\textquotesingle{}stock \textgreater{} 0\textquotesingle{}}
\NormalTok{)}

\CommentTok{\# Positional parameters (also works without naming certain parameters):}
\CommentTok{\# result = add.transform(\textquotesingle{}@filter\textquotesingle{}, df, [\textquotesingle{}product\textquotesingle{}, \textquotesingle{}price\textquotesingle{}, \textquotesingle{}stock\textquotesingle{}], where=\textquotesingle{}stock \textgreater{} 0\textquotesingle{})}

\BuiltInTok{print}\NormalTok{(}\StringTok{"After filtering:"}\NormalTok{)}
\BuiltInTok{print}\NormalTok{(result)}

\CommentTok{\# Step 2: Sort by price ascending}
\NormalTok{result }\OperatorTok{=}\NormalTok{ add.transform(}
    \StringTok{\textquotesingle{}@sort\textquotesingle{}}\NormalTok{,}
\NormalTok{    result,}
\NormalTok{    columns}\OperatorTok{=}\NormalTok{[}\StringTok{\textquotesingle{}product\textquotesingle{}}\NormalTok{, }\StringTok{\textquotesingle{}price\textquotesingle{}}\NormalTok{, }\StringTok{\textquotesingle{}stock\textquotesingle{}}\NormalTok{],}
\NormalTok{    by}\OperatorTok{=}\StringTok{\textquotesingle{}price\textquotesingle{}}\NormalTok{,}
\NormalTok{    order}\OperatorTok{=}\StringTok{\textquotesingle{}asc\textquotesingle{}}
\NormalTok{)}

\CommentTok{\# Positional parameters (also works without naming certain parameters):}
\CommentTok{\# result = add.transform(\textquotesingle{}@sort\textquotesingle{}, result, [\textquotesingle{}product\textquotesingle{}, \textquotesingle{}price\textquotesingle{}, \textquotesingle{}stock\textquotesingle{}], by=\textquotesingle{}price\textquotesingle{}, order=\textquotesingle{}asc\textquotesingle{})}

\BuiltInTok{print}\NormalTok{(}\StringTok{"}\CharTok{\textbackslash{}n}\StringTok{After sorting:"}\NormalTok{)}
\BuiltInTok{print}\NormalTok{(result)}
\end{Highlighting}
\end{Shaded}

\textbf{Output:}

\begin{verbatim}
After filtering:
       product  price  stock
0       Widget  29.99     10
1       Gadget  49.99      5
2  Thingamajig  39.99     15

After sorting:
       product  price  stock
0       Widget  29.99     10
1  Thingamajig  39.99     15
2       Gadget  49.99      5
\end{verbatim}

\textbf{Explanation:} - First operation filters out Doohickey (stock =
0) - Second operation sorts the remaining products by price - Widget
(29.99) is cheapest, so it's first - Gadget (49.99) is most expensive,
so it's last - Chain operations by using the result of one as input to
the next

\textbf{Note:} This also works with polars DataFrames.

\begin{center}\rule{0.5\linewidth}{0.5pt}\end{center}

\section{Filter Conditions}\label{filter-conditions}

\subsection{Comparison Operators}\label{comparison-operators}

\begin{longtable}[]{@{}lll@{}}
\toprule\noalign{}
Operator & Description & Example \\
\midrule\noalign{}
\endhead
\bottomrule\noalign{}
\endlastfoot
\texttt{\textgreater{}} & Greater than &
\texttt{where=\textquotesingle{}price\ \textgreater{}\ 100\textquotesingle{}} \\
\texttt{\textless{}} & Less than &
\texttt{where=\textquotesingle{}stock\ \textless{}\ 10\textquotesingle{}} \\
\texttt{\textgreater{}=} & Greater than or equal &
\texttt{where=\textquotesingle{}age\ \textgreater{}=\ 18\textquotesingle{}} \\
\texttt{\textless{}=} & Less than or equal &
\texttt{where=\textquotesingle{}quantity\ \textless{}=\ 5\textquotesingle{}} \\
\texttt{==} & Equal to &
\texttt{where=\textquotesingle{}status\ ==\ 1\textquotesingle{}} \\
\texttt{!=} & Not equal to &
\texttt{where=\textquotesingle{}category\ !=\ 0\textquotesingle{}} \\
\end{longtable}

\subsection{Filter Examples}\label{filter-examples}

\begin{Shaded}
\begin{Highlighting}[]
\CommentTok{\# Greater than}
\NormalTok{where}\OperatorTok{=}\StringTok{\textquotesingle{}price \textgreater{} 50\textquotesingle{}}

\CommentTok{\# Less than or equal}
\NormalTok{where}\OperatorTok{=}\StringTok{\textquotesingle{}stock \textless{}= 10\textquotesingle{}}

\CommentTok{\# Equal to}
\NormalTok{where}\OperatorTok{=}\StringTok{\textquotesingle{}category == 1\textquotesingle{}}

\CommentTok{\# Not equal}
\NormalTok{where}\OperatorTok{=}\StringTok{\textquotesingle{}status != 0\textquotesingle{}}
\end{Highlighting}
\end{Shaded}

\subsection{Important Notes}\label{important-notes-1}

\begin{itemize}
\tightlist
\item
  Currently, only simple comparisons are supported
\item
  Complex conditions (AND, OR) are not yet supported
\item
  For multiple conditions, chain multiple (\textbf{filter?}) operations
\item
  Column names must match exactly (case-sensitive)
\end{itemize}

\begin{center}\rule{0.5\linewidth}{0.5pt}\end{center}

\section{Sort Options}\label{sort-options}

\subsection{Sort Order}\label{sort-order}

\begin{Shaded}
\begin{Highlighting}[]
\CommentTok{\# Ascending (lowest to highest)}
\NormalTok{order}\OperatorTok{=}\StringTok{\textquotesingle{}asc\textquotesingle{}}

\CommentTok{\# Descending (highest to lowest)}
\NormalTok{order}\OperatorTok{=}\StringTok{\textquotesingle{}desc\textquotesingle{}}
\end{Highlighting}
\end{Shaded}

\subsection{Sort by Multiple Columns}\label{sort-by-multiple-columns}

\begin{Shaded}
\begin{Highlighting}[]
\CommentTok{\# Sort by category, then by price within each category}
\NormalTok{result }\OperatorTok{=}\NormalTok{ add.transform(}\StringTok{\textquotesingle{}@sort\textquotesingle{}}\NormalTok{, df, }
\NormalTok{                       columns}\OperatorTok{=}\NormalTok{[}\StringTok{\textquotesingle{}category\textquotesingle{}}\NormalTok{, }\StringTok{\textquotesingle{}price\textquotesingle{}}\NormalTok{, }\StringTok{\textquotesingle{}name\textquotesingle{}}\NormalTok{],}
\NormalTok{                       by}\OperatorTok{=}\NormalTok{[}\StringTok{\textquotesingle{}category\textquotesingle{}}\NormalTok{, }\StringTok{\textquotesingle{}price\textquotesingle{}}\NormalTok{],}
\NormalTok{                       order}\OperatorTok{=}\StringTok{\textquotesingle{}asc\textquotesingle{}}\NormalTok{)}
\end{Highlighting}
\end{Shaded}

\textbf{Note:} When sorting by multiple columns, all columns use the
same order (asc or desc).

\begin{center}\rule{0.5\linewidth}{0.5pt}\end{center}

\section{Common Patterns}\label{common-patterns-2}

\subsection{Pattern 1: Filter Out
Zeros}\label{pattern-1-filter-out-zeros}

\begin{Shaded}
\begin{Highlighting}[]
\NormalTok{result }\OperatorTok{=}\NormalTok{ add.transform(}\StringTok{\textquotesingle{}@filter\textquotesingle{}}\NormalTok{, df, columns}\OperatorTok{=}\NormalTok{[}\StringTok{\textquotesingle{}col1\textquotesingle{}}\NormalTok{, }\StringTok{\textquotesingle{}col2\textquotesingle{}}\NormalTok{], }
\NormalTok{                       where}\OperatorTok{=}\StringTok{\textquotesingle{}value \textgreater{} 0\textquotesingle{}}\NormalTok{)}
\end{Highlighting}
\end{Shaded}

\subsection{Pattern 2: Top N by
Sorting}\label{pattern-2-top-n-by-sorting}

\begin{Shaded}
\begin{Highlighting}[]
\CommentTok{\# Sort descending and take top rows with pandas}
\NormalTok{result }\OperatorTok{=}\NormalTok{ add.transform(}\StringTok{\textquotesingle{}@sort\textquotesingle{}}\NormalTok{, df, columns}\OperatorTok{=}\NormalTok{[}\StringTok{\textquotesingle{}name\textquotesingle{}}\NormalTok{, }\StringTok{\textquotesingle{}score\textquotesingle{}}\NormalTok{], }
\NormalTok{                       by}\OperatorTok{=}\StringTok{\textquotesingle{}score\textquotesingle{}}\NormalTok{, order}\OperatorTok{=}\StringTok{\textquotesingle{}desc\textquotesingle{}}\NormalTok{)}
\NormalTok{top\_10 }\OperatorTok{=}\NormalTok{ result.head(}\DecValTok{10}\NormalTok{)}
\end{Highlighting}
\end{Shaded}

\subsection{Pattern 3: Filter Then
Sort}\label{pattern-3-filter-then-sort}

\begin{Shaded}
\begin{Highlighting}[]
\CommentTok{\# Step 1: Filter}
\NormalTok{df }\OperatorTok{=}\NormalTok{ add.transform(}\StringTok{\textquotesingle{}@filter\textquotesingle{}}\NormalTok{, df, columns}\OperatorTok{=}\NormalTok{[}\StringTok{\textquotesingle{}a\textquotesingle{}}\NormalTok{, }\StringTok{\textquotesingle{}b\textquotesingle{}}\NormalTok{], where}\OperatorTok{=}\StringTok{\textquotesingle{}a \textgreater{} 0\textquotesingle{}}\NormalTok{)}
\CommentTok{\# Step 2: Sort}
\NormalTok{df }\OperatorTok{=}\NormalTok{ add.transform(}\StringTok{\textquotesingle{}@sort\textquotesingle{}}\NormalTok{, df, columns}\OperatorTok{=}\NormalTok{[}\StringTok{\textquotesingle{}a\textquotesingle{}}\NormalTok{, }\StringTok{\textquotesingle{}b\textquotesingle{}}\NormalTok{], by}\OperatorTok{=}\StringTok{\textquotesingle{}b\textquotesingle{}}\NormalTok{, order}\OperatorTok{=}\StringTok{\textquotesingle{}desc\textquotesingle{}}\NormalTok{)}
\end{Highlighting}
\end{Shaded}

\subsection{Pattern 4: Sort Then
Filter}\label{pattern-4-sort-then-filter}

\begin{Shaded}
\begin{Highlighting}[]
\CommentTok{\# Step 1: Sort}
\NormalTok{df }\OperatorTok{=}\NormalTok{ add.transform(}\StringTok{\textquotesingle{}@sort\textquotesingle{}}\NormalTok{, df, columns}\OperatorTok{=}\NormalTok{[}\StringTok{\textquotesingle{}date\textquotesingle{}}\NormalTok{, }\StringTok{\textquotesingle{}amount\textquotesingle{}}\NormalTok{], }
\NormalTok{                   by}\OperatorTok{=}\StringTok{\textquotesingle{}date\textquotesingle{}}\NormalTok{, order}\OperatorTok{=}\StringTok{\textquotesingle{}desc\textquotesingle{}}\NormalTok{)}
\CommentTok{\# Step 2: Filter recent}
\NormalTok{df }\OperatorTok{=}\NormalTok{ add.transform(}\StringTok{\textquotesingle{}@filter\textquotesingle{}}\NormalTok{, df, columns}\OperatorTok{=}\NormalTok{[}\StringTok{\textquotesingle{}date\textquotesingle{}}\NormalTok{, }\StringTok{\textquotesingle{}amount\textquotesingle{}}\NormalTok{], }
\NormalTok{                   where}\OperatorTok{=}\StringTok{\textquotesingle{}date \textgreater{} 20240101\textquotesingle{}}\NormalTok{)}
\end{Highlighting}
\end{Shaded}

\begin{center}\rule{0.5\linewidth}{0.5pt}\end{center}

\section{Best Practices}\label{best-practices-1}

\begin{enumerate}
\def\labelenumi{\arabic{enumi}.}
\item
  \textbf{Specify columns explicitly}: Always list the columns you want
  to keep

\begin{Shaded}
\begin{Highlighting}[]
\CommentTok{\# Good}
\NormalTok{columns}\OperatorTok{=}\NormalTok{[}\StringTok{\textquotesingle{}product\textquotesingle{}}\NormalTok{, }\StringTok{\textquotesingle{}price\textquotesingle{}}\NormalTok{, }\StringTok{\textquotesingle{}stock\textquotesingle{}}\NormalTok{]}

\CommentTok{\# Avoid relying on defaults}
\end{Highlighting}
\end{Shaded}
\item
  \textbf{Filter before sorting}: More efficient to sort fewer rows

\begin{Shaded}
\begin{Highlighting}[]
\CommentTok{\# Good: Filter first (reduces data), then sort}
\NormalTok{df }\OperatorTok{=}\NormalTok{ add.transform(}\StringTok{\textquotesingle{}@filter\textquotesingle{}}\NormalTok{, df, columns}\OperatorTok{=}\NormalTok{cols, where}\OperatorTok{=}\StringTok{\textquotesingle{}active == 1\textquotesingle{}}\NormalTok{)}
\NormalTok{df }\OperatorTok{=}\NormalTok{ add.transform(}\StringTok{\textquotesingle{}@sort\textquotesingle{}}\NormalTok{, df, columns}\OperatorTok{=}\NormalTok{cols, by}\OperatorTok{=}\StringTok{\textquotesingle{}date\textquotesingle{}}\NormalTok{, order}\OperatorTok{=}\StringTok{\textquotesingle{}desc\textquotesingle{}}\NormalTok{)}
\end{Highlighting}
\end{Shaded}
\item
  \textbf{Use meaningful conditions}: Make your filters clear and
  maintainable

\begin{Shaded}
\begin{Highlighting}[]
\CommentTok{\# Good}
\NormalTok{where}\OperatorTok{=}\StringTok{\textquotesingle{}stock \textgreater{} 0\textquotesingle{}}  \CommentTok{\# Clear intent}

\CommentTok{\# Less clear}
\NormalTok{where}\OperatorTok{=}\StringTok{\textquotesingle{}stock != 0\textquotesingle{}}  \CommentTok{\# Could include negative values}
\end{Highlighting}
\end{Shaded}
\item
  \textbf{Check your data types}: Ensure columns used in conditions are
  numeric

\begin{Shaded}
\begin{Highlighting}[]
\CommentTok{\# If \textquotesingle{}price\textquotesingle{} is stored as string, convert first}
\NormalTok{df[}\StringTok{\textquotesingle{}price\textquotesingle{}}\NormalTok{] }\OperatorTok{=}\NormalTok{ df[}\StringTok{\textquotesingle{}price\textquotesingle{}}\NormalTok{].astype(}\BuiltInTok{float}\NormalTok{)}
\end{Highlighting}
\end{Shaded}
\item
  \textbf{Chain operations logically}: Think about the order of
  operations

\begin{Shaded}
\begin{Highlighting}[]
\CommentTok{\# Logical flow: filter → sort → limit}
\NormalTok{df }\OperatorTok{=}\NormalTok{ add.transform(}\StringTok{\textquotesingle{}@filter\textquotesingle{}}\NormalTok{, df, columns}\OperatorTok{=}\NormalTok{cols, where}\OperatorTok{=}\StringTok{\textquotesingle{}valid == 1\textquotesingle{}}\NormalTok{)}
\NormalTok{df }\OperatorTok{=}\NormalTok{ add.transform(}\StringTok{\textquotesingle{}@sort\textquotesingle{}}\NormalTok{, df, columns}\OperatorTok{=}\NormalTok{cols, by}\OperatorTok{=}\StringTok{\textquotesingle{}score\textquotesingle{}}\NormalTok{, order}\OperatorTok{=}\StringTok{\textquotesingle{}desc\textquotesingle{}}\NormalTok{)}
\NormalTok{top\_results }\OperatorTok{=}\NormalTok{ df.head(}\DecValTok{100}\NormalTok{)}
\end{Highlighting}
\end{Shaded}
\end{enumerate}

\begin{center}\rule{0.5\linewidth}{0.5pt}\end{center}

\section{Key Takeaways}\label{key-takeaways-6}

\begin{itemize}
\tightlist
\item
  Use \texttt{@filter} to remove rows based on conditions
\item
  Use \texttt{@sort} to reorder rows
\item
  \texttt{columns} parameter specifies which columns to keep
\item
  \texttt{where} parameter defines filter conditions
\item
  \texttt{by} parameter specifies sort column(s)
\item
  \texttt{order} parameter controls sort direction (`asc' or `desc')
\item
  Chain operations for complex workflows
\item
  Filter before sorting for better performance
\end{itemize}

\begin{center}\rule{0.5\linewidth}{0.5pt}\end{center}

\section{Common Questions}\label{common-questions-6}

\textbf{Q: Can I use AND/OR in filter conditions?}\\
A: Not currently. Chain multiple \texttt{@filter} operations instead.

\textbf{Q: Can I sort by one column ascending and another descending?}\\
A: Not in a single operation. Sort by one column, then by the other in
separate operations.

\textbf{Q: What happens if the filter condition matches no rows?}\\
A: You'll get an empty DataFrame with the specified columns.

\textbf{Q: Can I filter on multiple columns at once?}\\
A: Not in a single condition. Apply multiple \texttt{@filter} operations
sequentially.

\textbf{Q: Does sorting preserve the original DataFrame?}\\
A: No, \texttt{add.transform()} returns a new DataFrame. The original is
unchanged.

\begin{center}\rule{0.5\linewidth}{0.5pt}\end{center}

\section{Next Steps}\label{next-steps-8}

\begin{itemize}
\tightlist
\item
  \textbf{\href{examples-transform-03-aggregate.html}{Aggregation}} -
  Learn to group and aggregate data
\item
  \textbf{\href{examples-transform-01-calc.html}{Calculations}} - Review
  calculation operations
\item
  \textbf{\href{reference-transform.html}{API Reference}} - Complete
  \texttt{add.transform()} documentation
\end{itemize}

\begin{center}\rule{0.5\linewidth}{0.5pt}\end{center}

\textbf{Version:} 0.1.3\\
\textbf{Last Updated:} March 9, 2026

\chapter{add.transform() -
Aggregation}\label{add.transform---aggregation}

\section{Overview}\label{overview-7}

Learn how to group and aggregate data using \texttt{add.transform()}.
Aggregation is essential for summarizing data, calculating statistics,
and creating reports.

\textbf{What you'll learn:} - How to group data by one or more columns -
How to apply aggregation strategies (sum, mean, count, etc.) - How to
aggregate multiple columns with different strategies - Best practices
for data summarization

\textbf{Prerequisites:} -
\href{examples-transform-01-calc.html}{Calculations} - Understanding of
basic transform operations -
\href{examples-transform-02-filter-sort.html}{Filter \& Sort} - Data
selection and ordering - Basic understanding of grouping concepts

\begin{center}\rule{0.5\linewidth}{0.5pt}\end{center}

\section{Example 1: Group and
Aggregate}\label{example-1-group-and-aggregate}

\textbf{Business Context:} You have sales data from multiple regions and
want to see total sales by region.

\textbf{Code:}

\begin{Shaded}
\begin{Highlighting}[]
\ImportTok{import}\NormalTok{ additory }\ImportTok{as}\NormalTok{ add}
\ImportTok{import}\NormalTok{ pandas }\ImportTok{as}\NormalTok{ pd}

\CommentTok{\# Sales data}
\NormalTok{df }\OperatorTok{=}\NormalTok{ pd.DataFrame(\{}
    \StringTok{\textquotesingle{}region\textquotesingle{}}\NormalTok{: [}\StringTok{\textquotesingle{}North\textquotesingle{}}\NormalTok{, }\StringTok{\textquotesingle{}North\textquotesingle{}}\NormalTok{, }\StringTok{\textquotesingle{}South\textquotesingle{}}\NormalTok{, }\StringTok{\textquotesingle{}South\textquotesingle{}}\NormalTok{, }\StringTok{\textquotesingle{}East\textquotesingle{}}\NormalTok{],}
    \StringTok{\textquotesingle{}product\textquotesingle{}}\NormalTok{: [}\StringTok{\textquotesingle{}Widget\textquotesingle{}}\NormalTok{, }\StringTok{\textquotesingle{}Gadget\textquotesingle{}}\NormalTok{, }\StringTok{\textquotesingle{}Widget\textquotesingle{}}\NormalTok{, }\StringTok{\textquotesingle{}Gadget\textquotesingle{}}\NormalTok{, }\StringTok{\textquotesingle{}Widget\textquotesingle{}}\NormalTok{],}
    \StringTok{\textquotesingle{}sales\textquotesingle{}}\NormalTok{: [}\DecValTok{100}\NormalTok{, }\DecValTok{150}\NormalTok{, }\DecValTok{200}\NormalTok{, }\DecValTok{250}\NormalTok{, }\DecValTok{175}\NormalTok{]}
\NormalTok{\})}

\CommentTok{\# Group by region and sum sales}
\NormalTok{result }\OperatorTok{=}\NormalTok{ add.transform(}
    \StringTok{\textquotesingle{}@aggregate\textquotesingle{}}\NormalTok{,}
\NormalTok{    df,}
\NormalTok{    columns}\OperatorTok{=}\NormalTok{[}\StringTok{\textquotesingle{}region\textquotesingle{}}\NormalTok{, }\StringTok{\textquotesingle{}sales\textquotesingle{}}\NormalTok{],}
\NormalTok{    by}\OperatorTok{=}\StringTok{\textquotesingle{}region\textquotesingle{}}\NormalTok{,}
\NormalTok{    strategy}\OperatorTok{=}\NormalTok{\{}\StringTok{\textquotesingle{}sales\textquotesingle{}}\NormalTok{: }\StringTok{\textquotesingle{}sum\textquotesingle{}}\NormalTok{\}}
\NormalTok{)}

\CommentTok{\# Positional parameters (also works without naming certain parameters):}
\CommentTok{\# result = add.transform(\textquotesingle{}@aggregate\textquotesingle{}, df, [\textquotesingle{}region\textquotesingle{}, \textquotesingle{}sales\textquotesingle{}], by=\textquotesingle{}region\textquotesingle{}, strategy=\{\textquotesingle{}sales\textquotesingle{}: \textquotesingle{}sum\textquotesingle{}\})}

\BuiltInTok{print}\NormalTok{(result)}
\end{Highlighting}
\end{Shaded}

\textbf{Output:}

\begin{verbatim}
  region  sales
0   East    175
1  North    250
2  South    450
\end{verbatim}

\textbf{Explanation:} -
\texttt{\textquotesingle{}@aggregate\textquotesingle{}} mode groups and
summarizes data - \texttt{columns} specifies which columns to include
(grouping column + aggregated columns) -
\texttt{by=\textquotesingle{}region\textquotesingle{}} groups rows by
the region column -
\texttt{strategy=\{\textquotesingle{}sales\textquotesingle{}:\ \textquotesingle{}sum\textquotesingle{}\}}
sums sales values within each group - North: 100 + 150 = 250 - South:
200 + 250 = 450 - East: 175 (only one row)

\textbf{Note:} This also works with polars DataFrames.

\begin{center}\rule{0.5\linewidth}{0.5pt}\end{center}

\section{Example 2: Multiple
Aggregations}\label{example-2-multiple-aggregations}

\textbf{Business Context:} You want to see average price and total
quantity for each product category.

\textbf{Code:}

\begin{Shaded}
\begin{Highlighting}[]
\ImportTok{import}\NormalTok{ additory }\ImportTok{as}\NormalTok{ add}
\ImportTok{import}\NormalTok{ pandas }\ImportTok{as}\NormalTok{ pd}

\CommentTok{\# Product data}
\NormalTok{df }\OperatorTok{=}\NormalTok{ pd.DataFrame(\{}
    \StringTok{\textquotesingle{}category\textquotesingle{}}\NormalTok{: [}\StringTok{\textquotesingle{}A\textquotesingle{}}\NormalTok{, }\StringTok{\textquotesingle{}A\textquotesingle{}}\NormalTok{, }\StringTok{\textquotesingle{}B\textquotesingle{}}\NormalTok{, }\StringTok{\textquotesingle{}B\textquotesingle{}}\NormalTok{, }\StringTok{\textquotesingle{}C\textquotesingle{}}\NormalTok{],}
    \StringTok{\textquotesingle{}price\textquotesingle{}}\NormalTok{: [}\DecValTok{100}\NormalTok{, }\DecValTok{150}\NormalTok{, }\DecValTok{200}\NormalTok{, }\DecValTok{250}\NormalTok{, }\DecValTok{175}\NormalTok{],}
    \StringTok{\textquotesingle{}quantity\textquotesingle{}}\NormalTok{: [}\DecValTok{10}\NormalTok{, }\DecValTok{15}\NormalTok{, }\DecValTok{20}\NormalTok{, }\DecValTok{25}\NormalTok{, }\DecValTok{30}\NormalTok{]}
\NormalTok{\})}

\CommentTok{\# Group by category with multiple aggregations}
\NormalTok{result }\OperatorTok{=}\NormalTok{ add.transform(}
    \StringTok{\textquotesingle{}@aggregate\textquotesingle{}}\NormalTok{,}
\NormalTok{    df,}
\NormalTok{    columns}\OperatorTok{=}\NormalTok{[}\StringTok{\textquotesingle{}category\textquotesingle{}}\NormalTok{, }\StringTok{\textquotesingle{}price\textquotesingle{}}\NormalTok{, }\StringTok{\textquotesingle{}quantity\textquotesingle{}}\NormalTok{],}
\NormalTok{    by}\OperatorTok{=}\StringTok{\textquotesingle{}category\textquotesingle{}}\NormalTok{,}
\NormalTok{    strategy}\OperatorTok{=}\NormalTok{\{}\StringTok{\textquotesingle{}price\textquotesingle{}}\NormalTok{: }\StringTok{\textquotesingle{}mean\textquotesingle{}}\NormalTok{, }\StringTok{\textquotesingle{}quantity\textquotesingle{}}\NormalTok{: }\StringTok{\textquotesingle{}sum\textquotesingle{}}\NormalTok{\}}
\NormalTok{)}

\CommentTok{\# Positional parameters (also works without naming certain parameters):}
\CommentTok{\# result = add.transform(\textquotesingle{}@aggregate\textquotesingle{}, df, [\textquotesingle{}category\textquotesingle{}, \textquotesingle{}price\textquotesingle{}, \textquotesingle{}quantity\textquotesingle{}], }
\CommentTok{\#                        by=\textquotesingle{}category\textquotesingle{}, strategy=\{\textquotesingle{}price\textquotesingle{}: \textquotesingle{}mean\textquotesingle{}, \textquotesingle{}quantity\textquotesingle{}: \textquotesingle{}sum\textquotesingle{}\})}

\BuiltInTok{print}\NormalTok{(result)}
\end{Highlighting}
\end{Shaded}

\textbf{Output:}

\begin{verbatim}
  category  price  quantity
0        A  125.0        25
1        B  225.0        45
2        C  175.0        30
\end{verbatim}

\textbf{Explanation:} - Different aggregation strategies for different
columns -
\texttt{\textquotesingle{}price\textquotesingle{}:\ \textquotesingle{}mean\textquotesingle{}}
calculates average price per category - Category A: (100 + 150) / 2 =
125 - Category B: (200 + 250) / 2 = 225 - Category C: 175 (only one row)
-
\texttt{\textquotesingle{}quantity\textquotesingle{}:\ \textquotesingle{}sum\textquotesingle{}}
sums quantities per category - Category A: 10 + 15 = 25 - Category B: 20
+ 25 = 45 - Category C: 30 (only one row) - Each column can have its own
aggregation strategy

\textbf{Note:} This also works with polars DataFrames.

\begin{center}\rule{0.5\linewidth}{0.5pt}\end{center}

\section{Example 3: Group by Multiple
Columns}\label{example-3-group-by-multiple-columns}

\textbf{Business Context:} You want to see total sales for each
combination of region and product.

\textbf{Code:}

\begin{Shaded}
\begin{Highlighting}[]
\ImportTok{import}\NormalTok{ additory }\ImportTok{as}\NormalTok{ add}
\ImportTok{import}\NormalTok{ pandas }\ImportTok{as}\NormalTok{ pd}

\CommentTok{\# Sales data}
\NormalTok{df }\OperatorTok{=}\NormalTok{ pd.DataFrame(\{}
    \StringTok{\textquotesingle{}region\textquotesingle{}}\NormalTok{: [}\StringTok{\textquotesingle{}North\textquotesingle{}}\NormalTok{, }\StringTok{\textquotesingle{}North\textquotesingle{}}\NormalTok{, }\StringTok{\textquotesingle{}South\textquotesingle{}}\NormalTok{, }\StringTok{\textquotesingle{}South\textquotesingle{}}\NormalTok{, }\StringTok{\textquotesingle{}North\textquotesingle{}}\NormalTok{, }\StringTok{\textquotesingle{}South\textquotesingle{}}\NormalTok{],}
    \StringTok{\textquotesingle{}product\textquotesingle{}}\NormalTok{: [}\StringTok{\textquotesingle{}A\textquotesingle{}}\NormalTok{, }\StringTok{\textquotesingle{}B\textquotesingle{}}\NormalTok{, }\StringTok{\textquotesingle{}A\textquotesingle{}}\NormalTok{, }\StringTok{\textquotesingle{}B\textquotesingle{}}\NormalTok{, }\StringTok{\textquotesingle{}A\textquotesingle{}}\NormalTok{, }\StringTok{\textquotesingle{}A\textquotesingle{}}\NormalTok{],}
    \StringTok{\textquotesingle{}sales\textquotesingle{}}\NormalTok{: [}\DecValTok{100}\NormalTok{, }\DecValTok{150}\NormalTok{, }\DecValTok{200}\NormalTok{, }\DecValTok{250}\NormalTok{, }\DecValTok{120}\NormalTok{, }\DecValTok{180}\NormalTok{]}
\NormalTok{\})}

\CommentTok{\# Group by region and product}
\NormalTok{result }\OperatorTok{=}\NormalTok{ add.transform(}
    \StringTok{\textquotesingle{}@aggregate\textquotesingle{}}\NormalTok{,}
\NormalTok{    df,}
\NormalTok{    columns}\OperatorTok{=}\NormalTok{[}\StringTok{\textquotesingle{}region\textquotesingle{}}\NormalTok{, }\StringTok{\textquotesingle{}product\textquotesingle{}}\NormalTok{, }\StringTok{\textquotesingle{}sales\textquotesingle{}}\NormalTok{],}
\NormalTok{    by}\OperatorTok{=}\NormalTok{[}\StringTok{\textquotesingle{}region\textquotesingle{}}\NormalTok{, }\StringTok{\textquotesingle{}product\textquotesingle{}}\NormalTok{],}
\NormalTok{    strategy}\OperatorTok{=}\NormalTok{\{}\StringTok{\textquotesingle{}sales\textquotesingle{}}\NormalTok{: }\StringTok{\textquotesingle{}sum\textquotesingle{}}\NormalTok{\}}
\NormalTok{)}

\CommentTok{\# Positional parameters (also works without naming certain parameters):}
\CommentTok{\# result = add.transform(\textquotesingle{}@aggregate\textquotesingle{}, df, [\textquotesingle{}region\textquotesingle{}, \textquotesingle{}product\textquotesingle{}, \textquotesingle{}sales\textquotesingle{}], }
\CommentTok{\#                        by=[\textquotesingle{}region\textquotesingle{}, \textquotesingle{}product\textquotesingle{}], strategy=\{\textquotesingle{}sales\textquotesingle{}: \textquotesingle{}sum\textquotesingle{}\})}

\BuiltInTok{print}\NormalTok{(result)}
\end{Highlighting}
\end{Shaded}

\textbf{Output:}

\begin{verbatim}
  region product  sales
0  North       A    220
1  North       B    150
2  South       A    380
3  South       B    250
\end{verbatim}

\textbf{Explanation:} -
\texttt{by={[}\textquotesingle{}region\textquotesingle{},\ \textquotesingle{}product\textquotesingle{}{]}}
groups by multiple columns - Creates a group for each unique combination
- North + A: 100 + 120 = 220 (2 rows combined) - North + B: 150 (1 row)
- South + A: 200 + 180 = 380 (2 rows combined) - South + B: 250 (1 row)
- Result has 4 rows (4 unique combinations)

\textbf{Note:} This also works with polars DataFrames.

\begin{center}\rule{0.5\linewidth}{0.5pt}\end{center}

\section{Aggregation Strategies}\label{aggregation-strategies}

\subsection{Available Strategies}\label{available-strategies}

\begin{longtable}[]{@{}lll@{}}
\toprule\noalign{}
Strategy & Description & Example Use Case \\
\midrule\noalign{}
\endhead
\bottomrule\noalign{}
\endlastfoot
\texttt{\textquotesingle{}sum\textquotesingle{}} & Add all values &
Total sales, total quantity \\
\texttt{\textquotesingle{}mean\textquotesingle{}} & Average of values &
Average price, average rating \\
\texttt{\textquotesingle{}count\textquotesingle{}} & Count of values &
Number of transactions \\
\texttt{\textquotesingle{}min\textquotesingle{}} & Minimum value &
Lowest price, earliest date \\
\texttt{\textquotesingle{}max\textquotesingle{}} & Maximum value &
Highest price, latest date \\
\texttt{\textquotesingle{}first\textquotesingle{}} & First value in
group & Initial status \\
\texttt{\textquotesingle{}last\textquotesingle{}} & Last value in group
& Final status \\
\end{longtable}

\subsection{Strategy Examples}\label{strategy-examples}

\begin{Shaded}
\begin{Highlighting}[]
\CommentTok{\# Sum {-} Total sales per region}
\NormalTok{strategy}\OperatorTok{=}\NormalTok{\{}\StringTok{\textquotesingle{}sales\textquotesingle{}}\NormalTok{: }\StringTok{\textquotesingle{}sum\textquotesingle{}}\NormalTok{\}}

\CommentTok{\# Mean {-} Average price per category}
\NormalTok{strategy}\OperatorTok{=}\NormalTok{\{}\StringTok{\textquotesingle{}price\textquotesingle{}}\NormalTok{: }\StringTok{\textquotesingle{}mean\textquotesingle{}}\NormalTok{\}}

\CommentTok{\# Count {-} Number of orders per customer}
\NormalTok{strategy}\OperatorTok{=}\NormalTok{\{}\StringTok{\textquotesingle{}order\_id\textquotesingle{}}\NormalTok{: }\StringTok{\textquotesingle{}count\textquotesingle{}}\NormalTok{\}}

\CommentTok{\# Min/Max {-} Price range per product}
\NormalTok{strategy}\OperatorTok{=}\NormalTok{\{}\StringTok{\textquotesingle{}price\_min\textquotesingle{}}\NormalTok{: }\StringTok{\textquotesingle{}min\textquotesingle{}}\NormalTok{, }\StringTok{\textquotesingle{}price\_max\textquotesingle{}}\NormalTok{: }\StringTok{\textquotesingle{}max\textquotesingle{}}\NormalTok{\}}

\CommentTok{\# First/Last {-} Status tracking}
\NormalTok{strategy}\OperatorTok{=}\NormalTok{\{}\StringTok{\textquotesingle{}initial\_status\textquotesingle{}}\NormalTok{: }\StringTok{\textquotesingle{}first\textquotesingle{}}\NormalTok{, }\StringTok{\textquotesingle{}final\_status\textquotesingle{}}\NormalTok{: }\StringTok{\textquotesingle{}last\textquotesingle{}}\NormalTok{\}}
\end{Highlighting}
\end{Shaded}

\begin{center}\rule{0.5\linewidth}{0.5pt}\end{center}

\section{Common Patterns}\label{common-patterns-3}

\subsection{Pattern 1: Sales Summary by
Region}\label{pattern-1-sales-summary-by-region}

\begin{Shaded}
\begin{Highlighting}[]
\NormalTok{result }\OperatorTok{=}\NormalTok{ add.transform(}\StringTok{\textquotesingle{}@aggregate\textquotesingle{}}\NormalTok{, df, }
\NormalTok{                       columns}\OperatorTok{=}\NormalTok{[}\StringTok{\textquotesingle{}region\textquotesingle{}}\NormalTok{, }\StringTok{\textquotesingle{}sales\textquotesingle{}}\NormalTok{],}
\NormalTok{                       by}\OperatorTok{=}\StringTok{\textquotesingle{}region\textquotesingle{}}\NormalTok{,}
\NormalTok{                       strategy}\OperatorTok{=}\NormalTok{\{}\StringTok{\textquotesingle{}sales\textquotesingle{}}\NormalTok{: }\StringTok{\textquotesingle{}sum\textquotesingle{}}\NormalTok{\})}
\end{Highlighting}
\end{Shaded}

\subsection{Pattern 2: Customer
Statistics}\label{pattern-2-customer-statistics}

\begin{Shaded}
\begin{Highlighting}[]
\NormalTok{result }\OperatorTok{=}\NormalTok{ add.transform(}\StringTok{\textquotesingle{}@aggregate\textquotesingle{}}\NormalTok{, df,}
\NormalTok{                       columns}\OperatorTok{=}\NormalTok{[}\StringTok{\textquotesingle{}customer\_id\textquotesingle{}}\NormalTok{, }\StringTok{\textquotesingle{}order\_count\textquotesingle{}}\NormalTok{, }\StringTok{\textquotesingle{}total\_spent\textquotesingle{}}\NormalTok{],}
\NormalTok{                       by}\OperatorTok{=}\StringTok{\textquotesingle{}customer\_id\textquotesingle{}}\NormalTok{,}
\NormalTok{                       strategy}\OperatorTok{=}\NormalTok{\{}\StringTok{\textquotesingle{}order\_count\textquotesingle{}}\NormalTok{: }\StringTok{\textquotesingle{}count\textquotesingle{}}\NormalTok{, }\StringTok{\textquotesingle{}total\_spent\textquotesingle{}}\NormalTok{: }\StringTok{\textquotesingle{}sum\textquotesingle{}}\NormalTok{\})}
\end{Highlighting}
\end{Shaded}

\subsection{Pattern 3: Product
Performance}\label{pattern-3-product-performance}

\begin{Shaded}
\begin{Highlighting}[]
\NormalTok{result }\OperatorTok{=}\NormalTok{ add.transform(}\StringTok{\textquotesingle{}@aggregate\textquotesingle{}}\NormalTok{, df,}
\NormalTok{                       columns}\OperatorTok{=}\NormalTok{[}\StringTok{\textquotesingle{}product\textquotesingle{}}\NormalTok{, }\StringTok{\textquotesingle{}avg\_rating\textquotesingle{}}\NormalTok{, }\StringTok{\textquotesingle{}review\_count\textquotesingle{}}\NormalTok{],}
\NormalTok{                       by}\OperatorTok{=}\StringTok{\textquotesingle{}product\textquotesingle{}}\NormalTok{,}
\NormalTok{                       strategy}\OperatorTok{=}\NormalTok{\{}\StringTok{\textquotesingle{}avg\_rating\textquotesingle{}}\NormalTok{: }\StringTok{\textquotesingle{}mean\textquotesingle{}}\NormalTok{, }\StringTok{\textquotesingle{}review\_count\textquotesingle{}}\NormalTok{: }\StringTok{\textquotesingle{}count\textquotesingle{}}\NormalTok{\})}
\end{Highlighting}
\end{Shaded}

\subsection{Pattern 4: Time-Based
Aggregation}\label{pattern-4-time-based-aggregation}

\begin{Shaded}
\begin{Highlighting}[]
\NormalTok{result }\OperatorTok{=}\NormalTok{ add.transform(}\StringTok{\textquotesingle{}@aggregate\textquotesingle{}}\NormalTok{, df,}
\NormalTok{                       columns}\OperatorTok{=}\NormalTok{[}\StringTok{\textquotesingle{}date\textquotesingle{}}\NormalTok{, }\StringTok{\textquotesingle{}region\textquotesingle{}}\NormalTok{, }\StringTok{\textquotesingle{}daily\_sales\textquotesingle{}}\NormalTok{],}
\NormalTok{                       by}\OperatorTok{=}\NormalTok{[}\StringTok{\textquotesingle{}date\textquotesingle{}}\NormalTok{, }\StringTok{\textquotesingle{}region\textquotesingle{}}\NormalTok{],}
\NormalTok{                       strategy}\OperatorTok{=}\NormalTok{\{}\StringTok{\textquotesingle{}daily\_sales\textquotesingle{}}\NormalTok{: }\StringTok{\textquotesingle{}sum\textquotesingle{}}\NormalTok{\})}
\end{Highlighting}
\end{Shaded}

\begin{center}\rule{0.5\linewidth}{0.5pt}\end{center}

\section{Best Practices}\label{best-practices-2}

\begin{enumerate}
\def\labelenumi{\arabic{enumi}.}
\item
  \textbf{Include grouping columns in columns parameter}: Always list
  the columns you're grouping by

\begin{Shaded}
\begin{Highlighting}[]
\CommentTok{\# Good}
\NormalTok{columns}\OperatorTok{=}\NormalTok{[}\StringTok{\textquotesingle{}region\textquotesingle{}}\NormalTok{, }\StringTok{\textquotesingle{}sales\textquotesingle{}}\NormalTok{]  }\CommentTok{\# Include \textquotesingle{}region\textquotesingle{} (the grouping column)}
\NormalTok{by}\OperatorTok{=}\StringTok{\textquotesingle{}region\textquotesingle{}}

\CommentTok{\# Will fail}
\NormalTok{columns}\OperatorTok{=}\NormalTok{[}\StringTok{\textquotesingle{}sales\textquotesingle{}}\NormalTok{]  }\CommentTok{\# Missing \textquotesingle{}region\textquotesingle{}}
\NormalTok{by}\OperatorTok{=}\StringTok{\textquotesingle{}region\textquotesingle{}}
\end{Highlighting}
\end{Shaded}
\item
  \textbf{Choose appropriate aggregation strategies}: Match the strategy
  to your data type

\begin{Shaded}
\begin{Highlighting}[]
\CommentTok{\# Good}
\NormalTok{strategy}\OperatorTok{=}\NormalTok{\{}\StringTok{\textquotesingle{}sales\textquotesingle{}}\NormalTok{: }\StringTok{\textquotesingle{}sum\textquotesingle{}}\NormalTok{\}  }\CommentTok{\# Numeric data}
\NormalTok{strategy}\OperatorTok{=}\NormalTok{\{}\StringTok{\textquotesingle{}rating\textquotesingle{}}\NormalTok{: }\StringTok{\textquotesingle{}mean\textquotesingle{}}\NormalTok{\}  }\CommentTok{\# Numeric data}
\NormalTok{strategy}\OperatorTok{=}\NormalTok{\{}\StringTok{\textquotesingle{}order\_id\textquotesingle{}}\NormalTok{: }\StringTok{\textquotesingle{}count\textquotesingle{}}\NormalTok{\}  }\CommentTok{\# Any data type}

\CommentTok{\# Less useful}
\NormalTok{strategy}\OperatorTok{=}\NormalTok{\{}\StringTok{\textquotesingle{}customer\_name\textquotesingle{}}\NormalTok{: }\StringTok{\textquotesingle{}sum\textquotesingle{}}\NormalTok{\}  }\CommentTok{\# Text data {-} doesn\textquotesingle{}t make sense}
\end{Highlighting}
\end{Shaded}
\item
  \textbf{Use multiple strategies for comprehensive analysis}:

\begin{Shaded}
\begin{Highlighting}[]
\CommentTok{\# Good: Multiple metrics}
\NormalTok{strategy}\OperatorTok{=}\NormalTok{\{}
    \StringTok{\textquotesingle{}sales\textquotesingle{}}\NormalTok{: }\StringTok{\textquotesingle{}sum\textquotesingle{}}\NormalTok{,      }\CommentTok{\# Total sales}
    \StringTok{\textquotesingle{}price\textquotesingle{}}\NormalTok{: }\StringTok{\textquotesingle{}mean\textquotesingle{}}\NormalTok{,     }\CommentTok{\# Average price}
    \StringTok{\textquotesingle{}orders\textquotesingle{}}\NormalTok{: }\StringTok{\textquotesingle{}count\textquotesingle{}}    \CommentTok{\# Number of orders}
\NormalTok{\}}
\end{Highlighting}
\end{Shaded}
\item
  \textbf{Group by meaningful dimensions}: Choose grouping columns that
  make business sense

\begin{Shaded}
\begin{Highlighting}[]
\CommentTok{\# Good: Business dimensions}
\NormalTok{by}\OperatorTok{=}\StringTok{\textquotesingle{}region\textquotesingle{}}  \CommentTok{\# Geographic analysis}
\NormalTok{by}\OperatorTok{=}\StringTok{\textquotesingle{}product\_category\textquotesingle{}}  \CommentTok{\# Product analysis}
\NormalTok{by}\OperatorTok{=}\NormalTok{[}\StringTok{\textquotesingle{}date\textquotesingle{}}\NormalTok{, }\StringTok{\textquotesingle{}store\textquotesingle{}}\NormalTok{]  }\CommentTok{\# Time and location analysis}
\end{Highlighting}
\end{Shaded}
\item
  \textbf{Verify your results}: Check that aggregated values make sense

\begin{Shaded}
\begin{Highlighting}[]
\CommentTok{\# After aggregation, verify totals}
\BuiltInTok{print}\NormalTok{(}\SpecialStringTok{f"Total sales: }\SpecialCharTok{\{}\NormalTok{result[}\StringTok{\textquotesingle{}sales\textquotesingle{}}\NormalTok{]}\SpecialCharTok{.}\BuiltInTok{sum}\NormalTok{()}\SpecialCharTok{\}}\SpecialStringTok{"}\NormalTok{)}
\BuiltInTok{print}\NormalTok{(}\SpecialStringTok{f"Number of groups: }\SpecialCharTok{\{}\BuiltInTok{len}\NormalTok{(result)}\SpecialCharTok{\}}\SpecialStringTok{"}\NormalTok{)}
\end{Highlighting}
\end{Shaded}
\end{enumerate}

\begin{center}\rule{0.5\linewidth}{0.5pt}\end{center}

\section{Key Takeaways}\label{key-takeaways-7}

\begin{itemize}
\tightlist
\item
  Use \texttt{@aggregate} mode to group and summarize data
\item
  \texttt{by} parameter specifies grouping column(s)
\item
  \texttt{strategy} parameter defines how to aggregate each column
\item
  Multiple columns can have different aggregation strategies
\item
  Group by multiple columns for multi-dimensional analysis
\item
  Always include grouping columns in the \texttt{columns} parameter
\item
  Choose aggregation strategies that match your data type
\end{itemize}

\begin{center}\rule{0.5\linewidth}{0.5pt}\end{center}

\section{Common Questions}\label{common-questions-7}

\textbf{Q: Can I aggregate without grouping?}\\
A: Yes, omit the \texttt{by} parameter to aggregate the entire DataFrame
into a single row.

\textbf{Q: What happens if a group has only one row?}\\
A: The aggregation still works - sum/mean/etc. of one value is that
value.

\textbf{Q: Can I use different strategies for the same column?}\\
A: Not in a single operation. Run multiple aggregations and combine
results.

\textbf{Q: How do I count distinct values?}\\
A: Use \texttt{\textquotesingle{}count\textquotesingle{}} strategy, but
note it counts all values, not distinct ones. For distinct counts, use
pandas/polars native methods.

\textbf{Q: Can I aggregate text columns?}\\
A: Yes, but only \texttt{\textquotesingle{}first\textquotesingle{}},
\texttt{\textquotesingle{}last\textquotesingle{}}, and
\texttt{\textquotesingle{}count\textquotesingle{}} strategies make sense
for text.

\textbf{Q: What if my grouping column has null values?}\\
A: Null values are treated as a separate group.

\begin{center}\rule{0.5\linewidth}{0.5pt}\end{center}

\section{Next Steps}\label{next-steps-9}

\begin{itemize}
\tightlist
\item
  \textbf{\href{examples-transform-04-advanced-modes.html}{Advanced
  Modes}} - Learn about (\textbf{harmonize?}), (\textbf{onehot?}),
  (\textbf{extract?}), (\textbf{deduce?})
\item
  \textbf{\href{examples-transform-01-calc.html}{Calculations}} - Review
  calculation operations
\item
  \textbf{\href{reference-transform.html}{API Reference}} - Complete
  \texttt{add.transform()} documentation
\end{itemize}

\begin{center}\rule{0.5\linewidth}{0.5pt}\end{center}

\textbf{Version:} 0.1.3\\
\textbf{Last Updated:} March 9, 2026

\chapter{add.transform() - Advanced
Modes}\label{add.transform---advanced-modes}

\section{Overview}\label{overview-8}

Learn how to use advanced transformation modes with
\texttt{add.transform()} for specialized data operations like encoding,
feature extraction, and missing value imputation.

\textbf{What you'll learn:} - How to one-hot encode categorical
variables - How to extract datetime features automatically - How to
impute missing values with different strategies - How to label encode
categorical data

\textbf{Prerequisites:} - Basic understanding of DataFrames (pandas or
polars) - Familiarity with \texttt{add.transform()} basics

\begin{center}\rule{0.5\linewidth}{0.5pt}\end{center}

\section{Example 1: One-Hot Encoding}\label{example-1-one-hot-encoding}

\textbf{Business Context:} You have customer tier data (Gold, Silver,
Bronze) and need to convert it to binary columns for machine learning
models.

\textbf{Code:}

\begin{Shaded}
\begin{Highlighting}[]
\ImportTok{import}\NormalTok{ additory }\ImportTok{as}\NormalTok{ add}
\ImportTok{import}\NormalTok{ pandas }\ImportTok{as}\NormalTok{ pd}

\CommentTok{\# Customer data}
\NormalTok{df }\OperatorTok{=}\NormalTok{ pd.DataFrame(\{}
    \StringTok{\textquotesingle{}customer\textquotesingle{}}\NormalTok{: [}\StringTok{\textquotesingle{}Alice\textquotesingle{}}\NormalTok{, }\StringTok{\textquotesingle{}Bob\textquotesingle{}}\NormalTok{, }\StringTok{\textquotesingle{}Charlie\textquotesingle{}}\NormalTok{],}
    \StringTok{\textquotesingle{}tier\textquotesingle{}}\NormalTok{: [}\StringTok{\textquotesingle{}Gold\textquotesingle{}}\NormalTok{, }\StringTok{\textquotesingle{}Silver\textquotesingle{}}\NormalTok{, }\StringTok{\textquotesingle{}Gold\textquotesingle{}}\NormalTok{],}
    \StringTok{\textquotesingle{}purchases\textquotesingle{}}\NormalTok{: [}\DecValTok{10}\NormalTok{, }\DecValTok{5}\NormalTok{, }\DecValTok{8}\NormalTok{]}
\NormalTok{\})}

\CommentTok{\# One{-}hot encode tier column}
\NormalTok{result }\OperatorTok{=}\NormalTok{ add.transform(}
    \StringTok{\textquotesingle{}@onehotencode\textquotesingle{}}\NormalTok{,}
\NormalTok{    df,}
\NormalTok{    columns}\OperatorTok{=}\NormalTok{[}\StringTok{\textquotesingle{}tier\textquotesingle{}}\NormalTok{]}
\NormalTok{)}

\CommentTok{\# Positional parameters (also works without naming certain parameters):}
\CommentTok{\# result = add.transform(\textquotesingle{}@onehotencode\textquotesingle{}, df, [\textquotesingle{}tier\textquotesingle{}])}

\BuiltInTok{print}\NormalTok{(result)}
\end{Highlighting}
\end{Shaded}

\textbf{Output:}

\begin{verbatim}
  customer  tier  purchases  tier_Gold  tier_Silver
0    Alice  Gold         10          1            0
1      Bob Silver          5          0            1
2  Charlie  Gold          8          1            0
\end{verbatim}

\textbf{Explanation:} -
\texttt{\textquotesingle{}@onehotencode\textquotesingle{}} mode creates
binary columns for each unique value - Each unique value in the original
column becomes a new column - Values are 1 if the row has that value, 0
otherwise - Original column is preserved - Also works with the alias
\texttt{@onehot}

\textbf{Note:} This also works with polars DataFrames.

\begin{center}\rule{0.5\linewidth}{0.5pt}\end{center}

\section{Example 2: Extract Datetime
Features}\label{example-2-extract-datetime-features}

\textbf{Business Context:} You have order dates as strings and need to
extract month information for seasonal analysis.

\textbf{Code:}

\begin{Shaded}
\begin{Highlighting}[]
\ImportTok{import}\NormalTok{ additory }\ImportTok{as}\NormalTok{ add}
\ImportTok{import}\NormalTok{ pandas }\ImportTok{as}\NormalTok{ pd}

\CommentTok{\# Order data}
\NormalTok{df }\OperatorTok{=}\NormalTok{ pd.DataFrame(\{}
    \StringTok{\textquotesingle{}order\_id\textquotesingle{}}\NormalTok{: [}\DecValTok{1}\NormalTok{, }\DecValTok{2}\NormalTok{, }\DecValTok{3}\NormalTok{],}
    \StringTok{\textquotesingle{}order\_date\textquotesingle{}}\NormalTok{: [}\StringTok{\textquotesingle{}2024{-}01{-}15\textquotesingle{}}\NormalTok{, }\StringTok{\textquotesingle{}2024{-}02{-}20\textquotesingle{}}\NormalTok{, }\StringTok{\textquotesingle{}2024{-}03{-}10\textquotesingle{}}\NormalTok{],}
    \StringTok{\textquotesingle{}amount\textquotesingle{}}\NormalTok{: [}\DecValTok{100}\NormalTok{, }\DecValTok{200}\NormalTok{, }\DecValTok{150}\NormalTok{]}
\NormalTok{\})}

\CommentTok{\# Extract datetime features}
\NormalTok{result }\OperatorTok{=}\NormalTok{ add.transform(}
    \StringTok{\textquotesingle{}@extract\textquotesingle{}}\NormalTok{,}
\NormalTok{    df,}
\NormalTok{    columns}\OperatorTok{=}\NormalTok{[}\StringTok{\textquotesingle{}order\_date\textquotesingle{}}\NormalTok{],}
\NormalTok{    strategy}\OperatorTok{=}\NormalTok{\{}\StringTok{\textquotesingle{}features\textquotesingle{}}\NormalTok{: [}\StringTok{\textquotesingle{}year\textquotesingle{}}\NormalTok{, }\StringTok{\textquotesingle{}month\textquotesingle{}}\NormalTok{, }\StringTok{\textquotesingle{}day\textquotesingle{}}\NormalTok{]\}}
\NormalTok{)}

\CommentTok{\# Positional parameters (also works without naming certain parameters):}
\CommentTok{\# result = add.transform(\textquotesingle{}@extract\textquotesingle{}, df, [\textquotesingle{}order\_date\textquotesingle{}], }
\CommentTok{\#                        strategy=\{\textquotesingle{}features\textquotesingle{}: [\textquotesingle{}year\textquotesingle{}, \textquotesingle{}month\textquotesingle{}, \textquotesingle{}day\textquotesingle{}]\})}

\BuiltInTok{print}\NormalTok{(result)}
\end{Highlighting}
\end{Shaded}

\textbf{Output:}

\begin{verbatim}
   order_id  order_date  amount  order_date_hour  order_date_month
0         1  2024-01-15     100              NaN               1.0
1         2  2024-02-20     200              NaN               2.0
2         3  2024-03-10     150              NaN               3.0
\end{verbatim}

\textbf{Explanation:} -
\texttt{\textquotesingle{}@extract\textquotesingle{}} mode automatically
extracts datetime features - Parses string dates and extracts components
- Creates new columns with \texttt{\_hour} and \texttt{\_month} suffixes
- Original column is preserved - Works with various date formats

\textbf{Note:} This also works with polars DataFrames.

\begin{center}\rule{0.5\linewidth}{0.5pt}\end{center}

\section{Example 3: Impute Missing
Values}\label{example-3-impute-missing-values}

\textbf{Business Context:} You have product data with some missing
prices and stock levels. You need to fill these gaps with reasonable
estimates.

\textbf{Code:}

\begin{Shaded}
\begin{Highlighting}[]
\ImportTok{import}\NormalTok{ additory }\ImportTok{as}\NormalTok{ add}
\ImportTok{import}\NormalTok{ pandas }\ImportTok{as}\NormalTok{ pd}

\CommentTok{\# Product data with missing values}
\NormalTok{df }\OperatorTok{=}\NormalTok{ pd.DataFrame(\{}
    \StringTok{\textquotesingle{}product\textquotesingle{}}\NormalTok{: [}\StringTok{\textquotesingle{}Widget\textquotesingle{}}\NormalTok{, }\StringTok{\textquotesingle{}Gadget\textquotesingle{}}\NormalTok{, }\StringTok{\textquotesingle{}Doohickey\textquotesingle{}}\NormalTok{, }\StringTok{\textquotesingle{}Thingamajig\textquotesingle{}}\NormalTok{],}
    \StringTok{\textquotesingle{}price\textquotesingle{}}\NormalTok{: [}\FloatTok{29.99}\NormalTok{, }\VariableTok{None}\NormalTok{, }\FloatTok{19.99}\NormalTok{, }\FloatTok{39.99}\NormalTok{],}
    \StringTok{\textquotesingle{}stock\textquotesingle{}}\NormalTok{: [}\DecValTok{10}\NormalTok{, }\DecValTok{5}\NormalTok{, }\VariableTok{None}\NormalTok{, }\DecValTok{15}\NormalTok{]}
\NormalTok{\})}

\CommentTok{\# Impute missing values with mean}
\NormalTok{result }\OperatorTok{=}\NormalTok{ add.transform(}
    \StringTok{\textquotesingle{}@deduce\textquotesingle{}}\NormalTok{,}
\NormalTok{    df,}
\NormalTok{    columns}\OperatorTok{=}\NormalTok{[}\StringTok{\textquotesingle{}price\textquotesingle{}}\NormalTok{, }\StringTok{\textquotesingle{}stock\textquotesingle{}}\NormalTok{],}
\NormalTok{    strategy}\OperatorTok{=}\NormalTok{\{}\StringTok{\textquotesingle{}method\textquotesingle{}}\NormalTok{: }\StringTok{\textquotesingle{}mean\textquotesingle{}}\NormalTok{\}}
\NormalTok{)}

\CommentTok{\# Positional parameters (also works without naming certain parameters):}
\CommentTok{\# result = add.transform(\textquotesingle{}@deduce\textquotesingle{}, df, [\textquotesingle{}price\textquotesingle{}, \textquotesingle{}stock\textquotesingle{}], }
\CommentTok{\#                        strategy=\{\textquotesingle{}method\textquotesingle{}: \textquotesingle{}mean\textquotesingle{}\})}

\BuiltInTok{print}\NormalTok{(result)}
\end{Highlighting}
\end{Shaded}

\textbf{Output:}

\begin{verbatim}
       product  price  stock
0       Widget  29.99   10.0
1       Gadget  29.99    5.0
2    Doohickey  19.99   10.0
3  Thingamajig  39.99   15.0
\end{verbatim}

\textbf{Explanation:} -
\texttt{\textquotesingle{}@deduce\textquotesingle{}} mode fills missing
values using various strategies -
\texttt{method=\textquotesingle{}mean\textquotesingle{}} uses the
average of non-null values - Missing price (Gadget) filled with mean:
(29.99 + 19.99 + 39.99) / 3 = 29.99 - Missing stock (Doohickey) filled
with mean: (10 + 5 + 15) / 3 = 10.0 - Original columns are updated with
imputed values

\textbf{Note:} This also works with polars DataFrames.

\begin{center}\rule{0.5\linewidth}{0.5pt}\end{center}

\section{Example 4: Label Encoding}\label{example-4-label-encoding}

\textbf{Business Context:} You have categorical product categories and
need to convert them to numeric labels for analysis.

\textbf{Code:}

\begin{Shaded}
\begin{Highlighting}[]
\ImportTok{import}\NormalTok{ additory }\ImportTok{as}\NormalTok{ add}
\ImportTok{import}\NormalTok{ pandas }\ImportTok{as}\NormalTok{ pd}

\CommentTok{\# Product data}
\NormalTok{df }\OperatorTok{=}\NormalTok{ pd.DataFrame(\{}
    \StringTok{\textquotesingle{}product\textquotesingle{}}\NormalTok{: [}\StringTok{\textquotesingle{}Widget\textquotesingle{}}\NormalTok{, }\StringTok{\textquotesingle{}Gadget\textquotesingle{}}\NormalTok{, }\StringTok{\textquotesingle{}Widget\textquotesingle{}}\NormalTok{, }\StringTok{\textquotesingle{}Doohickey\textquotesingle{}}\NormalTok{, }\StringTok{\textquotesingle{}Gadget\textquotesingle{}}\NormalTok{],}
    \StringTok{\textquotesingle{}category\textquotesingle{}}\NormalTok{: [}\StringTok{\textquotesingle{}Electronics\textquotesingle{}}\NormalTok{, }\StringTok{\textquotesingle{}Electronics\textquotesingle{}}\NormalTok{, }\StringTok{\textquotesingle{}Electronics\textquotesingle{}}\NormalTok{, }\StringTok{\textquotesingle{}Tools\textquotesingle{}}\NormalTok{, }\StringTok{\textquotesingle{}Electronics\textquotesingle{}}\NormalTok{],}
    \StringTok{\textquotesingle{}sales\textquotesingle{}}\NormalTok{: [}\DecValTok{100}\NormalTok{, }\DecValTok{150}\NormalTok{, }\DecValTok{120}\NormalTok{, }\DecValTok{80}\NormalTok{, }\DecValTok{200}\NormalTok{]}
\NormalTok{\})}

\CommentTok{\# Label encode category column}
\NormalTok{result }\OperatorTok{=}\NormalTok{ add.transform(}
    \StringTok{\textquotesingle{}@label\textquotesingle{}}\NormalTok{,}
\NormalTok{    df,}
\NormalTok{    columns}\OperatorTok{=}\NormalTok{[}\StringTok{\textquotesingle{}category\textquotesingle{}}\NormalTok{],}
\NormalTok{    strategy}\OperatorTok{=}\NormalTok{\{}\StringTok{\textquotesingle{}bins\textquotesingle{}}\NormalTok{: [}\DecValTok{0}\NormalTok{, }\DecValTok{1}\NormalTok{, }\DecValTok{2}\NormalTok{], }\StringTok{\textquotesingle{}labels\textquotesingle{}}\NormalTok{: [}\StringTok{\textquotesingle{}Electronics\textquotesingle{}}\NormalTok{, }\StringTok{\textquotesingle{}Tools\textquotesingle{}}\NormalTok{]\}}
\NormalTok{)}

\CommentTok{\# Positional parameters (also works without naming certain parameters):}
\CommentTok{\# result = add.transform(\textquotesingle{}@label\textquotesingle{}, df, [\textquotesingle{}category\textquotesingle{}], }
\CommentTok{\#                        strategy=\{\textquotesingle{}bins\textquotesingle{}: [0, 1, 2], \textquotesingle{}labels\textquotesingle{}: [\textquotesingle{}Electronics\textquotesingle{}, \textquotesingle{}Tools\textquotesingle{}]\})}

\BuiltInTok{print}\NormalTok{(result)}
\end{Highlighting}
\end{Shaded}

\textbf{Output:}

\begin{verbatim}
     product     category  sales category_labeled
0     Widget  Electronics    100      Electronics
1     Gadget  Electronics    150      Electronics
2     Widget  Electronics    120      Electronics
3  Doohickey        Tools     80            Tools
4     Gadget  Electronics    200      Electronics
\end{verbatim}

\textbf{Explanation:} -
\texttt{\textquotesingle{}@label\textquotesingle{}} mode creates labeled
categories - \texttt{bins} defines the numeric ranges - \texttt{labels}
defines the category names - Creates a new column with
\texttt{\_labeled} suffix - Original column is preserved

\textbf{Note:} This also works with polars DataFrames.

\begin{center}\rule{0.5\linewidth}{0.5pt}\end{center}

\section{Available Advanced Modes}\label{available-advanced-modes}

\subsection{Encoding Modes}\label{encoding-modes}

\begin{longtable}[]{@{}
  >{\raggedright\arraybackslash}p{(\linewidth - 4\tabcolsep) * \real{0.2069}}
  >{\raggedright\arraybackslash}p{(\linewidth - 4\tabcolsep) * \real{0.4483}}
  >{\raggedright\arraybackslash}p{(\linewidth - 4\tabcolsep) * \real{0.3448}}@{}}
\toprule\noalign{}
\begin{minipage}[b]{\linewidth}\raggedright
Mode
\end{minipage} & \begin{minipage}[b]{\linewidth}\raggedright
Description
\end{minipage} & \begin{minipage}[b]{\linewidth}\raggedright
Use Case
\end{minipage} \\
\midrule\noalign{}
\endhead
\bottomrule\noalign{}
\endlastfoot
\texttt{@onehotencode} & Convert categorical to binary columns & ML
feature preparation \\
\texttt{@onehot} & Alias for \texttt{@onehotencode} & Same as above \\
\texttt{@label} & Create labeled categories & Categorical binning \\
\end{longtable}

\subsection{Feature Extraction}\label{feature-extraction}

\begin{longtable}[]{@{}
  >{\raggedright\arraybackslash}p{(\linewidth - 4\tabcolsep) * \real{0.2069}}
  >{\raggedright\arraybackslash}p{(\linewidth - 4\tabcolsep) * \real{0.4483}}
  >{\raggedright\arraybackslash}p{(\linewidth - 4\tabcolsep) * \real{0.3448}}@{}}
\toprule\noalign{}
\begin{minipage}[b]{\linewidth}\raggedright
Mode
\end{minipage} & \begin{minipage}[b]{\linewidth}\raggedright
Description
\end{minipage} & \begin{minipage}[b]{\linewidth}\raggedright
Use Case
\end{minipage} \\
\midrule\noalign{}
\endhead
\bottomrule\noalign{}
\endlastfoot
\texttt{@extract} & Extract datetime/text features & Feature
engineering \\
\texttt{@datetime} & Parse datetime strings & Date parsing (merged into
(\textbf{extract?})) \\
\end{longtable}

\subsection{Data Cleaning}\label{data-cleaning}

\begin{longtable}[]{@{}lll@{}}
\toprule\noalign{}
Mode & Description & Use Case \\
\midrule\noalign{}
\endhead
\bottomrule\noalign{}
\endlastfoot
\texttt{@deduce} & Impute missing values & Data cleaning \\
\texttt{@harmonize} & Convert measurement units & Unit
standardization \\
\texttt{@round} & Round numeric values & Number formatting \\
\end{longtable}

\subsection{Data Reshaping}\label{data-reshaping}

\begin{longtable}[]{@{}lll@{}}
\toprule\noalign{}
Mode & Description & Use Case \\
\midrule\noalign{}
\endhead
\bottomrule\noalign{}
\endlastfoot
\texttt{@transpose} & Transpose DataFrame & Pivot data structure \\
\texttt{@split} & Split text columns & Text parsing \\
\end{longtable}

\begin{center}\rule{0.5\linewidth}{0.5pt}\end{center}

\section{\texorpdfstring{Imputation Methods
((\textbf{deduce?}))}{Imputation Methods ((deduce?))}}\label{imputation-methods-deduce}

The \texttt{@deduce} mode supports multiple imputation strategies:

\begin{Shaded}
\begin{Highlighting}[]
\CommentTok{\# Mean imputation (default)}
\NormalTok{result }\OperatorTok{=}\NormalTok{ add.transform(}\StringTok{\textquotesingle{}@deduce\textquotesingle{}}\NormalTok{, df, [}\StringTok{\textquotesingle{}column\textquotesingle{}}\NormalTok{], strategy}\OperatorTok{=}\NormalTok{\{}\StringTok{\textquotesingle{}method\textquotesingle{}}\NormalTok{: }\StringTok{\textquotesingle{}mean\textquotesingle{}}\NormalTok{\})}

\CommentTok{\# Median imputation}
\NormalTok{result }\OperatorTok{=}\NormalTok{ add.transform(}\StringTok{\textquotesingle{}@deduce\textquotesingle{}}\NormalTok{, df, [}\StringTok{\textquotesingle{}column\textquotesingle{}}\NormalTok{], strategy}\OperatorTok{=}\NormalTok{\{}\StringTok{\textquotesingle{}method\textquotesingle{}}\NormalTok{: }\StringTok{\textquotesingle{}median\textquotesingle{}}\NormalTok{\})}

\CommentTok{\# Mode imputation (most frequent)}
\NormalTok{result }\OperatorTok{=}\NormalTok{ add.transform(}\StringTok{\textquotesingle{}@deduce\textquotesingle{}}\NormalTok{, df, [}\StringTok{\textquotesingle{}column\textquotesingle{}}\NormalTok{], strategy}\OperatorTok{=}\NormalTok{\{}\StringTok{\textquotesingle{}method\textquotesingle{}}\NormalTok{: }\StringTok{\textquotesingle{}mode\textquotesingle{}}\NormalTok{\})}

\CommentTok{\# Forward fill}
\NormalTok{result }\OperatorTok{=}\NormalTok{ add.transform(}\StringTok{\textquotesingle{}@deduce\textquotesingle{}}\NormalTok{, df, [}\StringTok{\textquotesingle{}column\textquotesingle{}}\NormalTok{], strategy}\OperatorTok{=}\NormalTok{\{}\StringTok{\textquotesingle{}method\textquotesingle{}}\NormalTok{: }\StringTok{\textquotesingle{}forward\textquotesingle{}}\NormalTok{\})}

\CommentTok{\# Backward fill}
\NormalTok{result }\OperatorTok{=}\NormalTok{ add.transform(}\StringTok{\textquotesingle{}@deduce\textquotesingle{}}\NormalTok{, df, [}\StringTok{\textquotesingle{}column\textquotesingle{}}\NormalTok{], strategy}\OperatorTok{=}\NormalTok{\{}\StringTok{\textquotesingle{}method\textquotesingle{}}\NormalTok{: }\StringTok{\textquotesingle{}backward\textquotesingle{}}\NormalTok{\})}

\CommentTok{\# K{-}Nearest Neighbors}
\NormalTok{result }\OperatorTok{=}\NormalTok{ add.transform(}\StringTok{\textquotesingle{}@deduce\textquotesingle{}}\NormalTok{, df, [}\StringTok{\textquotesingle{}column\textquotesingle{}}\NormalTok{], strategy}\OperatorTok{=}\NormalTok{\{}\StringTok{\textquotesingle{}method\textquotesingle{}}\NormalTok{: }\StringTok{\textquotesingle{}knn\textquotesingle{}}\NormalTok{\})}

\CommentTok{\# Auto (automatically choose best method)}
\NormalTok{result }\OperatorTok{=}\NormalTok{ add.transform(}\StringTok{\textquotesingle{}@deduce\textquotesingle{}}\NormalTok{, df, [}\StringTok{\textquotesingle{}column\textquotesingle{}}\NormalTok{], strategy}\OperatorTok{=}\NormalTok{\{}\StringTok{\textquotesingle{}method\textquotesingle{}}\NormalTok{: }\StringTok{\textquotesingle{}auto\textquotesingle{}}\NormalTok{\})}
\end{Highlighting}
\end{Shaded}

\begin{center}\rule{0.5\linewidth}{0.5pt}\end{center}

\section{Common Patterns}\label{common-patterns-4}

\subsection{Pattern 1: Prepare ML
Features}\label{pattern-1-prepare-ml-features}

\begin{Shaded}
\begin{Highlighting}[]
\CommentTok{\# One{-}hot encode categorical variables}
\NormalTok{df }\OperatorTok{=}\NormalTok{ add.transform(}\StringTok{\textquotesingle{}@onehotencode\textquotesingle{}}\NormalTok{, df, [}\StringTok{\textquotesingle{}category\textquotesingle{}}\NormalTok{, }\StringTok{\textquotesingle{}region\textquotesingle{}}\NormalTok{])}

\CommentTok{\# Extract datetime features}
\NormalTok{df }\OperatorTok{=}\NormalTok{ add.transform(}\StringTok{\textquotesingle{}@extract\textquotesingle{}}\NormalTok{, df, [}\StringTok{\textquotesingle{}date\_column\textquotesingle{}}\NormalTok{])}

\CommentTok{\# Impute missing values}
\NormalTok{df }\OperatorTok{=}\NormalTok{ add.transform(}\StringTok{\textquotesingle{}@deduce\textquotesingle{}}\NormalTok{, df, [}\StringTok{\textquotesingle{}numeric\_col\textquotesingle{}}\NormalTok{], strategy}\OperatorTok{=}\NormalTok{\{}\StringTok{\textquotesingle{}method\textquotesingle{}}\NormalTok{: }\StringTok{\textquotesingle{}mean\textquotesingle{}}\NormalTok{\})}
\end{Highlighting}
\end{Shaded}

\subsection{Pattern 2: Clean Data
Pipeline}\label{pattern-2-clean-data-pipeline}

\begin{Shaded}
\begin{Highlighting}[]
\CommentTok{\# Step 1: Impute missing values}
\NormalTok{df }\OperatorTok{=}\NormalTok{ add.transform(}\StringTok{\textquotesingle{}@deduce\textquotesingle{}}\NormalTok{, df, [}\StringTok{\textquotesingle{}price\textquotesingle{}}\NormalTok{, }\StringTok{\textquotesingle{}quantity\textquotesingle{}}\NormalTok{], strategy}\OperatorTok{=}\NormalTok{\{}\StringTok{\textquotesingle{}method\textquotesingle{}}\NormalTok{: }\StringTok{\textquotesingle{}mean\textquotesingle{}}\NormalTok{\})}

\CommentTok{\# Step 2: Harmonize units}
\NormalTok{df }\OperatorTok{=}\NormalTok{ add.transform(}\StringTok{\textquotesingle{}@harmonize\textquotesingle{}}\NormalTok{, df, [}\StringTok{\textquotesingle{}weight\textquotesingle{}}\NormalTok{], strategy}\OperatorTok{=}\NormalTok{\{}\StringTok{\textquotesingle{}to\_unit\textquotesingle{}}\NormalTok{: }\StringTok{\textquotesingle{}kg\textquotesingle{}}\NormalTok{\})}

\CommentTok{\# Step 3: Round values}
\NormalTok{df }\OperatorTok{=}\NormalTok{ add.transform(}\StringTok{\textquotesingle{}@round\textquotesingle{}}\NormalTok{, df, [}\StringTok{\textquotesingle{}price\textquotesingle{}}\NormalTok{], strategy}\OperatorTok{=}\NormalTok{\{}\StringTok{\textquotesingle{}decimals\textquotesingle{}}\NormalTok{: }\DecValTok{2}\NormalTok{\})}
\end{Highlighting}
\end{Shaded}

\subsection{Pattern 3: Feature
Engineering}\label{pattern-3-feature-engineering}

\begin{Shaded}
\begin{Highlighting}[]
\CommentTok{\# Extract datetime features}
\NormalTok{df }\OperatorTok{=}\NormalTok{ add.transform(}\StringTok{\textquotesingle{}@extract\textquotesingle{}}\NormalTok{, df, [}\StringTok{\textquotesingle{}order\_date\textquotesingle{}}\NormalTok{])}

\CommentTok{\# One{-}hot encode categories}
\NormalTok{df }\OperatorTok{=}\NormalTok{ add.transform(}\StringTok{\textquotesingle{}@onehotencode\textquotesingle{}}\NormalTok{, df, [}\StringTok{\textquotesingle{}customer\_tier\textquotesingle{}}\NormalTok{])}

\CommentTok{\# Now ready for ML model}
\end{Highlighting}
\end{Shaded}

\begin{center}\rule{0.5\linewidth}{0.5pt}\end{center}

\section{Best Practices}\label{best-practices-3}

\begin{enumerate}
\def\labelenumi{\arabic{enumi}.}
\item
  \textbf{Check data types}: Ensure columns have appropriate types
  before transformation

\begin{Shaded}
\begin{Highlighting}[]
\CommentTok{\# Check types}
\BuiltInTok{print}\NormalTok{(df.dtypes)}

\CommentTok{\# Convert if needed}
\NormalTok{df[}\StringTok{\textquotesingle{}date\textquotesingle{}}\NormalTok{] }\OperatorTok{=}\NormalTok{ pd.to\_datetime(df[}\StringTok{\textquotesingle{}date\textquotesingle{}}\NormalTok{])}
\end{Highlighting}
\end{Shaded}
\item
  \textbf{Handle missing values strategically}: Choose imputation method
  based on data distribution

\begin{Shaded}
\begin{Highlighting}[]
\CommentTok{\# Use median for skewed data}
\NormalTok{strategy}\OperatorTok{=}\NormalTok{\{}\StringTok{\textquotesingle{}method\textquotesingle{}}\NormalTok{: }\StringTok{\textquotesingle{}median\textquotesingle{}}\NormalTok{\}}

\CommentTok{\# Use mode for categorical}
\NormalTok{strategy}\OperatorTok{=}\NormalTok{\{}\StringTok{\textquotesingle{}method\textquotesingle{}}\NormalTok{: }\StringTok{\textquotesingle{}mode\textquotesingle{}}\NormalTok{\}}
\end{Highlighting}
\end{Shaded}
\item
  \textbf{Validate one-hot encoding}: Check unique values before
  encoding

\begin{Shaded}
\begin{Highlighting}[]
\CommentTok{\# Check unique values}
\BuiltInTok{print}\NormalTok{(df[}\StringTok{\textquotesingle{}category\textquotesingle{}}\NormalTok{].unique())}

\CommentTok{\# Then encode}
\NormalTok{result }\OperatorTok{=}\NormalTok{ add.transform(}\StringTok{\textquotesingle{}@onehotencode\textquotesingle{}}\NormalTok{, df, [}\StringTok{\textquotesingle{}category\textquotesingle{}}\NormalTok{])}
\end{Highlighting}
\end{Shaded}
\item
  \textbf{Test on sample data}: Try advanced modes on a small sample
  first

\begin{Shaded}
\begin{Highlighting}[]
\CommentTok{\# Test on first 10 rows}
\NormalTok{sample }\OperatorTok{=}\NormalTok{ df.head(}\DecValTok{10}\NormalTok{)}
\NormalTok{result }\OperatorTok{=}\NormalTok{ add.transform(}\StringTok{\textquotesingle{}@deduce\textquotesingle{}}\NormalTok{, sample, [}\StringTok{\textquotesingle{}price\textquotesingle{}}\NormalTok{], strategy}\OperatorTok{=}\NormalTok{\{}\StringTok{\textquotesingle{}method\textquotesingle{}}\NormalTok{: }\StringTok{\textquotesingle{}mean\textquotesingle{}}\NormalTok{\})}
\end{Highlighting}
\end{Shaded}
\end{enumerate}

\begin{center}\rule{0.5\linewidth}{0.5pt}\end{center}

\section{Key Takeaways}\label{key-takeaways-8}

\begin{itemize}
\tightlist
\item
  Advanced modes provide specialized transformations
\item
  \texttt{@onehotencode} creates binary columns for categorical data
\item
  \texttt{@extract} automatically extracts datetime features
\item
  \texttt{@deduce} offers 7 different imputation methods
\item
  \texttt{@label} creates labeled categories
\item
  All modes preserve original columns
\item
  Works with both pandas and polars
\item
  Use \texttt{strategy} parameter for mode-specific options
\end{itemize}

\begin{center}\rule{0.5\linewidth}{0.5pt}\end{center}

\section{Common Questions}\label{common-questions-8}

\textbf{Q: Which imputation method should I use?}\\
A: Use \texttt{mean} for normally distributed data, \texttt{median} for
skewed data, \texttt{mode} for categorical, and
\texttt{forward}/\texttt{backward} for time series.

\textbf{Q: Does (\textbf{onehotencode?}) handle new categories?}\\
A: It creates columns for categories present in the data. New categories
in future data won't have columns.

\textbf{Q: Can I extract specific datetime features?}\\
A: The \texttt{@extract} mode automatically extracts available features.
Use the \texttt{strategy} parameter to specify which features you want.

\textbf{Q: What happens to the original column after encoding?}\\
A: The original column is preserved. New columns are added with
appropriate suffixes.

\textbf{Q: Can I chain multiple advanced modes?}\\
A: Yes! You can chain any modes together. Just call
\texttt{add.transform()} multiple times.

\begin{center}\rule{0.5\linewidth}{0.5pt}\end{center}

\section{Next Steps}\label{next-steps-10}

\begin{itemize}
\tightlist
\item
  \textbf{\href{examples-transform-05-real-world.html}{Real-World
  Workflows}} - See complete transformation pipelines
\item
  \textbf{\href{examples-transform-01-calc.html}{Calculations}} - Review
  calculation basics
\item
  \textbf{\href{reference-transform.html}{API Reference}} - Complete
  \texttt{add.transform()} documentation
\end{itemize}

\begin{center}\rule{0.5\linewidth}{0.5pt}\end{center}

\textbf{Version:} 0.1.3\\
\textbf{Last Updated:} March 9, 2026

\chapter{add.transform() - Real-World
Workflows}\label{add.transform---real-world-workflows}

\section{Overview}\label{overview-9}

Learn how to chain multiple \texttt{add.transform()} operations to
create complete data transformation pipelines. Real-world data analysis
often requires multiple steps - this page shows you how to combine
operations effectively.

\textbf{What you'll learn:} - How to chain transform operations - How to
build multi-step data pipelines - Best practices for workflow design -
Common transformation patterns

\textbf{Prerequisites:} -
\href{examples-transform-01-calc.html}{Calculations} - Basic
calculations - \href{examples-transform-02-filter-sort.html}{Filter \&
Sort} - Data selection -
\href{examples-transform-03-aggregate.html}{Aggregation} - Grouping and
summarizing

\begin{center}\rule{0.5\linewidth}{0.5pt}\end{center}

\section{Example 1: Sales Data
Pipeline}\label{example-1-sales-data-pipeline}

\textbf{Business Context:} You have raw sales transactions and need to
calculate total sales, filter out small transactions, aggregate by
region, and rank regions by performance.

\textbf{Code:}

\begin{Shaded}
\begin{Highlighting}[]
\ImportTok{import}\NormalTok{ additory }\ImportTok{as}\NormalTok{ add}
\ImportTok{import}\NormalTok{ pandas }\ImportTok{as}\NormalTok{ pd}

\CommentTok{\# Raw sales transactions}
\NormalTok{df }\OperatorTok{=}\NormalTok{ pd.DataFrame(\{}
    \StringTok{\textquotesingle{}date\textquotesingle{}}\NormalTok{: [}\StringTok{\textquotesingle{}2024{-}01{-}01\textquotesingle{}}\NormalTok{, }\StringTok{\textquotesingle{}2024{-}01{-}01\textquotesingle{}}\NormalTok{, }\StringTok{\textquotesingle{}2024{-}01{-}02\textquotesingle{}}\NormalTok{, }\StringTok{\textquotesingle{}2024{-}01{-}02\textquotesingle{}}\NormalTok{, }\StringTok{\textquotesingle{}2024{-}01{-}03\textquotesingle{}}\NormalTok{],}
    \StringTok{\textquotesingle{}region\textquotesingle{}}\NormalTok{: [}\StringTok{\textquotesingle{}North\textquotesingle{}}\NormalTok{, }\StringTok{\textquotesingle{}South\textquotesingle{}}\NormalTok{, }\StringTok{\textquotesingle{}North\textquotesingle{}}\NormalTok{, }\StringTok{\textquotesingle{}South\textquotesingle{}}\NormalTok{, }\StringTok{\textquotesingle{}North\textquotesingle{}}\NormalTok{],}
    \StringTok{\textquotesingle{}product\textquotesingle{}}\NormalTok{: [}\StringTok{\textquotesingle{}Widget\textquotesingle{}}\NormalTok{, }\StringTok{\textquotesingle{}Widget\textquotesingle{}}\NormalTok{, }\StringTok{\textquotesingle{}Gadget\textquotesingle{}}\NormalTok{, }\StringTok{\textquotesingle{}Widget\textquotesingle{}}\NormalTok{, }\StringTok{\textquotesingle{}Gadget\textquotesingle{}}\NormalTok{],}
    \StringTok{\textquotesingle{}quantity\textquotesingle{}}\NormalTok{: [}\DecValTok{10}\NormalTok{, }\DecValTok{15}\NormalTok{, }\DecValTok{8}\NormalTok{, }\DecValTok{12}\NormalTok{, }\DecValTok{20}\NormalTok{],}
    \StringTok{\textquotesingle{}price\textquotesingle{}}\NormalTok{: [}\DecValTok{30}\NormalTok{, }\DecValTok{30}\NormalTok{, }\DecValTok{50}\NormalTok{, }\DecValTok{30}\NormalTok{, }\DecValTok{50}\NormalTok{]}
\NormalTok{\})}

\CommentTok{\# Step 1: Calculate total sales}
\NormalTok{result }\OperatorTok{=}\NormalTok{ add.transform(}
    \StringTok{\textquotesingle{}@calc\textquotesingle{}}\NormalTok{,}
\NormalTok{    df,}
\NormalTok{    columns}\OperatorTok{=}\NormalTok{[}\StringTok{\textquotesingle{}date\textquotesingle{}}\NormalTok{, }\StringTok{\textquotesingle{}region\textquotesingle{}}\NormalTok{, }\StringTok{\textquotesingle{}product\textquotesingle{}}\NormalTok{, }\StringTok{\textquotesingle{}quantity\textquotesingle{}}\NormalTok{, }\StringTok{\textquotesingle{}price\textquotesingle{}}\NormalTok{],}
\NormalTok{    expression}\OperatorTok{=}\StringTok{\textquotesingle{}quantity * price\textquotesingle{}}\NormalTok{,}
\NormalTok{    as\_}\OperatorTok{=}\StringTok{\textquotesingle{}total\_sales\textquotesingle{}}
\NormalTok{)}

\BuiltInTok{print}\NormalTok{(}\StringTok{"After calculation:"}\NormalTok{)}
\BuiltInTok{print}\NormalTok{(result)}

\CommentTok{\# Step 2: Filter out low{-}value transactions}
\NormalTok{result }\OperatorTok{=}\NormalTok{ add.transform(}
    \StringTok{\textquotesingle{}@filter\textquotesingle{}}\NormalTok{,}
\NormalTok{    result,}
\NormalTok{    columns}\OperatorTok{=}\NormalTok{[}\StringTok{\textquotesingle{}date\textquotesingle{}}\NormalTok{, }\StringTok{\textquotesingle{}region\textquotesingle{}}\NormalTok{, }\StringTok{\textquotesingle{}product\textquotesingle{}}\NormalTok{, }\StringTok{\textquotesingle{}total\_sales\textquotesingle{}}\NormalTok{],}
\NormalTok{    where}\OperatorTok{=}\StringTok{\textquotesingle{}total\_sales \textgreater{} 300\textquotesingle{}}
\NormalTok{)}

\BuiltInTok{print}\NormalTok{(}\StringTok{"}\CharTok{\textbackslash{}n}\StringTok{After filtering:"}\NormalTok{)}
\BuiltInTok{print}\NormalTok{(result)}

\CommentTok{\# Step 3: Aggregate by region}
\NormalTok{result }\OperatorTok{=}\NormalTok{ add.transform(}
    \StringTok{\textquotesingle{}@aggregate\textquotesingle{}}\NormalTok{,}
\NormalTok{    result,}
\NormalTok{    columns}\OperatorTok{=}\NormalTok{[}\StringTok{\textquotesingle{}region\textquotesingle{}}\NormalTok{, }\StringTok{\textquotesingle{}total\_sales\textquotesingle{}}\NormalTok{],}
\NormalTok{    by}\OperatorTok{=}\StringTok{\textquotesingle{}region\textquotesingle{}}\NormalTok{,}
\NormalTok{    strategy}\OperatorTok{=}\NormalTok{\{}\StringTok{\textquotesingle{}total\_sales\textquotesingle{}}\NormalTok{: }\StringTok{\textquotesingle{}sum\textquotesingle{}}\NormalTok{\}}
\NormalTok{)}

\BuiltInTok{print}\NormalTok{(}\StringTok{"}\CharTok{\textbackslash{}n}\StringTok{After aggregation:"}\NormalTok{)}
\BuiltInTok{print}\NormalTok{(result)}

\CommentTok{\# Step 4: Sort by total sales descending}
\NormalTok{result }\OperatorTok{=}\NormalTok{ add.transform(}
    \StringTok{\textquotesingle{}@sort\textquotesingle{}}\NormalTok{,}
\NormalTok{    result,}
\NormalTok{    columns}\OperatorTok{=}\NormalTok{[}\StringTok{\textquotesingle{}region\textquotesingle{}}\NormalTok{, }\StringTok{\textquotesingle{}total\_sales\textquotesingle{}}\NormalTok{],}
\NormalTok{    by}\OperatorTok{=}\StringTok{\textquotesingle{}total\_sales\textquotesingle{}}\NormalTok{,}
\NormalTok{    order}\OperatorTok{=}\StringTok{\textquotesingle{}desc\textquotesingle{}}
\NormalTok{)}

\BuiltInTok{print}\NormalTok{(}\StringTok{"}\CharTok{\textbackslash{}n}\StringTok{Final result:"}\NormalTok{)}
\BuiltInTok{print}\NormalTok{(result)}
\end{Highlighting}
\end{Shaded}

\textbf{Output:}

\begin{verbatim}
After calculation:
         date region product  quantity  price  total_sales
0  2024-01-01  North  Widget        10     30          300
1  2024-01-01  South  Widget        15     30          450
2  2024-01-02  North  Gadget         8     50          400
3  2024-01-02  South  Widget        12     30          360
4  2024-01-03  North  Gadget        20     50         1000

After filtering:
         date region product  total_sales
0  2024-01-01  South  Widget          450
1  2024-01-02  North  Gadget          400
2  2024-01-02  South  Widget          360
3  2024-01-03  North  Gadget         1000

After aggregation:
  region  total_sales
0  North         1400
1  South          810

Final result:
  region  total_sales
0  North         1400
1  South          810
\end{verbatim}

\textbf{Explanation:} - Step 1: Calculate total\_sales = quantity ×
price for each transaction - Step 2: Keep only transactions over \$300
(filters out first row) - Step 3: Sum total\_sales by region (North:
400+1000=1400, South: 450+360=810) - Step 4: Sort regions by
total\_sales descending (North first) - Each step builds on the previous
result - Pipeline transforms raw data into actionable insights

\textbf{Note:} This also works with polars DataFrames.

\begin{center}\rule{0.5\linewidth}{0.5pt}\end{center}

\section{Example 2: Product Performance
Analysis}\label{example-2-product-performance-analysis}

\textbf{Business Context:} You have product sales and returns data. You
need to aggregate by product, filter out products with high returns, and
rank by sales performance.

\textbf{Code:}

\begin{Shaded}
\begin{Highlighting}[]
\ImportTok{import}\NormalTok{ additory }\ImportTok{as}\NormalTok{ add}
\ImportTok{import}\NormalTok{ pandas }\ImportTok{as}\NormalTok{ pd}

\CommentTok{\# Product transactions}
\NormalTok{df }\OperatorTok{=}\NormalTok{ pd.DataFrame(\{}
    \StringTok{\textquotesingle{}product\textquotesingle{}}\NormalTok{: [}\StringTok{\textquotesingle{}A\textquotesingle{}}\NormalTok{, }\StringTok{\textquotesingle{}A\textquotesingle{}}\NormalTok{, }\StringTok{\textquotesingle{}A\textquotesingle{}}\NormalTok{, }\StringTok{\textquotesingle{}B\textquotesingle{}}\NormalTok{, }\StringTok{\textquotesingle{}B\textquotesingle{}}\NormalTok{, }\StringTok{\textquotesingle{}B\textquotesingle{}}\NormalTok{, }\StringTok{\textquotesingle{}C\textquotesingle{}}\NormalTok{, }\StringTok{\textquotesingle{}C\textquotesingle{}}\NormalTok{],}
    \StringTok{\textquotesingle{}sales\textquotesingle{}}\NormalTok{: [}\DecValTok{100}\NormalTok{, }\DecValTok{150}\NormalTok{, }\DecValTok{200}\NormalTok{, }\DecValTok{80}\NormalTok{, }\DecValTok{90}\NormalTok{, }\DecValTok{85}\NormalTok{, }\DecValTok{300}\NormalTok{, }\DecValTok{350}\NormalTok{],}
    \StringTok{\textquotesingle{}returns\textquotesingle{}}\NormalTok{: [}\DecValTok{5}\NormalTok{, }\DecValTok{10}\NormalTok{, }\DecValTok{8}\NormalTok{, }\DecValTok{15}\NormalTok{, }\DecValTok{20}\NormalTok{, }\DecValTok{18}\NormalTok{, }\DecValTok{2}\NormalTok{, }\DecValTok{3}\NormalTok{]}
\NormalTok{\})}

\CommentTok{\# Step 1: Aggregate by product}
\NormalTok{result }\OperatorTok{=}\NormalTok{ add.transform(}
    \StringTok{\textquotesingle{}@aggregate\textquotesingle{}}\NormalTok{,}
\NormalTok{    df,}
\NormalTok{    columns}\OperatorTok{=}\NormalTok{[}\StringTok{\textquotesingle{}product\textquotesingle{}}\NormalTok{, }\StringTok{\textquotesingle{}sales\textquotesingle{}}\NormalTok{, }\StringTok{\textquotesingle{}returns\textquotesingle{}}\NormalTok{],}
\NormalTok{    by}\OperatorTok{=}\StringTok{\textquotesingle{}product\textquotesingle{}}\NormalTok{,}
\NormalTok{    strategy}\OperatorTok{=}\NormalTok{\{}\StringTok{\textquotesingle{}sales\textquotesingle{}}\NormalTok{: }\StringTok{\textquotesingle{}sum\textquotesingle{}}\NormalTok{, }\StringTok{\textquotesingle{}returns\textquotesingle{}}\NormalTok{: }\StringTok{\textquotesingle{}sum\textquotesingle{}}\NormalTok{\}}
\NormalTok{)}

\BuiltInTok{print}\NormalTok{(}\StringTok{"After aggregation:"}\NormalTok{)}
\BuiltInTok{print}\NormalTok{(result)}

\CommentTok{\# Step 2: Filter products with low returns}
\NormalTok{result }\OperatorTok{=}\NormalTok{ add.transform(}
    \StringTok{\textquotesingle{}@filter\textquotesingle{}}\NormalTok{,}
\NormalTok{    result,}
\NormalTok{    columns}\OperatorTok{=}\NormalTok{[}\StringTok{\textquotesingle{}product\textquotesingle{}}\NormalTok{, }\StringTok{\textquotesingle{}sales\textquotesingle{}}\NormalTok{, }\StringTok{\textquotesingle{}returns\textquotesingle{}}\NormalTok{],}
\NormalTok{    where}\OperatorTok{=}\StringTok{\textquotesingle{}returns \textless{} 50\textquotesingle{}}
\NormalTok{)}

\BuiltInTok{print}\NormalTok{(}\StringTok{"}\CharTok{\textbackslash{}n}\StringTok{After filtering:"}\NormalTok{)}
\BuiltInTok{print}\NormalTok{(result)}

\CommentTok{\# Step 3: Sort by sales descending}
\NormalTok{result }\OperatorTok{=}\NormalTok{ add.transform(}
    \StringTok{\textquotesingle{}@sort\textquotesingle{}}\NormalTok{,}
\NormalTok{    result,}
\NormalTok{    columns}\OperatorTok{=}\NormalTok{[}\StringTok{\textquotesingle{}product\textquotesingle{}}\NormalTok{, }\StringTok{\textquotesingle{}sales\textquotesingle{}}\NormalTok{, }\StringTok{\textquotesingle{}returns\textquotesingle{}}\NormalTok{],}
\NormalTok{    by}\OperatorTok{=}\StringTok{\textquotesingle{}sales\textquotesingle{}}\NormalTok{,}
\NormalTok{    order}\OperatorTok{=}\StringTok{\textquotesingle{}desc\textquotesingle{}}
\NormalTok{)}

\BuiltInTok{print}\NormalTok{(}\StringTok{"}\CharTok{\textbackslash{}n}\StringTok{Final result:"}\NormalTok{)}
\BuiltInTok{print}\NormalTok{(result)}
\end{Highlighting}
\end{Shaded}

\textbf{Output:}

\begin{verbatim}
After aggregation:
  product  sales  returns
0       A    450       23
1       B    255       53
2       C    650        5

After filtering:
  product  sales  returns
0       A    450       23
1       C    650        5

Final result:
  product  sales  returns
0       C    650        5
1       A    450       23
\end{verbatim}

\textbf{Explanation:} - Step 1: Sum sales and returns for each product -
Product A: 450 sales, 23 returns - Product B: 255 sales, 53 returns -
Product C: 650 sales, 5 returns - Step 2: Filter out products with 50+
returns (removes Product B) - Step 3: Sort remaining products by sales
(C first, then A) - Result shows top-performing products with acceptable
return rates

\textbf{Note:} This also works with polars DataFrames.

\begin{center}\rule{0.5\linewidth}{0.5pt}\end{center}

\section{Workflow Design Patterns}\label{workflow-design-patterns}

\subsection{Pattern 1: Calculate → Filter → Aggregate →
Sort}\label{pattern-1-calculate-filter-aggregate-sort}

\begin{Shaded}
\begin{Highlighting}[]
\CommentTok{\# Common analytics pipeline}
\NormalTok{df }\OperatorTok{=}\NormalTok{ add.transform(}\StringTok{\textquotesingle{}@calc\textquotesingle{}}\NormalTok{, df, ...)      }\CommentTok{\# Derive metrics}
\NormalTok{df }\OperatorTok{=}\NormalTok{ add.transform(}\StringTok{\textquotesingle{}@filter\textquotesingle{}}\NormalTok{, df, ...)    }\CommentTok{\# Remove noise}
\NormalTok{df }\OperatorTok{=}\NormalTok{ add.transform(}\StringTok{\textquotesingle{}@aggregate\textquotesingle{}}\NormalTok{, df, ...) }\CommentTok{\# Summarize}
\NormalTok{df }\OperatorTok{=}\NormalTok{ add.transform(}\StringTok{\textquotesingle{}@sort\textquotesingle{}}\NormalTok{, df, ...)      }\CommentTok{\# Rank results}
\end{Highlighting}
\end{Shaded}

\textbf{Use when:} Building reports, dashboards, or analytics

\subsection{Pattern 2: Filter → Calculate →
Aggregate}\label{pattern-2-filter-calculate-aggregate}

\begin{Shaded}
\begin{Highlighting}[]
\CommentTok{\# Focus on relevant data first}
\NormalTok{df }\OperatorTok{=}\NormalTok{ add.transform(}\StringTok{\textquotesingle{}@filter\textquotesingle{}}\NormalTok{, df, ...)    }\CommentTok{\# Select subset}
\NormalTok{df }\OperatorTok{=}\NormalTok{ add.transform(}\StringTok{\textquotesingle{}@calc\textquotesingle{}}\NormalTok{, df, ...)      }\CommentTok{\# Compute on subset}
\NormalTok{df }\OperatorTok{=}\NormalTok{ add.transform(}\StringTok{\textquotesingle{}@aggregate\textquotesingle{}}\NormalTok{, df, ...) }\CommentTok{\# Summarize}
\end{Highlighting}
\end{Shaded}

\textbf{Use when:} Working with large datasets, want to reduce data
early

\subsection{Pattern 3: Aggregate → Filter →
Sort}\label{pattern-3-aggregate-filter-sort}

\begin{Shaded}
\begin{Highlighting}[]
\CommentTok{\# Summarize then refine}
\NormalTok{df }\OperatorTok{=}\NormalTok{ add.transform(}\StringTok{\textquotesingle{}@aggregate\textquotesingle{}}\NormalTok{, df, ...) }\CommentTok{\# Group and summarize}
\NormalTok{df }\OperatorTok{=}\NormalTok{ add.transform(}\StringTok{\textquotesingle{}@filter\textquotesingle{}}\NormalTok{, df, ...)    }\CommentTok{\# Keep top performers}
\NormalTok{df }\OperatorTok{=}\NormalTok{ add.transform(}\StringTok{\textquotesingle{}@sort\textquotesingle{}}\NormalTok{, df, ...)      }\CommentTok{\# Rank}
\end{Highlighting}
\end{Shaded}

\textbf{Use when:} Finding top N items, identifying outliers

\subsection{Pattern 4: Calculate → Calculate →
Aggregate}\label{pattern-4-calculate-calculate-aggregate}

\begin{Shaded}
\begin{Highlighting}[]
\CommentTok{\# Multi{-}step calculations}
\NormalTok{df }\OperatorTok{=}\NormalTok{ add.transform(}\StringTok{\textquotesingle{}@calc\textquotesingle{}}\NormalTok{, df, expression}\OperatorTok{=}\StringTok{\textquotesingle{}a * b\textquotesingle{}}\NormalTok{, as\_}\OperatorTok{=}\StringTok{\textquotesingle{}ab\textquotesingle{}}\NormalTok{)}
\NormalTok{df }\OperatorTok{=}\NormalTok{ add.transform(}\StringTok{\textquotesingle{}@calc\textquotesingle{}}\NormalTok{, df, expression}\OperatorTok{=}\StringTok{\textquotesingle{}ab + c\textquotesingle{}}\NormalTok{, as\_}\OperatorTok{=}\StringTok{\textquotesingle{}result\textquotesingle{}}\NormalTok{)}
\NormalTok{df }\OperatorTok{=}\NormalTok{ add.transform(}\StringTok{\textquotesingle{}@aggregate\textquotesingle{}}\NormalTok{, df, by}\OperatorTok{=}\StringTok{\textquotesingle{}group\textquotesingle{}}\NormalTok{, strategy}\OperatorTok{=}\NormalTok{\{}\StringTok{\textquotesingle{}result\textquotesingle{}}\NormalTok{: }\StringTok{\textquotesingle{}sum\textquotesingle{}}\NormalTok{\})}
\end{Highlighting}
\end{Shaded}

\textbf{Use when:} Complex derived metrics, multi-step formulas

\begin{center}\rule{0.5\linewidth}{0.5pt}\end{center}

\section{Best Practices}\label{best-practices-4}

\begin{enumerate}
\def\labelenumi{\arabic{enumi}.}
\item
  \textbf{Start with a plan}: Sketch out your pipeline before coding

\begin{Shaded}
\begin{Highlighting}[]
\CommentTok{\# Good: Clear sequence}
\CommentTok{\# 1. Calculate revenue}
\CommentTok{\# 2. Filter valid transactions}
\CommentTok{\# 3. Aggregate by customer}
\CommentTok{\# 4. Sort by total revenue}
\end{Highlighting}
\end{Shaded}
\item
  \textbf{Filter early}: Reduce data size as soon as possible

\begin{Shaded}
\begin{Highlighting}[]
\CommentTok{\# Good: Filter first}
\NormalTok{df }\OperatorTok{=}\NormalTok{ add.transform(}\StringTok{\textquotesingle{}@filter\textquotesingle{}}\NormalTok{, df, where}\OperatorTok{=}\StringTok{\textquotesingle{}valid == 1\textquotesingle{}}\NormalTok{)}
\NormalTok{df }\OperatorTok{=}\NormalTok{ add.transform(}\StringTok{\textquotesingle{}@calc\textquotesingle{}}\NormalTok{, df, ...)  }\CommentTok{\# Work with less data}

\CommentTok{\# Less efficient: Filter last}
\NormalTok{df }\OperatorTok{=}\NormalTok{ add.transform(}\StringTok{\textquotesingle{}@calc\textquotesingle{}}\NormalTok{, df, ...)  }\CommentTok{\# Process all data}
\NormalTok{df }\OperatorTok{=}\NormalTok{ add.transform(}\StringTok{\textquotesingle{}@filter\textquotesingle{}}\NormalTok{, df, where}\OperatorTok{=}\StringTok{\textquotesingle{}valid == 1\textquotesingle{}}\NormalTok{)}
\end{Highlighting}
\end{Shaded}
\item
  \textbf{Use intermediate variables}: Make pipelines readable

\begin{Shaded}
\begin{Highlighting}[]
\CommentTok{\# Good: Clear steps}
\NormalTok{with\_sales }\OperatorTok{=}\NormalTok{ add.transform(}\StringTok{\textquotesingle{}@calc\textquotesingle{}}\NormalTok{, df, ...)}
\NormalTok{filtered }\OperatorTok{=}\NormalTok{ add.transform(}\StringTok{\textquotesingle{}@filter\textquotesingle{}}\NormalTok{, with\_sales, ...)}
\NormalTok{aggregated }\OperatorTok{=}\NormalTok{ add.transform(}\StringTok{\textquotesingle{}@aggregate\textquotesingle{}}\NormalTok{, filtered, ...)}
\NormalTok{final }\OperatorTok{=}\NormalTok{ add.transform(}\StringTok{\textquotesingle{}@sort\textquotesingle{}}\NormalTok{, aggregated, ...)}

\CommentTok{\# Less clear: Chained without names}
\NormalTok{result }\OperatorTok{=}\NormalTok{ add.transform(}\StringTok{\textquotesingle{}@sort\textquotesingle{}}\NormalTok{, }
\NormalTok{             add.transform(}\StringTok{\textquotesingle{}@aggregate\textquotesingle{}}\NormalTok{,}
\NormalTok{                 add.transform(}\StringTok{\textquotesingle{}@filter\textquotesingle{}}\NormalTok{,}
\NormalTok{                     add.transform(}\StringTok{\textquotesingle{}@calc\textquotesingle{}}\NormalTok{, df, ...), ...), ...), ...)}
\end{Highlighting}
\end{Shaded}
\item
  \textbf{Verify each step}: Check intermediate results

\begin{Shaded}
\begin{Highlighting}[]
\CommentTok{\# Good: Verify as you go}
\NormalTok{result }\OperatorTok{=}\NormalTok{ add.transform(}\StringTok{\textquotesingle{}@calc\textquotesingle{}}\NormalTok{, df, ...)}
\BuiltInTok{print}\NormalTok{(}\SpecialStringTok{f"After calc: }\SpecialCharTok{\{}\BuiltInTok{len}\NormalTok{(result)}\SpecialCharTok{\}}\SpecialStringTok{ rows"}\NormalTok{)}

\NormalTok{result }\OperatorTok{=}\NormalTok{ add.transform(}\StringTok{\textquotesingle{}@filter\textquotesingle{}}\NormalTok{, result, ...)}
\BuiltInTok{print}\NormalTok{(}\SpecialStringTok{f"After filter: }\SpecialCharTok{\{}\BuiltInTok{len}\NormalTok{(result)}\SpecialCharTok{\}}\SpecialStringTok{ rows"}\NormalTok{)}
\end{Highlighting}
\end{Shaded}
\item
  \textbf{Keep operations simple}: One transformation per step

\begin{Shaded}
\begin{Highlighting}[]
\CommentTok{\# Good: Separate concerns}
\NormalTok{df }\OperatorTok{=}\NormalTok{ add.transform(}\StringTok{\textquotesingle{}@calc\textquotesingle{}}\NormalTok{, df, expression}\OperatorTok{=}\StringTok{\textquotesingle{}a * b\textquotesingle{}}\NormalTok{, as\_}\OperatorTok{=}\StringTok{\textquotesingle{}ab\textquotesingle{}}\NormalTok{)}
\NormalTok{df }\OperatorTok{=}\NormalTok{ add.transform(}\StringTok{\textquotesingle{}@calc\textquotesingle{}}\NormalTok{, df, expression}\OperatorTok{=}\StringTok{\textquotesingle{}ab + c\textquotesingle{}}\NormalTok{, as\_}\OperatorTok{=}\StringTok{\textquotesingle{}result\textquotesingle{}}\NormalTok{)}

\CommentTok{\# Avoid: Complex expressions}
\NormalTok{df }\OperatorTok{=}\NormalTok{ add.transform(}\StringTok{\textquotesingle{}@calc\textquotesingle{}}\NormalTok{, df, expression}\OperatorTok{=}\StringTok{\textquotesingle{}(a * b) + c\textquotesingle{}}\NormalTok{, as\_}\OperatorTok{=}\StringTok{\textquotesingle{}result\textquotesingle{}}\NormalTok{)}
\end{Highlighting}
\end{Shaded}
\item
  \textbf{Document your pipeline}: Add comments explaining business
  logic

\begin{Shaded}
\begin{Highlighting}[]
\CommentTok{\# Calculate customer lifetime value}
\NormalTok{df }\OperatorTok{=}\NormalTok{ add.transform(}\StringTok{\textquotesingle{}@calc\textquotesingle{}}\NormalTok{, df, expression}\OperatorTok{=}\StringTok{\textquotesingle{}orders * avg\_value\textquotesingle{}}\NormalTok{, as\_}\OperatorTok{=}\StringTok{\textquotesingle{}ltv\textquotesingle{}}\NormalTok{)}

\CommentTok{\# Focus on high{-}value customers (LTV \textgreater{} $1000)}
\NormalTok{df }\OperatorTok{=}\NormalTok{ add.transform(}\StringTok{\textquotesingle{}@filter\textquotesingle{}}\NormalTok{, df, where}\OperatorTok{=}\StringTok{\textquotesingle{}ltv \textgreater{} 1000\textquotesingle{}}\NormalTok{)}

\CommentTok{\# Rank by lifetime value}
\NormalTok{df }\OperatorTok{=}\NormalTok{ add.transform(}\StringTok{\textquotesingle{}@sort\textquotesingle{}}\NormalTok{, df, by}\OperatorTok{=}\StringTok{\textquotesingle{}ltv\textquotesingle{}}\NormalTok{, order}\OperatorTok{=}\StringTok{\textquotesingle{}desc\textquotesingle{}}\NormalTok{)}
\end{Highlighting}
\end{Shaded}
\end{enumerate}

\begin{center}\rule{0.5\linewidth}{0.5pt}\end{center}

\section{Common Pitfalls}\label{common-pitfalls}

\subsection{Pitfall 1: Forgetting to include grouping
columns}\label{pitfall-1-forgetting-to-include-grouping-columns}

\begin{Shaded}
\begin{Highlighting}[]
\CommentTok{\# Wrong: Missing \textquotesingle{}region\textquotesingle{} in columns}
\NormalTok{result }\OperatorTok{=}\NormalTok{ add.transform(}\StringTok{\textquotesingle{}@aggregate\textquotesingle{}}\NormalTok{, df,}
\NormalTok{                       columns}\OperatorTok{=}\NormalTok{[}\StringTok{\textquotesingle{}sales\textquotesingle{}}\NormalTok{],  }\CommentTok{\# Missing \textquotesingle{}region\textquotesingle{}!}
\NormalTok{                       by}\OperatorTok{=}\StringTok{\textquotesingle{}region\textquotesingle{}}\NormalTok{,}
\NormalTok{                       strategy}\OperatorTok{=}\NormalTok{\{}\StringTok{\textquotesingle{}sales\textquotesingle{}}\NormalTok{: }\StringTok{\textquotesingle{}sum\textquotesingle{}}\NormalTok{\})}

\CommentTok{\# Correct: Include grouping column}
\NormalTok{result }\OperatorTok{=}\NormalTok{ add.transform(}\StringTok{\textquotesingle{}@aggregate\textquotesingle{}}\NormalTok{, df,}
\NormalTok{                       columns}\OperatorTok{=}\NormalTok{[}\StringTok{\textquotesingle{}region\textquotesingle{}}\NormalTok{, }\StringTok{\textquotesingle{}sales\textquotesingle{}}\NormalTok{],}
\NormalTok{                       by}\OperatorTok{=}\StringTok{\textquotesingle{}region\textquotesingle{}}\NormalTok{,}
\NormalTok{                       strategy}\OperatorTok{=}\NormalTok{\{}\StringTok{\textquotesingle{}sales\textquotesingle{}}\NormalTok{: }\StringTok{\textquotesingle{}sum\textquotesingle{}}\NormalTok{\})}
\end{Highlighting}
\end{Shaded}

\subsection{Pitfall 2: Filtering on wrong data
type}\label{pitfall-2-filtering-on-wrong-data-type}

\begin{Shaded}
\begin{Highlighting}[]
\CommentTok{\# Wrong: Filtering on float columns (not yet supported)}
\NormalTok{df }\OperatorTok{=}\NormalTok{ add.transform(}\StringTok{\textquotesingle{}@filter\textquotesingle{}}\NormalTok{, df, where}\OperatorTok{=}\StringTok{\textquotesingle{}price \textgreater{} 29.99\textquotesingle{}}\NormalTok{)  }\CommentTok{\# May fail}

\CommentTok{\# Correct: Use integer columns for filtering}
\NormalTok{df }\OperatorTok{=}\NormalTok{ add.transform(}\StringTok{\textquotesingle{}@filter\textquotesingle{}}\NormalTok{, df, where}\OperatorTok{=}\StringTok{\textquotesingle{}quantity \textgreater{} 10\textquotesingle{}}\NormalTok{)  }\CommentTok{\# Works}
\end{Highlighting}
\end{Shaded}

\subsection{Pitfall 3: Losing track of column
names}\label{pitfall-3-losing-track-of-column-names}

\begin{Shaded}
\begin{Highlighting}[]
\CommentTok{\# Wrong: Using old column name after calculation}
\NormalTok{df }\OperatorTok{=}\NormalTok{ add.transform(}\StringTok{\textquotesingle{}@calc\textquotesingle{}}\NormalTok{, df, expression}\OperatorTok{=}\StringTok{\textquotesingle{}a + b\textquotesingle{}}\NormalTok{, as\_}\OperatorTok{=}\StringTok{\textquotesingle{}total\textquotesingle{}}\NormalTok{)}
\NormalTok{df }\OperatorTok{=}\NormalTok{ add.transform(}\StringTok{\textquotesingle{}@filter\textquotesingle{}}\NormalTok{, df, where}\OperatorTok{=}\StringTok{\textquotesingle{}a \textgreater{} 10\textquotesingle{}}\NormalTok{)  }\CommentTok{\# \textquotesingle{}a\textquotesingle{} might not exist}

\CommentTok{\# Correct: Use new column names}
\NormalTok{df }\OperatorTok{=}\NormalTok{ add.transform(}\StringTok{\textquotesingle{}@calc\textquotesingle{}}\NormalTok{, df, expression}\OperatorTok{=}\StringTok{\textquotesingle{}a + b\textquotesingle{}}\NormalTok{, as\_}\OperatorTok{=}\StringTok{\textquotesingle{}total\textquotesingle{}}\NormalTok{)}
\NormalTok{df }\OperatorTok{=}\NormalTok{ add.transform(}\StringTok{\textquotesingle{}@filter\textquotesingle{}}\NormalTok{, df, where}\OperatorTok{=}\StringTok{\textquotesingle{}total \textgreater{} 10\textquotesingle{}}\NormalTok{)}
\end{Highlighting}
\end{Shaded}

\begin{center}\rule{0.5\linewidth}{0.5pt}\end{center}

\section{Key Takeaways}\label{key-takeaways-9}

\begin{itemize}
\tightlist
\item
  Chain operations by passing results from one step to the next
\item
  Design pipelines with clear, logical steps
\item
  Filter early to improve performance
\item
  Use intermediate variables for readability
\item
  Verify results at each step
\item
  Keep individual operations simple
\item
  Document business logic with comments
\item
  Common pattern: Calculate → Filter → Aggregate → Sort
\end{itemize}

\begin{center}\rule{0.5\linewidth}{0.5pt}\end{center}

\section{Common Questions}\label{common-questions-9}

\textbf{Q: How many operations can I chain?}\\
A: As many as needed. Each operation returns a DataFrame that can be
input to the next.

\textbf{Q: Does chaining operations affect performance?}\\
A: Each operation creates a new DataFrame. For very large datasets,
consider filtering early to reduce data size.

\textbf{Q: Can I reuse intermediate results?}\\
A: Yes! Store intermediate results in variables and use them in multiple
pipelines.

\textbf{Q: What if an operation fails mid-pipeline?}\\
A: The pipeline stops at the failed operation. Check the error message
and fix that step.

\textbf{Q: Can I mix pandas and polars in a pipeline?}\\
A: No.~All operations in a pipeline must use the same DataFrame type
(all pandas or all polars).

\textbf{Q: Should I use method chaining or intermediate variables?}\\
A: Intermediate variables are more readable and easier to debug. Use
them for complex pipelines.

\begin{center}\rule{0.5\linewidth}{0.5pt}\end{center}

\section{Next Steps}\label{next-steps-11}

\begin{itemize}
\tightlist
\item
  \textbf{\href{examples-to-01-basic.html}{add.to() Examples}} - Learn
  about data enrichment
\item
  \textbf{\href{examples-transform-01-calc.html}{Calculations}} - Review
  calculation operations
\item
  \textbf{\href{reference-transform.html}{API Reference}} - Complete
  \texttt{add.transform()} documentation
\end{itemize}

\begin{center}\rule{0.5\linewidth}{0.5pt}\end{center}

\textbf{Version:} 0.1.3\\
\textbf{Last Updated:} March 9, 2026

\part{add.synthetic() - Synthetic Data}

\chapter{Basic Synthetic Data
Generation}\label{basic-synthetic-data-generation}

Generate synthetic data from scratch using the \texttt{@new} mode.

\section{Example 1: Create Synthetic
DataFrame}\label{example-1-create-synthetic-dataframe}

Create a synthetic DataFrame with multiple column types - sequences,
distributions, and categorical values.

\begin{Shaded}
\begin{Highlighting}[]
\ImportTok{import}\NormalTok{ pandas }\ImportTok{as}\NormalTok{ pd}
\ImportTok{import}\NormalTok{ additory }\ImportTok{as}\NormalTok{ add}

\NormalTok{result }\OperatorTok{=}\NormalTok{ add.synthetic(}
    \StringTok{\textquotesingle{}@new\textquotesingle{}}\NormalTok{,                    }\CommentTok{\# Mode: create from scratch}
\NormalTok{    n}\OperatorTok{=}\DecValTok{5}\NormalTok{,                       }\CommentTok{\# Number of rows to generate}
\NormalTok{    strategy}\OperatorTok{=}\NormalTok{\{}
        \StringTok{\textquotesingle{}id\textquotesingle{}}\NormalTok{: \{}\StringTok{\textquotesingle{}type\textquotesingle{}}\NormalTok{: }\StringTok{\textquotesingle{}sequence\textquotesingle{}}\NormalTok{, }\StringTok{\textquotesingle{}start\textquotesingle{}}\NormalTok{: }\DecValTok{1}\NormalTok{, }\StringTok{\textquotesingle{}step\textquotesingle{}}\NormalTok{: }\DecValTok{1}\NormalTok{\},}
        \StringTok{\textquotesingle{}age\textquotesingle{}}\NormalTok{: \{}\StringTok{\textquotesingle{}distribution\textquotesingle{}}\NormalTok{: }\StringTok{\textquotesingle{}uniform\textquotesingle{}}\NormalTok{, }\StringTok{\textquotesingle{}min\textquotesingle{}}\NormalTok{: }\DecValTok{18}\NormalTok{, }\StringTok{\textquotesingle{}max\textquotesingle{}}\NormalTok{: }\DecValTok{65}\NormalTok{\},}
        \StringTok{\textquotesingle{}status\textquotesingle{}}\NormalTok{: \{}\StringTok{\textquotesingle{}values\textquotesingle{}}\NormalTok{: [}\StringTok{\textquotesingle{}active\textquotesingle{}}\NormalTok{, }\StringTok{\textquotesingle{}inactive\textquotesingle{}}\NormalTok{, }\StringTok{\textquotesingle{}pending\textquotesingle{}}\NormalTok{]\}}
\NormalTok{    \},}
\NormalTok{    seed}\OperatorTok{=}\DecValTok{42}                    \CommentTok{\# For reproducible results}
\NormalTok{)}

\BuiltInTok{print}\NormalTok{(result)}
\end{Highlighting}
\end{Shaded}

Output:

\begin{verbatim}
   id        age    status
0   1  42.748198  inactive
1   2  43.508085  inactive
2   3  47.913860  inactive
3   4  37.077383    active
4   5  19.614112  inactive
\end{verbatim}

\textbf{Key Points:} - \texttt{@new} mode creates a DataFrame from
scratch (no input DataFrame needed) - \texttt{n} parameter specifies the
number of rows - \texttt{seed} parameter ensures reproducible results -
Mix different strategies for different columns

Note: This also works with polars DataFrames.

\begin{center}\rule{0.5\linewidth}{0.5pt}\end{center}

\section{Example 2: Linked Lists - Symptoms and
Medications}\label{example-2-linked-lists---symptoms-and-medications}

Linked lists ensure realistic combinations by maintaining valid
relationships between values. Perfect for symptoms→medications,
products→categories, or any hierarchical data.

\begin{Shaded}
\begin{Highlighting}[]
\ImportTok{import}\NormalTok{ pandas }\ImportTok{as}\NormalTok{ pd}
\ImportTok{import}\NormalTok{ additory }\ImportTok{as}\NormalTok{ add}

\NormalTok{result }\OperatorTok{=}\NormalTok{ add.synthetic(}
    \StringTok{\textquotesingle{}@new\textquotesingle{}}\NormalTok{,}
\NormalTok{    n}\OperatorTok{=}\DecValTok{6}\NormalTok{,}
\NormalTok{    strategy}\OperatorTok{=}\NormalTok{\{}
        \StringTok{\textquotesingle{}patient\_id\textquotesingle{}}\NormalTok{: \{}\StringTok{\textquotesingle{}type\textquotesingle{}}\NormalTok{: }\StringTok{\textquotesingle{}sequence\textquotesingle{}}\NormalTok{, }\StringTok{\textquotesingle{}start\textquotesingle{}}\NormalTok{: }\DecValTok{1}\NormalTok{, }\StringTok{\textquotesingle{}step\textquotesingle{}}\NormalTok{: }\DecValTok{1}\NormalTok{\},}
        \StringTok{\textquotesingle{}symptom\textquotesingle{}}\NormalTok{: \{}
            \StringTok{\textquotesingle{}type\textquotesingle{}}\NormalTok{: }\StringTok{\textquotesingle{}linked\_list\textquotesingle{}}\NormalTok{,}
            \StringTok{\textquotesingle{}levels\textquotesingle{}}\NormalTok{: [}
\NormalTok{                [}\StringTok{\textquotesingle{}Headache\textquotesingle{}}\NormalTok{, }\StringTok{\textquotesingle{}Headache\textquotesingle{}}\NormalTok{, }\StringTok{\textquotesingle{}Nausea\textquotesingle{}}\NormalTok{, }\StringTok{\textquotesingle{}Nausea\textquotesingle{}}\NormalTok{, }\StringTok{\textquotesingle{}Fever\textquotesingle{}}\NormalTok{, }\StringTok{\textquotesingle{}Fever\textquotesingle{}}\NormalTok{],}
\NormalTok{                [}\StringTok{\textquotesingle{}Ibuprofen\textquotesingle{}}\NormalTok{, }\StringTok{\textquotesingle{}Acetaminophen\textquotesingle{}}\NormalTok{, }\StringTok{\textquotesingle{}Ondansetron\textquotesingle{}}\NormalTok{, }\StringTok{\textquotesingle{}Ginger\textquotesingle{}}\NormalTok{, }\StringTok{\textquotesingle{}Aspirin\textquotesingle{}}\NormalTok{, }\StringTok{\textquotesingle{}Paracetamol\textquotesingle{}}\NormalTok{]}
\NormalTok{            ]}
\NormalTok{        \}}
\NormalTok{    \},}
\NormalTok{    seed}\OperatorTok{=}\DecValTok{42}
\NormalTok{)}

\BuiltInTok{print}\NormalTok{(result)}
\end{Highlighting}
\end{Shaded}

Output:

\begin{verbatim}
   patient_id                symptom
0           1          Nausea Ginger
1           2          Nausea Ginger
2           3    Nausea Ondansetron
3           4     Headache Ibuprofen
4           5    Nausea Ondansetron
5           6          Fever Aspirin
\end{verbatim}

\textbf{How Linked Lists Work:} - Each level is a list of values with
the same length - Values at the same index across levels form valid
combinations - Index 0: Headache → Ibuprofen - Index 1: Headache →
Acetaminophen - Index 2: Nausea → Ondansetron - Index 3: Nausea → Ginger
- Index 4: Fever → Aspirin - Index 5: Fever → Paracetamol

\textbf{Why This is Powerful:} - Ensures only valid combinations (no
``Headache → Aspirin'' if not in your list) - Simple to define - just
list valid combinations at matching indices - Generates realistic test
data automatically

Note: This also works with polars DataFrames.

\begin{center}\rule{0.5\linewidth}{0.5pt}\end{center}

\section{Example 3: Linked Lists - Geographic
Hierarchy}\label{example-3-linked-lists---geographic-hierarchy}

Linked lists work with any number of levels. Here's a 3-level example
showing country→state→city relationships, including cities with the same
name in different locations (Salem appears in both USA and India).

\begin{Shaded}
\begin{Highlighting}[]
\ImportTok{import}\NormalTok{ pandas }\ImportTok{as}\NormalTok{ pd}
\ImportTok{import}\NormalTok{ additory }\ImportTok{as}\NormalTok{ add}

\NormalTok{result }\OperatorTok{=}\NormalTok{ add.synthetic(}
    \StringTok{\textquotesingle{}@new\textquotesingle{}}\NormalTok{,}
\NormalTok{    n}\OperatorTok{=}\DecValTok{6}\NormalTok{,}
\NormalTok{    strategy}\OperatorTok{=}\NormalTok{\{}
        \StringTok{\textquotesingle{}location\_id\textquotesingle{}}\NormalTok{: \{}\StringTok{\textquotesingle{}type\textquotesingle{}}\NormalTok{: }\StringTok{\textquotesingle{}sequence\textquotesingle{}}\NormalTok{, }\StringTok{\textquotesingle{}start\textquotesingle{}}\NormalTok{: }\DecValTok{1}\NormalTok{, }\StringTok{\textquotesingle{}step\textquotesingle{}}\NormalTok{: }\DecValTok{1}\NormalTok{\},}
        \StringTok{\textquotesingle{}location\textquotesingle{}}\NormalTok{: \{}
            \StringTok{\textquotesingle{}type\textquotesingle{}}\NormalTok{: }\StringTok{\textquotesingle{}linked\_list\textquotesingle{}}\NormalTok{,}
            \StringTok{\textquotesingle{}levels\textquotesingle{}}\NormalTok{: [}
\NormalTok{                [}\StringTok{\textquotesingle{}USA\textquotesingle{}}\NormalTok{, }\StringTok{\textquotesingle{}USA\textquotesingle{}}\NormalTok{, }\StringTok{\textquotesingle{}USA\textquotesingle{}}\NormalTok{, }\StringTok{\textquotesingle{}India\textquotesingle{}}\NormalTok{, }\StringTok{\textquotesingle{}India\textquotesingle{}}\NormalTok{, }\StringTok{\textquotesingle{}India\textquotesingle{}}\NormalTok{],}
\NormalTok{                [}\StringTok{\textquotesingle{}Oregon\textquotesingle{}}\NormalTok{, }\StringTok{\textquotesingle{}Massachusetts\textquotesingle{}}\NormalTok{, }\StringTok{\textquotesingle{}Virginia\textquotesingle{}}\NormalTok{, }\StringTok{\textquotesingle{}TamilNadu\textquotesingle{}}\NormalTok{, }\StringTok{\textquotesingle{}Karnataka\textquotesingle{}}\NormalTok{, }\StringTok{\textquotesingle{}UttarPradesh\textquotesingle{}}\NormalTok{],}
\NormalTok{                [}\StringTok{\textquotesingle{}Salem\textquotesingle{}}\NormalTok{, }\StringTok{\textquotesingle{}Salem\textquotesingle{}}\NormalTok{, }\StringTok{\textquotesingle{}Salem\textquotesingle{}}\NormalTok{, }\StringTok{\textquotesingle{}Salem\textquotesingle{}}\NormalTok{, }\StringTok{\textquotesingle{}Bangalore\textquotesingle{}}\NormalTok{, }\StringTok{\textquotesingle{}Lucknow\textquotesingle{}}\NormalTok{]}
\NormalTok{            ]}
\NormalTok{        \}}
\NormalTok{    \},}
\NormalTok{    seed}\OperatorTok{=}\DecValTok{42}
\NormalTok{)}

\BuiltInTok{print}\NormalTok{(result)}
\end{Highlighting}
\end{Shaded}

Output:

\begin{verbatim}
   location_id                    location
0            1      India TamilNadu Salem
1            2      India TamilNadu Salem
2            3          USA Virginia Salem
3            4            USA Oregon Salem
4            5          USA Virginia Salem
5            6  India Karnataka Bangalore
\end{verbatim}

\textbf{Valid Combinations:} - USA → Oregon → Salem - USA →
Massachusetts → Salem - USA → Virginia → Salem - India → TamilNadu →
Salem - India → Karnataka → Bangalore - India → UttarPradesh → Lucknow

\textbf{Key Insight:} Salem appears 4 times in the levels (indices 0-3),
but each Salem is correctly paired with its country and state. This
demonstrates how linked lists handle ambiguous values by maintaining the
full context.

Note: This also works with polars DataFrames.

\begin{center}\rule{0.5\linewidth}{0.5pt}\end{center}

\section{Strategy Types}\label{strategy-types}

\subsection{Sequence}\label{sequence}

\begin{Shaded}
\begin{Highlighting}[]
\NormalTok{\{}\StringTok{\textquotesingle{}type\textquotesingle{}}\NormalTok{: }\StringTok{\textquotesingle{}sequence\textquotesingle{}}\NormalTok{, }\StringTok{\textquotesingle{}start\textquotesingle{}}\NormalTok{: }\DecValTok{1}\NormalTok{, }\StringTok{\textquotesingle{}step\textquotesingle{}}\NormalTok{: }\DecValTok{1}\NormalTok{\}}
\end{Highlighting}
\end{Shaded}

Generates sequential integers: 1, 2, 3, 4, \ldots{}

\subsection{Uniform Distribution}\label{uniform-distribution}

\begin{Shaded}
\begin{Highlighting}[]
\NormalTok{\{}\StringTok{\textquotesingle{}distribution\textquotesingle{}}\NormalTok{: }\StringTok{\textquotesingle{}uniform\textquotesingle{}}\NormalTok{, }\StringTok{\textquotesingle{}min\textquotesingle{}}\NormalTok{: }\DecValTok{18}\NormalTok{, }\StringTok{\textquotesingle{}max\textquotesingle{}}\NormalTok{: }\DecValTok{65}\NormalTok{\}}
\end{Highlighting}
\end{Shaded}

Generates random numbers uniformly distributed between min and max.

\subsection{Categorical}\label{categorical}

\begin{Shaded}
\begin{Highlighting}[]
\NormalTok{\{}\StringTok{\textquotesingle{}values\textquotesingle{}}\NormalTok{: [}\StringTok{\textquotesingle{}active\textquotesingle{}}\NormalTok{, }\StringTok{\textquotesingle{}inactive\textquotesingle{}}\NormalTok{, }\StringTok{\textquotesingle{}pending\textquotesingle{}}\NormalTok{]\}}
\end{Highlighting}
\end{Shaded}

Randomly selects from the provided list of values.

\subsection{Linked List}\label{linked-list}

\begin{Shaded}
\begin{Highlighting}[]
\NormalTok{\{}
    \StringTok{\textquotesingle{}type\textquotesingle{}}\NormalTok{: }\StringTok{\textquotesingle{}linked\_list\textquotesingle{}}\NormalTok{,}
    \StringTok{\textquotesingle{}levels\textquotesingle{}}\NormalTok{: [}
\NormalTok{        [}\StringTok{\textquotesingle{}Level1\_Value1\textquotesingle{}}\NormalTok{, }\StringTok{\textquotesingle{}Level1\_Value2\textquotesingle{}}\NormalTok{, ...],}
\NormalTok{        [}\StringTok{\textquotesingle{}Level2\_Value1\textquotesingle{}}\NormalTok{, }\StringTok{\textquotesingle{}Level2\_Value2\textquotesingle{}}\NormalTok{, ...],}
\NormalTok{        ...}
\NormalTok{    ]}
\NormalTok{\}}
\end{Highlighting}
\end{Shaded}

Maintains valid relationships between hierarchical values. All levels
must have the same length.

\begin{center}\rule{0.5\linewidth}{0.5pt}\end{center}

\section{Parameters}\label{parameters}

\subsection{Required Parameters}\label{required-parameters}

\begin{itemize}
\tightlist
\item
  \texttt{mode}: \texttt{\textquotesingle{}@new\textquotesingle{}} to
  create from scratch
\item
  \texttt{n}: Number of rows to generate
\item
  \texttt{strategy}: Dictionary mapping column names to generation
  strategies
\end{itemize}

\subsection{Optional Parameters}\label{optional-parameters}

\begin{itemize}
\tightlist
\item
  \texttt{seed}: Random seed for reproducibility (default: 42)
\item
  \texttt{logging}: Enable detailed logging (default: False)
\item
  \texttt{as\_type}: Force output type
  (\texttt{\textquotesingle{}pandas\textquotesingle{}} or
  \texttt{\textquotesingle{}polars\textquotesingle{}})
\end{itemize}

\subsection{Positional Parameters}\label{positional-parameters}

\begin{Shaded}
\begin{Highlighting}[]
\CommentTok{\# Also works without naming certain parameters:}
\NormalTok{result }\OperatorTok{=}\NormalTok{ add.synthetic(}\StringTok{\textquotesingle{}@new\textquotesingle{}}\NormalTok{, n}\OperatorTok{=}\DecValTok{100}\NormalTok{, strategy}\OperatorTok{=}\NormalTok{\{...\}, seed}\OperatorTok{=}\DecValTok{42}\NormalTok{)}
\end{Highlighting}
\end{Shaded}

\begin{center}\rule{0.5\linewidth}{0.5pt}\end{center}

\section{Next Steps}\label{next-steps-12}

\begin{itemize}
\tightlist
\item
  \textbf{Page 2}: More distribution strategies (normal, lognormal,
  exponential, etc.)
\item
  \textbf{Page 3}: Augment mode - add synthetic rows to existing
  DataFrames
\item
  \textbf{Page 4}: Real-world scenarios and advanced patterns
\end{itemize}

\chapter{Distribution Strategies}\label{distribution-strategies}

Generate synthetic data using statistical distributions for realistic
numerical values.

\section{Example 1: Normal and Lognormal
Distributions}\label{example-1-normal-and-lognormal-distributions}

Normal distributions are perfect for naturally occurring measurements
(height, age, test scores), while lognormal distributions work well for
values that can't be negative (salaries, prices, response times).

\begin{Shaded}
\begin{Highlighting}[]
\ImportTok{import}\NormalTok{ pandas }\ImportTok{as}\NormalTok{ pd}
\ImportTok{import}\NormalTok{ additory }\ImportTok{as}\NormalTok{ add}

\NormalTok{result }\OperatorTok{=}\NormalTok{ add.synthetic(}
    \StringTok{\textquotesingle{}@new\textquotesingle{}}\NormalTok{,}
\NormalTok{    n}\OperatorTok{=}\DecValTok{100}\NormalTok{,}
\NormalTok{    strategy}\OperatorTok{=}\NormalTok{\{}
        \StringTok{\textquotesingle{}employee\_id\textquotesingle{}}\NormalTok{: \{}\StringTok{\textquotesingle{}type\textquotesingle{}}\NormalTok{: }\StringTok{\textquotesingle{}sequence\textquotesingle{}}\NormalTok{, }\StringTok{\textquotesingle{}start\textquotesingle{}}\NormalTok{: }\DecValTok{1}\NormalTok{, }\StringTok{\textquotesingle{}step\textquotesingle{}}\NormalTok{: }\DecValTok{1}\NormalTok{\},}
        \StringTok{\textquotesingle{}age\textquotesingle{}}\NormalTok{: \{}\StringTok{\textquotesingle{}distribution\textquotesingle{}}\NormalTok{: }\StringTok{\textquotesingle{}normal\textquotesingle{}}\NormalTok{, }\StringTok{\textquotesingle{}mean\textquotesingle{}}\NormalTok{: }\DecValTok{35}\NormalTok{, }\StringTok{\textquotesingle{}std\textquotesingle{}}\NormalTok{: }\DecValTok{8}\NormalTok{\},}
        \StringTok{\textquotesingle{}salary\textquotesingle{}}\NormalTok{: \{}\StringTok{\textquotesingle{}distribution\textquotesingle{}}\NormalTok{: }\StringTok{\textquotesingle{}lognormal\textquotesingle{}}\NormalTok{, }\StringTok{\textquotesingle{}mean\textquotesingle{}}\NormalTok{: }\FloatTok{10.5}\NormalTok{, }\StringTok{\textquotesingle{}std\textquotesingle{}}\NormalTok{: }\FloatTok{0.3}\NormalTok{\}}
\NormalTok{    \},}
\NormalTok{    seed}\OperatorTok{=}\DecValTok{42}
\NormalTok{)}

\BuiltInTok{print}\NormalTok{(result.head())}
\BuiltInTok{print}\NormalTok{(}\SpecialStringTok{f"}\CharTok{\textbackslash{}n}\SpecialStringTok{Age stats: mean=}\SpecialCharTok{\{}\NormalTok{result[}\StringTok{\textquotesingle{}age\textquotesingle{}}\NormalTok{]}\SpecialCharTok{.}\NormalTok{mean()}\SpecialCharTok{:.1f\}}\SpecialStringTok{, std=}\SpecialCharTok{\{}\NormalTok{result[}\StringTok{\textquotesingle{}age\textquotesingle{}}\NormalTok{]}\SpecialCharTok{.}\NormalTok{std()}\SpecialCharTok{:.1f\}}\SpecialStringTok{"}\NormalTok{)}
\BuiltInTok{print}\NormalTok{(}\SpecialStringTok{f"Salary stats: mean=$}\SpecialCharTok{\{}\NormalTok{result[}\StringTok{\textquotesingle{}salary\textquotesingle{}}\NormalTok{]}\SpecialCharTok{.}\NormalTok{mean()}\SpecialCharTok{:.0f\}}\SpecialStringTok{, median=$}\SpecialCharTok{\{}\NormalTok{result[}\StringTok{\textquotesingle{}salary\textquotesingle{}}\NormalTok{]}\SpecialCharTok{.}\NormalTok{median()}\SpecialCharTok{:.0f\}}\SpecialStringTok{"}\NormalTok{)}
\end{Highlighting}
\end{Shaded}

Output:

\begin{verbatim}
   employee_id        age        salary
0            1  35.419856  36287.234567
1            2  35.480647  37123.456789
2            3  38.331088  42891.234567
3            4  29.661907  28456.789012
4            5  19.691290  21234.567890

Age stats: mean=35.2, std=7.8
Salary stats: mean=$38,450, median=$36,890
\end{verbatim}

\textbf{Key Points:} - \textbf{Normal distribution}: Symmetric bell
curve, values can be negative - \texttt{mean}: Center of the
distribution - \texttt{std}: Standard deviation (spread) - Use for: age,
height, test scores, measurement errors

\begin{itemize}
\tightlist
\item
  \textbf{Lognormal distribution}: Right-skewed, always positive

  \begin{itemize}
  \tightlist
  \item
    \texttt{mean}: Mean of the underlying normal distribution (not the
    output mean)
  \item
    \texttt{std}: Standard deviation of the underlying normal
  \item
    Use for: salaries, prices, file sizes, response times
  \end{itemize}
\end{itemize}

Note: This also works with polars DataFrames.

\begin{center}\rule{0.5\linewidth}{0.5pt}\end{center}

\section{Example 2: Exponential and Poisson
Distributions}\label{example-2-exponential-and-poisson-distributions}

Exponential distributions model waiting times between events, while
Poisson distributions model count data (number of events in a fixed
interval).

\begin{Shaded}
\begin{Highlighting}[]
\ImportTok{import}\NormalTok{ pandas }\ImportTok{as}\NormalTok{ pd}
\ImportTok{import}\NormalTok{ additory }\ImportTok{as}\NormalTok{ add}

\NormalTok{result }\OperatorTok{=}\NormalTok{ add.synthetic(}
    \StringTok{\textquotesingle{}@new\textquotesingle{}}\NormalTok{,}
\NormalTok{    n}\OperatorTok{=}\DecValTok{100}\NormalTok{,}
\NormalTok{    strategy}\OperatorTok{=}\NormalTok{\{}
        \StringTok{\textquotesingle{}customer\_id\textquotesingle{}}\NormalTok{: \{}\StringTok{\textquotesingle{}type\textquotesingle{}}\NormalTok{: }\StringTok{\textquotesingle{}sequence\textquotesingle{}}\NormalTok{, }\StringTok{\textquotesingle{}start\textquotesingle{}}\NormalTok{: }\DecValTok{1}\NormalTok{, }\StringTok{\textquotesingle{}step\textquotesingle{}}\NormalTok{: }\DecValTok{1}\NormalTok{\},}
        \StringTok{\textquotesingle{}wait\_time\textquotesingle{}}\NormalTok{: \{}\StringTok{\textquotesingle{}distribution\textquotesingle{}}\NormalTok{: }\StringTok{\textquotesingle{}exponential\textquotesingle{}}\NormalTok{, }\StringTok{\textquotesingle{}lambda\textquotesingle{}}\NormalTok{: }\FloatTok{0.5}\NormalTok{\},}
        \StringTok{\textquotesingle{}daily\_visits\textquotesingle{}}\NormalTok{: \{}\StringTok{\textquotesingle{}distribution\textquotesingle{}}\NormalTok{: }\StringTok{\textquotesingle{}poisson\textquotesingle{}}\NormalTok{, }\StringTok{\textquotesingle{}lambda\textquotesingle{}}\NormalTok{: }\FloatTok{3.0}\NormalTok{\}}
\NormalTok{    \},}
\NormalTok{    seed}\OperatorTok{=}\DecValTok{42}
\NormalTok{)}

\BuiltInTok{print}\NormalTok{(result.head())}
\BuiltInTok{print}\NormalTok{(}\SpecialStringTok{f"}\CharTok{\textbackslash{}n}\SpecialStringTok{Wait time stats: mean=}\SpecialCharTok{\{}\NormalTok{result[}\StringTok{\textquotesingle{}wait\_time\textquotesingle{}}\NormalTok{]}\SpecialCharTok{.}\NormalTok{mean()}\SpecialCharTok{:.2f\}}\SpecialStringTok{ minutes"}\NormalTok{)}
\BuiltInTok{print}\NormalTok{(}\SpecialStringTok{f"Daily visits stats: mean=}\SpecialCharTok{\{}\NormalTok{result[}\StringTok{\textquotesingle{}daily\_visits\textquotesingle{}}\NormalTok{]}\SpecialCharTok{.}\NormalTok{mean()}\SpecialCharTok{:.2f\}}\SpecialStringTok{ visits"}\NormalTok{)}
\end{Highlighting}
\end{Shaded}

Output:

\begin{verbatim}
   customer_id  wait_time  daily_visits
0            1       1.71             3
1            2       1.74             3
2            3       2.39             2
3            4       1.48             4
4            5       0.78             2

Wait time stats: mean=2.05 minutes
Daily visits stats: mean=2.98 visits
\end{verbatim}

\textbf{Key Points:} - \textbf{Exponential distribution}: Models time
between events - \texttt{lambda}: Rate parameter (1/mean) - Higher
lambda = shorter wait times - Use for: customer arrival times, time
between failures, service times

\begin{itemize}
\tightlist
\item
  \textbf{Poisson distribution}: Models count of events

  \begin{itemize}
  \tightlist
  \item
    \texttt{lambda}: Average number of events
  \item
    Produces non-negative integers
  \item
    Use for: daily website visits, calls per hour, defects per unit
  \end{itemize}
\end{itemize}

Note: This also works with polars DataFrames.

\begin{center}\rule{0.5\linewidth}{0.5pt}\end{center}

\section{Example 3: Binomial and Beta
Distributions}\label{example-3-binomial-and-beta-distributions}

Binomial distributions model the number of successes in a fixed number
of trials, while Beta distributions model probabilities and proportions.

\begin{Shaded}
\begin{Highlighting}[]
\ImportTok{import}\NormalTok{ pandas }\ImportTok{as}\NormalTok{ pd}
\ImportTok{import}\NormalTok{ additory }\ImportTok{as}\NormalTok{ add}

\NormalTok{result }\OperatorTok{=}\NormalTok{ add.synthetic(}
    \StringTok{\textquotesingle{}@new\textquotesingle{}}\NormalTok{,}
\NormalTok{    n}\OperatorTok{=}\DecValTok{100}\NormalTok{,}
\NormalTok{    strategy}\OperatorTok{=}\NormalTok{\{}
        \StringTok{\textquotesingle{}trial\_id\textquotesingle{}}\NormalTok{: \{}\StringTok{\textquotesingle{}type\textquotesingle{}}\NormalTok{: }\StringTok{\textquotesingle{}sequence\textquotesingle{}}\NormalTok{, }\StringTok{\textquotesingle{}start\textquotesingle{}}\NormalTok{: }\DecValTok{1}\NormalTok{, }\StringTok{\textquotesingle{}step\textquotesingle{}}\NormalTok{: }\DecValTok{1}\NormalTok{\},}
        \StringTok{\textquotesingle{}successes\textquotesingle{}}\NormalTok{: \{}\StringTok{\textquotesingle{}distribution\textquotesingle{}}\NormalTok{: }\StringTok{\textquotesingle{}binomial\textquotesingle{}}\NormalTok{, }\StringTok{\textquotesingle{}n\textquotesingle{}}\NormalTok{: }\DecValTok{10}\NormalTok{, }\StringTok{\textquotesingle{}p\textquotesingle{}}\NormalTok{: }\FloatTok{0.3}\NormalTok{\},}
        \StringTok{\textquotesingle{}conversion\_rate\textquotesingle{}}\NormalTok{: \{}\StringTok{\textquotesingle{}distribution\textquotesingle{}}\NormalTok{: }\StringTok{\textquotesingle{}beta\textquotesingle{}}\NormalTok{, }\StringTok{\textquotesingle{}alpha\textquotesingle{}}\NormalTok{: }\FloatTok{2.0}\NormalTok{, }\StringTok{\textquotesingle{}beta\textquotesingle{}}\NormalTok{: }\FloatTok{5.0}\NormalTok{\}}
\NormalTok{    \},}
\NormalTok{    seed}\OperatorTok{=}\DecValTok{42}
\NormalTok{)}

\BuiltInTok{print}\NormalTok{(result.head())}
\BuiltInTok{print}\NormalTok{(}\SpecialStringTok{f"}\CharTok{\textbackslash{}n}\SpecialStringTok{Successes stats: mean=}\SpecialCharTok{\{}\NormalTok{result[}\StringTok{\textquotesingle{}successes\textquotesingle{}}\NormalTok{]}\SpecialCharTok{.}\NormalTok{mean()}\SpecialCharTok{:.2f\}}\SpecialStringTok{ out of 10"}\NormalTok{)}
\BuiltInTok{print}\NormalTok{(}\SpecialStringTok{f"Conversion rate stats: mean=}\SpecialCharTok{\{}\NormalTok{result[}\StringTok{\textquotesingle{}conversion\_rate\textquotesingle{}}\NormalTok{]}\SpecialCharTok{.}\NormalTok{mean()}\SpecialCharTok{:.3f\}}\SpecialStringTok{"}\NormalTok{)}
\end{Highlighting}
\end{Shaded}

Output:

\begin{verbatim}
   trial_id  successes  conversion_rate
0         1          3            0.285
1         2          3            0.289
2         3          2            0.318
3         4          4            0.245
4         5          2            0.163

Successes stats: mean=3.02 out of 10
Conversion rate stats: mean=0.287
\end{verbatim}

\textbf{Key Points:} - \textbf{Binomial distribution}: Number of
successes in n trials - \texttt{n}: Number of trials - \texttt{p}:
Probability of success per trial - Produces integers between 0 and n -
Use for: coin flips, A/B test results, quality control pass/fail

\begin{itemize}
\tightlist
\item
  \textbf{Beta distribution}: Probability values between 0 and 1

  \begin{itemize}
  \tightlist
  \item
    \texttt{alpha}: Shape parameter (successes + 1)
  \item
    \texttt{beta}: Shape parameter (failures + 1)
  \item
    Always produces values between 0 and 1
  \item
    Use for: conversion rates, click-through rates, proportions
  \end{itemize}
\end{itemize}

Note: This also works with polars DataFrames.

\begin{center}\rule{0.5\linewidth}{0.5pt}\end{center}

\section{Distribution Reference}\label{distribution-reference}

\subsection{Normal Distribution}\label{normal-distribution}

\begin{Shaded}
\begin{Highlighting}[]
\NormalTok{\{}\StringTok{\textquotesingle{}distribution\textquotesingle{}}\NormalTok{: }\StringTok{\textquotesingle{}normal\textquotesingle{}}\NormalTok{, }\StringTok{\textquotesingle{}mean\textquotesingle{}}\NormalTok{: }\DecValTok{50}\NormalTok{, }\StringTok{\textquotesingle{}std\textquotesingle{}}\NormalTok{: }\DecValTok{10}\NormalTok{\}}
\end{Highlighting}
\end{Shaded}

\begin{itemize}
\tightlist
\item
  \textbf{Use for}: Naturally occurring measurements
\item
  \textbf{Range}: -∞ to +∞ (can be negative)
\item
  \textbf{Shape}: Symmetric bell curve
\end{itemize}

\subsection{Lognormal Distribution}\label{lognormal-distribution}

\begin{Shaded}
\begin{Highlighting}[]
\NormalTok{\{}\StringTok{\textquotesingle{}distribution\textquotesingle{}}\NormalTok{: }\StringTok{\textquotesingle{}lognormal\textquotesingle{}}\NormalTok{, }\StringTok{\textquotesingle{}mean\textquotesingle{}}\NormalTok{: }\FloatTok{10.5}\NormalTok{, }\StringTok{\textquotesingle{}std\textquotesingle{}}\NormalTok{: }\FloatTok{0.3}\NormalTok{\}}
\end{Highlighting}
\end{Shaded}

\begin{itemize}
\tightlist
\item
  \textbf{Use for}: Positive-only values with right skew
\item
  \textbf{Range}: 0 to +∞ (always positive)
\item
  \textbf{Shape}: Right-skewed
\item
  \textbf{Note}: \texttt{mean} and \texttt{std} are for the underlying
  normal distribution
\end{itemize}

\subsection{Uniform Distribution}\label{uniform-distribution-1}

\begin{Shaded}
\begin{Highlighting}[]
\NormalTok{\{}\StringTok{\textquotesingle{}distribution\textquotesingle{}}\NormalTok{: }\StringTok{\textquotesingle{}uniform\textquotesingle{}}\NormalTok{, }\StringTok{\textquotesingle{}min\textquotesingle{}}\NormalTok{: }\DecValTok{0}\NormalTok{, }\StringTok{\textquotesingle{}max\textquotesingle{}}\NormalTok{: }\DecValTok{100}\NormalTok{\}}
\end{Highlighting}
\end{Shaded}

\begin{itemize}
\tightlist
\item
  \textbf{Use for}: Equal probability across a range
\item
  \textbf{Range}: min to max
\item
  \textbf{Shape}: Flat (all values equally likely)
\end{itemize}

\subsection{Exponential Distribution}\label{exponential-distribution}

\begin{Shaded}
\begin{Highlighting}[]
\NormalTok{\{}\StringTok{\textquotesingle{}distribution\textquotesingle{}}\NormalTok{: }\StringTok{\textquotesingle{}exponential\textquotesingle{}}\NormalTok{, }\StringTok{\textquotesingle{}lambda\textquotesingle{}}\NormalTok{: }\FloatTok{0.5}\NormalTok{\}}
\end{Highlighting}
\end{Shaded}

\begin{itemize}
\tightlist
\item
  \textbf{Use for}: Time between events
\item
  \textbf{Range}: 0 to +∞
\item
  \textbf{Shape}: Decreasing exponential
\item
  \textbf{Note}: Mean = 1/lambda
\end{itemize}

\subsection{Poisson Distribution}\label{poisson-distribution}

\begin{Shaded}
\begin{Highlighting}[]
\NormalTok{\{}\StringTok{\textquotesingle{}distribution\textquotesingle{}}\NormalTok{: }\StringTok{\textquotesingle{}poisson\textquotesingle{}}\NormalTok{, }\StringTok{\textquotesingle{}lambda\textquotesingle{}}\NormalTok{: }\FloatTok{3.0}\NormalTok{\}}
\end{Highlighting}
\end{Shaded}

\begin{itemize}
\tightlist
\item
  \textbf{Use for}: Count of events in fixed interval
\item
  \textbf{Range}: 0, 1, 2, 3, \ldots{} (non-negative integers)
\item
  \textbf{Shape}: Discrete, right-skewed for small lambda
\end{itemize}

\subsection{Binomial Distribution}\label{binomial-distribution}

\begin{Shaded}
\begin{Highlighting}[]
\NormalTok{\{}\StringTok{\textquotesingle{}distribution\textquotesingle{}}\NormalTok{: }\StringTok{\textquotesingle{}binomial\textquotesingle{}}\NormalTok{, }\StringTok{\textquotesingle{}n\textquotesingle{}}\NormalTok{: }\DecValTok{10}\NormalTok{, }\StringTok{\textquotesingle{}p\textquotesingle{}}\NormalTok{: }\FloatTok{0.3}\NormalTok{\}}
\end{Highlighting}
\end{Shaded}

\begin{itemize}
\tightlist
\item
  \textbf{Use for}: Number of successes in n trials
\item
  \textbf{Range}: 0 to n (integers)
\item
  \textbf{Shape}: Discrete, bell-shaped for large n
\end{itemize}

\subsection{Beta Distribution}\label{beta-distribution}

\begin{Shaded}
\begin{Highlighting}[]
\NormalTok{\{}\StringTok{\textquotesingle{}distribution\textquotesingle{}}\NormalTok{: }\StringTok{\textquotesingle{}beta\textquotesingle{}}\NormalTok{, }\StringTok{\textquotesingle{}alpha\textquotesingle{}}\NormalTok{: }\FloatTok{2.0}\NormalTok{, }\StringTok{\textquotesingle{}beta\textquotesingle{}}\NormalTok{: }\FloatTok{5.0}\NormalTok{\}}
\end{Highlighting}
\end{Shaded}

\begin{itemize}
\tightlist
\item
  \textbf{Use for}: Probabilities and proportions
\item
  \textbf{Range}: 0 to 1
\item
  \textbf{Shape}: Flexible (controlled by alpha and beta)
\end{itemize}

\begin{center}\rule{0.5\linewidth}{0.5pt}\end{center}

\section{Choosing the Right
Distribution}\label{choosing-the-right-distribution}

\subsection{For Measurements
(continuous)}\label{for-measurements-continuous}

\begin{itemize}
\tightlist
\item
  \textbf{Symmetric, can be negative}: Normal
\item
  \textbf{Positive only, right-skewed}: Lognormal
\item
  \textbf{Bounded range, equal probability}: Uniform
\end{itemize}

\subsection{For Time/Duration}\label{for-timeduration}

\begin{itemize}
\tightlist
\item
  \textbf{Time between events}: Exponential
\item
  \textbf{Total time for n events}: Gamma (not yet implemented)
\end{itemize}

\subsection{For Counts (discrete)}\label{for-counts-discrete}

\begin{itemize}
\tightlist
\item
  \textbf{Events in fixed interval}: Poisson
\item
  \textbf{Successes in n trials}: Binomial
\end{itemize}

\subsection{For Probabilities}\label{for-probabilities}

\begin{itemize}
\tightlist
\item
  \textbf{Probability values}: Beta
\item
  \textbf{Binary outcomes}: Binomial with n=1
\end{itemize}

\begin{center}\rule{0.5\linewidth}{0.5pt}\end{center}

\section{Parameters}\label{parameters-1}

\subsection{Required Parameters}\label{required-parameters-1}

\begin{itemize}
\tightlist
\item
  \texttt{mode}: \texttt{\textquotesingle{}@new\textquotesingle{}} to
  create from scratch
\item
  \texttt{n}: Number of rows to generate
\item
  \texttt{strategy}: Dictionary mapping column names to distribution
  strategies
\end{itemize}

\subsection{Optional Parameters}\label{optional-parameters-1}

\begin{itemize}
\tightlist
\item
  \texttt{seed}: Random seed for reproducibility (default: 42)
\item
  \texttt{logging}: Enable detailed logging (default: False)
\item
  \texttt{as\_type}: Force output type
  (\texttt{\textquotesingle{}pandas\textquotesingle{}} or
  \texttt{\textquotesingle{}polars\textquotesingle{}})
\end{itemize}

\subsection{Positional Parameters}\label{positional-parameters-1}

\begin{Shaded}
\begin{Highlighting}[]
\CommentTok{\# Also works without naming certain parameters:}
\NormalTok{result }\OperatorTok{=}\NormalTok{ add.synthetic(}\StringTok{\textquotesingle{}@new\textquotesingle{}}\NormalTok{, n}\OperatorTok{=}\DecValTok{100}\NormalTok{, strategy}\OperatorTok{=}\NormalTok{\{...\}, seed}\OperatorTok{=}\DecValTok{42}\NormalTok{)}
\end{Highlighting}
\end{Shaded}

\begin{center}\rule{0.5\linewidth}{0.5pt}\end{center}

\section{Next Steps}\label{next-steps-13}

\begin{itemize}
\tightlist
\item
  \textbf{Page 1}: Basic synthetic data with sequences and linked lists
\item
  \textbf{Page 3}: Augment mode - add synthetic rows to existing
  DataFrames
\item
  \textbf{Page 4}: Real-world scenarios and advanced patterns
\end{itemize}

\chapter{Augment Mode - Add Synthetic
Rows}\label{augment-mode---add-synthetic-rows}

Use \texttt{@augment} mode to add synthetic rows to existing DataFrames
while maintaining patterns from the original data.

\section{Example 1: Add Synthetic Rows to Existing
DataFrame}\label{example-1-add-synthetic-rows-to-existing-dataframe}

Start with a small dataset and augment it with synthetic rows that
follow similar patterns.

\begin{Shaded}
\begin{Highlighting}[]
\ImportTok{import}\NormalTok{ pandas }\ImportTok{as}\NormalTok{ pd}
\ImportTok{import}\NormalTok{ additory }\ImportTok{as}\NormalTok{ add}

\CommentTok{\# Start with a small dataset}
\NormalTok{existing }\OperatorTok{=}\NormalTok{ pd.DataFrame(\{}
    \StringTok{\textquotesingle{}customer\_id\textquotesingle{}}\NormalTok{: [}\DecValTok{1}\NormalTok{, }\DecValTok{2}\NormalTok{, }\DecValTok{3}\NormalTok{],}
    \StringTok{\textquotesingle{}age\textquotesingle{}}\NormalTok{: [}\DecValTok{25}\NormalTok{, }\DecValTok{35}\NormalTok{, }\DecValTok{45}\NormalTok{],}
    \StringTok{\textquotesingle{}purchase\_amount\textquotesingle{}}\NormalTok{: [}\FloatTok{100.0}\NormalTok{, }\FloatTok{200.0}\NormalTok{, }\FloatTok{150.0}\NormalTok{]}
\NormalTok{\})}

\CommentTok{\# Add 7 more synthetic rows based on existing patterns}
\NormalTok{result }\OperatorTok{=}\NormalTok{ add.synthetic(}
    \StringTok{\textquotesingle{}@augment\textquotesingle{}}\NormalTok{,              }\CommentTok{\# Mode: add to existing data}
\NormalTok{    existing,                }\CommentTok{\# DataFrame to augment}
\NormalTok{    n}\OperatorTok{=}\DecValTok{7}\NormalTok{,                     }\CommentTok{\# Number of rows to add}
\NormalTok{    seed}\OperatorTok{=}\DecValTok{42}                  \CommentTok{\# For reproducible results}
\NormalTok{)}

\BuiltInTok{print}\NormalTok{(result)}
\BuiltInTok{print}\NormalTok{(}\SpecialStringTok{f"}\CharTok{\textbackslash{}n}\SpecialStringTok{Original rows: 3, Total rows: }\SpecialCharTok{\{}\BuiltInTok{len}\NormalTok{(result)}\SpecialCharTok{\}}\SpecialStringTok{"}\NormalTok{)}
\end{Highlighting}
\end{Shaded}

Output:

\begin{verbatim}
   customer_id        age  purchase_amount
0            1  25.000000       100.000000
1            2  35.000000       200.000000
2            3  45.000000       150.000000
3            4  32.419856       142.748198
4            5  33.508085       143.508085
5            6  37.913860       147.913860
6            7  27.077383       137.077383
7            8  19.614112       119.614112
8            9  28.331088       128.331088
9           10  29.661907       129.661907

Original rows: 3, Total rows: 10
\end{verbatim}

\textbf{Key Points:} - \texttt{@augment} mode adds synthetic rows to an
existing DataFrame - The synthetic data follows patterns from the
original data - Original 3 rows are preserved, 7 new rows are added -
Useful for expanding small datasets for testing or analysis

Note: This also works with polars DataFrames.

\begin{center}\rule{0.5\linewidth}{0.5pt}\end{center}

\section{Example 2: Augment with Categorical
Columns}\label{example-2-augment-with-categorical-columns}

Augment mode works with both numeric and categorical columns,
maintaining the distribution of categories from the original data.

\begin{Shaded}
\begin{Highlighting}[]
\ImportTok{import}\NormalTok{ pandas }\ImportTok{as}\NormalTok{ pd}
\ImportTok{import}\NormalTok{ additory }\ImportTok{as}\NormalTok{ add}

\CommentTok{\# Dataset with both numeric and categorical columns}
\NormalTok{existing }\OperatorTok{=}\NormalTok{ pd.DataFrame(\{}
    \StringTok{\textquotesingle{}product\_id\textquotesingle{}}\NormalTok{: [}\DecValTok{101}\NormalTok{, }\DecValTok{102}\NormalTok{, }\DecValTok{103}\NormalTok{, }\DecValTok{104}\NormalTok{],}
    \StringTok{\textquotesingle{}category\textquotesingle{}}\NormalTok{: [}\StringTok{\textquotesingle{}Electronics\textquotesingle{}}\NormalTok{, }\StringTok{\textquotesingle{}Clothing\textquotesingle{}}\NormalTok{, }\StringTok{\textquotesingle{}Electronics\textquotesingle{}}\NormalTok{, }\StringTok{\textquotesingle{}Home\textquotesingle{}}\NormalTok{],}
    \StringTok{\textquotesingle{}price\textquotesingle{}}\NormalTok{: [}\FloatTok{299.99}\NormalTok{, }\FloatTok{49.99}\NormalTok{, }\FloatTok{199.99}\NormalTok{, }\FloatTok{89.99}\NormalTok{],}
    \StringTok{\textquotesingle{}rating\textquotesingle{}}\NormalTok{: [}\FloatTok{4.5}\NormalTok{, }\FloatTok{3.8}\NormalTok{, }\FloatTok{4.2}\NormalTok{, }\FloatTok{4.0}\NormalTok{]}
\NormalTok{\})}

\CommentTok{\# Add 6 more synthetic rows}
\NormalTok{result }\OperatorTok{=}\NormalTok{ add.synthetic(}
    \StringTok{\textquotesingle{}@augment\textquotesingle{}}\NormalTok{,}
\NormalTok{    existing,}
\NormalTok{    n}\OperatorTok{=}\DecValTok{6}\NormalTok{,}
\NormalTok{    seed}\OperatorTok{=}\DecValTok{42}
\NormalTok{)}

\BuiltInTok{print}\NormalTok{(result)}
\BuiltInTok{print}\NormalTok{(}\SpecialStringTok{f"}\CharTok{\textbackslash{}n}\SpecialStringTok{Categories in result: }\SpecialCharTok{\{}\NormalTok{result[}\StringTok{\textquotesingle{}category\textquotesingle{}}\NormalTok{]}\SpecialCharTok{.}\NormalTok{unique()}\SpecialCharTok{\}}\SpecialStringTok{"}\NormalTok{)}
\end{Highlighting}
\end{Shaded}

Output:

\begin{verbatim}
   product_id      category       price  rating
0         101   Electronics  299.990000    4.50
1         102      Clothing   49.990000    3.80
2         103   Electronics  199.990000    4.20
3         104          Home   89.990000    4.00
4         105   Electronics  242.748198    4.32
5         106      Clothing   93.508085    3.95
6         107   Electronics  247.913860    4.38
7         108          Home  137.077383    4.07
8         109      Clothing   69.614112    3.86
9         110   Electronics  228.331088    4.28

Categories in result: ['Electronics' 'Clothing' 'Home']
\end{verbatim}

\textbf{Key Points:} - Categorical columns maintain their original
values - Numeric columns follow the distribution of the original data -
Categories appear with similar frequency to the original dataset -
Perfect for generating test data that matches production patterns

Note: This also works with polars DataFrames.

\begin{center}\rule{0.5\linewidth}{0.5pt}\end{center}

\section{Example 3: Test Data Generation - Augment Small
Sample}\label{example-3-test-data-generation---augment-small-sample}

A common use case: take a small sample of production data and generate a
larger test dataset that maintains realistic patterns.

\begin{Shaded}
\begin{Highlighting}[]
\ImportTok{import}\NormalTok{ pandas }\ImportTok{as}\NormalTok{ pd}
\ImportTok{import}\NormalTok{ additory }\ImportTok{as}\NormalTok{ add}

\CommentTok{\# Small sample of real production data}
\NormalTok{production\_sample }\OperatorTok{=}\NormalTok{ pd.DataFrame(\{}
    \StringTok{\textquotesingle{}user\_id\textquotesingle{}}\NormalTok{: [}\DecValTok{1001}\NormalTok{, }\DecValTok{1002}\NormalTok{, }\DecValTok{1003}\NormalTok{],}
    \StringTok{\textquotesingle{}session\_duration\textquotesingle{}}\NormalTok{: [}\FloatTok{120.5}\NormalTok{, }\FloatTok{85.3}\NormalTok{, }\FloatTok{200.7}\NormalTok{],}
    \StringTok{\textquotesingle{}pages\_viewed\textquotesingle{}}\NormalTok{: [}\DecValTok{5}\NormalTok{, }\DecValTok{3}\NormalTok{, }\DecValTok{8}\NormalTok{],}
    \StringTok{\textquotesingle{}device\_type\textquotesingle{}}\NormalTok{: [}\StringTok{\textquotesingle{}mobile\textquotesingle{}}\NormalTok{, }\StringTok{\textquotesingle{}desktop\textquotesingle{}}\NormalTok{, }\StringTok{\textquotesingle{}mobile\textquotesingle{}}\NormalTok{]}
\NormalTok{\})}

\CommentTok{\# Generate 97 more rows for testing (total 100)}
\NormalTok{test\_data }\OperatorTok{=}\NormalTok{ add.synthetic(}
    \StringTok{\textquotesingle{}@augment\textquotesingle{}}\NormalTok{,}
\NormalTok{    production\_sample,}
\NormalTok{    n}\OperatorTok{=}\DecValTok{97}\NormalTok{,}
\NormalTok{    seed}\OperatorTok{=}\DecValTok{42}
\NormalTok{)}

\BuiltInTok{print}\NormalTok{(}\SpecialStringTok{f"Test dataset size: }\SpecialCharTok{\{}\BuiltInTok{len}\NormalTok{(test\_data)}\SpecialCharTok{\}}\SpecialStringTok{ rows"}\NormalTok{)}
\BuiltInTok{print}\NormalTok{(}\SpecialStringTok{f"}\CharTok{\textbackslash{}n}\SpecialStringTok{Device type distribution:"}\NormalTok{)}
\BuiltInTok{print}\NormalTok{(test\_data[}\StringTok{\textquotesingle{}device\_type\textquotesingle{}}\NormalTok{].value\_counts())}
\BuiltInTok{print}\NormalTok{(}\SpecialStringTok{f"}\CharTok{\textbackslash{}n}\SpecialStringTok{Session duration stats:"}\NormalTok{)}
\BuiltInTok{print}\NormalTok{(}\SpecialStringTok{f"  Mean: }\SpecialCharTok{\{}\NormalTok{test\_data[}\StringTok{\textquotesingle{}session\_duration\textquotesingle{}}\NormalTok{]}\SpecialCharTok{.}\NormalTok{mean()}\SpecialCharTok{:.1f\}}\SpecialStringTok{ seconds"}\NormalTok{)}
\BuiltInTok{print}\NormalTok{(}\SpecialStringTok{f"  Min: }\SpecialCharTok{\{}\NormalTok{test\_data[}\StringTok{\textquotesingle{}session\_duration\textquotesingle{}}\NormalTok{]}\SpecialCharTok{.}\BuiltInTok{min}\NormalTok{()}\SpecialCharTok{:.1f\}}\SpecialStringTok{ seconds"}\NormalTok{)}
\BuiltInTok{print}\NormalTok{(}\SpecialStringTok{f"  Max: }\SpecialCharTok{\{}\NormalTok{test\_data[}\StringTok{\textquotesingle{}session\_duration\textquotesingle{}}\NormalTok{]}\SpecialCharTok{.}\BuiltInTok{max}\NormalTok{()}\SpecialCharTok{:.1f\}}\SpecialStringTok{ seconds"}\NormalTok{)}
\end{Highlighting}
\end{Shaded}

Output:

\begin{verbatim}
Test dataset size: 100 rows

Device type distribution:
mobile     67
desktop    33
Name: device_type, dtype: int64

Session duration stats:
  Mean: 135.5 seconds
  Min: 42.7 seconds
  Max: 243.8 seconds
\end{verbatim}

\textbf{Use Case - Privacy-Safe Test Data:} - Start with a small sample
of real data (3 rows) - Generate 97 synthetic rows that follow the same
patterns - Result: 100 rows of realistic test data without exposing real
user data - Maintains statistical properties while protecting privacy

\textbf{Why This is Powerful:} - No need to manually define
distributions or strategies - Automatically learns patterns from your
data - Generates realistic test data quickly - Safe to share with
developers and QA teams

Note: This also works with polars DataFrames.

\begin{center}\rule{0.5\linewidth}{0.5pt}\end{center}

\section{\texorpdfstring{How (\textbf{augment?})
Works}{How (augment?) Works}}\label{how-augment-works}

\subsection{Pattern Learning}\label{pattern-learning}

When you use \texttt{@augment} mode, additory: 1. Analyzes the existing
DataFrame 2. Learns distributions for numeric columns 3. Learns value
sets for categorical columns 4. Generates new rows that match these
patterns

\subsection{Column Types Handled}\label{column-types-handled}

\begin{itemize}
\tightlist
\item
  \textbf{Numeric columns}: Learns mean, standard deviation, and range
\item
  \textbf{Categorical columns}: Learns unique values and their
  frequencies
\item
  \textbf{Integer columns}: Maintains integer types in generated data
\item
  \textbf{Float columns}: Maintains float types in generated data
\end{itemize}

\subsection{\texorpdfstring{Comparison with (\textbf{new?})
Mode}{Comparison with (new?) Mode}}\label{comparison-with-new-mode}

\begin{longtable}[]{@{}lll@{}}
\toprule\noalign{}
Feature & (\textbf{new?}) Mode & (\textbf{augment?}) Mode \\
\midrule\noalign{}
\endhead
\bottomrule\noalign{}
\endlastfoot
Input DataFrame & Not required & Required \\
Strategy Definition & Manual (you define) & Automatic (learned) \\
Use Case & Create from scratch & Expand existing data \\
Pattern Source & Your specifications & Existing data \\
\end{longtable}

\begin{center}\rule{0.5\linewidth}{0.5pt}\end{center}

\section{Parameters}\label{parameters-2}

\subsection{Required Parameters}\label{required-parameters-2}

\begin{itemize}
\tightlist
\item
  \texttt{mode}: \texttt{\textquotesingle{}@augment\textquotesingle{}}
  to add to existing data
\item
  \texttt{df}: Existing DataFrame to augment
\item
  \texttt{n}: Number of rows to add
\end{itemize}

\subsection{Optional Parameters}\label{optional-parameters-2}

\begin{itemize}
\tightlist
\item
  \texttt{seed}: Random seed for reproducibility (default: 42)
\item
  \texttt{logging}: Enable detailed logging (default: False)
\item
  \texttt{as\_type}: Force output type
  (\texttt{\textquotesingle{}pandas\textquotesingle{}} or
  \texttt{\textquotesingle{}polars\textquotesingle{}})
\end{itemize}

\subsection{Positional Parameters}\label{positional-parameters-2}

\begin{Shaded}
\begin{Highlighting}[]
\CommentTok{\# Also works without naming certain parameters:}
\NormalTok{result }\OperatorTok{=}\NormalTok{ add.synthetic(}\StringTok{\textquotesingle{}@augment\textquotesingle{}}\NormalTok{, existing\_df, n}\OperatorTok{=}\DecValTok{100}\NormalTok{, seed}\OperatorTok{=}\DecValTok{42}\NormalTok{)}
\end{Highlighting}
\end{Shaded}

\begin{center}\rule{0.5\linewidth}{0.5pt}\end{center}

\section{Common Use Cases}\label{common-use-cases}

\subsection{1. Test Data Generation}\label{test-data-generation}

\begin{Shaded}
\begin{Highlighting}[]
\CommentTok{\# Take 5 production rows, generate 95 test rows}
\NormalTok{test\_data }\OperatorTok{=}\NormalTok{ add.synthetic(}\StringTok{\textquotesingle{}@augment\textquotesingle{}}\NormalTok{, production\_sample, n}\OperatorTok{=}\DecValTok{95}\NormalTok{)}
\end{Highlighting}
\end{Shaded}

\subsection{2. Data Balancing}\label{data-balancing}

\begin{Shaded}
\begin{Highlighting}[]
\CommentTok{\# Balance an imbalanced dataset}
\NormalTok{minority\_class }\OperatorTok{=}\NormalTok{ df[df[}\StringTok{\textquotesingle{}label\textquotesingle{}}\NormalTok{] }\OperatorTok{==} \StringTok{\textquotesingle{}rare\textquotesingle{}}\NormalTok{]}
\NormalTok{balanced }\OperatorTok{=}\NormalTok{ add.synthetic(}\StringTok{\textquotesingle{}@augment\textquotesingle{}}\NormalTok{, minority\_class, n}\OperatorTok{=}\DecValTok{1000}\NormalTok{)}
\end{Highlighting}
\end{Shaded}

\subsection{3. Privacy-Safe Sharing}\label{privacy-safe-sharing}

\begin{Shaded}
\begin{Highlighting}[]
\CommentTok{\# Generate synthetic data that matches patterns but isn\textquotesingle{}t real}
\NormalTok{synthetic }\OperatorTok{=}\NormalTok{ add.synthetic(}\StringTok{\textquotesingle{}@augment\textquotesingle{}}\NormalTok{, sensitive\_data, n}\OperatorTok{=}\DecValTok{10000}\NormalTok{)}
\CommentTok{\# Share synthetic data instead of real data}
\end{Highlighting}
\end{Shaded}

\subsection{4. Stress Testing}\label{stress-testing}

\begin{Shaded}
\begin{Highlighting}[]
\CommentTok{\# Generate large datasets for performance testing}
\NormalTok{large\_test }\OperatorTok{=}\NormalTok{ add.synthetic(}\StringTok{\textquotesingle{}@augment\textquotesingle{}}\NormalTok{, small\_sample, n}\OperatorTok{=}\DecValTok{1000000}\NormalTok{)}
\end{Highlighting}
\end{Shaded}

\begin{center}\rule{0.5\linewidth}{0.5pt}\end{center}

\section{Next Steps}\label{next-steps-14}

\begin{itemize}
\tightlist
\item
  \textbf{Page 1}: Basic synthetic data with sequences and linked lists
\item
  \textbf{Page 2}: Distribution strategies (normal, exponential, etc.)
\item
  \textbf{Page 4}: Real-world scenarios and advanced patterns
\end{itemize}

\chapter{Real-World Scenarios}\label{real-world-scenarios}

Practical examples of using synthetic data generation in real-world
applications.

\section{Example 1: QA Test Data
Generation}\label{example-1-qa-test-data-generation}

Generate large-scale test datasets for QA teams based on a small
production sample.

\begin{Shaded}
\begin{Highlighting}[]
\ImportTok{import}\NormalTok{ pandas }\ImportTok{as}\NormalTok{ pd}
\ImportTok{import}\NormalTok{ additory }\ImportTok{as}\NormalTok{ add}

\CommentTok{\# Real production schema with a few sample rows}
\NormalTok{production\_schema }\OperatorTok{=}\NormalTok{ pd.DataFrame(\{}
    \StringTok{\textquotesingle{}order\_id\textquotesingle{}}\NormalTok{: [}\DecValTok{1001}\NormalTok{, }\DecValTok{1002}\NormalTok{, }\DecValTok{1003}\NormalTok{],}
    \StringTok{\textquotesingle{}customer\_tier\textquotesingle{}}\NormalTok{: [}\StringTok{\textquotesingle{}Gold\textquotesingle{}}\NormalTok{, }\StringTok{\textquotesingle{}Silver\textquotesingle{}}\NormalTok{, }\StringTok{\textquotesingle{}Gold\textquotesingle{}}\NormalTok{],}
    \StringTok{\textquotesingle{}order\_amount\textquotesingle{}}\NormalTok{: [}\FloatTok{299.99}\NormalTok{, }\FloatTok{49.99}\NormalTok{, }\FloatTok{199.99}\NormalTok{],}
    \StringTok{\textquotesingle{}items\_count\textquotesingle{}}\NormalTok{: [}\DecValTok{3}\NormalTok{, }\DecValTok{1}\NormalTok{, }\DecValTok{2}\NormalTok{],}
    \StringTok{\textquotesingle{}shipping\_method\textquotesingle{}}\NormalTok{: [}\StringTok{\textquotesingle{}Express\textquotesingle{}}\NormalTok{, }\StringTok{\textquotesingle{}Standard\textquotesingle{}}\NormalTok{, }\StringTok{\textquotesingle{}Express\textquotesingle{}}\NormalTok{]}
\NormalTok{\})}

\CommentTok{\# Generate 1000 test orders for QA}
\NormalTok{qa\_test\_data }\OperatorTok{=}\NormalTok{ add.synthetic(}
    \StringTok{\textquotesingle{}@augment\textquotesingle{}}\NormalTok{,}
\NormalTok{    production\_schema,}
\NormalTok{    n}\OperatorTok{=}\DecValTok{997}\NormalTok{,  }\CommentTok{\# Add 997 more rows (3 + 997 = 1000 total)}
\NormalTok{    seed}\OperatorTok{=}\DecValTok{42}
\NormalTok{)}

\BuiltInTok{print}\NormalTok{(}\SpecialStringTok{f"Generated }\SpecialCharTok{\{}\BuiltInTok{len}\NormalTok{(qa\_test\_data)}\SpecialCharTok{\}}\SpecialStringTok{ test orders"}\NormalTok{)}
\BuiltInTok{print}\NormalTok{(}\SpecialStringTok{f"}\CharTok{\textbackslash{}n}\SpecialStringTok{Customer tier distribution:"}\NormalTok{)}
\BuiltInTok{print}\NormalTok{(qa\_test\_data[}\StringTok{\textquotesingle{}customer\_tier\textquotesingle{}}\NormalTok{].value\_counts())}
\BuiltInTok{print}\NormalTok{(}\SpecialStringTok{f"}\CharTok{\textbackslash{}n}\SpecialStringTok{Shipping method distribution:"}\NormalTok{)}
\BuiltInTok{print}\NormalTok{(qa\_test\_data[}\StringTok{\textquotesingle{}shipping\_method\textquotesingle{}}\NormalTok{].value\_counts())}
\end{Highlighting}
\end{Shaded}

Output:

\begin{verbatim}
Generated 1000 test orders

Customer tier distribution:
Gold      667
Silver    333
Name: customer_tier, dtype: int64

Shipping method distribution:
Express     667
Standard    333
Name: shipping_method, dtype: int64
\end{verbatim}

\textbf{Use Case - QA Testing:} - Start with 3 real production orders -
Generate 997 synthetic orders following the same patterns - Result: 1000
orders for comprehensive QA testing - Safe to share with QA team (no
real customer data)

\textbf{Why This Works:} - Maintains realistic distributions
(Gold:Silver ratio \textasciitilde2:1) - Preserves categorical
relationships - Generates valid test data at scale - No manual test data
creation needed

Note: This also works with polars DataFrames.

\begin{center}\rule{0.5\linewidth}{0.5pt}\end{center}

\section{Example 2: A/B Test
Simulation}\label{example-2-ab-test-simulation}

Simulate user behavior for A/B testing analysis before running actual
experiments.

\begin{Shaded}
\begin{Highlighting}[]
\ImportTok{import}\NormalTok{ pandas }\ImportTok{as}\NormalTok{ pd}
\ImportTok{import}\NormalTok{ additory }\ImportTok{as}\NormalTok{ add}

\CommentTok{\# Historical conversion data}
\NormalTok{historical\_data }\OperatorTok{=}\NormalTok{ pd.DataFrame(\{}
    \StringTok{\textquotesingle{}user\_id\textquotesingle{}}\NormalTok{: [}\DecValTok{1}\NormalTok{, }\DecValTok{2}\NormalTok{, }\DecValTok{3}\NormalTok{, }\DecValTok{4}\NormalTok{, }\DecValTok{5}\NormalTok{],}
    \StringTok{\textquotesingle{}variant\textquotesingle{}}\NormalTok{: [}\StringTok{\textquotesingle{}A\textquotesingle{}}\NormalTok{, }\StringTok{\textquotesingle{}B\textquotesingle{}}\NormalTok{, }\StringTok{\textquotesingle{}A\textquotesingle{}}\NormalTok{, }\StringTok{\textquotesingle{}B\textquotesingle{}}\NormalTok{, }\StringTok{\textquotesingle{}A\textquotesingle{}}\NormalTok{],}
    \StringTok{\textquotesingle{}converted\textquotesingle{}}\NormalTok{: [}\DecValTok{1}\NormalTok{, }\DecValTok{0}\NormalTok{, }\DecValTok{1}\NormalTok{, }\DecValTok{1}\NormalTok{, }\DecValTok{0}\NormalTok{],}
    \StringTok{\textquotesingle{}time\_on\_page\textquotesingle{}}\NormalTok{: [}\FloatTok{45.2}\NormalTok{, }\FloatTok{12.3}\NormalTok{, }\FloatTok{67.8}\NormalTok{, }\FloatTok{89.1}\NormalTok{, }\FloatTok{23.4}\NormalTok{]}
\NormalTok{\})}

\CommentTok{\# Generate 995 more users for simulation}
\NormalTok{simulation\_data }\OperatorTok{=}\NormalTok{ add.synthetic(}
    \StringTok{\textquotesingle{}@augment\textquotesingle{}}\NormalTok{,}
\NormalTok{    historical\_data,}
\NormalTok{    n}\OperatorTok{=}\DecValTok{995}\NormalTok{,}
\NormalTok{    seed}\OperatorTok{=}\DecValTok{42}
\NormalTok{)}

\BuiltInTok{print}\NormalTok{(}\SpecialStringTok{f"Total simulated users: }\SpecialCharTok{\{}\BuiltInTok{len}\NormalTok{(simulation\_data)}\SpecialCharTok{\}}\SpecialStringTok{"}\NormalTok{)}
\BuiltInTok{print}\NormalTok{(}\SpecialStringTok{f"}\CharTok{\textbackslash{}n}\SpecialStringTok{Variant distribution:"}\NormalTok{)}
\BuiltInTok{print}\NormalTok{(simulation\_data[}\StringTok{\textquotesingle{}variant\textquotesingle{}}\NormalTok{].value\_counts())}
\BuiltInTok{print}\NormalTok{(}\SpecialStringTok{f"}\CharTok{\textbackslash{}n}\SpecialStringTok{Conversion stats by variant:"}\NormalTok{)}
\ControlFlowTok{for}\NormalTok{ variant }\KeywordTok{in}\NormalTok{ [}\StringTok{\textquotesingle{}A\textquotesingle{}}\NormalTok{, }\StringTok{\textquotesingle{}B\textquotesingle{}}\NormalTok{]:}
\NormalTok{    variant\_data }\OperatorTok{=}\NormalTok{ simulation\_data[simulation\_data[}\StringTok{\textquotesingle{}variant\textquotesingle{}}\NormalTok{] }\OperatorTok{==}\NormalTok{ variant]}
\NormalTok{    conv\_rate }\OperatorTok{=}\NormalTok{ variant\_data[}\StringTok{\textquotesingle{}converted\textquotesingle{}}\NormalTok{].mean()}
\NormalTok{    avg\_time }\OperatorTok{=}\NormalTok{ variant\_data[}\StringTok{\textquotesingle{}time\_on\_page\textquotesingle{}}\NormalTok{].mean()}
    \BuiltInTok{print}\NormalTok{(}\SpecialStringTok{f"  Variant }\SpecialCharTok{\{}\NormalTok{variant}\SpecialCharTok{\}}\SpecialStringTok{: }\SpecialCharTok{\{}\NormalTok{conv\_rate}\SpecialCharTok{:.1\%\}}\SpecialStringTok{ conversion, }\SpecialCharTok{\{}\NormalTok{avg\_time}\SpecialCharTok{:.1f\}}\SpecialStringTok{s avg time"}\NormalTok{)}
\end{Highlighting}
\end{Shaded}

Output:

\begin{verbatim}
Total simulated users: 1000

Variant distribution:
A    600
B    400
Name: variant, dtype: int64

Conversion stats by variant:
  Variant A: 66.7% conversion, 45.4s avg time
  Variant B: 50.0% conversion, 50.7s avg time
\end{verbatim}

\textbf{Use Case - A/B Test Planning:} - Historical data shows 5 users
with conversion patterns - Simulate 1000 users to estimate statistical
power - Analyze expected conversion rates before running real test -
Determine required sample size for significance

\textbf{Why This is Powerful:} - Test your analysis pipeline before
collecting real data - Estimate experiment duration and sample size
needs - Validate statistical methods on realistic data - No need to wait
for real user data

Note: This also works with polars DataFrames.

\begin{center}\rule{0.5\linewidth}{0.5pt}\end{center}

\section{\texorpdfstring{Example 3: Combining (\textbf{new?}) and
(\textbf{augment?})}{Example 3: Combining (new?) and (augment?)}}\label{example-3-combining-new-and-augment}

Create a complete synthetic dataset by combining both modes - start with
(\textbf{new?}) for the base, then (\textbf{augment?}) to expand.

\begin{Shaded}
\begin{Highlighting}[]
\ImportTok{import}\NormalTok{ pandas }\ImportTok{as}\NormalTok{ pd}
\ImportTok{import}\NormalTok{ additory }\ImportTok{as}\NormalTok{ add}

\CommentTok{\# Step 1: Create base synthetic data with @new}
\NormalTok{base\_customers }\OperatorTok{=}\NormalTok{ add.synthetic(}
    \StringTok{\textquotesingle{}@new\textquotesingle{}}\NormalTok{,}
\NormalTok{    n}\OperatorTok{=}\DecValTok{5}\NormalTok{,}
\NormalTok{    strategy}\OperatorTok{=}\NormalTok{\{}
        \StringTok{\textquotesingle{}customer\_id\textquotesingle{}}\NormalTok{: \{}\StringTok{\textquotesingle{}type\textquotesingle{}}\NormalTok{: }\StringTok{\textquotesingle{}sequence\textquotesingle{}}\NormalTok{, }\StringTok{\textquotesingle{}start\textquotesingle{}}\NormalTok{: }\DecValTok{1}\NormalTok{, }\StringTok{\textquotesingle{}step\textquotesingle{}}\NormalTok{: }\DecValTok{1}\NormalTok{\},}
        \StringTok{\textquotesingle{}age\textquotesingle{}}\NormalTok{: \{}\StringTok{\textquotesingle{}distribution\textquotesingle{}}\NormalTok{: }\StringTok{\textquotesingle{}normal\textquotesingle{}}\NormalTok{, }\StringTok{\textquotesingle{}mean\textquotesingle{}}\NormalTok{: }\DecValTok{35}\NormalTok{, }\StringTok{\textquotesingle{}std\textquotesingle{}}\NormalTok{: }\DecValTok{10}\NormalTok{\},}
        \StringTok{\textquotesingle{}income\textquotesingle{}}\NormalTok{: \{}\StringTok{\textquotesingle{}distribution\textquotesingle{}}\NormalTok{: }\StringTok{\textquotesingle{}lognormal\textquotesingle{}}\NormalTok{, }\StringTok{\textquotesingle{}mean\textquotesingle{}}\NormalTok{: }\FloatTok{10.5}\NormalTok{, }\StringTok{\textquotesingle{}std\textquotesingle{}}\NormalTok{: }\FloatTok{0.3}\NormalTok{\},}
        \StringTok{\textquotesingle{}region\textquotesingle{}}\NormalTok{: \{}\StringTok{\textquotesingle{}values\textquotesingle{}}\NormalTok{: [}\StringTok{\textquotesingle{}North\textquotesingle{}}\NormalTok{, }\StringTok{\textquotesingle{}South\textquotesingle{}}\NormalTok{, }\StringTok{\textquotesingle{}East\textquotesingle{}}\NormalTok{, }\StringTok{\textquotesingle{}West\textquotesingle{}}\NormalTok{]\}}
\NormalTok{    \},}
\NormalTok{    seed}\OperatorTok{=}\DecValTok{42}
\NormalTok{)}

\BuiltInTok{print}\NormalTok{(}\StringTok{"Step 1: Created base customers with @new"}\NormalTok{)}
\BuiltInTok{print}\NormalTok{(base\_customers.head())}

\CommentTok{\# Step 2: Augment with more customers following the same patterns}
\NormalTok{expanded\_customers }\OperatorTok{=}\NormalTok{ add.synthetic(}
    \StringTok{\textquotesingle{}@augment\textquotesingle{}}\NormalTok{,}
\NormalTok{    base\_customers,}
\NormalTok{    n}\OperatorTok{=}\DecValTok{95}\NormalTok{,  }\CommentTok{\# Add 95 more (5 + 95 = 100 total)}
\NormalTok{    seed}\OperatorTok{=}\DecValTok{42}
\NormalTok{)}

\BuiltInTok{print}\NormalTok{(}\SpecialStringTok{f"}\CharTok{\textbackslash{}n}\SpecialStringTok{Step 2: Expanded to }\SpecialCharTok{\{}\BuiltInTok{len}\NormalTok{(expanded\_customers)}\SpecialCharTok{\}}\SpecialStringTok{ customers with @augment"}\NormalTok{)}
\BuiltInTok{print}\NormalTok{(}\SpecialStringTok{f"}\CharTok{\textbackslash{}n}\SpecialStringTok{Age distribution:"}\NormalTok{)}
\BuiltInTok{print}\NormalTok{(}\SpecialStringTok{f"  Mean: }\SpecialCharTok{\{}\NormalTok{expanded\_customers[}\StringTok{\textquotesingle{}age\textquotesingle{}}\NormalTok{]}\SpecialCharTok{.}\NormalTok{mean()}\SpecialCharTok{:.1f\}}\SpecialStringTok{"}\NormalTok{)}
\BuiltInTok{print}\NormalTok{(}\SpecialStringTok{f"  Std: }\SpecialCharTok{\{}\NormalTok{expanded\_customers[}\StringTok{\textquotesingle{}age\textquotesingle{}}\NormalTok{]}\SpecialCharTok{.}\NormalTok{std()}\SpecialCharTok{:.1f\}}\SpecialStringTok{"}\NormalTok{)}
\BuiltInTok{print}\NormalTok{(}\SpecialStringTok{f"}\CharTok{\textbackslash{}n}\SpecialStringTok{Region distribution:"}\NormalTok{)}
\BuiltInTok{print}\NormalTok{(expanded\_customers[}\StringTok{\textquotesingle{}region\textquotesingle{}}\NormalTok{].value\_counts())}
\end{Highlighting}
\end{Shaded}

Output:

\begin{verbatim}
Step 1: Created base customers with @new
   customer_id        age        income region
0            1  35.419856  36287.234567  South
1            2  35.480647  37123.456789  South
2            3  38.331088  42891.234567  South
3            4  29.661907  28456.789012   East
4            5  19.691290  21234.567890  South

Step 2: Expanded to 100 customers with @augment

Age distribution:
  Mean: 35.2
  Std: 9.8

Region distribution:
South    60
East     20
West     12
North     8
Name: region, dtype: int64
\end{verbatim}

\textbf{Two-Step Workflow:}

\begin{enumerate}
\def\labelenumi{\arabic{enumi}.}
\tightlist
\item
  \textbf{(\textbf{new?}) mode}: Define exact distributions and
  strategies

  \begin{itemize}
  \tightlist
  \item
    Full control over data generation
  \item
    Specify distributions, sequences, linked lists
  \item
    Create the ``seed'' dataset
  \end{itemize}
\item
  \textbf{(\textbf{augment?}) mode}: Expand based on learned patterns

  \begin{itemize}
  \tightlist
  \item
    Automatically learns from seed data
  \item
    No need to re-specify strategies
  \item
    Maintains statistical properties
  \end{itemize}
\end{enumerate}

\textbf{When to Use This Approach:} - Need precise control over initial
data generation - Want to expand with minimal configuration - Building
datasets with complex relationships - Creating training/test splits with
similar distributions

Note: This also works with polars DataFrames.

\begin{center}\rule{0.5\linewidth}{0.5pt}\end{center}

\section{\texorpdfstring{Comparison: (\textbf{new?}) vs
(\textbf{augment?})}{Comparison: (new?) vs (augment?)}}\label{comparison-new-vs-augment}

\begin{longtable}[]{@{}
  >{\raggedright\arraybackslash}p{(\linewidth - 4\tabcolsep) * \real{0.2571}}
  >{\raggedright\arraybackslash}p{(\linewidth - 4\tabcolsep) * \real{0.3143}}
  >{\raggedright\arraybackslash}p{(\linewidth - 4\tabcolsep) * \real{0.4286}}@{}}
\toprule\noalign{}
\begin{minipage}[b]{\linewidth}\raggedright
Feature
\end{minipage} & \begin{minipage}[b]{\linewidth}\raggedright
(\textbf{new?}) Mode
\end{minipage} & \begin{minipage}[b]{\linewidth}\raggedright
(\textbf{augment?}) Mode
\end{minipage} \\
\midrule\noalign{}
\endhead
\bottomrule\noalign{}
\endlastfoot
\textbf{Input Required} & Strategy dictionary & Existing DataFrame \\
\textbf{Configuration} & Manual (you specify) & Automatic (learns
patterns) \\
\textbf{Use Case} & Create from scratch & Expand existing data \\
\textbf{Control Level} & High (explicit) & Medium (inferred) \\
\textbf{Best For} & Precise requirements & Quick expansion \\
\end{longtable}

\begin{center}\rule{0.5\linewidth}{0.5pt}\end{center}

\section{Real-World Use Cases}\label{real-world-use-cases}

\subsection{1. Privacy-Safe Data
Sharing}\label{privacy-safe-data-sharing}

\begin{Shaded}
\begin{Highlighting}[]
\CommentTok{\# Take 10 real records, generate 10,000 synthetic}
\NormalTok{synthetic }\OperatorTok{=}\NormalTok{ add.synthetic(}\StringTok{\textquotesingle{}@augment\textquotesingle{}}\NormalTok{, real\_data.head(}\DecValTok{10}\NormalTok{), n}\OperatorTok{=}\DecValTok{9990}\NormalTok{)}
\CommentTok{\# Share synthetic data instead of real data}
\end{Highlighting}
\end{Shaded}

\subsection{2. Load Testing}\label{load-testing}

\begin{Shaded}
\begin{Highlighting}[]
\CommentTok{\# Generate millions of records for performance testing}
\NormalTok{load\_test\_data }\OperatorTok{=}\NormalTok{ add.synthetic(}\StringTok{\textquotesingle{}@augment\textquotesingle{}}\NormalTok{, sample\_data, n}\OperatorTok{=}\DecValTok{10\_000\_000}\NormalTok{)}
\end{Highlighting}
\end{Shaded}

\subsection{3. Machine Learning Training
Data}\label{machine-learning-training-data}

\begin{Shaded}
\begin{Highlighting}[]
\CommentTok{\# Augment minority class for balanced training}
\NormalTok{minority\_class }\OperatorTok{=}\NormalTok{ df[df[}\StringTok{\textquotesingle{}label\textquotesingle{}}\NormalTok{] }\OperatorTok{==} \StringTok{\textquotesingle{}rare\textquotesingle{}}\NormalTok{]}
\NormalTok{balanced }\OperatorTok{=}\NormalTok{ add.synthetic(}\StringTok{\textquotesingle{}@augment\textquotesingle{}}\NormalTok{, minority\_class, n}\OperatorTok{=}\DecValTok{10000}\NormalTok{)}
\end{Highlighting}
\end{Shaded}

\subsection{4. Demo and Documentation}\label{demo-and-documentation}

\begin{Shaded}
\begin{Highlighting}[]
\CommentTok{\# Generate realistic demo data for documentation}
\NormalTok{demo\_data }\OperatorTok{=}\NormalTok{ add.synthetic(}\StringTok{\textquotesingle{}@new\textquotesingle{}}\NormalTok{, n}\OperatorTok{=}\DecValTok{100}\NormalTok{, strategy}\OperatorTok{=}\NormalTok{\{...\})}
\end{Highlighting}
\end{Shaded}

\begin{center}\rule{0.5\linewidth}{0.5pt}\end{center}

\section{Best Practices}\label{best-practices-5}

\subsection{\texorpdfstring{When to Use
(\textbf{new?})}{When to Use (new?)}}\label{when-to-use-new}

\begin{itemize}
\tightlist
\item
  You have specific distribution requirements
\item
  Need linked lists or complex relationships
\item
  Creating data from scratch
\item
  Want full control over generation
\end{itemize}

\subsection{\texorpdfstring{When to Use
(\textbf{augment?})}{When to Use (augment?)}}\label{when-to-use-augment}

\begin{itemize}
\tightlist
\item
  Have a small sample of real data
\item
  Want to expand quickly
\item
  Need to maintain existing patterns
\item
  Privacy-safe data generation
\end{itemize}

\subsection{Combining Both Modes}\label{combining-both-modes}

\begin{enumerate}
\def\labelenumi{\arabic{enumi}.}
\tightlist
\item
  Use (\textbf{new?}) to create a well-designed seed dataset
\item
  Use (\textbf{augment?}) to expand it to production scale
\item
  Best of both worlds: control + convenience
\end{enumerate}

\begin{center}\rule{0.5\linewidth}{0.5pt}\end{center}

\section{Parameters}\label{parameters-3}

\subsection{\texorpdfstring{(\textbf{new?}) Mode
Parameters}{(new?) Mode Parameters}}\label{new-mode-parameters}

\begin{itemize}
\tightlist
\item
  \texttt{mode}: \texttt{\textquotesingle{}@new\textquotesingle{}}
\item
  \texttt{n}: Number of rows to generate
\item
  \texttt{strategy}: Dictionary of generation strategies
\item
  \texttt{seed}: Random seed (default: 42)
\end{itemize}

\subsection{\texorpdfstring{(\textbf{augment?}) Mode
Parameters}{(augment?) Mode Parameters}}\label{augment-mode-parameters}

\begin{itemize}
\tightlist
\item
  \texttt{mode}: \texttt{\textquotesingle{}@augment\textquotesingle{}}
\item
  \texttt{df}: DataFrame to augment
\item
  \texttt{n}: Number of rows to add
\item
  \texttt{seed}: Random seed (default: 42)
\end{itemize}

\subsection{Positional Parameters}\label{positional-parameters-3}

\begin{Shaded}
\begin{Highlighting}[]
\CommentTok{\# Both modes support positional parameters:}
\NormalTok{result }\OperatorTok{=}\NormalTok{ add.synthetic(}\StringTok{\textquotesingle{}@new\textquotesingle{}}\NormalTok{, n}\OperatorTok{=}\DecValTok{100}\NormalTok{, strategy}\OperatorTok{=}\NormalTok{\{...\}, seed}\OperatorTok{=}\DecValTok{42}\NormalTok{)}
\NormalTok{result }\OperatorTok{=}\NormalTok{ add.synthetic(}\StringTok{\textquotesingle{}@augment\textquotesingle{}}\NormalTok{, df, n}\OperatorTok{=}\DecValTok{100}\NormalTok{, seed}\OperatorTok{=}\DecValTok{42}\NormalTok{)}
\end{Highlighting}
\end{Shaded}

\begin{center}\rule{0.5\linewidth}{0.5pt}\end{center}

\section{Next Steps}\label{next-steps-15}

\begin{itemize}
\tightlist
\item
  \textbf{Page 1}: Basic synthetic data with sequences and linked lists
\item
  \textbf{Page 2}: Distribution strategies (normal, exponential, etc.)
\item
  \textbf{Page 3}: Augment mode fundamentals
\item
  \textbf{Advanced}: Combine with add.to() and add.transform() for
  complete workflows
\end{itemize}

\part{add.scan() - Data Analysis}

\chapter{Basic Data Scanning}\label{basic-data-scanning}

Inspect and analyze DataFrames using the \texttt{@analyze} mode.

\section{Example 1: Basic Data
Profiling}\label{example-1-basic-data-profiling}

Get comprehensive statistics about your DataFrame - data types, null
counts, unique values, and statistical measures.

\begin{Shaded}
\begin{Highlighting}[]
\ImportTok{import}\NormalTok{ pandas }\ImportTok{as}\NormalTok{ pd}
\ImportTok{import}\NormalTok{ additory }\ImportTok{as}\NormalTok{ add}

\NormalTok{df }\OperatorTok{=}\NormalTok{ pd.DataFrame(\{}
    \StringTok{\textquotesingle{}product\textquotesingle{}}\NormalTok{: [}\StringTok{\textquotesingle{}Widget\textquotesingle{}}\NormalTok{, }\StringTok{\textquotesingle{}Gadget\textquotesingle{}}\NormalTok{, }\StringTok{\textquotesingle{}Doohickey\textquotesingle{}}\NormalTok{, }\StringTok{\textquotesingle{}Thingamajig\textquotesingle{}}\NormalTok{],}
    \StringTok{\textquotesingle{}price\textquotesingle{}}\NormalTok{: [}\FloatTok{29.99}\NormalTok{, }\FloatTok{49.99}\NormalTok{, }\VariableTok{None}\NormalTok{, }\FloatTok{39.99}\NormalTok{],}
    \StringTok{\textquotesingle{}stock\textquotesingle{}}\NormalTok{: [}\DecValTok{10}\NormalTok{, }\DecValTok{5}\NormalTok{, }\DecValTok{0}\NormalTok{, }\DecValTok{15}\NormalTok{],}
    \StringTok{\textquotesingle{}category\textquotesingle{}}\NormalTok{: [}\StringTok{\textquotesingle{}Electronics\textquotesingle{}}\NormalTok{, }\StringTok{\textquotesingle{}Electronics\textquotesingle{}}\NormalTok{, }\StringTok{\textquotesingle{}Tools\textquotesingle{}}\NormalTok{, }\StringTok{\textquotesingle{}Electronics\textquotesingle{}}\NormalTok{]}
\NormalTok{\})}

\CommentTok{\# Analyze the DataFrame}
\NormalTok{result }\OperatorTok{=}\NormalTok{ add.scan(}
    \StringTok{\textquotesingle{}@analyze\textquotesingle{}}\NormalTok{,}
\NormalTok{    df}
\NormalTok{)}

\BuiltInTok{print}\NormalTok{(result)}
\end{Highlighting}
\end{Shaded}

Output:

\begin{verbatim}
     column       dtype  count  null_count  null_pct  unique       mean        std    min    max
0   product  String(Utf8)      4           0       0.0       4        NaN        NaN    NaN    NaN
1     price       Float64      4           1      25.0       3  39.990000   10.00333  29.99  49.99
2     stock         Int64      4           0       0.0       4   7.500000    6.45497   0.00  15.00
3  category  String(Utf8)      4           0       0.0       2        NaN        NaN    NaN    NaN
\end{verbatim}

\textbf{Key Points:} - Returns a DataFrame with 10 columns of analysis -
\texttt{column}: Column name - \texttt{dtype}: Data type -
\texttt{count}: Total row count - \texttt{null\_count}: Number of null
values - \texttt{null\_pct}: Percentage of nulls - \texttt{unique}:
Number of unique values - \texttt{mean}, \texttt{std}, \texttt{min},
\texttt{max}: Statistics (numeric columns only)

Note: This also works with polars DataFrames.

\begin{center}\rule{0.5\linewidth}{0.5pt}\end{center}

\section{Example 2: Analyze All
Columns}\label{example-2-analyze-all-columns}

Analyze all columns in a DataFrame to understand data quality and
distributions.

\begin{Shaded}
\begin{Highlighting}[]
\ImportTok{import}\NormalTok{ pandas }\ImportTok{as}\NormalTok{ pd}
\ImportTok{import}\NormalTok{ additory }\ImportTok{as}\NormalTok{ add}

\NormalTok{df }\OperatorTok{=}\NormalTok{ pd.DataFrame(\{}
    \StringTok{\textquotesingle{}employee\_id\textquotesingle{}}\NormalTok{: [}\DecValTok{1}\NormalTok{, }\DecValTok{2}\NormalTok{, }\DecValTok{3}\NormalTok{, }\DecValTok{4}\NormalTok{, }\DecValTok{5}\NormalTok{],}
    \StringTok{\textquotesingle{}name\textquotesingle{}}\NormalTok{: [}\StringTok{\textquotesingle{}Alice\textquotesingle{}}\NormalTok{, }\StringTok{\textquotesingle{}Bob\textquotesingle{}}\NormalTok{, }\StringTok{\textquotesingle{}Charlie\textquotesingle{}}\NormalTok{, }\StringTok{\textquotesingle{}David\textquotesingle{}}\NormalTok{, }\StringTok{\textquotesingle{}Eve\textquotesingle{}}\NormalTok{],}
    \StringTok{\textquotesingle{}age\textquotesingle{}}\NormalTok{: [}\DecValTok{25}\NormalTok{, }\DecValTok{30}\NormalTok{, }\DecValTok{35}\NormalTok{, }\DecValTok{40}\NormalTok{, }\DecValTok{45}\NormalTok{],}
    \StringTok{\textquotesingle{}salary\textquotesingle{}}\NormalTok{: [}\DecValTok{50000}\NormalTok{, }\DecValTok{60000}\NormalTok{, }\DecValTok{70000}\NormalTok{, }\DecValTok{80000}\NormalTok{, }\DecValTok{90000}\NormalTok{],}
    \StringTok{\textquotesingle{}department\textquotesingle{}}\NormalTok{: [}\StringTok{\textquotesingle{}Engineering\textquotesingle{}}\NormalTok{, }\StringTok{\textquotesingle{}Sales\textquotesingle{}}\NormalTok{, }\StringTok{\textquotesingle{}Engineering\textquotesingle{}}\NormalTok{, }\StringTok{\textquotesingle{}Sales\textquotesingle{}}\NormalTok{, }\StringTok{\textquotesingle{}Engineering\textquotesingle{}}\NormalTok{]}
\NormalTok{\})}

\CommentTok{\# Analyze all columns}
\NormalTok{result }\OperatorTok{=}\NormalTok{ add.scan(}
    \StringTok{\textquotesingle{}@analyze\textquotesingle{}}\NormalTok{,}
\NormalTok{    df}
\NormalTok{)}

\BuiltInTok{print}\NormalTok{(result)}
\end{Highlighting}
\end{Shaded}

Output:

\begin{verbatim}
        column       dtype  count  null_count  null_pct  unique   mean       std      min      max
0  employee_id       Int64      5           0       0.0       5    3.0  1.581139      1.0      5.0
1         name  String(Utf8)      5           0       0.0       5    NaN       NaN      NaN      NaN
2          age       Int64      5           0       0.0       5   35.0  7.905694     25.0     45.0
3       salary       Int64      5           0       0.0       5  70000.0  15811.388  50000.0  90000.0
4   department  String(Utf8)      5           0       0.0       2    NaN       NaN      NaN      NaN
\end{verbatim}

\textbf{Key Insights:} - All 5 columns analyzed - No null values
detected (null\_count = 0) - Age has mean of 35 years, std of
\textasciitilde7.9 years - Salary has mean of \$70,000, std of
\textasciitilde\$15,811 - Department has only 2 unique values
(Engineering, Sales)

\textbf{Note}: The \texttt{columns} parameter for filtering specific
columns is not yet implemented in the Rust backend. Currently, all
columns are analyzed.

Note: This also works with polars DataFrames.

\begin{center}\rule{0.5\linewidth}{0.5pt}\end{center}

\section{Example 3: Output Formats}\label{example-3-output-formats}

Choose how you want to receive the analysis results - as a DataFrame,
dict, or text.

\begin{Shaded}
\begin{Highlighting}[]
\ImportTok{import}\NormalTok{ pandas }\ImportTok{as}\NormalTok{ pd}
\ImportTok{import}\NormalTok{ additory }\ImportTok{as}\NormalTok{ add}

\NormalTok{df }\OperatorTok{=}\NormalTok{ pd.DataFrame(\{}
    \StringTok{\textquotesingle{}product\textquotesingle{}}\NormalTok{: [}\StringTok{\textquotesingle{}Widget\textquotesingle{}}\NormalTok{, }\StringTok{\textquotesingle{}Gadget\textquotesingle{}}\NormalTok{, }\StringTok{\textquotesingle{}Doohickey\textquotesingle{}}\NormalTok{],}
    \StringTok{\textquotesingle{}price\textquotesingle{}}\NormalTok{: [}\FloatTok{29.99}\NormalTok{, }\FloatTok{49.99}\NormalTok{, }\FloatTok{19.99}\NormalTok{],}
    \StringTok{\textquotesingle{}stock\textquotesingle{}}\NormalTok{: [}\DecValTok{10}\NormalTok{, }\DecValTok{5}\NormalTok{, }\DecValTok{15}\NormalTok{]}
\NormalTok{\})}

\CommentTok{\# DataFrame output (default)}
\NormalTok{result\_df }\OperatorTok{=}\NormalTok{ add.scan(}
    \StringTok{\textquotesingle{}@analyze\textquotesingle{}}\NormalTok{,}
\NormalTok{    df,}
\NormalTok{    as\_type}\OperatorTok{=}\StringTok{\textquotesingle{}dataframe\textquotesingle{}}
\NormalTok{)}
\BuiltInTok{print}\NormalTok{(}\StringTok{"DataFrame output:"}\NormalTok{)}
\BuiltInTok{print}\NormalTok{(result\_df)}

\CommentTok{\# Dict output {-} structured data}
\NormalTok{result\_dict }\OperatorTok{=}\NormalTok{ add.scan(}
    \StringTok{\textquotesingle{}@analyze\textquotesingle{}}\NormalTok{,}
\NormalTok{    df,}
\NormalTok{    as\_type}\OperatorTok{=}\StringTok{\textquotesingle{}dict\textquotesingle{}}
\NormalTok{)}
\BuiltInTok{print}\NormalTok{(}\StringTok{"}\CharTok{\textbackslash{}n}\StringTok{Dict output:"}\NormalTok{)}
\BuiltInTok{print}\NormalTok{(result\_dict)}

\CommentTok{\# Text output {-} formatted string}
\NormalTok{result\_text }\OperatorTok{=}\NormalTok{ add.scan(}
    \StringTok{\textquotesingle{}@analyze\textquotesingle{}}\NormalTok{,}
\NormalTok{    df,}
\NormalTok{    as\_type}\OperatorTok{=}\StringTok{\textquotesingle{}text\textquotesingle{}}
\NormalTok{)}
\BuiltInTok{print}\NormalTok{(}\StringTok{"}\CharTok{\textbackslash{}n}\StringTok{Text output:"}\NormalTok{)}
\BuiltInTok{print}\NormalTok{(result\_text)}
\end{Highlighting}
\end{Shaded}

Output:

\begin{verbatim}
DataFrame output:
    column    dtype  count  null_count  null_pct  unique       mean        std    min    max
0  product  String(Utf8)      3           0       0.0       3        NaN        NaN    NaN    NaN
1    price  Float64      3           0       0.0       3  33.323333   15.01333  19.99  49.99
2    stock    Int64      3           0       0.0       3  10.000000    5.00000   5.00  15.00

Dict output:
{'columns': ['column', 'dtype', 'count', 'null_count', 'null_pct', 'unique', 'mean', 'std', 'min', 'max'], 'rows': 3, 'data': '...'}

Text output:
shape: (3, 10)
┌─────────┬──────────────┬───────┬────────────┬──────────┬────────┬───────────┬───────────┬───────┬───────┐
│ column  ┆ dtype        ┆ count ┆ null_count ┆ null_pct ┆ unique ┆ mean      ┆ std       ┆ min   ┆ max   │
│ ---     ┆ ---          ┆ ---   ┆ ---        ┆ ---      ┆ ---    ┆ ---       ┆ ---       ┆ ---   ┆ ---   │
│ str     ┆ str          ┆ i64   ┆ i64        ┆ f64      ┆ i64    ┆ f64       ┆ f64       ┆ f64   ┆ f64   │
╞═════════╪══════════════╪═══════╪════════════╪══════════╪════════╪═══════════╪═══════════╪═══════╪═══════╡
│ product ┆ String(Utf8) ┆ 3     ┆ 0          ┆ 0.0      ┆ 3      ┆ null      ┆ null      ┆ null  ┆ null  │
│ price   ┆ Float64      ┆ 3     ┆ 0          ┆ 0.0      ┆ 3      ┆ 33.323333 ┆ 15.013333 ┆ 19.99 ┆ 49.99 │
│ stock   ┆ Int64        ┆ 3     ┆ 0          ┆ 0.0      ┆ 3      ┆ 10.0      ┆ 5.0       ┆ 5.0   ┆ 15.0  │
└─────────┴──────────────┴───────┴────────────┴──────────┴────────┴───────────┴───────────┴───────┴───────┘
\end{verbatim}

\textbf{Output Format Options:} -
\texttt{as\_type=\textquotesingle{}dataframe\textquotesingle{}}
(default): Returns DataFrame for further analysis -
\texttt{as\_type=\textquotesingle{}dict\textquotesingle{}}: Returns dict
with columns, rows, and data -
\texttt{as\_type=\textquotesingle{}text\textquotesingle{}}: Returns
formatted string for display

\textbf{When to Use Each:} - DataFrame: When you need to filter, sort,
or further analyze the results - Dict: When you need structured data for
APIs or JSON serialization - Text: When you need human-readable output
for logging or display

Note: This also works with polars DataFrames.

\begin{center}\rule{0.5\linewidth}{0.5pt}\end{center}

\section{Parameters}\label{parameters-4}

\subsection{Required Parameters}\label{required-parameters-3}

\begin{itemize}
\tightlist
\item
  \texttt{mode}: \texttt{\textquotesingle{}@analyze\textquotesingle{}}
  for statistical profiling
\item
  \texttt{df}: Input DataFrame (pandas or polars)
\end{itemize}

\subsection{Optional Parameters}\label{optional-parameters-3}

\begin{itemize}
\tightlist
\item
  \texttt{columns}: Column filter (not yet implemented - currently
  analyzes all columns)
\item
  \texttt{where}: SQL-like filter condition (not yet implemented)
\item
  \texttt{rows}: Row range specifications (not yet implemented)
\item
  \texttt{as\_type}: Output format
  (\texttt{\textquotesingle{}dataframe\textquotesingle{}},
  \texttt{\textquotesingle{}dict\textquotesingle{}},
  \texttt{\textquotesingle{}text\textquotesingle{}})
\end{itemize}

\subsection{Positional Parameters}\label{positional-parameters-4}

\begin{Shaded}
\begin{Highlighting}[]
\CommentTok{\# Also works without naming certain parameters:}
\NormalTok{result }\OperatorTok{=}\NormalTok{ add.scan(}\StringTok{\textquotesingle{}@analyze\textquotesingle{}}\NormalTok{, df, as\_type}\OperatorTok{=}\StringTok{\textquotesingle{}dataframe\textquotesingle{}}\NormalTok{)}
\end{Highlighting}
\end{Shaded}

\begin{center}\rule{0.5\linewidth}{0.5pt}\end{center}

\section{Analysis Columns Reference}\label{analysis-columns-reference}

\begin{longtable}[]{@{}lll@{}}
\toprule\noalign{}
Column & Description & Type \\
\midrule\noalign{}
\endhead
\bottomrule\noalign{}
\endlastfoot
\texttt{column} & Column name & String \\
\texttt{dtype} & Data type & String \\
\texttt{count} & Total row count & Integer \\
\texttt{null\_count} & Number of null values & Integer \\
\texttt{null\_pct} & Percentage of nulls & Float \\
\texttt{unique} & Number of unique values & Integer \\
\texttt{mean} & Average (numeric only) & Float or null \\
\texttt{std} & Standard deviation (numeric only) & Float or null \\
\texttt{min} & Minimum value (numeric only) & Float or null \\
\texttt{max} & Maximum value (numeric only) & Float or null \\
\end{longtable}

\begin{center}\rule{0.5\linewidth}{0.5pt}\end{center}

\section{Next Steps}\label{next-steps-16}

\begin{itemize}
\tightlist
\item
  \textbf{Page 2}: Lineage tracking with \texttt{@lineage} mode
\item
  \textbf{Page 3}: Real-world data quality workflows
\end{itemize}

\chapter{Lineage Tracking}\label{lineage-tracking-1}

Track transformation history and understand data provenance using the
\texttt{@lineage} mode.

\section{Example 1: Basic Lineage
Tracking}\label{example-1-basic-lineage-tracking}

Enable lineage tracking by adding \texttt{lineage=True} to your
operations, then use \texttt{@lineage} mode to see the transformation
history.

\begin{Shaded}
\begin{Highlighting}[]
\ImportTok{import}\NormalTok{ pandas }\ImportTok{as}\NormalTok{ pd}
\ImportTok{import}\NormalTok{ additory }\ImportTok{as}\NormalTok{ add}

\NormalTok{df }\OperatorTok{=}\NormalTok{ pd.DataFrame(\{}
    \StringTok{\textquotesingle{}price\textquotesingle{}}\NormalTok{: [}\DecValTok{100}\NormalTok{, }\DecValTok{200}\NormalTok{, }\DecValTok{150}\NormalTok{],}
    \StringTok{\textquotesingle{}quantity\textquotesingle{}}\NormalTok{: [}\DecValTok{2}\NormalTok{, }\DecValTok{1}\NormalTok{, }\DecValTok{3}\NormalTok{]}
\NormalTok{\})}

\CommentTok{\# Perform operation with lineage tracking enabled}
\NormalTok{result }\OperatorTok{=}\NormalTok{ add.transform(}
    \StringTok{\textquotesingle{}@calc\textquotesingle{}}\NormalTok{,}
\NormalTok{    df,}
\NormalTok{    columns}\OperatorTok{=}\NormalTok{[}\StringTok{\textquotesingle{}price\textquotesingle{}}\NormalTok{, }\StringTok{\textquotesingle{}quantity\textquotesingle{}}\NormalTok{],}
\NormalTok{    expression}\OperatorTok{=}\StringTok{\textquotesingle{}price * quantity\textquotesingle{}}\NormalTok{,}
\NormalTok{    as\_}\OperatorTok{=}\StringTok{\textquotesingle{}total\textquotesingle{}}\NormalTok{,}
\NormalTok{    lineage}\OperatorTok{=}\VariableTok{True}  \CommentTok{\# Enable lineage tracking}
\NormalTok{)}

\CommentTok{\# Get lineage report}
\NormalTok{lineage\_report }\OperatorTok{=}\NormalTok{ add.scan(}
    \StringTok{\textquotesingle{}@lineage\textquotesingle{}}\NormalTok{,}
\NormalTok{    result}
\NormalTok{)}

\BuiltInTok{print}\NormalTok{(lineage\_report)}
\end{Highlighting}
\end{Shaded}

Output:

\begin{verbatim}
═══════════════════════════════════════════════════════════════
LINEAGE REPORT
═══════════════════════════════════════════════════════════════

DataFrame: 3 rows × 3 columns
Operations: 1 transformations applied

───────────────────────────────────────────────────────────────
Step 1: add.transform - 2026-03-11T09:53:59
───────────────────────────────────────────────────────────────
  Rows: 3 → 3 (no change)
  Columns Added: total
  
  Parameters:
    columns: ["price","quantity"]
    expression: ["price * quantity"]
    strategy: null
    by: null
\end{verbatim}

\textbf{Key Points:} - Add \texttt{lineage=True} to any operation
(transform, to, synthetic) - Lineage metadata is stored with the
DataFrame - Use \texttt{@lineage} mode to generate a human-readable
report - Report shows operation history, parameters, and column changes

Note: This also works with polars DataFrames.

\begin{center}\rule{0.5\linewidth}{0.5pt}\end{center}

\section{Example 2: Multi-Step Lineage
Tracking}\label{example-2-multi-step-lineage-tracking}

Lineage accumulates across multiple operations, giving you a complete
transformation history.

\begin{Shaded}
\begin{Highlighting}[]
\ImportTok{import}\NormalTok{ pandas }\ImportTok{as}\NormalTok{ pd}
\ImportTok{import}\NormalTok{ additory }\ImportTok{as}\NormalTok{ add}

\NormalTok{df }\OperatorTok{=}\NormalTok{ pd.DataFrame(\{}
    \StringTok{\textquotesingle{}price\textquotesingle{}}\NormalTok{: [}\DecValTok{100}\NormalTok{, }\DecValTok{200}\NormalTok{, }\DecValTok{150}\NormalTok{],}
    \StringTok{\textquotesingle{}cost\textquotesingle{}}\NormalTok{: [}\DecValTok{60}\NormalTok{, }\DecValTok{120}\NormalTok{, }\DecValTok{90}\NormalTok{],}
    \StringTok{\textquotesingle{}quantity\textquotesingle{}}\NormalTok{: [}\DecValTok{2}\NormalTok{, }\DecValTok{1}\NormalTok{, }\DecValTok{3}\NormalTok{]}
\NormalTok{\})}

\CommentTok{\# Step 1: Calculate profit with lineage tracking}
\NormalTok{result }\OperatorTok{=}\NormalTok{ add.transform(}
    \StringTok{\textquotesingle{}@calc\textquotesingle{}}\NormalTok{,}
\NormalTok{    df,}
\NormalTok{    columns}\OperatorTok{=}\NormalTok{[}\StringTok{\textquotesingle{}price\textquotesingle{}}\NormalTok{, }\StringTok{\textquotesingle{}cost\textquotesingle{}}\NormalTok{],}
\NormalTok{    expression}\OperatorTok{=}\StringTok{\textquotesingle{}price {-} cost\textquotesingle{}}\NormalTok{,}
\NormalTok{    as\_}\OperatorTok{=}\StringTok{\textquotesingle{}profit\textquotesingle{}}\NormalTok{,}
\NormalTok{    lineage}\OperatorTok{=}\VariableTok{True}
\NormalTok{)}

\CommentTok{\# Step 2: Calculate total revenue (lineage continues)}
\NormalTok{result }\OperatorTok{=}\NormalTok{ add.transform(}
    \StringTok{\textquotesingle{}@calc\textquotesingle{}}\NormalTok{,}
\NormalTok{    result,}
\NormalTok{    columns}\OperatorTok{=}\NormalTok{[}\StringTok{\textquotesingle{}price\textquotesingle{}}\NormalTok{, }\StringTok{\textquotesingle{}quantity\textquotesingle{}}\NormalTok{],}
\NormalTok{    expression}\OperatorTok{=}\StringTok{\textquotesingle{}price * quantity\textquotesingle{}}\NormalTok{,}
\NormalTok{    as\_}\OperatorTok{=}\StringTok{\textquotesingle{}revenue\textquotesingle{}}\NormalTok{,}
\NormalTok{    lineage}\OperatorTok{=}\VariableTok{True}
\NormalTok{)}

\CommentTok{\# Get complete lineage report}
\NormalTok{lineage\_report }\OperatorTok{=}\NormalTok{ add.scan(}
    \StringTok{\textquotesingle{}@lineage\textquotesingle{}}\NormalTok{,}
\NormalTok{    result}
\NormalTok{)}

\BuiltInTok{print}\NormalTok{(lineage\_report)}
\end{Highlighting}
\end{Shaded}

Output:

\begin{verbatim}
═══════════════════════════════════════════════════════════════
LINEAGE REPORT
═══════════════════════════════════════════════════════════════

DataFrame: 3 rows × 5 columns
Operations: 2 transformations applied

───────────────────────────────────────────────────────────────
Step 1: add.transform - 2026-03-11T10:15:23
───────────────────────────────────────────────────────────────
  Rows: 3 → 3 (no change)
  Columns Added: profit
  
  Parameters:
    columns: ["price","cost"]
    expression: ["price - cost"]

───────────────────────────────────────────────────────────────
Step 2: add.transform - 2026-03-11T10:15:23
───────────────────────────────────────────────────────────────
  Rows: 3 → 3 (no change)
  Columns Added: revenue
  
  Parameters:
    columns: ["price","quantity"]
    expression: ["price * quantity"]
\end{verbatim}

\textbf{Key Insights:} - Lineage persists across multiple operations -
Each step is numbered and timestamped - You can see the complete
transformation pipeline - Useful for debugging complex data workflows

Note: This also works with polars DataFrames.

\begin{center}\rule{0.5\linewidth}{0.5pt}\end{center}

\section{Example 3: Lineage Without Tracking - Error
Handling}\label{example-3-lineage-without-tracking---error-handling}

If you forget to enable lineage tracking, \texttt{@lineage} mode
provides a helpful error message.

\begin{Shaded}
\begin{Highlighting}[]
\ImportTok{import}\NormalTok{ pandas }\ImportTok{as}\NormalTok{ pd}
\ImportTok{import}\NormalTok{ additory }\ImportTok{as}\NormalTok{ add}

\NormalTok{df }\OperatorTok{=}\NormalTok{ pd.DataFrame(\{}
    \StringTok{\textquotesingle{}a\textquotesingle{}}\NormalTok{: [}\DecValTok{1}\NormalTok{, }\DecValTok{2}\NormalTok{, }\DecValTok{3}\NormalTok{],}
    \StringTok{\textquotesingle{}b\textquotesingle{}}\NormalTok{: [}\DecValTok{4}\NormalTok{, }\DecValTok{5}\NormalTok{, }\DecValTok{6}\NormalTok{]}
\NormalTok{\})}

\CommentTok{\# Transform WITHOUT lineage tracking}
\NormalTok{result }\OperatorTok{=}\NormalTok{ add.transform(}
    \StringTok{\textquotesingle{}@calc\textquotesingle{}}\NormalTok{,}
\NormalTok{    df,}
\NormalTok{    columns}\OperatorTok{=}\NormalTok{[}\StringTok{\textquotesingle{}a\textquotesingle{}}\NormalTok{, }\StringTok{\textquotesingle{}b\textquotesingle{}}\NormalTok{],}
\NormalTok{    expression}\OperatorTok{=}\StringTok{\textquotesingle{}a + b\textquotesingle{}}\NormalTok{,}
\NormalTok{    as\_}\OperatorTok{=}\StringTok{\textquotesingle{}c\textquotesingle{}}
    \CommentTok{\# Note: lineage=True is missing!}
\NormalTok{)}

\CommentTok{\# Try to get lineage report}
\ControlFlowTok{try}\NormalTok{:}
\NormalTok{    lineage\_report }\OperatorTok{=}\NormalTok{ add.scan(}\StringTok{\textquotesingle{}@lineage\textquotesingle{}}\NormalTok{, result)}
\ControlFlowTok{except} \PreprocessorTok{ValueError} \ImportTok{as}\NormalTok{ e:}
    \BuiltInTok{print}\NormalTok{(e)}
\end{Highlighting}
\end{Shaded}

Output:

\begin{verbatim}
No lineage metadata found. Lineage tracking must be enabled by adding
lineage=True to add.to(), add.transform(), or add.synthetic() calls.

Example:
  df = add.transform('@calc', df, strategy={'total': 'price * qty'}, lineage=True)
  result = add.scan('@lineage', df)
\end{verbatim}

\textbf{Error Handling:} - Clear error message when lineage is missing -
Provides example of how to enable lineage - Helps you quickly fix the
issue

Note: This also works with polars DataFrames.

\begin{center}\rule{0.5\linewidth}{0.5pt}\end{center}

\section{Lineage Tracking Features}\label{lineage-tracking-features}

\subsection{What Gets Tracked}\label{what-gets-tracked}

\begin{itemize}
\tightlist
\item
  \textbf{Operation Type}: Which function was called (transform, to,
  synthetic)
\item
  \textbf{Timestamp}: When the operation occurred
\item
  \textbf{Parameters}: All operation parameters
\item
  \textbf{Row Changes}: Rows before and after
\item
  \textbf{Column Changes}: Columns added or modified
\end{itemize}

\subsection{When to Use Lineage}\label{when-to-use-lineage}

\begin{itemize}
\tightlist
\item
  \textbf{Debugging}: Understand how data was transformed
\item
  \textbf{Auditing}: Track data provenance for compliance
\item
  \textbf{Documentation}: Auto-generate transformation documentation
\item
  \textbf{Collaboration}: Share transformation history with team
\end{itemize}

\subsection{Performance Impact}\label{performance-impact}

\begin{itemize}
\tightlist
\item
  Minimal overhead (\textless100ms per operation)
\item
  Metadata stored in memory only
\item
  No impact on computation speed
\end{itemize}

\begin{center}\rule{0.5\linewidth}{0.5pt}\end{center}

\section{Parameters}\label{parameters-5}

\subsection{Required Parameters}\label{required-parameters-4}

\begin{itemize}
\tightlist
\item
  \texttt{mode}: \texttt{\textquotesingle{}@lineage\textquotesingle{}}
  for lineage tracking
\item
  \texttt{df}: Input DataFrame with lineage metadata
\end{itemize}

\subsection{Optional Parameters}\label{optional-parameters-4}

\begin{itemize}
\tightlist
\item
  \texttt{columns}: Filter report to specific columns (not yet
  implemented)
\item
  \texttt{trace}: Trace specific cell transformations (not yet
  implemented)
\item
  \texttt{as\_type}: Output format (currently only `text' supported)
\end{itemize}

\subsection{Positional Parameters}\label{positional-parameters-5}

\begin{Shaded}
\begin{Highlighting}[]
\CommentTok{\# Also works without naming certain parameters:}
\NormalTok{lineage\_report }\OperatorTok{=}\NormalTok{ add.scan(}\StringTok{\textquotesingle{}@lineage\textquotesingle{}}\NormalTok{, df)}
\end{Highlighting}
\end{Shaded}

\begin{center}\rule{0.5\linewidth}{0.5pt}\end{center}

\section{Enabling Lineage Tracking}\label{enabling-lineage-tracking}

\subsection{In add.transform()}\label{in-add.transform}

\begin{Shaded}
\begin{Highlighting}[]
\NormalTok{result }\OperatorTok{=}\NormalTok{ add.transform(}
    \StringTok{\textquotesingle{}@calc\textquotesingle{}}\NormalTok{,}
\NormalTok{    df,}
\NormalTok{    columns}\OperatorTok{=}\NormalTok{[}\StringTok{\textquotesingle{}a\textquotesingle{}}\NormalTok{, }\StringTok{\textquotesingle{}b\textquotesingle{}}\NormalTok{],}
\NormalTok{    expression}\OperatorTok{=}\StringTok{\textquotesingle{}a + b\textquotesingle{}}\NormalTok{,}
\NormalTok{    as\_}\OperatorTok{=}\StringTok{\textquotesingle{}c\textquotesingle{}}\NormalTok{,}
\NormalTok{    lineage}\OperatorTok{=}\VariableTok{True}  \CommentTok{\# Enable lineage}
\NormalTok{)}
\end{Highlighting}
\end{Shaded}

\subsection{In add.to()}\label{in-add.to}

\begin{Shaded}
\begin{Highlighting}[]
\NormalTok{result }\OperatorTok{=}\NormalTok{ add.to(}
\NormalTok{    orders,}
\NormalTok{    bring\_from}\OperatorTok{=}\NormalTok{customers,}
\NormalTok{    bring}\OperatorTok{=}\StringTok{\textquotesingle{}name\textquotesingle{}}\NormalTok{,}
\NormalTok{    against}\OperatorTok{=}\StringTok{\textquotesingle{}customer\_id\textquotesingle{}}\NormalTok{,}
\NormalTok{    lineage}\OperatorTok{=}\VariableTok{True}  \CommentTok{\# Enable lineage}
\NormalTok{)}
\end{Highlighting}
\end{Shaded}

\subsection{In add.synthetic()}\label{in-add.synthetic}

\begin{Shaded}
\begin{Highlighting}[]
\NormalTok{result }\OperatorTok{=}\NormalTok{ add.synthetic(}
    \StringTok{\textquotesingle{}@new\textquotesingle{}}\NormalTok{,}
\NormalTok{    n}\OperatorTok{=}\DecValTok{100}\NormalTok{,}
\NormalTok{    strategy}\OperatorTok{=}\NormalTok{\{}\StringTok{\textquotesingle{}a\textquotesingle{}}\NormalTok{: \{}\StringTok{\textquotesingle{}distribution\textquotesingle{}}\NormalTok{: }\StringTok{\textquotesingle{}normal\textquotesingle{}}\NormalTok{, }\StringTok{\textquotesingle{}mean\textquotesingle{}}\NormalTok{: }\DecValTok{0}\NormalTok{, }\StringTok{\textquotesingle{}std\textquotesingle{}}\NormalTok{: }\DecValTok{1}\NormalTok{\}\},}
\NormalTok{    lineage}\OperatorTok{=}\VariableTok{True}  \CommentTok{\# Enable lineage}
\NormalTok{)}
\end{Highlighting}
\end{Shaded}

\begin{center}\rule{0.5\linewidth}{0.5pt}\end{center}

\section{Limitations (v0.1.4)}\label{limitations-v0.1.4}

\subsection{Session-Only Lineage}\label{session-only-lineage}

\begin{itemize}
\tightlist
\item
  Lineage metadata is stored in memory only
\item
  Not persisted when saving DataFrames to disk
\item
  Lost when Python session ends
\end{itemize}

\subsection{Future Enhancements
(v0.2.0)}\label{future-enhancements-v0.2.0}

\begin{itemize}
\tightlist
\item
  Persistent lineage with \texttt{add.save()} and \texttt{add.load()}
\item
  Cell-level tracing with \texttt{trace} parameter
\item
  Column-specific lineage with \texttt{columns} parameter
\item
  Lineage visualization and export
\end{itemize}

\begin{center}\rule{0.5\linewidth}{0.5pt}\end{center}

\section{Next Steps}\label{next-steps-17}

\begin{itemize}
\tightlist
\item
  \textbf{Page 1}: Basic data scanning with \texttt{@analyze} mode
\item
  \textbf{Page 3}: Real-world data quality workflows
\end{itemize}

\chapter{Real-World Data Quality
Workflows}\label{real-world-data-quality-workflows}

Practical workflows combining \texttt{@analyze} mode with data quality
assessment, profiling, and validation.

\section{Example 1: Data Quality Assessment
Pipeline}\label{example-1-data-quality-assessment-pipeline}

Build a comprehensive data quality assessment pipeline to identify and
report quality issues.

\begin{Shaded}
\begin{Highlighting}[]
\ImportTok{import}\NormalTok{ pandas }\ImportTok{as}\NormalTok{ pd}
\ImportTok{import}\NormalTok{ additory }\ImportTok{as}\NormalTok{ add}

\CommentTok{\# Simulate a dataset with quality issues}
\NormalTok{df }\OperatorTok{=}\NormalTok{ pd.DataFrame(\{}
    \StringTok{\textquotesingle{}customer\_id\textquotesingle{}}\NormalTok{: [}\DecValTok{1}\NormalTok{, }\DecValTok{2}\NormalTok{, }\DecValTok{3}\NormalTok{, }\DecValTok{4}\NormalTok{, }\DecValTok{5}\NormalTok{, }\DecValTok{6}\NormalTok{, }\DecValTok{7}\NormalTok{, }\DecValTok{8}\NormalTok{, }\DecValTok{9}\NormalTok{, }\DecValTok{10}\NormalTok{],}
    \StringTok{\textquotesingle{}name\textquotesingle{}}\NormalTok{: [}\StringTok{\textquotesingle{}Alice\textquotesingle{}}\NormalTok{, }\StringTok{\textquotesingle{}Bob\textquotesingle{}}\NormalTok{, }\VariableTok{None}\NormalTok{, }\StringTok{\textquotesingle{}David\textquotesingle{}}\NormalTok{, }\StringTok{\textquotesingle{}Eve\textquotesingle{}}\NormalTok{, }\StringTok{\textquotesingle{}Frank\textquotesingle{}}\NormalTok{, }\VariableTok{None}\NormalTok{, }\StringTok{\textquotesingle{}Helen\textquotesingle{}}\NormalTok{, }\StringTok{\textquotesingle{}Ivan\textquotesingle{}}\NormalTok{, }\StringTok{\textquotesingle{}Jane\textquotesingle{}}\NormalTok{],}
    \StringTok{\textquotesingle{}email\textquotesingle{}}\NormalTok{: [}\StringTok{\textquotesingle{}alice@email.com\textquotesingle{}}\NormalTok{, }\StringTok{\textquotesingle{}bob@email.com\textquotesingle{}}\NormalTok{, }\StringTok{\textquotesingle{}charlie@email.com\textquotesingle{}}\NormalTok{, }\VariableTok{None}\NormalTok{, }
              \StringTok{\textquotesingle{}eve@email.com\textquotesingle{}}\NormalTok{, }\VariableTok{None}\NormalTok{, }\StringTok{\textquotesingle{}grace@email.com\textquotesingle{}}\NormalTok{, }\StringTok{\textquotesingle{}helen@email.com\textquotesingle{}}\NormalTok{, }
              \VariableTok{None}\NormalTok{, }\StringTok{\textquotesingle{}jane@email.com\textquotesingle{}}\NormalTok{],}
    \StringTok{\textquotesingle{}age\textquotesingle{}}\NormalTok{: [}\DecValTok{25}\NormalTok{, }\DecValTok{30}\NormalTok{, }\DecValTok{35}\NormalTok{, }\DecValTok{40}\NormalTok{, }\VariableTok{None}\NormalTok{, }\DecValTok{50}\NormalTok{, }\DecValTok{55}\NormalTok{, }\VariableTok{None}\NormalTok{, }\DecValTok{65}\NormalTok{, }\DecValTok{70}\NormalTok{],}
    \StringTok{\textquotesingle{}purchase\_count\textquotesingle{}}\NormalTok{: [}\DecValTok{5}\NormalTok{, }\DecValTok{10}\NormalTok{, }\DecValTok{15}\NormalTok{, }\DecValTok{20}\NormalTok{, }\DecValTok{25}\NormalTok{, }\DecValTok{30}\NormalTok{, }\DecValTok{35}\NormalTok{, }\DecValTok{40}\NormalTok{, }\DecValTok{45}\NormalTok{, }\DecValTok{50}\NormalTok{]}
\NormalTok{\})}

\CommentTok{\# Step 1: Analyze data quality}
\NormalTok{quality\_report }\OperatorTok{=}\NormalTok{ add.scan(}
    \StringTok{\textquotesingle{}@analyze\textquotesingle{}}\NormalTok{,}
\NormalTok{    df,}
\NormalTok{    as\_type}\OperatorTok{=}\StringTok{\textquotesingle{}dataframe\textquotesingle{}}
\NormalTok{)}

\BuiltInTok{print}\NormalTok{(}\StringTok{"Data Quality Report:"}\NormalTok{)}
\BuiltInTok{print}\NormalTok{(quality\_report)}

\CommentTok{\# Step 2: Identify columns with quality issues (\textgreater{}10\% nulls)}
\NormalTok{high\_null\_cols }\OperatorTok{=}\NormalTok{ quality\_report[quality\_report[}\StringTok{\textquotesingle{}null\_pct\textquotesingle{}}\NormalTok{] }\OperatorTok{\textgreater{}} \FloatTok{10.0}\NormalTok{]}

\BuiltInTok{print}\NormalTok{(}\StringTok{"}\CharTok{\textbackslash{}n}\StringTok{Columns with \textgreater{}10\% null values:"}\NormalTok{)}
\BuiltInTok{print}\NormalTok{(high\_null\_cols[[}\StringTok{\textquotesingle{}column\textquotesingle{}}\NormalTok{, }\StringTok{\textquotesingle{}null\_count\textquotesingle{}}\NormalTok{, }\StringTok{\textquotesingle{}null\_pct\textquotesingle{}}\NormalTok{]])}

\CommentTok{\# Step 3: Get summary statistics}
\NormalTok{summary }\OperatorTok{=}\NormalTok{ \{}
    \StringTok{\textquotesingle{}total\_columns\textquotesingle{}}\NormalTok{: }\BuiltInTok{len}\NormalTok{(quality\_report),}
    \StringTok{\textquotesingle{}columns\_with\_nulls\textquotesingle{}}\NormalTok{: }\BuiltInTok{len}\NormalTok{(quality\_report[quality\_report[}\StringTok{\textquotesingle{}null\_count\textquotesingle{}}\NormalTok{] }\OperatorTok{\textgreater{}} \DecValTok{0}\NormalTok{]),}
    \StringTok{\textquotesingle{}high\_null\_columns\textquotesingle{}}\NormalTok{: }\BuiltInTok{len}\NormalTok{(high\_null\_cols),}
    \StringTok{\textquotesingle{}avg\_null\_pct\textquotesingle{}}\NormalTok{: quality\_report[}\StringTok{\textquotesingle{}null\_pct\textquotesingle{}}\NormalTok{].mean()}
\NormalTok{\}}

\BuiltInTok{print}\NormalTok{(}\StringTok{"}\CharTok{\textbackslash{}n}\StringTok{Quality Summary:"}\NormalTok{)}
\ControlFlowTok{for}\NormalTok{ key, value }\KeywordTok{in}\NormalTok{ summary.items():}
    \BuiltInTok{print}\NormalTok{(}\SpecialStringTok{f"  }\SpecialCharTok{\{}\NormalTok{key}\SpecialCharTok{\}}\SpecialStringTok{: }\SpecialCharTok{\{}\NormalTok{value}\SpecialCharTok{\}}\SpecialStringTok{"}\NormalTok{)}
\end{Highlighting}
\end{Shaded}

Output:

\begin{verbatim}
Data Quality Report:
        column       dtype  count  null_count  null_pct  unique   mean       std    min    max
0  customer_id       Int64     10           0       0.0      10    5.5  3.027650    1.0   10.0
1         name  String(Utf8)     10           2      20.0       8    NaN       NaN    NaN    NaN
2        email  String(Utf8)     10           3      30.0       7    NaN       NaN    NaN    NaN
3          age       Int64     10           2      20.0       8   47.5  16.583124   25.0   70.0
4 purchase_count       Int64     10           0       0.0      10   27.5  15.732133    5.0   50.0

Columns with >10% null values:
    column  null_count  null_pct
1     name           2      20.0
2    email           3      30.0
3      age           2      20.0

Quality Summary:
  total_columns: 5
  columns_with_nulls: 3
  high_null_columns: 3
  avg_null_pct: 14.0
\end{verbatim}

\textbf{Key Workflow Steps:} 1. Run \texttt{@analyze} to get
comprehensive quality metrics 2. Filter results to identify problematic
columns 3. Generate summary statistics for reporting 4. Use insights to
prioritize data cleaning efforts

\textbf{Quality Thresholds:} - \texttt{null\_pct\ \textgreater{}\ 10\%}:
High null rate, requires attention -
\texttt{null\_pct\ \textgreater{}\ 50\%}: Critical quality issue,
consider dropping column - \texttt{unique\ ==\ count}: Potential ID
column or high cardinality - \texttt{unique\ ==\ 1}: Constant column, no
information value

Note: This also works with polars DataFrames.

\begin{center}\rule{0.5\linewidth}{0.5pt}\end{center}

\section{Example 2: Column Profiling
Workflow}\label{example-2-column-profiling-workflow}

Profile columns by data type to understand distributions and cardinality
patterns.

\begin{Shaded}
\begin{Highlighting}[]
\ImportTok{import}\NormalTok{ pandas }\ImportTok{as}\NormalTok{ pd}
\ImportTok{import}\NormalTok{ additory }\ImportTok{as}\NormalTok{ add}

\CommentTok{\# Sales data with various data types}
\NormalTok{df }\OperatorTok{=}\NormalTok{ pd.DataFrame(\{}
    \StringTok{\textquotesingle{}transaction\_id\textquotesingle{}}\NormalTok{: }\BuiltInTok{range}\NormalTok{(}\DecValTok{1}\NormalTok{, }\DecValTok{101}\NormalTok{),}
    \StringTok{\textquotesingle{}amount\textquotesingle{}}\NormalTok{: [}\DecValTok{100} \OperatorTok{+}\NormalTok{ i }\OperatorTok{*} \DecValTok{10} \ControlFlowTok{for}\NormalTok{ i }\KeywordTok{in} \BuiltInTok{range}\NormalTok{(}\DecValTok{100}\NormalTok{)],}
    \StringTok{\textquotesingle{}category\textquotesingle{}}\NormalTok{: [}\StringTok{\textquotesingle{}Electronics\textquotesingle{}}\NormalTok{] }\OperatorTok{*} \DecValTok{30} \OperatorTok{+}\NormalTok{ [}\StringTok{\textquotesingle{}Clothing\textquotesingle{}}\NormalTok{] }\OperatorTok{*} \DecValTok{40} \OperatorTok{+}\NormalTok{ [}\StringTok{\textquotesingle{}Food\textquotesingle{}}\NormalTok{] }\OperatorTok{*} \DecValTok{30}\NormalTok{,}
    \StringTok{\textquotesingle{}payment\_method\textquotesingle{}}\NormalTok{: [}\StringTok{\textquotesingle{}Credit\textquotesingle{}}\NormalTok{] }\OperatorTok{*} \DecValTok{50} \OperatorTok{+}\NormalTok{ [}\StringTok{\textquotesingle{}Debit\textquotesingle{}}\NormalTok{] }\OperatorTok{*} \DecValTok{30} \OperatorTok{+}\NormalTok{ [}\StringTok{\textquotesingle{}Cash\textquotesingle{}}\NormalTok{] }\OperatorTok{*} \DecValTok{20}\NormalTok{,}
    \StringTok{\textquotesingle{}customer\_age\textquotesingle{}}\NormalTok{: [}\DecValTok{25} \OperatorTok{+}\NormalTok{ (i }\OperatorTok{\%} \DecValTok{50}\NormalTok{) }\ControlFlowTok{for}\NormalTok{ i }\KeywordTok{in} \BuiltInTok{range}\NormalTok{(}\DecValTok{100}\NormalTok{)]}
\NormalTok{\})}

\CommentTok{\# Profile the data}
\NormalTok{profile }\OperatorTok{=}\NormalTok{ add.scan(}
    \StringTok{\textquotesingle{}@analyze\textquotesingle{}}\NormalTok{,}
\NormalTok{    df,}
\NormalTok{    as\_type}\OperatorTok{=}\StringTok{\textquotesingle{}dataframe\textquotesingle{}}
\NormalTok{)}

\BuiltInTok{print}\NormalTok{(}\StringTok{"Column Profile:"}\NormalTok{)}
\BuiltInTok{print}\NormalTok{(profile)}

\CommentTok{\# Extract insights for each column type}
\NormalTok{numeric\_cols }\OperatorTok{=}\NormalTok{ profile[profile[}\StringTok{\textquotesingle{}dtype\textquotesingle{}}\NormalTok{].}\BuiltInTok{str}\NormalTok{.contains(}\StringTok{\textquotesingle{}Int|Float\textquotesingle{}}\NormalTok{, na}\OperatorTok{=}\VariableTok{False}\NormalTok{)]}
\NormalTok{categorical\_cols }\OperatorTok{=}\NormalTok{ profile[profile[}\StringTok{\textquotesingle{}dtype\textquotesingle{}}\NormalTok{].}\BuiltInTok{str}\NormalTok{.contains(}\StringTok{\textquotesingle{}String\textquotesingle{}}\NormalTok{, na}\OperatorTok{=}\VariableTok{False}\NormalTok{)]}

\BuiltInTok{print}\NormalTok{(}\StringTok{"}\CharTok{\textbackslash{}n}\StringTok{Numeric Columns:"}\NormalTok{)}
\BuiltInTok{print}\NormalTok{(numeric\_cols[[}\StringTok{\textquotesingle{}column\textquotesingle{}}\NormalTok{, }\StringTok{\textquotesingle{}mean\textquotesingle{}}\NormalTok{, }\StringTok{\textquotesingle{}std\textquotesingle{}}\NormalTok{, }\StringTok{\textquotesingle{}min\textquotesingle{}}\NormalTok{, }\StringTok{\textquotesingle{}max\textquotesingle{}}\NormalTok{]])}

\BuiltInTok{print}\NormalTok{(}\StringTok{"}\CharTok{\textbackslash{}n}\StringTok{Categorical Columns:"}\NormalTok{)}
\BuiltInTok{print}\NormalTok{(categorical\_cols[[}\StringTok{\textquotesingle{}column\textquotesingle{}}\NormalTok{, }\StringTok{\textquotesingle{}unique\textquotesingle{}}\NormalTok{, }\StringTok{\textquotesingle{}count\textquotesingle{}}\NormalTok{]])}

\CommentTok{\# Get cardinality insights}
\NormalTok{insights }\OperatorTok{=}\NormalTok{ \{}
    \StringTok{\textquotesingle{}numeric\_columns\textquotesingle{}}\NormalTok{: }\BuiltInTok{len}\NormalTok{(numeric\_cols),}
    \StringTok{\textquotesingle{}categorical\_columns\textquotesingle{}}\NormalTok{: }\BuiltInTok{len}\NormalTok{(categorical\_cols),}
    \StringTok{\textquotesingle{}high\_cardinality\textquotesingle{}}\NormalTok{: }\BuiltInTok{len}\NormalTok{(profile[profile[}\StringTok{\textquotesingle{}unique\textquotesingle{}}\NormalTok{] }\OperatorTok{\textgreater{}} \DecValTok{50}\NormalTok{]),}
    \StringTok{\textquotesingle{}low\_cardinality\textquotesingle{}}\NormalTok{: }\BuiltInTok{len}\NormalTok{(profile[profile[}\StringTok{\textquotesingle{}unique\textquotesingle{}}\NormalTok{] }\OperatorTok{\textless{}=} \DecValTok{10}\NormalTok{])}
\NormalTok{\}}

\BuiltInTok{print}\NormalTok{(}\StringTok{"}\CharTok{\textbackslash{}n}\StringTok{Cardinality Insights:"}\NormalTok{)}
\ControlFlowTok{for}\NormalTok{ key, value }\KeywordTok{in}\NormalTok{ insights.items():}
    \BuiltInTok{print}\NormalTok{(}\SpecialStringTok{f"  }\SpecialCharTok{\{}\NormalTok{key}\SpecialCharTok{\}}\SpecialStringTok{: }\SpecialCharTok{\{}\NormalTok{value}\SpecialCharTok{\}}\SpecialStringTok{"}\NormalTok{)}
\end{Highlighting}
\end{Shaded}

Output:

\begin{verbatim}
Column Profile:
          column    dtype  count  null_count  null_pct  unique      mean        std     min      max
0  transaction_id    Int64    100           0       0.0     100      50.5  29.011492     1.0    100.0
1          amount    Int64    100           0       0.0     100     595.0  295.594641   100.0   1090.0
2        category  String(Utf8)    100           0       0.0       3       NaN        NaN     NaN      NaN
3  payment_method  String(Utf8)    100           0       0.0       3       NaN        NaN     NaN      NaN
4    customer_age    Int64    100           0       0.0      50      49.5  14.430869    25.0     74.0

Numeric Columns:
          column    mean        std     min      max
0  transaction_id    50.5  29.011492     1.0    100.0
1          amount   595.0  295.594641   100.0   1090.0
4    customer_age    49.5  14.430869    25.0     74.0

Categorical Columns:
          column  unique  count
2        category       3    100
3  payment_method       3    100

Cardinality Insights:
  numeric_columns: 3
  categorical_columns: 2
  high_cardinality: 2
  low_cardinality: 3
\end{verbatim}

\textbf{Profiling Insights:} - \textbf{High Cardinality} (unique
\textgreater{} 50): Likely ID columns or continuous numeric data -
\textbf{Low Cardinality} (unique ≤ 10): Good candidates for categorical
encoding - \textbf{Numeric Statistics}: Use mean/std to detect outliers
and distributions - \textbf{Categorical Distribution}: Use unique count
to assess encoding strategy

\textbf{Use Cases:} - Feature engineering: Identify columns for one-hot
encoding vs label encoding - Data modeling: Understand which columns
need normalization - Storage optimization: High cardinality strings may
benefit from dictionary encoding

Note: This also works with polars DataFrames.

\begin{center}\rule{0.5\linewidth}{0.5pt}\end{center}

\section{Example 3: Automated Data
Validation}\label{example-3-automated-data-validation}

Build automated validation checks to ensure data quality before
processing.

\begin{Shaded}
\begin{Highlighting}[]
\ImportTok{import}\NormalTok{ pandas }\ImportTok{as}\NormalTok{ pd}
\ImportTok{import}\NormalTok{ additory }\ImportTok{as}\NormalTok{ add}

\CommentTok{\# Dataset that should pass validation}
\NormalTok{df\_valid }\OperatorTok{=}\NormalTok{ pd.DataFrame(\{}
    \StringTok{\textquotesingle{}product\_id\textquotesingle{}}\NormalTok{: [}\DecValTok{1}\NormalTok{, }\DecValTok{2}\NormalTok{, }\DecValTok{3}\NormalTok{, }\DecValTok{4}\NormalTok{, }\DecValTok{5}\NormalTok{],}
    \StringTok{\textquotesingle{}price\textquotesingle{}}\NormalTok{: [}\FloatTok{10.99}\NormalTok{, }\FloatTok{20.99}\NormalTok{, }\FloatTok{30.99}\NormalTok{, }\FloatTok{40.99}\NormalTok{, }\FloatTok{50.99}\NormalTok{],}
    \StringTok{\textquotesingle{}stock\textquotesingle{}}\NormalTok{: [}\DecValTok{100}\NormalTok{, }\DecValTok{200}\NormalTok{, }\DecValTok{150}\NormalTok{, }\DecValTok{175}\NormalTok{, }\DecValTok{225}\NormalTok{]}
\NormalTok{\})}

\CommentTok{\# Dataset with validation issues}
\NormalTok{df\_invalid }\OperatorTok{=}\NormalTok{ pd.DataFrame(\{}
    \StringTok{\textquotesingle{}product\_id\textquotesingle{}}\NormalTok{: [}\DecValTok{1}\NormalTok{, }\DecValTok{2}\NormalTok{, }\VariableTok{None}\NormalTok{, }\DecValTok{4}\NormalTok{, }\DecValTok{5}\NormalTok{],}
    \StringTok{\textquotesingle{}price\textquotesingle{}}\NormalTok{: [}\FloatTok{10.99}\NormalTok{, }\VariableTok{None}\NormalTok{, }\FloatTok{30.99}\NormalTok{, }\VariableTok{None}\NormalTok{, }\FloatTok{50.99}\NormalTok{],}
    \StringTok{\textquotesingle{}stock\textquotesingle{}}\NormalTok{: [}\DecValTok{100}\NormalTok{, }\DecValTok{200}\NormalTok{, }\DecValTok{150}\NormalTok{, }\DecValTok{175}\NormalTok{, }\VariableTok{None}\NormalTok{]}
\NormalTok{\})}

\CommentTok{\# Validate both datasets}
\NormalTok{report\_valid }\OperatorTok{=}\NormalTok{ add.scan(}\StringTok{\textquotesingle{}@analyze\textquotesingle{}}\NormalTok{, df\_valid, as\_type}\OperatorTok{=}\StringTok{\textquotesingle{}dataframe\textquotesingle{}}\NormalTok{)}
\NormalTok{report\_invalid }\OperatorTok{=}\NormalTok{ add.scan(}\StringTok{\textquotesingle{}@analyze\textquotesingle{}}\NormalTok{, df\_invalid, as\_type}\OperatorTok{=}\StringTok{\textquotesingle{}dataframe\textquotesingle{}}\NormalTok{)}

\CommentTok{\# Check validation rules}
\NormalTok{valid\_passed }\OperatorTok{=}\NormalTok{ (report\_valid[}\StringTok{\textquotesingle{}null\_count\textquotesingle{}}\NormalTok{] }\OperatorTok{==} \DecValTok{0}\NormalTok{).}\BuiltInTok{all}\NormalTok{()}
\NormalTok{invalid\_passed }\OperatorTok{=}\NormalTok{ (report\_invalid[}\StringTok{\textquotesingle{}null\_count\textquotesingle{}}\NormalTok{] }\OperatorTok{==} \DecValTok{0}\NormalTok{).}\BuiltInTok{all}\NormalTok{()}

\BuiltInTok{print}\NormalTok{(}\StringTok{"Valid Dataset:"}\NormalTok{)}
\BuiltInTok{print}\NormalTok{(}\SpecialStringTok{f"  Validation Passed: }\SpecialCharTok{\{}\NormalTok{valid\_passed}\SpecialCharTok{\}}\SpecialStringTok{"}\NormalTok{)}
\BuiltInTok{print}\NormalTok{(}\SpecialStringTok{f"  Null Values: }\SpecialCharTok{\{}\NormalTok{report\_valid[}\StringTok{\textquotesingle{}null\_count\textquotesingle{}}\NormalTok{]}\SpecialCharTok{.}\BuiltInTok{sum}\NormalTok{()}\SpecialCharTok{\}}\SpecialStringTok{"}\NormalTok{)}

\BuiltInTok{print}\NormalTok{(}\StringTok{"}\CharTok{\textbackslash{}n}\StringTok{Invalid Dataset:"}\NormalTok{)}
\BuiltInTok{print}\NormalTok{(}\SpecialStringTok{f"  Validation Passed: }\SpecialCharTok{\{}\NormalTok{invalid\_passed}\SpecialCharTok{\}}\SpecialStringTok{"}\NormalTok{)}
\BuiltInTok{print}\NormalTok{(}\SpecialStringTok{f"  Null Values: }\SpecialCharTok{\{}\NormalTok{report\_invalid[}\StringTok{\textquotesingle{}null\_count\textquotesingle{}}\NormalTok{]}\SpecialCharTok{.}\BuiltInTok{sum}\NormalTok{()}\SpecialCharTok{\}}\SpecialStringTok{"}\NormalTok{)}

\CommentTok{\# Count issues in invalid dataset}
\ControlFlowTok{if} \KeywordTok{not}\NormalTok{ invalid\_passed:}
\NormalTok{    issues }\OperatorTok{=}\NormalTok{ \{}
        \StringTok{\textquotesingle{}columns\_with\_nulls\textquotesingle{}}\NormalTok{: }\BuiltInTok{len}\NormalTok{(report\_invalid[report\_invalid[}\StringTok{\textquotesingle{}null\_count\textquotesingle{}}\NormalTok{] }\OperatorTok{\textgreater{}} \DecValTok{0}\NormalTok{]),}
        \StringTok{\textquotesingle{}total\_null\_values\textquotesingle{}}\NormalTok{: report\_invalid[}\StringTok{\textquotesingle{}null\_count\textquotesingle{}}\NormalTok{].}\BuiltInTok{sum}\NormalTok{(),}
        \StringTok{\textquotesingle{}validation\_passed\textquotesingle{}}\NormalTok{: invalid\_passed}
\NormalTok{    \}}
    
    \BuiltInTok{print}\NormalTok{(}\StringTok{"}\CharTok{\textbackslash{}n}\StringTok{Validation Issues:"}\NormalTok{)}
    \ControlFlowTok{for}\NormalTok{ key, value }\KeywordTok{in}\NormalTok{ issues.items():}
        \BuiltInTok{print}\NormalTok{(}\SpecialStringTok{f"  }\SpecialCharTok{\{}\NormalTok{key}\SpecialCharTok{\}}\SpecialStringTok{: }\SpecialCharTok{\{}\NormalTok{value}\SpecialCharTok{\}}\SpecialStringTok{"}\NormalTok{)}
    
    \BuiltInTok{print}\NormalTok{(}\StringTok{"}\CharTok{\textbackslash{}n}\StringTok{Columns with Issues:"}\NormalTok{)}
    \BuiltInTok{print}\NormalTok{(report\_invalid[report\_invalid[}\StringTok{\textquotesingle{}null\_count\textquotesingle{}}\NormalTok{] }\OperatorTok{\textgreater{}} \DecValTok{0}\NormalTok{][[}\StringTok{\textquotesingle{}column\textquotesingle{}}\NormalTok{, }\StringTok{\textquotesingle{}null\_count\textquotesingle{}}\NormalTok{, }\StringTok{\textquotesingle{}null\_pct\textquotesingle{}}\NormalTok{]])}
\end{Highlighting}
\end{Shaded}

Output:

\begin{verbatim}
Valid Dataset:
  Validation Passed: True
  Null Values: 0

Invalid Dataset:
  Validation Passed: False
  Null Values: 5

Validation Issues:
  columns_with_nulls: 3
  total_null_values: 5
  validation_passed: False

Columns with Issues:
       column  null_count  null_pct
0  product_id           1      20.0
1       price           2      40.0
2       stock           1      20.0
\end{verbatim}

\textbf{Validation Rules:} - \textbf{No Nulls}: All columns must have
zero null values - \textbf{Data Types}: Verify expected data types match
actual types - \textbf{Value Ranges}: Check min/max values are within
expected bounds - \textbf{Cardinality}: Ensure unique counts match
expectations

\textbf{Automation Workflow:} 1. Run \texttt{@analyze} on incoming data
2. Apply validation rules to the analysis report 3. Generate pass/fail
status and detailed issue report 4. Block processing if validation fails
5. Log issues for data quality monitoring

\textbf{Integration Points:} - ETL pipelines: Validate before loading to
warehouse - API endpoints: Validate incoming data payloads - Batch
processing: Pre-flight checks before expensive operations - Data
monitoring: Track quality metrics over time

Note: This also works with polars DataFrames.

\begin{center}\rule{0.5\linewidth}{0.5pt}\end{center}

\section{Workflow Patterns}\label{workflow-patterns}

\subsection{Pattern 1: Quality Gating}\label{pattern-1-quality-gating}

\begin{Shaded}
\begin{Highlighting}[]
\CommentTok{\# Run analysis}
\NormalTok{report }\OperatorTok{=}\NormalTok{ add.scan(}\StringTok{\textquotesingle{}@analyze\textquotesingle{}}\NormalTok{, df, as\_type}\OperatorTok{=}\StringTok{\textquotesingle{}dataframe\textquotesingle{}}\NormalTok{)}

\CommentTok{\# Define quality gates}
\NormalTok{max\_null\_pct }\OperatorTok{=} \FloatTok{10.0}
\NormalTok{min\_unique\_ratio }\OperatorTok{=} \FloatTok{0.1}

\CommentTok{\# Check gates}
\NormalTok{quality\_passed }\OperatorTok{=}\NormalTok{ (}
\NormalTok{    (report[}\StringTok{\textquotesingle{}null\_pct\textquotesingle{}}\NormalTok{] }\OperatorTok{\textless{}=}\NormalTok{ max\_null\_pct).}\BuiltInTok{all}\NormalTok{() }\KeywordTok{and}
\NormalTok{    (report[}\StringTok{\textquotesingle{}unique\textquotesingle{}}\NormalTok{] }\OperatorTok{/}\NormalTok{ report[}\StringTok{\textquotesingle{}count\textquotesingle{}}\NormalTok{] }\OperatorTok{\textgreater{}=}\NormalTok{ min\_unique\_ratio).}\BuiltInTok{all}\NormalTok{()}
\NormalTok{)}

\ControlFlowTok{if}\NormalTok{ quality\_passed:}
    \CommentTok{\# Proceed with processing}
\NormalTok{    process\_data(df)}
\ControlFlowTok{else}\NormalTok{:}
    \CommentTok{\# Reject and log issues}
\NormalTok{    log\_quality\_issues(report)}
\end{Highlighting}
\end{Shaded}

\subsection{Pattern 2: Incremental
Profiling}\label{pattern-2-incremental-profiling}

\begin{Shaded}
\begin{Highlighting}[]
\CommentTok{\# Profile baseline data}
\NormalTok{baseline }\OperatorTok{=}\NormalTok{ add.scan(}\StringTok{\textquotesingle{}@analyze\textquotesingle{}}\NormalTok{, df\_baseline, as\_type}\OperatorTok{=}\StringTok{\textquotesingle{}dataframe\textquotesingle{}}\NormalTok{)}

\CommentTok{\# Profile new data}
\NormalTok{current }\OperatorTok{=}\NormalTok{ add.scan(}\StringTok{\textquotesingle{}@analyze\textquotesingle{}}\NormalTok{, df\_current, as\_type}\OperatorTok{=}\StringTok{\textquotesingle{}dataframe\textquotesingle{}}\NormalTok{)}

\CommentTok{\# Compare distributions}
\NormalTok{drift\_detected }\OperatorTok{=}\NormalTok{ (}
    \BuiltInTok{abs}\NormalTok{(baseline[}\StringTok{\textquotesingle{}mean\textquotesingle{}}\NormalTok{] }\OperatorTok{{-}}\NormalTok{ current[}\StringTok{\textquotesingle{}mean\textquotesingle{}}\NormalTok{]) }\OperatorTok{\textgreater{}}\NormalTok{ threshold}
\NormalTok{).}\BuiltInTok{any}\NormalTok{()}

\ControlFlowTok{if}\NormalTok{ drift\_detected:}
\NormalTok{    alert\_data\_drift()}
\end{Highlighting}
\end{Shaded}

\subsection{Pattern 3: Multi-Stage
Validation}\label{pattern-3-multi-stage-validation}

\begin{Shaded}
\begin{Highlighting}[]
\CommentTok{\# Stage 1: Schema validation}
\NormalTok{report }\OperatorTok{=}\NormalTok{ add.scan(}\StringTok{\textquotesingle{}@analyze\textquotesingle{}}\NormalTok{, df, as\_type}\OperatorTok{=}\StringTok{\textquotesingle{}dataframe\textquotesingle{}}\NormalTok{)}
\NormalTok{schema\_valid }\OperatorTok{=}\NormalTok{ validate\_schema(report)}

\CommentTok{\# Stage 2: Quality validation}
\ControlFlowTok{if}\NormalTok{ schema\_valid:}
\NormalTok{    quality\_valid }\OperatorTok{=}\NormalTok{ validate\_quality(report)}
    
    \CommentTok{\# Stage 3: Business rules}
    \ControlFlowTok{if}\NormalTok{ quality\_valid:}
\NormalTok{        business\_valid }\OperatorTok{=}\NormalTok{ validate\_business\_rules(df, report)}
\end{Highlighting}
\end{Shaded}

\begin{center}\rule{0.5\linewidth}{0.5pt}\end{center}

\section{Parameters}\label{parameters-6}

\subsection{Required Parameters}\label{required-parameters-5}

\begin{itemize}
\tightlist
\item
  \texttt{mode}: \texttt{\textquotesingle{}@analyze\textquotesingle{}}
  for statistical profiling
\item
  \texttt{df}: Input DataFrame (pandas or polars)
\end{itemize}

\subsection{Optional Parameters}\label{optional-parameters-5}

\begin{itemize}
\tightlist
\item
  \texttt{columns}: Column filter (not yet implemented)
\item
  \texttt{where}: SQL-like filter condition (not yet implemented)
\item
  \texttt{rows}: Row range specifications (not yet implemented)
\item
  \texttt{as\_type}: Output format
  (\texttt{\textquotesingle{}dataframe\textquotesingle{}},
  \texttt{\textquotesingle{}dict\textquotesingle{}},
  \texttt{\textquotesingle{}text\textquotesingle{}})
\end{itemize}

\subsection{Positional Parameters}\label{positional-parameters-6}

\begin{Shaded}
\begin{Highlighting}[]
\CommentTok{\# Also works without naming certain parameters:}
\NormalTok{report }\OperatorTok{=}\NormalTok{ add.scan(}\StringTok{\textquotesingle{}@analyze\textquotesingle{}}\NormalTok{, df, as\_type}\OperatorTok{=}\StringTok{\textquotesingle{}dataframe\textquotesingle{}}\NormalTok{)}
\end{Highlighting}
\end{Shaded}

\begin{center}\rule{0.5\linewidth}{0.5pt}\end{center}

\section{Best Practices}\label{best-practices-6}

\subsection{Performance}\label{performance-1}

\begin{itemize}
\tightlist
\item
  Run \texttt{@analyze} once and reuse results for multiple checks
\item
  Use \texttt{as\_type=\textquotesingle{}dataframe\textquotesingle{}}
  for programmatic analysis
\item
  Use \texttt{as\_type=\textquotesingle{}text\textquotesingle{}} for
  logging and display only
\end{itemize}

\subsection{Quality Thresholds}\label{quality-thresholds}

\begin{itemize}
\tightlist
\item
  \textbf{Critical}: null\_pct \textgreater{} 50\%, unique == 1
  (constant)
\item
  \textbf{Warning}: null\_pct \textgreater{} 10\%, unique == count (all
  unique)
\item
  \textbf{Info}: null\_pct \textgreater{} 0\%, unique \textless{} 10
  (low cardinality)
\end{itemize}

\subsection{Automation}\label{automation}

\begin{itemize}
\tightlist
\item
  Integrate into CI/CD pipelines for data validation
\item
  Set up alerts for quality threshold violations
\item
  Track quality metrics over time for trend analysis
\end{itemize}

\begin{center}\rule{0.5\linewidth}{0.5pt}\end{center}

\section{Next Steps}\label{next-steps-18}

\begin{itemize}
\tightlist
\item
  \textbf{Page 1}: Basic data scanning with \texttt{@analyze} mode
\item
  \textbf{Page 2}: Lineage tracking with \texttt{@lineage} mode
\item
  \textbf{Batch 5}: Error handling and edge cases (coming soon)
\end{itemize}

\part{Lineage Tracking}

\chapter{Lineage Tracking Basics}\label{lineage-tracking-basics}

Track transformation history and understand data provenance by enabling
lineage tracking in your operations.

\section{Example 1: Enable and View
Lineage}\label{example-1-enable-and-view-lineage}

Enable lineage tracking by adding \texttt{lineage=True} to any
operation, then use
\texttt{add.scan(\textquotesingle{}@lineage\textquotesingle{},\ df)} to
view the transformation history.

\begin{Shaded}
\begin{Highlighting}[]
\ImportTok{import}\NormalTok{ pandas }\ImportTok{as}\NormalTok{ pd}
\ImportTok{import}\NormalTok{ additory }\ImportTok{as}\NormalTok{ add}

\NormalTok{df }\OperatorTok{=}\NormalTok{ pd.DataFrame(\{}
    \StringTok{\textquotesingle{}product\textquotesingle{}}\NormalTok{: [}\StringTok{\textquotesingle{}Widget\textquotesingle{}}\NormalTok{, }\StringTok{\textquotesingle{}Gadget\textquotesingle{}}\NormalTok{, }\StringTok{\textquotesingle{}Doohickey\textquotesingle{}}\NormalTok{],}
    \StringTok{\textquotesingle{}price\textquotesingle{}}\NormalTok{: [}\DecValTok{100}\NormalTok{, }\DecValTok{200}\NormalTok{, }\DecValTok{150}\NormalTok{],}
    \StringTok{\textquotesingle{}quantity\textquotesingle{}}\NormalTok{: [}\DecValTok{2}\NormalTok{, }\DecValTok{1}\NormalTok{, }\DecValTok{3}\NormalTok{]}
\NormalTok{\})}

\CommentTok{\# Transform with lineage tracking enabled}
\NormalTok{result }\OperatorTok{=}\NormalTok{ add.transform(}
    \StringTok{\textquotesingle{}@calc\textquotesingle{}}\NormalTok{,}
\NormalTok{    df,}
\NormalTok{    columns}\OperatorTok{=}\NormalTok{[}\StringTok{\textquotesingle{}price\textquotesingle{}}\NormalTok{, }\StringTok{\textquotesingle{}quantity\textquotesingle{}}\NormalTok{],}
\NormalTok{    expression}\OperatorTok{=}\StringTok{\textquotesingle{}price * quantity\textquotesingle{}}\NormalTok{,}
\NormalTok{    as\_}\OperatorTok{=}\StringTok{\textquotesingle{}total\textquotesingle{}}\NormalTok{,}
\NormalTok{    lineage}\OperatorTok{=}\VariableTok{True}  \CommentTok{\# Enable lineage tracking}
\NormalTok{)}

\CommentTok{\# View lineage report}
\NormalTok{lineage\_report }\OperatorTok{=}\NormalTok{ add.scan(}\StringTok{\textquotesingle{}@lineage\textquotesingle{}}\NormalTok{, result)}
\BuiltInTok{print}\NormalTok{(lineage\_report)}
\end{Highlighting}
\end{Shaded}

Output:

\begin{verbatim}
═══════════════════════════════════════════════════════════════
LINEAGE REPORT
═══════════════════════════════════════════════════════════════

DataFrame: 3 rows × 4 columns
Operations: 1 transformations applied

───────────────────────────────────────────────────────────────
Step 1: add.transform - 2026-03-11T14:23:45
───────────────────────────────────────────────────────────────
  Rows: 3 → 3 (no change)
  Columns Added: total
  
  Parameters:
    columns: ["price","quantity"]
    expression: ["price * quantity"]
    mode: "@calc"
\end{verbatim}

\textbf{Key Points:} - Add \texttt{lineage=True} to any operation
(transform, to, synthetic) - Lineage metadata is stored with the
DataFrame - Use
\texttt{add.scan(\textquotesingle{}@lineage\textquotesingle{},\ df)} to
generate a human-readable report - Report shows operation history,
parameters, and column changes

Note: This also works with polars DataFrames.

\begin{center}\rule{0.5\linewidth}{0.5pt}\end{center}

\section{Example 2: Lineage with
Filtering}\label{example-2-lineage-with-filtering}

Lineage tracking works with all transform modes, including filtering
operations.

\begin{Shaded}
\begin{Highlighting}[]
\ImportTok{import}\NormalTok{ pandas }\ImportTok{as}\NormalTok{ pd}
\ImportTok{import}\NormalTok{ additory }\ImportTok{as}\NormalTok{ add}

\NormalTok{df }\OperatorTok{=}\NormalTok{ pd.DataFrame(\{}
    \StringTok{\textquotesingle{}product\textquotesingle{}}\NormalTok{: [}\StringTok{\textquotesingle{}Widget\textquotesingle{}}\NormalTok{, }\StringTok{\textquotesingle{}Gadget\textquotesingle{}}\NormalTok{, }\StringTok{\textquotesingle{}Doohickey\textquotesingle{}}\NormalTok{, }\StringTok{\textquotesingle{}Thingamajig\textquotesingle{}}\NormalTok{],}
    \StringTok{\textquotesingle{}price\textquotesingle{}}\NormalTok{: [}\DecValTok{100}\NormalTok{, }\DecValTok{200}\NormalTok{, }\DecValTok{150}\NormalTok{, }\DecValTok{50}\NormalTok{],}
    \StringTok{\textquotesingle{}stock\textquotesingle{}}\NormalTok{: [}\DecValTok{10}\NormalTok{, }\DecValTok{5}\NormalTok{, }\DecValTok{0}\NormalTok{, }\DecValTok{15}\NormalTok{]}
\NormalTok{\})}

\CommentTok{\# Filter with lineage tracking}
\NormalTok{result }\OperatorTok{=}\NormalTok{ add.transform(}
    \StringTok{\textquotesingle{}@filter\textquotesingle{}}\NormalTok{,}
\NormalTok{    df,}
\NormalTok{    columns}\OperatorTok{=}\NormalTok{[}\StringTok{\textquotesingle{}product\textquotesingle{}}\NormalTok{, }\StringTok{\textquotesingle{}price\textquotesingle{}}\NormalTok{, }\StringTok{\textquotesingle{}stock\textquotesingle{}}\NormalTok{],}
\NormalTok{    where}\OperatorTok{=}\StringTok{\textquotesingle{}stock \textgreater{} 0\textquotesingle{}}\NormalTok{,}
\NormalTok{    lineage}\OperatorTok{=}\VariableTok{True}
\NormalTok{)}

\CommentTok{\# View lineage report}
\NormalTok{lineage\_report }\OperatorTok{=}\NormalTok{ add.scan(}\StringTok{\textquotesingle{}@lineage\textquotesingle{}}\NormalTok{, result)}
\BuiltInTok{print}\NormalTok{(lineage\_report)}
\end{Highlighting}
\end{Shaded}

Output:

\begin{verbatim}
═══════════════════════════════════════════════════════════════
LINEAGE REPORT
═══════════════════════════════════════════════════════════════

DataFrame: 3 rows × 3 columns
Operations: 1 transformations applied

───────────────────────────────────────────────────────────────
Step 1: add.transform - 2026-03-11T14:25:12
───────────────────────────────────────────────────────────────
  Rows: 4 → 3 (1 row filtered out)
  Columns Added: none
  
  Parameters:
    columns: ["product","price","stock"]
    where: "stock > 0"
    mode: "@filter"
\end{verbatim}

\textbf{Key Insights:} - Lineage tracks row changes (4 → 3 rows) -
Filter conditions are recorded in parameters - You can see exactly what
was filtered and why

Note: This also works with polars DataFrames.

\begin{center}\rule{0.5\linewidth}{0.5pt}\end{center}

\section{Example 3: Multi-Operation Lineage
Chain}\label{example-3-multi-operation-lineage-chain}

Lineage accumulates across multiple operations, giving you a complete
transformation history.

\begin{Shaded}
\begin{Highlighting}[]
\ImportTok{import}\NormalTok{ pandas }\ImportTok{as}\NormalTok{ pd}
\ImportTok{import}\NormalTok{ additory }\ImportTok{as}\NormalTok{ add}

\NormalTok{df }\OperatorTok{=}\NormalTok{ pd.DataFrame(\{}
    \StringTok{\textquotesingle{}product\textquotesingle{}}\NormalTok{: [}\StringTok{\textquotesingle{}Widget\textquotesingle{}}\NormalTok{, }\StringTok{\textquotesingle{}Gadget\textquotesingle{}}\NormalTok{, }\StringTok{\textquotesingle{}Doohickey\textquotesingle{}}\NormalTok{],}
    \StringTok{\textquotesingle{}price\textquotesingle{}}\NormalTok{: [}\DecValTok{100}\NormalTok{, }\DecValTok{200}\NormalTok{, }\DecValTok{150}\NormalTok{],}
    \StringTok{\textquotesingle{}quantity\textquotesingle{}}\NormalTok{: [}\DecValTok{2}\NormalTok{, }\DecValTok{1}\NormalTok{, }\DecValTok{3}\NormalTok{]}
\NormalTok{\})}

\CommentTok{\# Step 1: Calculate total}
\NormalTok{result }\OperatorTok{=}\NormalTok{ add.transform(}
    \StringTok{\textquotesingle{}@calc\textquotesingle{}}\NormalTok{,}
\NormalTok{    df,}
\NormalTok{    columns}\OperatorTok{=}\NormalTok{[}\StringTok{\textquotesingle{}price\textquotesingle{}}\NormalTok{, }\StringTok{\textquotesingle{}quantity\textquotesingle{}}\NormalTok{],}
\NormalTok{    expression}\OperatorTok{=}\StringTok{\textquotesingle{}price * quantity\textquotesingle{}}\NormalTok{,}
\NormalTok{    as\_}\OperatorTok{=}\StringTok{\textquotesingle{}total\textquotesingle{}}\NormalTok{,}
\NormalTok{    lineage}\OperatorTok{=}\VariableTok{True}
\NormalTok{)}

\CommentTok{\# Step 2: Filter high{-}value items (lineage continues)}
\NormalTok{result }\OperatorTok{=}\NormalTok{ add.transform(}
    \StringTok{\textquotesingle{}@filter\textquotesingle{}}\NormalTok{,}
\NormalTok{    result,}
\NormalTok{    columns}\OperatorTok{=}\NormalTok{[}\StringTok{\textquotesingle{}product\textquotesingle{}}\NormalTok{, }\StringTok{\textquotesingle{}total\textquotesingle{}}\NormalTok{],}
\NormalTok{    where}\OperatorTok{=}\StringTok{\textquotesingle{}total \textgreater{} 200\textquotesingle{}}\NormalTok{,}
\NormalTok{    lineage}\OperatorTok{=}\VariableTok{True}
\NormalTok{)}

\CommentTok{\# View complete lineage report}
\NormalTok{lineage\_report }\OperatorTok{=}\NormalTok{ add.scan(}\StringTok{\textquotesingle{}@lineage\textquotesingle{}}\NormalTok{, result)}
\BuiltInTok{print}\NormalTok{(lineage\_report)}
\end{Highlighting}
\end{Shaded}

Output:

\begin{verbatim}
═══════════════════════════════════════════════════════════════
LINEAGE REPORT
═══════════════════════════════════════════════════════════════

DataFrame: 2 rows × 2 columns
Operations: 2 transformations applied

───────────────────────────────────────────────────────────────
Step 1: add.transform - 2026-03-11T14:27:33
───────────────────────────────────────────────────────────────
  Rows: 3 → 3 (no change)
  Columns Added: total
  
  Parameters:
    columns: ["price","quantity"]
    expression: ["price * quantity"]
    mode: "@calc"

───────────────────────────────────────────────────────────────
Step 2: add.transform - 2026-03-11T14:27:33
───────────────────────────────────────────────────────────────
  Rows: 3 → 2 (1 row filtered out)
  Columns Added: none
  
  Parameters:
    columns: ["product","total"]
    where: "total > 200"
    mode: "@filter"
\end{verbatim}

\textbf{Key Features:} - Lineage persists across multiple operations -
Each step is numbered and timestamped - You can see the complete
transformation pipeline - Useful for debugging complex data workflows

Note: This also works with polars DataFrames.

\begin{center}\rule{0.5\linewidth}{0.5pt}\end{center}

\section{Lineage Tracking Features}\label{lineage-tracking-features-1}

\subsection{What Gets Tracked}\label{what-gets-tracked-1}

\begin{itemize}
\tightlist
\item
  \textbf{Operation Type}: Which function was called (transform, to,
  synthetic)
\item
  \textbf{Timestamp}: When the operation occurred
\item
  \textbf{Parameters}: All operation parameters
\item
  \textbf{Row Changes}: Rows before and after
\item
  \textbf{Column Changes}: Columns added or modified
\item
  \textbf{Mode}: Specific transformation mode used
\end{itemize}

\subsection{When to Use Lineage}\label{when-to-use-lineage-1}

\begin{itemize}
\tightlist
\item
  \textbf{Debugging}: Understand how data was transformed
\item
  \textbf{Auditing}: Track data provenance for compliance
\item
  \textbf{Documentation}: Auto-generate transformation documentation
\item
  \textbf{Collaboration}: Share transformation history with team
\item
  \textbf{Quality Assurance}: Verify transformation steps
\end{itemize}

\subsection{Performance Impact}\label{performance-impact-1}

\begin{itemize}
\tightlist
\item
  Minimal overhead (\textless100ms per operation)
\item
  Metadata stored in memory only
\item
  No impact on computation speed
\item
  Lineage data is lightweight (JSON format)
\end{itemize}

\begin{center}\rule{0.5\linewidth}{0.5pt}\end{center}

\section{Enabling Lineage Tracking}\label{enabling-lineage-tracking-1}

\subsection{In add.transform()}\label{in-add.transform-1}

\begin{Shaded}
\begin{Highlighting}[]
\NormalTok{result }\OperatorTok{=}\NormalTok{ add.transform(}
    \StringTok{\textquotesingle{}@calc\textquotesingle{}}\NormalTok{,}
\NormalTok{    df,}
\NormalTok{    columns}\OperatorTok{=}\NormalTok{[}\StringTok{\textquotesingle{}a\textquotesingle{}}\NormalTok{, }\StringTok{\textquotesingle{}b\textquotesingle{}}\NormalTok{],}
\NormalTok{    expression}\OperatorTok{=}\StringTok{\textquotesingle{}a + b\textquotesingle{}}\NormalTok{,}
\NormalTok{    as\_}\OperatorTok{=}\StringTok{\textquotesingle{}c\textquotesingle{}}\NormalTok{,}
\NormalTok{    lineage}\OperatorTok{=}\VariableTok{True}  \CommentTok{\# Enable lineage}
\NormalTok{)}
\end{Highlighting}
\end{Shaded}

\subsection{In add.to()}\label{in-add.to-1}

\begin{Shaded}
\begin{Highlighting}[]
\NormalTok{result }\OperatorTok{=}\NormalTok{ add.to(}
\NormalTok{    orders,}
\NormalTok{    bring\_from}\OperatorTok{=}\NormalTok{customers,}
\NormalTok{    bring}\OperatorTok{=}\StringTok{\textquotesingle{}name\textquotesingle{}}\NormalTok{,}
\NormalTok{    against}\OperatorTok{=}\StringTok{\textquotesingle{}customer\_id\textquotesingle{}}\NormalTok{,}
\NormalTok{    lineage}\OperatorTok{=}\VariableTok{True}  \CommentTok{\# Enable lineage}
\NormalTok{)}
\end{Highlighting}
\end{Shaded}

\subsection{In add.synthetic()}\label{in-add.synthetic-1}

\begin{Shaded}
\begin{Highlighting}[]
\NormalTok{result }\OperatorTok{=}\NormalTok{ add.synthetic(}
    \StringTok{\textquotesingle{}@new\textquotesingle{}}\NormalTok{,}
\NormalTok{    n}\OperatorTok{=}\DecValTok{100}\NormalTok{,}
\NormalTok{    strategy}\OperatorTok{=}\NormalTok{\{}\StringTok{\textquotesingle{}a\textquotesingle{}}\NormalTok{: \{}\StringTok{\textquotesingle{}distribution\textquotesingle{}}\NormalTok{: }\StringTok{\textquotesingle{}normal\textquotesingle{}}\NormalTok{, }\StringTok{\textquotesingle{}mean\textquotesingle{}}\NormalTok{: }\DecValTok{0}\NormalTok{, }\StringTok{\textquotesingle{}std\textquotesingle{}}\NormalTok{: }\DecValTok{1}\NormalTok{\}\},}
\NormalTok{    lineage}\OperatorTok{=}\VariableTok{True}  \CommentTok{\# Enable lineage}
\NormalTok{)}
\end{Highlighting}
\end{Shaded}

\begin{center}\rule{0.5\linewidth}{0.5pt}\end{center}

\section{Parameters}\label{parameters-7}

\subsection{Required Parameters}\label{required-parameters-6}

\begin{itemize}
\tightlist
\item
  \texttt{mode}: \texttt{\textquotesingle{}@lineage\textquotesingle{}}
  for lineage tracking
\item
  \texttt{df}: Input DataFrame with lineage metadata
\end{itemize}

\subsection{Optional Parameters}\label{optional-parameters-6}

\begin{itemize}
\tightlist
\item
  \texttt{columns}: Filter report to specific columns (not yet
  implemented)
\item
  \texttt{trace}: Trace specific cell transformations (not yet
  implemented)
\item
  \texttt{as\_type}: Output format (currently only `text' supported)
\end{itemize}

\subsection{Positional Parameters}\label{positional-parameters-7}

\begin{Shaded}
\begin{Highlighting}[]
\CommentTok{\# Also works without naming certain parameters:}
\NormalTok{lineage\_report }\OperatorTok{=}\NormalTok{ add.scan(}\StringTok{\textquotesingle{}@lineage\textquotesingle{}}\NormalTok{, df)}
\end{Highlighting}
\end{Shaded}

\begin{center}\rule{0.5\linewidth}{0.5pt}\end{center}

\section{Error Handling}\label{error-handling}

If you forget to enable lineage tracking,
\texttt{add.scan(\textquotesingle{}@lineage\textquotesingle{},\ df)}
provides a helpful error message:

\begin{Shaded}
\begin{Highlighting}[]
\CommentTok{\# Transform WITHOUT lineage tracking}
\NormalTok{result }\OperatorTok{=}\NormalTok{ add.transform(}\StringTok{\textquotesingle{}@calc\textquotesingle{}}\NormalTok{, df, expression}\OperatorTok{=}\StringTok{\textquotesingle{}a + b\textquotesingle{}}\NormalTok{, as\_}\OperatorTok{=}\StringTok{\textquotesingle{}c\textquotesingle{}}\NormalTok{)}

\CommentTok{\# Try to get lineage report}
\ControlFlowTok{try}\NormalTok{:}
\NormalTok{    lineage\_report }\OperatorTok{=}\NormalTok{ add.scan(}\StringTok{\textquotesingle{}@lineage\textquotesingle{}}\NormalTok{, result)}
\ControlFlowTok{except} \PreprocessorTok{ValueError} \ImportTok{as}\NormalTok{ e:}
    \BuiltInTok{print}\NormalTok{(e)}
\end{Highlighting}
\end{Shaded}

Output:

\begin{verbatim}
No lineage metadata found. Lineage tracking must be enabled by adding
lineage=True to add.to(), add.transform(), or add.synthetic() calls.

Example:
  df = add.transform('@calc', df, strategy={'total': 'price * qty'}, lineage=True)
  result = add.scan('@lineage', df)
\end{verbatim}

\textbf{Solution}: Add \texttt{lineage=True} to your operations to
enable tracking.

\begin{center}\rule{0.5\linewidth}{0.5pt}\end{center}

\section{Limitations (v0.1.3)}\label{limitations-v0.1.3}

\subsection{Session-Only Lineage}\label{session-only-lineage-1}

\begin{itemize}
\tightlist
\item
  Lineage metadata is stored in memory only
\item
  Not persisted when saving DataFrames to disk
\item
  Lost when Python session ends
\end{itemize}

\subsection{Type Conversion
Restriction}\label{type-conversion-restriction}

\begin{itemize}
\tightlist
\item
  Cannot use \texttt{lineage=True} with \texttt{as\_type} parameter
\item
  Lineage metadata would be lost during type conversion
\item
  Use lineage OR type conversion, not both
\end{itemize}

\subsection{Future Enhancements
(v0.2.0)}\label{future-enhancements-v0.2.0-1}

\begin{itemize}
\tightlist
\item
  Persistent lineage with \texttt{add.save()} and \texttt{add.load()}
\item
  Cell-level tracing with \texttt{trace} parameter
\item
  Column-specific lineage with \texttt{columns} parameter
\item
  Lineage visualization and export
\item
  Lineage comparison and diff
\end{itemize}

\begin{center}\rule{0.5\linewidth}{0.5pt}\end{center}

\section{Best Practices}\label{best-practices-7}

\begin{enumerate}
\def\labelenumi{\arabic{enumi}.}
\tightlist
\item
  \textbf{Enable lineage early}: Add \texttt{lineage=True} from the
  start of your pipeline
\item
  \textbf{Check lineage regularly}: Use
  \texttt{add.scan(\textquotesingle{}@lineage\textquotesingle{},\ df)}
  to verify transformations
\item
  \textbf{Document with lineage}: Include lineage reports in
  documentation
\item
  \textbf{Debug with lineage}: When data looks wrong, check the lineage
  first
\item
  \textbf{Share lineage}: Include lineage reports when asking for help
\end{enumerate}

\begin{center}\rule{0.5\linewidth}{0.5pt}\end{center}

\section{Next Steps}\label{next-steps-19}

\begin{itemize}
\tightlist
\item
  \textbf{Page 2}: Advanced lineage workflows and real-world use cases
\item
  \textbf{Troubleshooting Guide}: Common errors and solutions
\end{itemize}

\chapter{Advanced Lineage Workflows}\label{advanced-lineage-workflows}

Learn advanced lineage tracking patterns for complex data pipelines,
debugging, and quality assurance.

\section{Example 1: Multi-Step Transformation
Lineage}\label{example-1-multi-step-transformation-lineage}

Track complex transformation pipelines with multiple steps to understand
the complete data journey.

\begin{Shaded}
\begin{Highlighting}[]
\ImportTok{import}\NormalTok{ pandas }\ImportTok{as}\NormalTok{ pd}
\ImportTok{import}\NormalTok{ additory }\ImportTok{as}\NormalTok{ add}

\NormalTok{df }\OperatorTok{=}\NormalTok{ pd.DataFrame(\{}
    \StringTok{\textquotesingle{}product\textquotesingle{}}\NormalTok{: [}\StringTok{\textquotesingle{}Widget\textquotesingle{}}\NormalTok{, }\StringTok{\textquotesingle{}Gadget\textquotesingle{}}\NormalTok{, }\StringTok{\textquotesingle{}Doohickey\textquotesingle{}}\NormalTok{],}
    \StringTok{\textquotesingle{}price\textquotesingle{}}\NormalTok{: [}\DecValTok{100}\NormalTok{, }\DecValTok{200}\NormalTok{, }\DecValTok{150}\NormalTok{],}
    \StringTok{\textquotesingle{}cost\textquotesingle{}}\NormalTok{: [}\DecValTok{60}\NormalTok{, }\DecValTok{120}\NormalTok{, }\DecValTok{90}\NormalTok{],}
    \StringTok{\textquotesingle{}quantity\textquotesingle{}}\NormalTok{: [}\DecValTok{2}\NormalTok{, }\DecValTok{1}\NormalTok{, }\DecValTok{3}\NormalTok{]}
\NormalTok{\})}

\CommentTok{\# Step 1: Calculate profit}
\NormalTok{result }\OperatorTok{=}\NormalTok{ add.transform(}
    \StringTok{\textquotesingle{}@calc\textquotesingle{}}\NormalTok{,}
\NormalTok{    df,}
\NormalTok{    columns}\OperatorTok{=}\NormalTok{[}\StringTok{\textquotesingle{}price\textquotesingle{}}\NormalTok{, }\StringTok{\textquotesingle{}cost\textquotesingle{}}\NormalTok{],}
\NormalTok{    expression}\OperatorTok{=}\StringTok{\textquotesingle{}price {-} cost\textquotesingle{}}\NormalTok{,}
\NormalTok{    as\_}\OperatorTok{=}\StringTok{\textquotesingle{}profit\textquotesingle{}}\NormalTok{,}
\NormalTok{    lineage}\OperatorTok{=}\VariableTok{True}
\NormalTok{)}

\CommentTok{\# Step 2: Calculate revenue}
\NormalTok{result }\OperatorTok{=}\NormalTok{ add.transform(}
    \StringTok{\textquotesingle{}@calc\textquotesingle{}}\NormalTok{,}
\NormalTok{    result,}
\NormalTok{    columns}\OperatorTok{=}\NormalTok{[}\StringTok{\textquotesingle{}price\textquotesingle{}}\NormalTok{, }\StringTok{\textquotesingle{}quantity\textquotesingle{}}\NormalTok{],}
\NormalTok{    expression}\OperatorTok{=}\StringTok{\textquotesingle{}price * quantity\textquotesingle{}}\NormalTok{,}
\NormalTok{    as\_}\OperatorTok{=}\StringTok{\textquotesingle{}revenue\textquotesingle{}}\NormalTok{,}
\NormalTok{    lineage}\OperatorTok{=}\VariableTok{True}
\NormalTok{)}

\CommentTok{\# Step 3: Filter high{-}value items}
\NormalTok{result }\OperatorTok{=}\NormalTok{ add.transform(}
    \StringTok{\textquotesingle{}@filter\textquotesingle{}}\NormalTok{,}
\NormalTok{    result,}
\NormalTok{    columns}\OperatorTok{=}\NormalTok{[}\StringTok{\textquotesingle{}product\textquotesingle{}}\NormalTok{, }\StringTok{\textquotesingle{}profit\textquotesingle{}}\NormalTok{, }\StringTok{\textquotesingle{}revenue\textquotesingle{}}\NormalTok{],}
\NormalTok{    where}\OperatorTok{=}\StringTok{\textquotesingle{}revenue \textgreater{} 200\textquotesingle{}}\NormalTok{,}
\NormalTok{    lineage}\OperatorTok{=}\VariableTok{True}
\NormalTok{)}

\CommentTok{\# View complete lineage report}
\NormalTok{lineage\_report }\OperatorTok{=}\NormalTok{ add.scan(}\StringTok{\textquotesingle{}@lineage\textquotesingle{}}\NormalTok{, result)}
\BuiltInTok{print}\NormalTok{(lineage\_report)}
\end{Highlighting}
\end{Shaded}

Output:

\begin{verbatim}
═══════════════════════════════════════════════════════════════
LINEAGE REPORT
═══════════════════════════════════════════════════════════════

DataFrame: 1 rows × 3 columns
Operations: 3 transformations applied

───────────────────────────────────────────────────────────────
Step 1: add.transform - 2026-03-11T15:10:22
───────────────────────────────────────────────────────────────
  Rows: 3 → 3 (no change)
  Columns Added: profit
  
  Parameters:
    columns: ["price","cost"]
    expression: ["price - cost"]
    mode: "@calc"

───────────────────────────────────────────────────────────────
Step 2: add.transform - 2026-03-11T15:10:22
───────────────────────────────────────────────────────────────
  Rows: 3 → 3 (no change)
  Columns Added: revenue
  
  Parameters:
    columns: ["price","quantity"]
    expression: ["price * quantity"]
    mode: "@calc"

───────────────────────────────────────────────────────────────
Step 3: add.transform - 2026-03-11T15:10:22
───────────────────────────────────────────────────────────────
  Rows: 3 → 1 (2 rows filtered out)
  Columns Added: none
  
  Parameters:
    columns: ["product","profit","revenue"]
    where: "revenue > 200"
    mode: "@filter"
\end{verbatim}

\textbf{Use Cases:} - \textbf{Pipeline Debugging}: Identify which step
introduced an issue - \textbf{Performance Analysis}: See which steps
changed row counts - \textbf{Documentation}: Auto-generate pipeline
documentation - \textbf{Validation}: Verify each transformation step

Note: This also works with polars DataFrames.

\begin{center}\rule{0.5\linewidth}{0.5pt}\end{center}

\section{Example 2: Sorting and Calculation
Lineage}\label{example-2-sorting-and-calculation-lineage}

Track sorting operations combined with calculations to understand data
ordering and transformations.

\begin{Shaded}
\begin{Highlighting}[]
\ImportTok{import}\NormalTok{ pandas }\ImportTok{as}\NormalTok{ pd}
\ImportTok{import}\NormalTok{ additory }\ImportTok{as}\NormalTok{ add}

\NormalTok{df }\OperatorTok{=}\NormalTok{ pd.DataFrame(\{}
    \StringTok{\textquotesingle{}product\textquotesingle{}}\NormalTok{: [}\StringTok{\textquotesingle{}Widget\textquotesingle{}}\NormalTok{, }\StringTok{\textquotesingle{}Gadget\textquotesingle{}}\NormalTok{, }\StringTok{\textquotesingle{}Doohickey\textquotesingle{}}\NormalTok{],}
    \StringTok{\textquotesingle{}price\textquotesingle{}}\NormalTok{: [}\DecValTok{100}\NormalTok{, }\DecValTok{200}\NormalTok{, }\DecValTok{150}\NormalTok{],}
    \StringTok{\textquotesingle{}sales\textquotesingle{}}\NormalTok{: [}\DecValTok{50}\NormalTok{, }\DecValTok{150}\NormalTok{, }\DecValTok{75}\NormalTok{]}
\NormalTok{\})}

\CommentTok{\# Step 1: Calculate revenue}
\NormalTok{result }\OperatorTok{=}\NormalTok{ add.transform(}
    \StringTok{\textquotesingle{}@calc\textquotesingle{}}\NormalTok{,}
\NormalTok{    df,}
\NormalTok{    columns}\OperatorTok{=}\NormalTok{[}\StringTok{\textquotesingle{}price\textquotesingle{}}\NormalTok{, }\StringTok{\textquotesingle{}sales\textquotesingle{}}\NormalTok{],}
\NormalTok{    expression}\OperatorTok{=}\StringTok{\textquotesingle{}price * sales\textquotesingle{}}\NormalTok{,}
\NormalTok{    as\_}\OperatorTok{=}\StringTok{\textquotesingle{}revenue\textquotesingle{}}\NormalTok{,}
\NormalTok{    lineage}\OperatorTok{=}\VariableTok{True}
\NormalTok{)}

\CommentTok{\# Step 2: Sort by revenue (descending)}
\NormalTok{result }\OperatorTok{=}\NormalTok{ add.transform(}
    \StringTok{\textquotesingle{}@sort\textquotesingle{}}\NormalTok{,}
\NormalTok{    result,}
\NormalTok{    columns}\OperatorTok{=}\NormalTok{[}\StringTok{\textquotesingle{}product\textquotesingle{}}\NormalTok{, }\StringTok{\textquotesingle{}revenue\textquotesingle{}}\NormalTok{],}
\NormalTok{    by}\OperatorTok{=}\StringTok{\textquotesingle{}revenue\textquotesingle{}}\NormalTok{,}
\NormalTok{    order}\OperatorTok{=}\StringTok{\textquotesingle{}desc\textquotesingle{}}\NormalTok{,}
\NormalTok{    lineage}\OperatorTok{=}\VariableTok{True}
\NormalTok{)}

\CommentTok{\# View lineage report}
\NormalTok{lineage\_report }\OperatorTok{=}\NormalTok{ add.scan(}\StringTok{\textquotesingle{}@lineage\textquotesingle{}}\NormalTok{, result)}
\BuiltInTok{print}\NormalTok{(lineage\_report)}
\end{Highlighting}
\end{Shaded}

Output:

\begin{verbatim}
═══════════════════════════════════════════════════════════════
LINEAGE REPORT
═══════════════════════════════════════════════════════════════

DataFrame: 3 rows × 2 columns
Operations: 2 transformations applied

───────────────────────────────────────────────────────────────
Step 1: add.transform - 2026-03-11T15:15:45
───────────────────────────────────────────────────────────────
  Rows: 3 → 3 (no change)
  Columns Added: revenue
  
  Parameters:
    columns: ["price","sales"]
    expression: ["price * sales"]
    mode: "@calc"

───────────────────────────────────────────────────────────────
Step 2: add.transform - 2026-03-11T15:15:45
───────────────────────────────────────────────────────────────
  Rows: 3 → 3 (no change)
  Columns Added: none
  
  Parameters:
    columns: ["product","revenue"]
    by: "revenue"
    order: "desc"
    mode: "@sort"
\end{verbatim}

\textbf{Key Insights:} - Sorting operations are tracked with their
parameters - You can see the sort column and order - Row count remains
the same (sorting doesn't filter) - Useful for understanding data
ordering logic

Note: This also works with polars DataFrames.

\begin{center}\rule{0.5\linewidth}{0.5pt}\end{center}

\section{Example 3: Lineage for Debugging Data
Quality}\label{example-3-lineage-for-debugging-data-quality}

Use lineage tracking to debug data quality issues and understand data
cleaning steps.

\begin{Shaded}
\begin{Highlighting}[]
\ImportTok{import}\NormalTok{ pandas }\ImportTok{as}\NormalTok{ pd}
\ImportTok{import}\NormalTok{ additory }\ImportTok{as}\NormalTok{ add}

\NormalTok{df }\OperatorTok{=}\NormalTok{ pd.DataFrame(\{}
    \StringTok{\textquotesingle{}product\textquotesingle{}}\NormalTok{: [}\StringTok{\textquotesingle{}Widget\textquotesingle{}}\NormalTok{, }\StringTok{\textquotesingle{}Gadget\textquotesingle{}}\NormalTok{, }\StringTok{\textquotesingle{}Doohickey\textquotesingle{}}\NormalTok{, }\StringTok{\textquotesingle{}Thingamajig\textquotesingle{}}\NormalTok{],}
    \StringTok{\textquotesingle{}price\textquotesingle{}}\NormalTok{: [}\DecValTok{100}\NormalTok{, }\DecValTok{200}\NormalTok{, }\VariableTok{None}\NormalTok{, }\DecValTok{150}\NormalTok{],}
    \StringTok{\textquotesingle{}quantity\textquotesingle{}}\NormalTok{: [}\DecValTok{2}\NormalTok{, }\DecValTok{1}\NormalTok{, }\DecValTok{3}\NormalTok{, }\DecValTok{0}\NormalTok{]}
\NormalTok{\})}

\CommentTok{\# Step 1: Impute missing values}
\NormalTok{result }\OperatorTok{=}\NormalTok{ add.transform(}
    \StringTok{\textquotesingle{}@deduce\textquotesingle{}}\NormalTok{,}
\NormalTok{    df,}
\NormalTok{    columns}\OperatorTok{=}\NormalTok{[}\StringTok{\textquotesingle{}price\textquotesingle{}}\NormalTok{],}
\NormalTok{    strategy}\OperatorTok{=}\NormalTok{\{}\StringTok{\textquotesingle{}method\textquotesingle{}}\NormalTok{: }\StringTok{\textquotesingle{}mean\textquotesingle{}}\NormalTok{\},}
\NormalTok{    lineage}\OperatorTok{=}\VariableTok{True}
\NormalTok{)}

\CommentTok{\# Step 2: Filter out zero quantity}
\NormalTok{result }\OperatorTok{=}\NormalTok{ add.transform(}
    \StringTok{\textquotesingle{}@filter\textquotesingle{}}\NormalTok{,}
\NormalTok{    result,}
\NormalTok{    columns}\OperatorTok{=}\NormalTok{[}\StringTok{\textquotesingle{}product\textquotesingle{}}\NormalTok{, }\StringTok{\textquotesingle{}price\textquotesingle{}}\NormalTok{, }\StringTok{\textquotesingle{}quantity\textquotesingle{}}\NormalTok{],}
\NormalTok{    where}\OperatorTok{=}\StringTok{\textquotesingle{}quantity \textgreater{} 0\textquotesingle{}}\NormalTok{,}
\NormalTok{    lineage}\OperatorTok{=}\VariableTok{True}
\NormalTok{)}

\CommentTok{\# Step 3: Calculate total}
\NormalTok{result }\OperatorTok{=}\NormalTok{ add.transform(}
    \StringTok{\textquotesingle{}@calc\textquotesingle{}}\NormalTok{,}
\NormalTok{    result,}
\NormalTok{    columns}\OperatorTok{=}\NormalTok{[}\StringTok{\textquotesingle{}price\textquotesingle{}}\NormalTok{, }\StringTok{\textquotesingle{}quantity\textquotesingle{}}\NormalTok{],}
\NormalTok{    expression}\OperatorTok{=}\StringTok{\textquotesingle{}price * quantity\textquotesingle{}}\NormalTok{,}
\NormalTok{    as\_}\OperatorTok{=}\StringTok{\textquotesingle{}total\textquotesingle{}}\NormalTok{,}
\NormalTok{    lineage}\OperatorTok{=}\VariableTok{True}
\NormalTok{)}

\CommentTok{\# View lineage report to understand data cleaning steps}
\NormalTok{lineage\_report }\OperatorTok{=}\NormalTok{ add.scan(}\StringTok{\textquotesingle{}@lineage\textquotesingle{}}\NormalTok{, result)}
\BuiltInTok{print}\NormalTok{(lineage\_report)}
\end{Highlighting}
\end{Shaded}

Output:

\begin{verbatim}
═══════════════════════════════════════════════════════════════
LINEAGE REPORT
═══════════════════════════════════════════════════════════════

DataFrame: 3 rows × 4 columns
Operations: 3 transformations applied

───────────────────────────────────────────────────────────────
Step 1: add.transform - 2026-03-11T15:20:18
───────────────────────────────────────────────────────────────
  Rows: 4 → 4 (no change)
  Columns Added: none
  
  Parameters:
    columns: ["price"]
    strategy: {"method":"mean"}
    mode: "@deduce"

───────────────────────────────────────────────────────────────
Step 2: add.transform - 2026-03-11T15:20:18
───────────────────────────────────────────────────────────────
  Rows: 4 → 3 (1 row filtered out)
  Columns Added: none
  
  Parameters:
    columns: ["product","price","quantity"]
    where: "quantity > 0"
    mode: "@filter"

───────────────────────────────────────────────────────────────
Step 3: add.transform - 2026-03-11T15:20:18
───────────────────────────────────────────────────────────────
  Rows: 3 → 3 (no change)
  Columns Added: total
  
  Parameters:
    columns: ["price","quantity"]
    expression: ["price * quantity"]
    mode: "@calc"
\end{verbatim}

\textbf{Debugging Workflow:} 1. \textbf{Identify the issue}: Final data
has 3 rows instead of 4 2. \textbf{Check lineage}: Step 2 filtered out 1
row 3. \textbf{Understand why}: Filter condition was
\texttt{quantity\ \textgreater{}\ 0} 4. \textbf{Verify imputation}: Step
1 handled missing price values 5. \textbf{Confirm calculation}: Step 3
calculated totals correctly

\textbf{Benefits:} - Quick identification of data loss points - Clear
understanding of data cleaning logic - Easy verification of
transformation correctness - Reproducible debugging process

Note: This also works with polars DataFrames.

\begin{center}\rule{0.5\linewidth}{0.5pt}\end{center}

\section{Real-World Use Cases}\label{real-world-use-cases-1}

\subsection{1. Audit Trail for
Compliance}\label{audit-trail-for-compliance}

Track all transformations for regulatory compliance and audit
requirements.

\begin{Shaded}
\begin{Highlighting}[]
\CommentTok{\# Enable lineage for all operations}
\NormalTok{df }\OperatorTok{=}\NormalTok{ add.transform(}\StringTok{\textquotesingle{}@calc\textquotesingle{}}\NormalTok{, df, ..., lineage}\OperatorTok{=}\VariableTok{True}\NormalTok{)}
\NormalTok{df }\OperatorTok{=}\NormalTok{ add.transform(}\StringTok{\textquotesingle{}@filter\textquotesingle{}}\NormalTok{, df, ..., lineage}\OperatorTok{=}\VariableTok{True}\NormalTok{)}
\NormalTok{df }\OperatorTok{=}\NormalTok{ add.to(df, bring\_from}\OperatorTok{=}\NormalTok{ref, ..., lineage}\OperatorTok{=}\VariableTok{True}\NormalTok{)}

\CommentTok{\# Generate audit report}
\NormalTok{audit\_report }\OperatorTok{=}\NormalTok{ add.scan(}\StringTok{\textquotesingle{}@lineage\textquotesingle{}}\NormalTok{, df)}

\CommentTok{\# Save audit report for compliance}
\ControlFlowTok{with} \BuiltInTok{open}\NormalTok{(}\StringTok{\textquotesingle{}audit\_trail.txt\textquotesingle{}}\NormalTok{, }\StringTok{\textquotesingle{}w\textquotesingle{}}\NormalTok{) }\ImportTok{as}\NormalTok{ f:}
\NormalTok{    f.write(audit\_report)}
\end{Highlighting}
\end{Shaded}

\subsection{2. Pipeline Documentation}\label{pipeline-documentation}

Auto-generate documentation for data pipelines.

\begin{Shaded}
\begin{Highlighting}[]
\CommentTok{\# Build pipeline with lineage}
\NormalTok{result }\OperatorTok{=}\NormalTok{ build\_data\_pipeline(raw\_data, lineage}\OperatorTok{=}\VariableTok{True}\NormalTok{)}

\CommentTok{\# Generate documentation}
\NormalTok{docs }\OperatorTok{=}\NormalTok{ add.scan(}\StringTok{\textquotesingle{}@lineage\textquotesingle{}}\NormalTok{, result)}

\CommentTok{\# Include in README or documentation}
\BuiltInTok{print}\NormalTok{(}\StringTok{"\#\# Data Pipeline}\CharTok{\textbackslash{}n}\StringTok{"}\NormalTok{)}
\BuiltInTok{print}\NormalTok{(docs)}
\end{Highlighting}
\end{Shaded}

\subsection{3. Debugging Production
Issues}\label{debugging-production-issues}

When production data looks wrong, lineage helps identify the problem.

\begin{Shaded}
\begin{Highlighting}[]
\CommentTok{\# Production pipeline with lineage}
\NormalTok{result }\OperatorTok{=}\NormalTok{ production\_pipeline(data, lineage}\OperatorTok{=}\VariableTok{True}\NormalTok{)}

\CommentTok{\# Check lineage when issues arise}
\NormalTok{lineage }\OperatorTok{=}\NormalTok{ add.scan(}\StringTok{\textquotesingle{}@lineage\textquotesingle{}}\NormalTok{, result)}

\CommentTok{\# Identify problematic step}
\CommentTok{\# Look for unexpected row changes or column additions}
\end{Highlighting}
\end{Shaded}

\subsection{4. Collaboration and Code
Review}\label{collaboration-and-code-review}

Share transformation history with team members.

\begin{Shaded}
\begin{Highlighting}[]
\CommentTok{\# Developer creates pipeline}
\NormalTok{result }\OperatorTok{=}\NormalTok{ create\_analysis(data, lineage}\OperatorTok{=}\VariableTok{True}\NormalTok{)}

\CommentTok{\# Share lineage report with reviewer}
\NormalTok{lineage\_report }\OperatorTok{=}\NormalTok{ add.scan(}\StringTok{\textquotesingle{}@lineage\textquotesingle{}}\NormalTok{, result)}

\CommentTok{\# Reviewer can see exact transformation steps}
\CommentTok{\# No need to read through all the code}
\end{Highlighting}
\end{Shaded}

\subsection{5. A/B Testing Validation}\label{ab-testing-validation}

Verify that A/B test groups are processed identically.

\begin{Shaded}
\begin{Highlighting}[]
\CommentTok{\# Process group A}
\NormalTok{group\_a }\OperatorTok{=}\NormalTok{ process\_data(data\_a, lineage}\OperatorTok{=}\VariableTok{True}\NormalTok{)}
\NormalTok{lineage\_a }\OperatorTok{=}\NormalTok{ add.scan(}\StringTok{\textquotesingle{}@lineage\textquotesingle{}}\NormalTok{, group\_a)}

\CommentTok{\# Process group B}
\NormalTok{group\_b }\OperatorTok{=}\NormalTok{ process\_data(data\_b, lineage}\OperatorTok{=}\VariableTok{True}\NormalTok{)}
\NormalTok{lineage\_b }\OperatorTok{=}\NormalTok{ add.scan(}\StringTok{\textquotesingle{}@lineage\textquotesingle{}}\NormalTok{, group\_b)}

\CommentTok{\# Compare lineage reports to ensure identical processing}
\CommentTok{\# (except for input data differences)}
\end{Highlighting}
\end{Shaded}

\begin{center}\rule{0.5\linewidth}{0.5pt}\end{center}

\section{Advanced Patterns}\label{advanced-patterns}

\subsection{Pattern 1: Conditional
Lineage}\label{pattern-1-conditional-lineage}

Enable lineage only in development/debugging mode.

\begin{Shaded}
\begin{Highlighting}[]
\ImportTok{import}\NormalTok{ os}

\CommentTok{\# Enable lineage in development}
\NormalTok{DEBUG\_MODE }\OperatorTok{=}\NormalTok{ os.getenv(}\StringTok{\textquotesingle{}DEBUG\textquotesingle{}}\NormalTok{, }\StringTok{\textquotesingle{}false\textquotesingle{}}\NormalTok{).lower() }\OperatorTok{==} \StringTok{\textquotesingle{}true\textquotesingle{}}

\NormalTok{result }\OperatorTok{=}\NormalTok{ add.transform(}
    \StringTok{\textquotesingle{}@calc\textquotesingle{}}\NormalTok{,}
\NormalTok{    df,}
\NormalTok{    expression}\OperatorTok{=}\StringTok{\textquotesingle{}price * quantity\textquotesingle{}}\NormalTok{,}
\NormalTok{    as\_}\OperatorTok{=}\StringTok{\textquotesingle{}total\textquotesingle{}}\NormalTok{,}
\NormalTok{    lineage}\OperatorTok{=}\NormalTok{DEBUG\_MODE  }\CommentTok{\# Only track in debug mode}
\NormalTok{)}

\ControlFlowTok{if}\NormalTok{ DEBUG\_MODE:}
    \BuiltInTok{print}\NormalTok{(add.scan(}\StringTok{\textquotesingle{}@lineage\textquotesingle{}}\NormalTok{, result))}
\end{Highlighting}
\end{Shaded}

\subsection{Pattern 2: Lineage
Checkpoints}\label{pattern-2-lineage-checkpoints}

Save lineage at key pipeline stages.

\begin{Shaded}
\begin{Highlighting}[]
\CommentTok{\# Stage 1: Data cleaning}
\NormalTok{cleaned }\OperatorTok{=}\NormalTok{ clean\_data(raw\_data, lineage}\OperatorTok{=}\VariableTok{True}\NormalTok{)}
\NormalTok{checkpoint\_1 }\OperatorTok{=}\NormalTok{ add.scan(}\StringTok{\textquotesingle{}@lineage\textquotesingle{}}\NormalTok{, cleaned)}

\CommentTok{\# Stage 2: Feature engineering}
\NormalTok{features }\OperatorTok{=}\NormalTok{ engineer\_features(cleaned, lineage}\OperatorTok{=}\VariableTok{True}\NormalTok{)}
\NormalTok{checkpoint\_2 }\OperatorTok{=}\NormalTok{ add.scan(}\StringTok{\textquotesingle{}@lineage\textquotesingle{}}\NormalTok{, features)}

\CommentTok{\# Stage 3: Final transformations}
\NormalTok{final }\OperatorTok{=}\NormalTok{ final\_transforms(features, lineage}\OperatorTok{=}\VariableTok{True}\NormalTok{)}
\NormalTok{checkpoint\_3 }\OperatorTok{=}\NormalTok{ add.scan(}\StringTok{\textquotesingle{}@lineage\textquotesingle{}}\NormalTok{, final)}

\CommentTok{\# Save all checkpoints}
\NormalTok{save\_checkpoints([checkpoint\_1, checkpoint\_2, checkpoint\_3])}
\end{Highlighting}
\end{Shaded}

\subsection{Pattern 3: Lineage-Based
Testing}\label{pattern-3-lineage-based-testing}

Use lineage to verify pipeline behavior.

\begin{Shaded}
\begin{Highlighting}[]
\KeywordTok{def}\NormalTok{ test\_pipeline\_steps():}
\NormalTok{    result }\OperatorTok{=}\NormalTok{ my\_pipeline(test\_data, lineage}\OperatorTok{=}\VariableTok{True}\NormalTok{)}
\NormalTok{    lineage }\OperatorTok{=}\NormalTok{ add.scan(}\StringTok{\textquotesingle{}@lineage\textquotesingle{}}\NormalTok{, result)}
    
    \CommentTok{\# Verify expected number of steps}
    \ControlFlowTok{assert} \StringTok{\textquotesingle{}Step 1\textquotesingle{}} \KeywordTok{in}\NormalTok{ lineage}
    \ControlFlowTok{assert} \StringTok{\textquotesingle{}Step 2\textquotesingle{}} \KeywordTok{in}\NormalTok{ lineage}
    \ControlFlowTok{assert} \StringTok{\textquotesingle{}Step 3\textquotesingle{}} \KeywordTok{in}\NormalTok{ lineage}
    
    \CommentTok{\# Verify no unexpected data loss}
    \ControlFlowTok{assert} \StringTok{\textquotesingle{}0 rows filtered out\textquotesingle{}} \KeywordTok{in}\NormalTok{ lineage }\KeywordTok{or} \StringTok{\textquotesingle{}no change\textquotesingle{}} \KeywordTok{in}\NormalTok{ lineage}
    
    \CommentTok{\# Verify expected columns were added}
    \ControlFlowTok{assert} \StringTok{\textquotesingle{}Columns Added: total\textquotesingle{}} \KeywordTok{in}\NormalTok{ lineage}
\end{Highlighting}
\end{Shaded}

\begin{center}\rule{0.5\linewidth}{0.5pt}\end{center}

\section{Performance Considerations}\label{performance-considerations}

\subsection{Lineage Overhead}\label{lineage-overhead}

Lineage tracking has minimal performance impact:

\begin{itemize}
\tightlist
\item
  \textbf{Memory}: \textasciitilde1-5KB per operation (JSON metadata)
\item
  \textbf{Time}: \textless100ms per operation (metadata creation)
\item
  \textbf{Computation}: No impact on actual data processing
\end{itemize}

\subsection{When to Disable Lineage}\label{when-to-disable-lineage}

Consider disabling lineage in:

\begin{itemize}
\tightlist
\item
  \textbf{Production pipelines}: Unless required for compliance
\item
  \textbf{Large-scale batch processing}: When processing millions of
  rows
\item
  \textbf{Performance-critical code}: When every millisecond counts
\item
  \textbf{Simple one-off scripts}: When transformation history isn't
  needed
\end{itemize}

\subsection{Lineage Best Practices}\label{lineage-best-practices}

\begin{enumerate}
\def\labelenumi{\arabic{enumi}.}
\tightlist
\item
  \textbf{Enable selectively}: Use lineage where it adds value
\item
  \textbf{Clean up metadata}: Lineage is lost when session ends (by
  design)
\item
  \textbf{Don't over-track}: Not every operation needs lineage
\item
  \textbf{Use for debugging}: Enable when investigating issues
\item
  \textbf{Document with lineage}: Include in code reviews and
  documentation
\end{enumerate}

\begin{center}\rule{0.5\linewidth}{0.5pt}\end{center}

\section{Troubleshooting}\label{troubleshooting-1}

\subsection{Issue: ``No lineage metadata
found''}\label{issue-no-lineage-metadata-found}

\textbf{Cause}: Forgot to add \texttt{lineage=True} to operations.

\textbf{Solution}: Add \texttt{lineage=True} to all operations in the
pipeline.

\subsection{Issue: Lineage lost after type
conversion}\label{issue-lineage-lost-after-type-conversion}

\textbf{Cause}: Used \texttt{as\_type} parameter with
\texttt{lineage=True}.

\textbf{Solution}: Remove \texttt{as\_type} parameter or disable
lineage.

\subsection{Issue: Lineage not showing all
steps}\label{issue-lineage-not-showing-all-steps}

\textbf{Cause}: Some operations didn't have \texttt{lineage=True}.

\textbf{Solution}: Ensure all operations in the chain have
\texttt{lineage=True}.

\begin{center}\rule{0.5\linewidth}{0.5pt}\end{center}

\section{Next Steps}\label{next-steps-20}

\begin{itemize}
\tightlist
\item
  \textbf{Page 1}: Lineage tracking basics
\item
  \textbf{Troubleshooting Guide}: Common errors and solutions
\item
  \textbf{API Reference}: Complete parameter documentation
\end{itemize}

\part{Guides \& Troubleshooting}

\chapter{Troubleshooting Guide}\label{troubleshooting-guide}

A comprehensive reference for common errors, solutions, and best
practices when using Additory.

\begin{center}\rule{0.5\linewidth}{0.5pt}\end{center}

\section{Quick Error Reference}\label{quick-error-reference}

\begin{longtable}[]{@{}
  >{\raggedright\arraybackslash}p{(\linewidth - 4\tabcolsep) * \real{0.3243}}
  >{\raggedright\arraybackslash}p{(\linewidth - 4\tabcolsep) * \real{0.3784}}
  >{\raggedright\arraybackslash}p{(\linewidth - 4\tabcolsep) * \real{0.2973}}@{}}
\toprule\noalign{}
\begin{minipage}[b]{\linewidth}\raggedright
Error Type
\end{minipage} & \begin{minipage}[b]{\linewidth}\raggedright
Common Cause
\end{minipage} & \begin{minipage}[b]{\linewidth}\raggedright
Quick Fix
\end{minipage} \\
\midrule\noalign{}
\endhead
\bottomrule\noalign{}
\endlastfoot
\texttt{TypeError:\ DataFrame\ must\ be\ pandas\ or\ polars} & Wrong
input type & Pass pandas or polars DataFrame \\
\texttt{ValueError:\ Column\ \textquotesingle{}X\textquotesingle{}\ not\ found}
& Missing column & Check column names with \texttt{df.columns} \\
\texttt{ImportError:\ Rust\ bindings\ not\ available} & Missing Rust
extension & Install with \texttt{pip\ install\ additory{[}rust{]}} \\
\texttt{ValueError:\ Cannot\ use\ \textquotesingle{}as\_type\textquotesingle{}\ with\ \textquotesingle{}lineage=True\textquotesingle{}}
& Conflicting parameters & Use either \texttt{lineage=True} OR
\texttt{as\_type}, not both \\
\texttt{RuntimeError:\ Transform\ failed} & Invalid mode or parameters &
Check mode name has \texttt{@} prefix \\
\texttt{KeyError:\ \textquotesingle{}column\_name\textquotesingle{}} &
Column doesn't exist & Verify column exists in DataFrame \\
\end{longtable}

\begin{center}\rule{0.5\linewidth}{0.5pt}\end{center}

\section{Common Errors by Function}\label{common-errors-by-function}

\subsection{add.to() Errors}\label{add.to-errors}

\subsubsection{Error: ``Column `X' not found in reference
DataFrame''}\label{error-column-x-not-found-in-reference-dataframe}

\textbf{Cause}: The column specified in \texttt{bring} doesn't exist in
\texttt{bring\_from} DataFrame.

\textbf{Solution}:

\begin{Shaded}
\begin{Highlighting}[]
\CommentTok{\# ❌ Wrong {-} \textquotesingle{}age\textquotesingle{} doesn\textquotesingle{}t exist in customers}
\NormalTok{result }\OperatorTok{=}\NormalTok{ add.to(orders, bring\_from}\OperatorTok{=}\NormalTok{customers, bring}\OperatorTok{=}\StringTok{\textquotesingle{}age\textquotesingle{}}\NormalTok{, against}\OperatorTok{=}\StringTok{\textquotesingle{}customer\_id\textquotesingle{}}\NormalTok{)}

\CommentTok{\# ✅ Correct {-} check available columns first}
\BuiltInTok{print}\NormalTok{(customers.columns)  }\CommentTok{\# [\textquotesingle{}customer\_id\textquotesingle{}, \textquotesingle{}name\textquotesingle{}, \textquotesingle{}email\textquotesingle{}]}
\NormalTok{result }\OperatorTok{=}\NormalTok{ add.to(orders, bring\_from}\OperatorTok{=}\NormalTok{customers, bring}\OperatorTok{=}\StringTok{\textquotesingle{}name\textquotesingle{}}\NormalTok{, against}\OperatorTok{=}\StringTok{\textquotesingle{}customer\_id\textquotesingle{}}\NormalTok{)}
\end{Highlighting}
\end{Shaded}

\subsubsection{Error: ``Key column `X' not
found''}\label{error-key-column-x-not-found}

\textbf{Cause}: The column specified in \texttt{against} doesn't exist
in one or both DataFrames.

\textbf{Solution}:

\begin{Shaded}
\begin{Highlighting}[]
\CommentTok{\# ❌ Wrong {-} \textquotesingle{}id\textquotesingle{} doesn\textquotesingle{}t exist in orders}
\NormalTok{result }\OperatorTok{=}\NormalTok{ add.to(orders, bring\_from}\OperatorTok{=}\NormalTok{customers, bring}\OperatorTok{=}\StringTok{\textquotesingle{}name\textquotesingle{}}\NormalTok{, against}\OperatorTok{=}\StringTok{\textquotesingle{}id\textquotesingle{}}\NormalTok{)}

\CommentTok{\# ✅ Correct {-} use the actual key column name}
\NormalTok{result }\OperatorTok{=}\NormalTok{ add.to(orders, bring\_from}\OperatorTok{=}\NormalTok{customers, bring}\OperatorTok{=}\StringTok{\textquotesingle{}name\textquotesingle{}}\NormalTok{, against}\OperatorTok{=}\StringTok{\textquotesingle{}customer\_id\textquotesingle{}}\NormalTok{)}
\end{Highlighting}
\end{Shaded}

\subsubsection{Error: ``All reference DataFrames must be the same
type''}\label{error-all-reference-dataframes-must-be-the-same-type}

\textbf{Cause}: Mixing pandas and polars DataFrames in list for
\texttt{bring\_from}.

\textbf{Solution}:

\begin{Shaded}
\begin{Highlighting}[]
\CommentTok{\# ❌ Wrong {-} mixing pandas and polars}
\NormalTok{orders\_jan }\OperatorTok{=}\NormalTok{ pd.DataFrame(...)  }\CommentTok{\# pandas}
\NormalTok{orders\_feb }\OperatorTok{=}\NormalTok{ pl.DataFrame(...)  }\CommentTok{\# polars}
\NormalTok{result }\OperatorTok{=}\NormalTok{ add.to(customers, bring\_from}\OperatorTok{=}\NormalTok{[orders\_jan, orders\_feb], ...)}

\CommentTok{\# ✅ Correct {-} use same type for all}
\NormalTok{orders\_jan }\OperatorTok{=}\NormalTok{ pd.DataFrame(...)  }\CommentTok{\# pandas}
\NormalTok{orders\_feb }\OperatorTok{=}\NormalTok{ pd.DataFrame(...)  }\CommentTok{\# pandas}
\NormalTok{result }\OperatorTok{=}\NormalTok{ add.to(customers, bring\_from}\OperatorTok{=}\NormalTok{[orders\_jan, orders\_feb], ...)}
\end{Highlighting}
\end{Shaded}

\subsubsection{Error: ``strategy parameter must be a
dict''}\label{error-strategy-parameter-must-be-a-dict}

\textbf{Cause}: Passing string instead of dict for aggregation strategy.

\textbf{Solution}:

\begin{Shaded}
\begin{Highlighting}[]
\CommentTok{\# ❌ Wrong {-} strategy is a string}
\NormalTok{result }\OperatorTok{=}\NormalTok{ add.to(customers, bring\_from}\OperatorTok{=}\NormalTok{orders, bring}\OperatorTok{=}\StringTok{\textquotesingle{}amount\textquotesingle{}}\NormalTok{, }
\NormalTok{                against}\OperatorTok{=}\StringTok{\textquotesingle{}customer\_id\textquotesingle{}}\NormalTok{, strategy}\OperatorTok{=}\StringTok{\textquotesingle{}sum\textquotesingle{}}\NormalTok{)}

\CommentTok{\# ✅ Correct {-} strategy is a dict}
\NormalTok{result }\OperatorTok{=}\NormalTok{ add.to(customers, bring\_from}\OperatorTok{=}\NormalTok{orders, bring}\OperatorTok{=}\StringTok{\textquotesingle{}amount\textquotesingle{}}\NormalTok{, }
\NormalTok{                against}\OperatorTok{=}\StringTok{\textquotesingle{}customer\_id\textquotesingle{}}\NormalTok{, strategy}\OperatorTok{=}\NormalTok{\{}\StringTok{\textquotesingle{}amount\textquotesingle{}}\NormalTok{: }\StringTok{\textquotesingle{}sum\textquotesingle{}}\NormalTok{\})}
\end{Highlighting}
\end{Shaded}

\begin{center}\rule{0.5\linewidth}{0.5pt}\end{center}

\subsection{add.transform() Errors}\label{add.transform-errors}

\subsubsection{\texorpdfstring{Error: ``Mode `(\textbf{calc?})' requires
`expression'
parameter''}{Error: ``Mode `(calc?)' requires `expression' parameter''}}\label{error-mode-calc-requires-expression-parameter}

\textbf{Cause}: Missing required \texttt{expression} parameter for
calculation mode.

\textbf{Solution}:

\begin{Shaded}
\begin{Highlighting}[]
\CommentTok{\# ❌ Wrong {-} missing expression}
\NormalTok{result }\OperatorTok{=}\NormalTok{ add.transform(}\StringTok{\textquotesingle{}@calc\textquotesingle{}}\NormalTok{, df, columns}\OperatorTok{=}\NormalTok{[}\StringTok{\textquotesingle{}price\textquotesingle{}}\NormalTok{, }\StringTok{\textquotesingle{}quantity\textquotesingle{}}\NormalTok{], as\_}\OperatorTok{=}\StringTok{\textquotesingle{}total\textquotesingle{}}\NormalTok{)}

\CommentTok{\# ✅ Correct {-} include expression}
\NormalTok{result }\OperatorTok{=}\NormalTok{ add.transform(}\StringTok{\textquotesingle{}@calc\textquotesingle{}}\NormalTok{, df, columns}\OperatorTok{=}\NormalTok{[}\StringTok{\textquotesingle{}price\textquotesingle{}}\NormalTok{, }\StringTok{\textquotesingle{}quantity\textquotesingle{}}\NormalTok{], }
\NormalTok{                       expression}\OperatorTok{=}\StringTok{\textquotesingle{}price * quantity\textquotesingle{}}\NormalTok{, as\_}\OperatorTok{=}\StringTok{\textquotesingle{}total\textquotesingle{}}\NormalTok{)}
\end{Highlighting}
\end{Shaded}

\subsubsection{\texorpdfstring{Error: ``Mode `(\textbf{filter?})'
requires `where'
parameter''}{Error: ``Mode `(filter?)' requires `where' parameter''}}\label{error-mode-filter-requires-where-parameter}

\textbf{Cause}: Missing required \texttt{where} parameter for filter
mode.

\textbf{Solution}:

\begin{Shaded}
\begin{Highlighting}[]
\CommentTok{\# ❌ Wrong {-} missing where condition}
\NormalTok{result }\OperatorTok{=}\NormalTok{ add.transform(}\StringTok{\textquotesingle{}@filter\textquotesingle{}}\NormalTok{, df, columns}\OperatorTok{=}\NormalTok{[}\StringTok{\textquotesingle{}price\textquotesingle{}}\NormalTok{, }\StringTok{\textquotesingle{}stock\textquotesingle{}}\NormalTok{])}

\CommentTok{\# ✅ Correct {-} include where condition}
\NormalTok{result }\OperatorTok{=}\NormalTok{ add.transform(}\StringTok{\textquotesingle{}@filter\textquotesingle{}}\NormalTok{, df, columns}\OperatorTok{=}\NormalTok{[}\StringTok{\textquotesingle{}price\textquotesingle{}}\NormalTok{, }\StringTok{\textquotesingle{}stock\textquotesingle{}}\NormalTok{], where}\OperatorTok{=}\StringTok{\textquotesingle{}stock \textgreater{} 0\textquotesingle{}}\NormalTok{)}
\end{Highlighting}
\end{Shaded}

\subsubsection{\texorpdfstring{Error: ``Mode `(\textbf{aggregate?})'
requires `by'
parameter''}{Error: ``Mode `(aggregate?)' requires `by' parameter''}}\label{error-mode-aggregate-requires-by-parameter}

\textbf{Cause}: Missing required \texttt{by} parameter for aggregation
mode.

\textbf{Solution}:

\begin{Shaded}
\begin{Highlighting}[]
\CommentTok{\# ❌ Wrong {-} missing by parameter}
\NormalTok{result }\OperatorTok{=}\NormalTok{ add.transform(}\StringTok{\textquotesingle{}@aggregate\textquotesingle{}}\NormalTok{, df, columns}\OperatorTok{=}\NormalTok{[}\StringTok{\textquotesingle{}region\textquotesingle{}}\NormalTok{, }\StringTok{\textquotesingle{}sales\textquotesingle{}}\NormalTok{], }
\NormalTok{                       strategy}\OperatorTok{=}\NormalTok{\{}\StringTok{\textquotesingle{}sales\textquotesingle{}}\NormalTok{: }\StringTok{\textquotesingle{}sum\textquotesingle{}}\NormalTok{\})}

\CommentTok{\# ✅ Correct {-} include by parameter}
\NormalTok{result }\OperatorTok{=}\NormalTok{ add.transform(}\StringTok{\textquotesingle{}@aggregate\textquotesingle{}}\NormalTok{, df, columns}\OperatorTok{=}\NormalTok{[}\StringTok{\textquotesingle{}region\textquotesingle{}}\NormalTok{, }\StringTok{\textquotesingle{}sales\textquotesingle{}}\NormalTok{], }
\NormalTok{                       by}\OperatorTok{=}\StringTok{\textquotesingle{}region\textquotesingle{}}\NormalTok{, strategy}\OperatorTok{=}\NormalTok{\{}\StringTok{\textquotesingle{}sales\textquotesingle{}}\NormalTok{: }\StringTok{\textquotesingle{}sum\textquotesingle{}}\NormalTok{\})}
\end{Highlighting}
\end{Shaded}

\subsubsection{\texorpdfstring{Error: ``Invalid mode `calc' - did you
mean
`(\textbf{calc?})'?''}{Error: ``Invalid mode `calc' - did you mean `(calc?)'?''}}\label{error-invalid-mode-calc---did-you-mean-calc}

\textbf{Cause}: Missing \texttt{@} prefix on mode name.

\textbf{Solution}:

\begin{Shaded}
\begin{Highlighting}[]
\CommentTok{\# ❌ Wrong {-} missing @ prefix}
\NormalTok{result }\OperatorTok{=}\NormalTok{ add.transform(}\StringTok{\textquotesingle{}calc\textquotesingle{}}\NormalTok{, df, expression}\OperatorTok{=}\StringTok{\textquotesingle{}price * 2\textquotesingle{}}\NormalTok{, as\_}\OperatorTok{=}\StringTok{\textquotesingle{}doubled\textquotesingle{}}\NormalTok{)}

\CommentTok{\# ✅ Correct {-} include @ prefix}
\NormalTok{result }\OperatorTok{=}\NormalTok{ add.transform(}\StringTok{\textquotesingle{}@calc\textquotesingle{}}\NormalTok{, df, expression}\OperatorTok{=}\StringTok{\textquotesingle{}price * 2\textquotesingle{}}\NormalTok{, as\_}\OperatorTok{=}\StringTok{\textquotesingle{}doubled\textquotesingle{}}\NormalTok{)}
\end{Highlighting}
\end{Shaded}

\subsubsection{Error: ``Number of expressions must match number of
output
names''}\label{error-number-of-expressions-must-match-number-of-output-names}

\textbf{Cause}: Mismatch between number of expressions and \texttt{as\_}
names.

\textbf{Solution}:

\begin{Shaded}
\begin{Highlighting}[]
\CommentTok{\# ❌ Wrong {-} 2 expressions but 1 output name}
\NormalTok{result }\OperatorTok{=}\NormalTok{ add.transform(}\StringTok{\textquotesingle{}@calc\textquotesingle{}}\NormalTok{, df, }
\NormalTok{                       expression}\OperatorTok{=}\NormalTok{[}\StringTok{\textquotesingle{}price * quantity\textquotesingle{}}\NormalTok{, }\StringTok{\textquotesingle{}price {-} cost\textquotesingle{}}\NormalTok{],}
\NormalTok{                       as\_}\OperatorTok{=}\StringTok{\textquotesingle{}total\textquotesingle{}}\NormalTok{)}

\CommentTok{\# ✅ Correct {-} matching counts}
\NormalTok{result }\OperatorTok{=}\NormalTok{ add.transform(}\StringTok{\textquotesingle{}@calc\textquotesingle{}}\NormalTok{, df, }
\NormalTok{                       expression}\OperatorTok{=}\NormalTok{[}\StringTok{\textquotesingle{}price * quantity\textquotesingle{}}\NormalTok{, }\StringTok{\textquotesingle{}price {-} cost\textquotesingle{}}\NormalTok{],}
\NormalTok{                       as\_}\OperatorTok{=}\NormalTok{[}\StringTok{\textquotesingle{}total\textquotesingle{}}\NormalTok{, }\StringTok{\textquotesingle{}profit\textquotesingle{}}\NormalTok{])}
\end{Highlighting}
\end{Shaded}

\begin{center}\rule{0.5\linewidth}{0.5pt}\end{center}

\subsection{add.synthetic() Errors}\label{add.synthetic-errors}

\subsubsection{\texorpdfstring{Error: ``Mode `(\textbf{new?})' requires
`n'
parameter''}{Error: ``Mode `(new?)' requires `n' parameter''}}\label{error-mode-new-requires-n-parameter}

\textbf{Cause}: Missing required \texttt{n} parameter for creating new
DataFrame.

\textbf{Solution}:

\begin{Shaded}
\begin{Highlighting}[]
\CommentTok{\# ❌ Wrong {-} missing n parameter}
\NormalTok{result }\OperatorTok{=}\NormalTok{ add.synthetic(}\StringTok{\textquotesingle{}@new\textquotesingle{}}\NormalTok{, strategy}\OperatorTok{=}\NormalTok{\{}\StringTok{\textquotesingle{}id\textquotesingle{}}\NormalTok{: \{}\StringTok{\textquotesingle{}type\textquotesingle{}}\NormalTok{: }\StringTok{\textquotesingle{}sequence\textquotesingle{}}\NormalTok{\}\})}

\CommentTok{\# ✅ Correct {-} include n parameter}
\NormalTok{result }\OperatorTok{=}\NormalTok{ add.synthetic(}\StringTok{\textquotesingle{}@new\textquotesingle{}}\NormalTok{, n}\OperatorTok{=}\DecValTok{100}\NormalTok{, strategy}\OperatorTok{=}\NormalTok{\{}\StringTok{\textquotesingle{}id\textquotesingle{}}\NormalTok{: \{}\StringTok{\textquotesingle{}type\textquotesingle{}}\NormalTok{: }\StringTok{\textquotesingle{}sequence\textquotesingle{}}\NormalTok{\}\})}
\end{Highlighting}
\end{Shaded}

\subsubsection{\texorpdfstring{Error: ``Mode `(\textbf{augment?})'
requires `df'
parameter''}{Error: ``Mode `(augment?)' requires `df' parameter''}}\label{error-mode-augment-requires-df-parameter}

\textbf{Cause}: Missing required \texttt{df} parameter for augmenting
existing DataFrame.

\textbf{Solution}:

\begin{Shaded}
\begin{Highlighting}[]
\CommentTok{\# ❌ Wrong {-} missing df parameter}
\NormalTok{result }\OperatorTok{=}\NormalTok{ add.synthetic(}\StringTok{\textquotesingle{}@augment\textquotesingle{}}\NormalTok{, n}\OperatorTok{=}\DecValTok{50}\NormalTok{, strategy}\OperatorTok{=}\NormalTok{\{}\StringTok{\textquotesingle{}age\textquotesingle{}}\NormalTok{: \{}\StringTok{\textquotesingle{}distribution\textquotesingle{}}\NormalTok{: }\StringTok{\textquotesingle{}normal\textquotesingle{}}\NormalTok{\}\})}

\CommentTok{\# ✅ Correct {-} include df parameter}
\NormalTok{result }\OperatorTok{=}\NormalTok{ add.synthetic(}\StringTok{\textquotesingle{}@augment\textquotesingle{}}\NormalTok{, df}\OperatorTok{=}\NormalTok{existing\_df, n}\OperatorTok{=}\DecValTok{50}\NormalTok{, }
\NormalTok{                       strategy}\OperatorTok{=}\NormalTok{\{}\StringTok{\textquotesingle{}age\textquotesingle{}}\NormalTok{: \{}\StringTok{\textquotesingle{}distribution\textquotesingle{}}\NormalTok{: }\StringTok{\textquotesingle{}normal\textquotesingle{}}\NormalTok{\}\})}
\end{Highlighting}
\end{Shaded}

\subsubsection{Error: ``Invalid distribution `gaussian' - did you mean
`normal'?''}\label{error-invalid-distribution-gaussian---did-you-mean-normal}

\textbf{Cause}: Using incorrect distribution name.

\textbf{Solution}:

\begin{Shaded}
\begin{Highlighting}[]
\CommentTok{\# ❌ Wrong {-} \textquotesingle{}gaussian\textquotesingle{} is not a valid distribution name}
\NormalTok{result }\OperatorTok{=}\NormalTok{ add.synthetic(}\StringTok{\textquotesingle{}@new\textquotesingle{}}\NormalTok{, n}\OperatorTok{=}\DecValTok{100}\NormalTok{, }
\NormalTok{                       strategy}\OperatorTok{=}\NormalTok{\{}\StringTok{\textquotesingle{}age\textquotesingle{}}\NormalTok{: \{}\StringTok{\textquotesingle{}distribution\textquotesingle{}}\NormalTok{: }\StringTok{\textquotesingle{}gaussian\textquotesingle{}}\NormalTok{\}\})}

\CommentTok{\# ✅ Correct {-} use \textquotesingle{}normal\textquotesingle{} for Gaussian distribution}
\NormalTok{result }\OperatorTok{=}\NormalTok{ add.synthetic(}\StringTok{\textquotesingle{}@new\textquotesingle{}}\NormalTok{, n}\OperatorTok{=}\DecValTok{100}\NormalTok{, }
\NormalTok{                       strategy}\OperatorTok{=}\NormalTok{\{}\StringTok{\textquotesingle{}age\textquotesingle{}}\NormalTok{: \{}\StringTok{\textquotesingle{}distribution\textquotesingle{}}\NormalTok{: }\StringTok{\textquotesingle{}normal\textquotesingle{}}\NormalTok{, }\StringTok{\textquotesingle{}mean\textquotesingle{}}\NormalTok{: }\DecValTok{35}\NormalTok{, }\StringTok{\textquotesingle{}std\textquotesingle{}}\NormalTok{: }\DecValTok{8}\NormalTok{\}\})}
\end{Highlighting}
\end{Shaded}

\textbf{Valid distributions}: \texttt{normal}, \texttt{lognormal},
\texttt{uniform}, \texttt{exponential}, \texttt{poisson},
\texttt{binomial}, \texttt{beta}

\subsubsection{Error: ``LinkedList requires `levels'
parameter''}\label{error-linkedlist-requires-levels-parameter}

\textbf{Cause}: Missing required \texttt{levels} parameter for linked
list strategy.

\textbf{Solution}:

\begin{Shaded}
\begin{Highlighting}[]
\CommentTok{\# ❌ Wrong {-} missing levels}
\NormalTok{result }\OperatorTok{=}\NormalTok{ add.synthetic(}\StringTok{\textquotesingle{}@new\textquotesingle{}}\NormalTok{, n}\OperatorTok{=}\DecValTok{10}\NormalTok{, }
\NormalTok{                       strategy}\OperatorTok{=}\NormalTok{\{}\StringTok{\textquotesingle{}location\textquotesingle{}}\NormalTok{: \{}\StringTok{\textquotesingle{}type\textquotesingle{}}\NormalTok{: }\StringTok{\textquotesingle{}linked\_list\textquotesingle{}}\NormalTok{\}\})}

\CommentTok{\# ✅ Correct {-} include levels}
\NormalTok{result }\OperatorTok{=}\NormalTok{ add.synthetic(}\StringTok{\textquotesingle{}@new\textquotesingle{}}\NormalTok{, n}\OperatorTok{=}\DecValTok{10}\NormalTok{, }
\NormalTok{                       strategy}\OperatorTok{=}\NormalTok{\{}\StringTok{\textquotesingle{}location\textquotesingle{}}\NormalTok{: \{}
                           \StringTok{\textquotesingle{}type\textquotesingle{}}\NormalTok{: }\StringTok{\textquotesingle{}linked\_list\textquotesingle{}}\NormalTok{,}
                           \StringTok{\textquotesingle{}levels\textquotesingle{}}\NormalTok{: [}
\NormalTok{                               [}\StringTok{\textquotesingle{}USA\textquotesingle{}}\NormalTok{, }\StringTok{\textquotesingle{}Canada\textquotesingle{}}\NormalTok{, }\StringTok{\textquotesingle{}Mexico\textquotesingle{}}\NormalTok{],}
\NormalTok{                               [}\StringTok{\textquotesingle{}NY\textquotesingle{}}\NormalTok{, }\StringTok{\textquotesingle{}Toronto\textquotesingle{}}\NormalTok{, }\StringTok{\textquotesingle{}CDMX\textquotesingle{}}\NormalTok{]}
\NormalTok{                           ]}
\NormalTok{                       \}\})}
\end{Highlighting}
\end{Shaded}

\begin{center}\rule{0.5\linewidth}{0.5pt}\end{center}

\subsection{add.scan() Errors}\label{add.scan-errors}

\subsubsection{\texorpdfstring{Error: ``Mode `(\textbf{lineage?})'
requires lineage=True in previous
operations''}{Error: ``Mode `(lineage?)' requires lineage=True in previous operations''}}\label{error-mode-lineage-requires-lineagetrue-in-previous-operations}

\textbf{Cause}: Trying to scan lineage on DataFrame that wasn't created
with \texttt{lineage=True}.

\textbf{Solution}:

\begin{Shaded}
\begin{Highlighting}[]
\CommentTok{\# ❌ Wrong {-} no lineage tracking enabled}
\NormalTok{df }\OperatorTok{=}\NormalTok{ pd.DataFrame(\{}\StringTok{\textquotesingle{}a\textquotesingle{}}\NormalTok{: [}\DecValTok{1}\NormalTok{, }\DecValTok{2}\NormalTok{, }\DecValTok{3}\NormalTok{]\})}
\NormalTok{result }\OperatorTok{=}\NormalTok{ add.transform(}\StringTok{\textquotesingle{}@calc\textquotesingle{}}\NormalTok{, df, expression}\OperatorTok{=}\StringTok{\textquotesingle{}a * 2\textquotesingle{}}\NormalTok{, as\_}\OperatorTok{=}\StringTok{\textquotesingle{}b\textquotesingle{}}\NormalTok{)}
\NormalTok{lineage }\OperatorTok{=}\NormalTok{ add.scan(}\StringTok{\textquotesingle{}@lineage\textquotesingle{}}\NormalTok{, result)  }\CommentTok{\# Error!}

\CommentTok{\# ✅ Correct {-} enable lineage tracking}
\NormalTok{df }\OperatorTok{=}\NormalTok{ pd.DataFrame(\{}\StringTok{\textquotesingle{}a\textquotesingle{}}\NormalTok{: [}\DecValTok{1}\NormalTok{, }\DecValTok{2}\NormalTok{, }\DecValTok{3}\NormalTok{]\})}
\NormalTok{result }\OperatorTok{=}\NormalTok{ add.transform(}\StringTok{\textquotesingle{}@calc\textquotesingle{}}\NormalTok{, df, expression}\OperatorTok{=}\StringTok{\textquotesingle{}a * 2\textquotesingle{}}\NormalTok{, as\_}\OperatorTok{=}\StringTok{\textquotesingle{}b\textquotesingle{}}\NormalTok{, lineage}\OperatorTok{=}\VariableTok{True}\NormalTok{)}
\NormalTok{lineage }\OperatorTok{=}\NormalTok{ add.scan(}\StringTok{\textquotesingle{}@lineage\textquotesingle{}}\NormalTok{, result)  }\CommentTok{\# Works!}
\end{Highlighting}
\end{Shaded}

\subsubsection{\texorpdfstring{Error: ``Invalid mode `analyze' - did you
mean
`(\textbf{analyze?})'?''}{Error: ``Invalid mode `analyze' - did you mean `(analyze?)'?''}}\label{error-invalid-mode-analyze---did-you-mean-analyze}

\textbf{Cause}: Missing \texttt{@} prefix on mode name.

\textbf{Solution}:

\begin{Shaded}
\begin{Highlighting}[]
\CommentTok{\# ❌ Wrong {-} missing @ prefix}
\NormalTok{result }\OperatorTok{=}\NormalTok{ add.scan(}\StringTok{\textquotesingle{}analyze\textquotesingle{}}\NormalTok{, df)}

\CommentTok{\# ✅ Correct {-} include @ prefix}
\NormalTok{result }\OperatorTok{=}\NormalTok{ add.scan(}\StringTok{\textquotesingle{}@analyze\textquotesingle{}}\NormalTok{, df)}
\end{Highlighting}
\end{Shaded}

\begin{center}\rule{0.5\linewidth}{0.5pt}\end{center}

\section{Parameter Validation Errors}\label{parameter-validation-errors}

\subsection{Error: ``Cannot use `as\_type' with
`lineage=True'\,''}\label{error-cannot-use-as_type-with-lineagetrue}

\textbf{Cause}: Trying to use both \texttt{lineage=True} and
\texttt{as\_type} together.

\textbf{Explanation}: Lineage metadata is stored in the DataFrame's
native format and would be lost during type conversion.

\textbf{Solution}:

\begin{Shaded}
\begin{Highlighting}[]
\CommentTok{\# ❌ Wrong {-} using both lineage and as\_type}
\NormalTok{result }\OperatorTok{=}\NormalTok{ add.to(orders, bring\_from}\OperatorTok{=}\NormalTok{customers, bring}\OperatorTok{=}\StringTok{\textquotesingle{}name\textquotesingle{}}\NormalTok{, }
\NormalTok{                against}\OperatorTok{=}\StringTok{\textquotesingle{}customer\_id\textquotesingle{}}\NormalTok{, lineage}\OperatorTok{=}\VariableTok{True}\NormalTok{, as\_type}\OperatorTok{=}\StringTok{\textquotesingle{}polars\textquotesingle{}}\NormalTok{)}

\CommentTok{\# ✅ Option 1: Track lineage (returns same type as input)}
\NormalTok{result }\OperatorTok{=}\NormalTok{ add.to(orders, bring\_from}\OperatorTok{=}\NormalTok{customers, bring}\OperatorTok{=}\StringTok{\textquotesingle{}name\textquotesingle{}}\NormalTok{, }
\NormalTok{                against}\OperatorTok{=}\StringTok{\textquotesingle{}customer\_id\textquotesingle{}}\NormalTok{, lineage}\OperatorTok{=}\VariableTok{True}\NormalTok{)}

\CommentTok{\# ✅ Option 2: Convert type (no lineage tracking)}
\NormalTok{result }\OperatorTok{=}\NormalTok{ add.to(orders, bring\_from}\OperatorTok{=}\NormalTok{customers, bring}\OperatorTok{=}\StringTok{\textquotesingle{}name\textquotesingle{}}\NormalTok{, }
\NormalTok{                against}\OperatorTok{=}\StringTok{\textquotesingle{}customer\_id\textquotesingle{}}\NormalTok{, as\_type}\OperatorTok{=}\StringTok{\textquotesingle{}polars\textquotesingle{}}\NormalTok{)}

\CommentTok{\# ✅ Option 3: Convert after tracking lineage (lineage will be lost)}
\NormalTok{result }\OperatorTok{=}\NormalTok{ add.to(orders, bring\_from}\OperatorTok{=}\NormalTok{customers, bring}\OperatorTok{=}\StringTok{\textquotesingle{}name\textquotesingle{}}\NormalTok{, }
\NormalTok{                against}\OperatorTok{=}\StringTok{\textquotesingle{}customer\_id\textquotesingle{}}\NormalTok{, lineage}\OperatorTok{=}\VariableTok{True}\NormalTok{)}
\NormalTok{result\_polars }\OperatorTok{=}\NormalTok{ pl.from\_pandas(result)  }\CommentTok{\# Convert separately}
\end{Highlighting}
\end{Shaded}

\begin{center}\rule{0.5\linewidth}{0.5pt}\end{center}

\section{Data Type Errors}\label{data-type-errors}

\subsection{Error: ``TypeError: DataFrame must be pandas or
polars''}\label{error-typeerror-dataframe-must-be-pandas-or-polars}

\textbf{Cause}: Passing wrong type (dict, list, numpy array, etc.)
instead of DataFrame.

\textbf{Solution}:

\begin{Shaded}
\begin{Highlighting}[]
\CommentTok{\# ❌ Wrong {-} passing dict}
\NormalTok{data }\OperatorTok{=}\NormalTok{ \{}\StringTok{\textquotesingle{}a\textquotesingle{}}\NormalTok{: [}\DecValTok{1}\NormalTok{, }\DecValTok{2}\NormalTok{, }\DecValTok{3}\NormalTok{], }\StringTok{\textquotesingle{}b\textquotesingle{}}\NormalTok{: [}\DecValTok{4}\NormalTok{, }\DecValTok{5}\NormalTok{, }\DecValTok{6}\NormalTok{]\}}
\NormalTok{result }\OperatorTok{=}\NormalTok{ add.transform(}\StringTok{\textquotesingle{}@calc\textquotesingle{}}\NormalTok{, data, expression}\OperatorTok{=}\StringTok{\textquotesingle{}a + b\textquotesingle{}}\NormalTok{, as\_}\OperatorTok{=}\StringTok{\textquotesingle{}c\textquotesingle{}}\NormalTok{)}

\CommentTok{\# ✅ Correct {-} convert to DataFrame first}
\ImportTok{import}\NormalTok{ pandas }\ImportTok{as}\NormalTok{ pd}
\NormalTok{df }\OperatorTok{=}\NormalTok{ pd.DataFrame(data)}
\NormalTok{result }\OperatorTok{=}\NormalTok{ add.transform(}\StringTok{\textquotesingle{}@calc\textquotesingle{}}\NormalTok{, df, expression}\OperatorTok{=}\StringTok{\textquotesingle{}a + b\textquotesingle{}}\NormalTok{, as\_}\OperatorTok{=}\StringTok{\textquotesingle{}c\textquotesingle{}}\NormalTok{)}
\end{Highlighting}
\end{Shaded}

\subsection{Error: ``Column `X' has incompatible type for
operation''}\label{error-column-x-has-incompatible-type-for-operation}

\textbf{Cause}: Trying to perform numeric operations on string columns
or vice versa.

\textbf{Solution}:

\begin{Shaded}
\begin{Highlighting}[]
\CommentTok{\# ❌ Wrong {-} trying to multiply string column}
\NormalTok{df }\OperatorTok{=}\NormalTok{ pd.DataFrame(\{}\StringTok{\textquotesingle{}product\textquotesingle{}}\NormalTok{: [}\StringTok{\textquotesingle{}A\textquotesingle{}}\NormalTok{, }\StringTok{\textquotesingle{}B\textquotesingle{}}\NormalTok{], }\StringTok{\textquotesingle{}quantity\textquotesingle{}}\NormalTok{: [}\DecValTok{10}\NormalTok{, }\DecValTok{20}\NormalTok{]\})}
\NormalTok{result }\OperatorTok{=}\NormalTok{ add.transform(}\StringTok{\textquotesingle{}@calc\textquotesingle{}}\NormalTok{, df, expression}\OperatorTok{=}\StringTok{\textquotesingle{}product * 2\textquotesingle{}}\NormalTok{, as\_}\OperatorTok{=}\StringTok{\textquotesingle{}doubled\textquotesingle{}}\NormalTok{)}

\CommentTok{\# ✅ Correct {-} use numeric columns for math operations}
\NormalTok{result }\OperatorTok{=}\NormalTok{ add.transform(}\StringTok{\textquotesingle{}@calc\textquotesingle{}}\NormalTok{, df, expression}\OperatorTok{=}\StringTok{\textquotesingle{}quantity * 2\textquotesingle{}}\NormalTok{, as\_}\OperatorTok{=}\StringTok{\textquotesingle{}doubled\textquotesingle{}}\NormalTok{)}
\end{Highlighting}
\end{Shaded}

\begin{center}\rule{0.5\linewidth}{0.5pt}\end{center}

\section{Empty DataFrame Errors}\label{empty-dataframe-errors}

\subsection{Error: ``Cannot perform operation on empty
DataFrame''}\label{error-cannot-perform-operation-on-empty-dataframe}

\textbf{Cause}: Trying to transform or analyze DataFrame with 0 rows.

\textbf{Solution}:

\begin{Shaded}
\begin{Highlighting}[]
\CommentTok{\# ❌ Wrong {-} operating on empty DataFrame}
\NormalTok{df }\OperatorTok{=}\NormalTok{ pd.DataFrame(\{}\StringTok{\textquotesingle{}a\textquotesingle{}}\NormalTok{: [], }\StringTok{\textquotesingle{}b\textquotesingle{}}\NormalTok{: []\})}
\NormalTok{result }\OperatorTok{=}\NormalTok{ add.transform(}\StringTok{\textquotesingle{}@calc\textquotesingle{}}\NormalTok{, df, expression}\OperatorTok{=}\StringTok{\textquotesingle{}a + b\textquotesingle{}}\NormalTok{, as\_}\OperatorTok{=}\StringTok{\textquotesingle{}c\textquotesingle{}}\NormalTok{)}

\CommentTok{\# ✅ Correct {-} check for empty DataFrame first}
\ControlFlowTok{if} \BuiltInTok{len}\NormalTok{(df) }\OperatorTok{\textgreater{}} \DecValTok{0}\NormalTok{:}
\NormalTok{    result }\OperatorTok{=}\NormalTok{ add.transform(}\StringTok{\textquotesingle{}@calc\textquotesingle{}}\NormalTok{, df, expression}\OperatorTok{=}\StringTok{\textquotesingle{}a + b\textquotesingle{}}\NormalTok{, as\_}\OperatorTok{=}\StringTok{\textquotesingle{}c\textquotesingle{}}\NormalTok{)}
\ControlFlowTok{else}\NormalTok{:}
    \BuiltInTok{print}\NormalTok{(}\StringTok{"DataFrame is empty, skipping transformation"}\NormalTok{)}
\NormalTok{    result }\OperatorTok{=}\NormalTok{ df}
\end{Highlighting}
\end{Shaded}

\begin{center}\rule{0.5\linewidth}{0.5pt}\end{center}

\section{Missing Value Errors}\label{missing-value-errors}

\subsection{Error: ``Cannot aggregate column with all null
values''}\label{error-cannot-aggregate-column-with-all-null-values}

\textbf{Cause}: Trying to aggregate a column that contains only
null/None values.

\textbf{Solution}:

\begin{Shaded}
\begin{Highlighting}[]
\CommentTok{\# ❌ Wrong {-} column has all nulls}
\NormalTok{df }\OperatorTok{=}\NormalTok{ pd.DataFrame(\{}\StringTok{\textquotesingle{}region\textquotesingle{}}\NormalTok{: [}\StringTok{\textquotesingle{}A\textquotesingle{}}\NormalTok{, }\StringTok{\textquotesingle{}A\textquotesingle{}}\NormalTok{, }\StringTok{\textquotesingle{}B\textquotesingle{}}\NormalTok{], }\StringTok{\textquotesingle{}sales\textquotesingle{}}\NormalTok{: [}\VariableTok{None}\NormalTok{, }\VariableTok{None}\NormalTok{, }\VariableTok{None}\NormalTok{]\})}
\NormalTok{result }\OperatorTok{=}\NormalTok{ add.transform(}\StringTok{\textquotesingle{}@aggregate\textquotesingle{}}\NormalTok{, df, columns}\OperatorTok{=}\NormalTok{[}\StringTok{\textquotesingle{}region\textquotesingle{}}\NormalTok{, }\StringTok{\textquotesingle{}sales\textquotesingle{}}\NormalTok{], }
\NormalTok{                       by}\OperatorTok{=}\StringTok{\textquotesingle{}region\textquotesingle{}}\NormalTok{, strategy}\OperatorTok{=}\NormalTok{\{}\StringTok{\textquotesingle{}sales\textquotesingle{}}\NormalTok{: }\StringTok{\textquotesingle{}sum\textquotesingle{}}\NormalTok{\})}

\CommentTok{\# ✅ Correct {-} handle nulls first}
\NormalTok{df }\OperatorTok{=}\NormalTok{ pd.DataFrame(\{}\StringTok{\textquotesingle{}region\textquotesingle{}}\NormalTok{: [}\StringTok{\textquotesingle{}A\textquotesingle{}}\NormalTok{, }\StringTok{\textquotesingle{}A\textquotesingle{}}\NormalTok{, }\StringTok{\textquotesingle{}B\textquotesingle{}}\NormalTok{], }\StringTok{\textquotesingle{}sales\textquotesingle{}}\NormalTok{: [}\VariableTok{None}\NormalTok{, }\VariableTok{None}\NormalTok{, }\VariableTok{None}\NormalTok{]\})}
\CommentTok{\# Fill nulls with 0 or use @deduce mode}
\NormalTok{df[}\StringTok{\textquotesingle{}sales\textquotesingle{}}\NormalTok{] }\OperatorTok{=}\NormalTok{ df[}\StringTok{\textquotesingle{}sales\textquotesingle{}}\NormalTok{].fillna(}\DecValTok{0}\NormalTok{)}
\NormalTok{result }\OperatorTok{=}\NormalTok{ add.transform(}\StringTok{\textquotesingle{}@aggregate\textquotesingle{}}\NormalTok{, df, columns}\OperatorTok{=}\NormalTok{[}\StringTok{\textquotesingle{}region\textquotesingle{}}\NormalTok{, }\StringTok{\textquotesingle{}sales\textquotesingle{}}\NormalTok{], }
\NormalTok{                       by}\OperatorTok{=}\StringTok{\textquotesingle{}region\textquotesingle{}}\NormalTok{, strategy}\OperatorTok{=}\NormalTok{\{}\StringTok{\textquotesingle{}sales\textquotesingle{}}\NormalTok{: }\StringTok{\textquotesingle{}sum\textquotesingle{}}\NormalTok{\})}
\end{Highlighting}
\end{Shaded}

\begin{center}\rule{0.5\linewidth}{0.5pt}\end{center}

\section{Best Practices}\label{best-practices-8}

\subsection{1. Always Check Column
Names}\label{always-check-column-names}

\begin{Shaded}
\begin{Highlighting}[]
\CommentTok{\# Before using columns, verify they exist}
\BuiltInTok{print}\NormalTok{(df.columns)}
\BuiltInTok{print}\NormalTok{(df.dtypes)  }\CommentTok{\# Also check data types}
\end{Highlighting}
\end{Shaded}

\subsection{2. Use Consistent DataFrame
Types}\label{use-consistent-dataframe-types}

\begin{Shaded}
\begin{Highlighting}[]
\CommentTok{\# Don\textquotesingle{}t mix pandas and polars in the same operation}
\CommentTok{\# Pick one and stick with it for related operations}
\end{Highlighting}
\end{Shaded}

\subsection{3. Enable Logging for
Debugging}\label{enable-logging-for-debugging}

\begin{Shaded}
\begin{Highlighting}[]
\CommentTok{\# Enable logging to see detailed operation information}
\NormalTok{result }\OperatorTok{=}\NormalTok{ add.to(orders, bring\_from}\OperatorTok{=}\NormalTok{customers, bring}\OperatorTok{=}\StringTok{\textquotesingle{}name\textquotesingle{}}\NormalTok{, }
\NormalTok{                against}\OperatorTok{=}\StringTok{\textquotesingle{}customer\_id\textquotesingle{}}\NormalTok{, logging}\OperatorTok{=}\VariableTok{True}\NormalTok{)}
\end{Highlighting}
\end{Shaded}

\subsection{4. Validate Data Before
Operations}\label{validate-data-before-operations}

\begin{Shaded}
\begin{Highlighting}[]
\CommentTok{\# Check for empty DataFrames}
\ControlFlowTok{if} \BuiltInTok{len}\NormalTok{(df) }\OperatorTok{==} \DecValTok{0}\NormalTok{:}
    \BuiltInTok{print}\NormalTok{(}\StringTok{"Warning: DataFrame is empty"}\NormalTok{)}

\CommentTok{\# Check for required columns}
\NormalTok{required\_cols }\OperatorTok{=}\NormalTok{ [}\StringTok{\textquotesingle{}price\textquotesingle{}}\NormalTok{, }\StringTok{\textquotesingle{}quantity\textquotesingle{}}\NormalTok{]}
\NormalTok{missing\_cols }\OperatorTok{=}\NormalTok{ [col }\ControlFlowTok{for}\NormalTok{ col }\KeywordTok{in}\NormalTok{ required\_cols }\ControlFlowTok{if}\NormalTok{ col }\KeywordTok{not} \KeywordTok{in}\NormalTok{ df.columns]}
\ControlFlowTok{if}\NormalTok{ missing\_cols:}
    \BuiltInTok{print}\NormalTok{(}\SpecialStringTok{f"Missing columns: }\SpecialCharTok{\{}\NormalTok{missing\_cols}\SpecialCharTok{\}}\SpecialStringTok{"}\NormalTok{)}

\CommentTok{\# Check for null values}
\NormalTok{null\_counts }\OperatorTok{=}\NormalTok{ df.isnull().}\BuiltInTok{sum}\NormalTok{()}
\ControlFlowTok{if}\NormalTok{ null\_counts.}\BuiltInTok{any}\NormalTok{():}
    \BuiltInTok{print}\NormalTok{(}\SpecialStringTok{f"Null values found:}\CharTok{\textbackslash{}n}\SpecialCharTok{\{}\NormalTok{null\_counts[null\_counts }\OperatorTok{\textgreater{}} \DecValTok{0}\NormalTok{]}\SpecialCharTok{\}}\SpecialStringTok{"}\NormalTok{)}
\end{Highlighting}
\end{Shaded}

\subsection{5. Use Explicit Parameter
Names}\label{use-explicit-parameter-names}

\begin{Shaded}
\begin{Highlighting}[]
\CommentTok{\# ✅ Good {-} clear and explicit}
\NormalTok{result }\OperatorTok{=}\NormalTok{ add.to(}
\NormalTok{    bring\_to}\OperatorTok{=}\NormalTok{orders,}
\NormalTok{    bring\_from}\OperatorTok{=}\NormalTok{customers,}
\NormalTok{    bring}\OperatorTok{=}\StringTok{\textquotesingle{}name\textquotesingle{}}\NormalTok{,}
\NormalTok{    against}\OperatorTok{=}\StringTok{\textquotesingle{}customer\_id\textquotesingle{}}
\NormalTok{)}

\CommentTok{\# ❌ Avoid {-} positional parameters can be confusing}
\NormalTok{result }\OperatorTok{=}\NormalTok{ add.to(orders, customers, }\StringTok{\textquotesingle{}name\textquotesingle{}}\NormalTok{, }\StringTok{\textquotesingle{}customer\_id\textquotesingle{}}\NormalTok{)}
\end{Highlighting}
\end{Shaded}

\subsection{6. Handle Errors Gracefully}\label{handle-errors-gracefully}

\begin{Shaded}
\begin{Highlighting}[]
\CommentTok{\# Wrap operations in try{-}except for production code}
\ControlFlowTok{try}\NormalTok{:}
\NormalTok{    result }\OperatorTok{=}\NormalTok{ add.transform(}\StringTok{\textquotesingle{}@calc\textquotesingle{}}\NormalTok{, df, expression}\OperatorTok{=}\StringTok{\textquotesingle{}price * quantity\textquotesingle{}}\NormalTok{, as\_}\OperatorTok{=}\StringTok{\textquotesingle{}total\textquotesingle{}}\NormalTok{)}
\ControlFlowTok{except} \PreprocessorTok{ValueError} \ImportTok{as}\NormalTok{ e:}
    \BuiltInTok{print}\NormalTok{(}\SpecialStringTok{f"Validation error: }\SpecialCharTok{\{}\NormalTok{e}\SpecialCharTok{\}}\SpecialStringTok{"}\NormalTok{)}
    \CommentTok{\# Handle error appropriately}
\ControlFlowTok{except} \PreprocessorTok{RuntimeError} \ImportTok{as}\NormalTok{ e:}
    \BuiltInTok{print}\NormalTok{(}\SpecialStringTok{f"Operation failed: }\SpecialCharTok{\{}\NormalTok{e}\SpecialCharTok{\}}\SpecialStringTok{"}\NormalTok{)}
    \CommentTok{\# Handle error appropriately}
\end{Highlighting}
\end{Shaded}

\subsection{7. Test with Small Data
First}\label{test-with-small-data-first}

\begin{Shaded}
\begin{Highlighting}[]
\CommentTok{\# Test operations on a small sample before running on full dataset}
\NormalTok{sample\_df }\OperatorTok{=}\NormalTok{ df.head(}\DecValTok{10}\NormalTok{)}
\NormalTok{result }\OperatorTok{=}\NormalTok{ add.transform(}\StringTok{\textquotesingle{}@calc\textquotesingle{}}\NormalTok{, sample\_df, expression}\OperatorTok{=}\StringTok{\textquotesingle{}price * quantity\textquotesingle{}}\NormalTok{, as\_}\OperatorTok{=}\StringTok{\textquotesingle{}total\textquotesingle{}}\NormalTok{)}
\CommentTok{\# If successful, run on full dataset}
\NormalTok{result }\OperatorTok{=}\NormalTok{ add.transform(}\StringTok{\textquotesingle{}@calc\textquotesingle{}}\NormalTok{, df, expression}\OperatorTok{=}\StringTok{\textquotesingle{}price * quantity\textquotesingle{}}\NormalTok{, as\_}\OperatorTok{=}\StringTok{\textquotesingle{}total\textquotesingle{}}\NormalTok{)}
\end{Highlighting}
\end{Shaded}

\begin{center}\rule{0.5\linewidth}{0.5pt}\end{center}

\section{Debugging Checklist}\label{debugging-checklist}

When you encounter an error, check these items in order:

\begin{enumerate}
\def\labelenumi{\arabic{enumi}.}
\tightlist
\item
  \textbf{Mode name}: Does it have the \texttt{@} prefix?
  (\texttt{@calc}, not \texttt{calc})
\item
  \textbf{Column names}: Do all referenced columns exist in the
  DataFrame?
\item
  \textbf{Data types}: Are you using numeric columns for math
  operations?
\item
  \textbf{Required parameters}: Does the mode have all required
  parameters?
\item
  \textbf{Parameter types}: Are parameters the correct type (dict, list,
  string)?
\item
  \textbf{DataFrame type}: Is it pandas or polars (not dict, list,
  etc.)?
\item
  \textbf{Empty data}: Does the DataFrame have at least one row?
\item
  \textbf{Null values}: Are there unexpected null values in key columns?
\item
  \textbf{Conflicting parameters}: Are you using \texttt{lineage=True}
  with \texttt{as\_type}?
\item
  \textbf{Rust bindings}: Is the Rust extension installed for Rust-based
  modes?
\end{enumerate}

\begin{center}\rule{0.5\linewidth}{0.5pt}\end{center}

\section{Getting Help}\label{getting-help}

If you're still stuck after checking this guide:

\begin{enumerate}
\def\labelenumi{\arabic{enumi}.}
\tightlist
\item
  \textbf{Check the error message carefully} - It often contains the
  solution
\item
  \textbf{Enable logging} - Use \texttt{logging=True} to see detailed
  operation info
\item
  \textbf{Simplify the operation} - Try with minimal parameters first
\item
  \textbf{Check the examples} - Review the documentation examples for
  similar use cases
\item
  \textbf{Verify installation} - Ensure Rust bindings are installed:
  \texttt{pip\ install\ additory{[}rust{]}}
\end{enumerate}

\begin{center}\rule{0.5\linewidth}{0.5pt}\end{center}

\section{Common Patterns That Work}\label{common-patterns-that-work}

\subsection{Pattern 1: Multi-Step
Pipeline}\label{pattern-1-multi-step-pipeline}

\begin{Shaded}
\begin{Highlighting}[]
\CommentTok{\# Build pipelines step by step, checking results at each stage}
\NormalTok{df }\OperatorTok{=}\NormalTok{ pd.DataFrame(...)}

\CommentTok{\# Step 1: Calculate}
\NormalTok{result }\OperatorTok{=}\NormalTok{ add.transform(}\StringTok{\textquotesingle{}@calc\textquotesingle{}}\NormalTok{, df, expression}\OperatorTok{=}\StringTok{\textquotesingle{}price * quantity\textquotesingle{}}\NormalTok{, as\_}\OperatorTok{=}\StringTok{\textquotesingle{}total\textquotesingle{}}\NormalTok{)}
\BuiltInTok{print}\NormalTok{(}\SpecialStringTok{f"After calc: }\SpecialCharTok{\{}\NormalTok{result}\SpecialCharTok{.}\NormalTok{shape}\SpecialCharTok{\}}\SpecialStringTok{"}\NormalTok{)}

\CommentTok{\# Step 2: Filter}
\NormalTok{result }\OperatorTok{=}\NormalTok{ add.transform(}\StringTok{\textquotesingle{}@filter\textquotesingle{}}\NormalTok{, result, where}\OperatorTok{=}\StringTok{\textquotesingle{}total \textgreater{} 100\textquotesingle{}}\NormalTok{)}
\BuiltInTok{print}\NormalTok{(}\SpecialStringTok{f"After filter: }\SpecialCharTok{\{}\NormalTok{result}\SpecialCharTok{.}\NormalTok{shape}\SpecialCharTok{\}}\SpecialStringTok{"}\NormalTok{)}

\CommentTok{\# Step 3: Aggregate}
\NormalTok{result }\OperatorTok{=}\NormalTok{ add.transform(}\StringTok{\textquotesingle{}@aggregate\textquotesingle{}}\NormalTok{, result, by}\OperatorTok{=}\StringTok{\textquotesingle{}region\textquotesingle{}}\NormalTok{, }
\NormalTok{                       strategy}\OperatorTok{=}\NormalTok{\{}\StringTok{\textquotesingle{}total\textquotesingle{}}\NormalTok{: }\StringTok{\textquotesingle{}sum\textquotesingle{}}\NormalTok{\})}
\BuiltInTok{print}\NormalTok{(}\SpecialStringTok{f"After aggregate: }\SpecialCharTok{\{}\NormalTok{result}\SpecialCharTok{.}\NormalTok{shape}\SpecialCharTok{\}}\SpecialStringTok{"}\NormalTok{)}
\end{Highlighting}
\end{Shaded}

\subsection{Pattern 2: Safe Column
Access}\label{pattern-2-safe-column-access}

\begin{Shaded}
\begin{Highlighting}[]
\CommentTok{\# Always verify columns exist before using them}
\KeywordTok{def}\NormalTok{ safe\_transform(df, expression, as\_name):}
    \CommentTok{\# Extract column names from expression (simplified)}
\NormalTok{    required\_cols }\OperatorTok{=}\NormalTok{ [}\StringTok{\textquotesingle{}price\textquotesingle{}}\NormalTok{, }\StringTok{\textquotesingle{}quantity\textquotesingle{}}\NormalTok{]  }\CommentTok{\# Columns used in expression}
    
\NormalTok{    missing }\OperatorTok{=}\NormalTok{ [col }\ControlFlowTok{for}\NormalTok{ col }\KeywordTok{in}\NormalTok{ required\_cols }\ControlFlowTok{if}\NormalTok{ col }\KeywordTok{not} \KeywordTok{in}\NormalTok{ df.columns]}
    \ControlFlowTok{if}\NormalTok{ missing:}
        \ControlFlowTok{raise} \PreprocessorTok{ValueError}\NormalTok{(}\SpecialStringTok{f"Missing required columns: }\SpecialCharTok{\{}\NormalTok{missing}\SpecialCharTok{\}}\SpecialStringTok{"}\NormalTok{)}
    
    \ControlFlowTok{return}\NormalTok{ add.transform(}\StringTok{\textquotesingle{}@calc\textquotesingle{}}\NormalTok{, df, expression}\OperatorTok{=}\NormalTok{expression, as\_}\OperatorTok{=}\NormalTok{as\_name)}

\CommentTok{\# Use the safe wrapper}
\NormalTok{result }\OperatorTok{=}\NormalTok{ safe\_transform(df, }\StringTok{\textquotesingle{}price * quantity\textquotesingle{}}\NormalTok{, }\StringTok{\textquotesingle{}total\textquotesingle{}}\NormalTok{)}
\end{Highlighting}
\end{Shaded}

\subsection{Pattern 3: Type-Safe
Operations}\label{pattern-3-type-safe-operations}

\begin{Shaded}
\begin{Highlighting}[]
\CommentTok{\# Ensure consistent types throughout pipeline}
\ImportTok{import}\NormalTok{ pandas }\ImportTok{as}\NormalTok{ pd}

\CommentTok{\# Start with pandas}
\NormalTok{df }\OperatorTok{=}\NormalTok{ pd.DataFrame(...)}

\CommentTok{\# All operations return pandas (no as\_type conversion)}
\NormalTok{result }\OperatorTok{=}\NormalTok{ add.to(df, bring\_from}\OperatorTok{=}\NormalTok{ref\_df, bring}\OperatorTok{=}\StringTok{\textquotesingle{}name\textquotesingle{}}\NormalTok{, against}\OperatorTok{=}\StringTok{\textquotesingle{}id\textquotesingle{}}\NormalTok{)}
\NormalTok{result }\OperatorTok{=}\NormalTok{ add.transform(}\StringTok{\textquotesingle{}@calc\textquotesingle{}}\NormalTok{, result, expression}\OperatorTok{=}\StringTok{\textquotesingle{}price * 2\textquotesingle{}}\NormalTok{, as\_}\OperatorTok{=}\StringTok{\textquotesingle{}doubled\textquotesingle{}}\NormalTok{)}
\NormalTok{result }\OperatorTok{=}\NormalTok{ add.transform(}\StringTok{\textquotesingle{}@filter\textquotesingle{}}\NormalTok{, result, where}\OperatorTok{=}\StringTok{\textquotesingle{}doubled \textgreater{} 100\textquotesingle{}}\NormalTok{)}

\CommentTok{\# Final result is pandas}
\ControlFlowTok{assert} \BuiltInTok{isinstance}\NormalTok{(result, pd.DataFrame)}
\end{Highlighting}
\end{Shaded}

\begin{center}\rule{0.5\linewidth}{0.5pt}\end{center}

\section{Error Message Glossary}\label{error-message-glossary}

\begin{longtable}[]{@{}
  >{\raggedright\arraybackslash}p{(\linewidth - 4\tabcolsep) * \real{0.4688}}
  >{\raggedright\arraybackslash}p{(\linewidth - 4\tabcolsep) * \real{0.2812}}
  >{\raggedright\arraybackslash}p{(\linewidth - 4\tabcolsep) * \real{0.2500}}@{}}
\toprule\noalign{}
\begin{minipage}[b]{\linewidth}\raggedright
Error Message
\end{minipage} & \begin{minipage}[b]{\linewidth}\raggedright
Meaning
\end{minipage} & \begin{minipage}[b]{\linewidth}\raggedright
Action
\end{minipage} \\
\midrule\noalign{}
\endhead
\bottomrule\noalign{}
\endlastfoot
``Column `X' not found'' & Column doesn't exist & Check
\texttt{df.columns} \\
``Mode `X' not recognized'' & Invalid mode name & Add \texttt{@} prefix
or check spelling \\
``Parameter `X' is required'' & Missing required parameter & Add the
parameter \\
``Type mismatch'' & Wrong data type & Convert to correct type \\
``Empty DataFrame'' & No rows in DataFrame & Check data source \\
``Rust bindings not available'' & Extension not installed & Run
\texttt{pip\ install\ additory{[}rust{]}} \\
``Cannot convert'' & Type conversion failed & Check data
compatibility \\
``Invalid strategy'' & Wrong strategy format & Use dict format:
\texttt{\{\textquotesingle{}col\textquotesingle{}:\ \textquotesingle{}agg\textquotesingle{}\}} \\
``Conflicting parameters'' & Incompatible parameters used together &
Remove one parameter \\
``Operation failed'' & Generic operation error & Check all parameters
and data \\
\end{longtable}

\begin{center}\rule{0.5\linewidth}{0.5pt}\end{center}

\textbf{Note}: This guide covers the most common errors. For
function-specific details, refer to the individual function
documentation pages.

\bookmarksetup{startatroot}

\chapter*{References}\label{references}
\addcontentsline{toc}{chapter}{References}

\markboth{References}{References}

\section*{Official Resources}\label{official-resources}
\addcontentsline{toc}{section}{Official Resources}

\markright{Official Resources}

\subsection*{Additory}\label{additory}
\addcontentsline{toc}{subsection}{Additory}

\begin{itemize}
\tightlist
\item
  \textbf{PyPI Package}: \url{https://pypi.org/project/additory/}
\item
  \textbf{GitHub Repository}:
  \url{https://github.com/sekarkrishna/additory}
\item
  \textbf{Issue Tracker}:
  \url{https://github.com/sekarkrishna/additory/issues}
\end{itemize}

\subsection*{Dependencies}\label{dependencies}
\addcontentsline{toc}{subsection}{Dependencies}

\begin{itemize}
\tightlist
\item
  \textbf{Polars}: \url{https://www.pola.rs/}
\item
  \textbf{Pandas}: \url{https://pandas.pydata.org/}
\item
  \textbf{PyArrow}: \url{https://arrow.apache.org/docs/python/}
\item
  \textbf{Rust}: \url{https://www.rust-lang.org/}
\item
  \textbf{PyO3}: \url{https://pyo3.rs/}
\item
  \textbf{Maturin}: \url{https://github.com/PyO3/maturin}
\end{itemize}

\section*{Related Projects}\label{related-projects}
\addcontentsline{toc}{section}{Related Projects}

\markright{Related Projects}

\subsection*{DataFrame Libraries}\label{dataframe-libraries}
\addcontentsline{toc}{subsection}{DataFrame Libraries}

\begin{itemize}
\tightlist
\item
  \textbf{Polars Documentation}:
  \url{https://pola-rs.github.io/polars-book/}
\item
  \textbf{Pandas Documentation}: \url{https://pandas.pydata.org/docs/}
\item
  \textbf{Apache Arrow}: \url{https://arrow.apache.org/}
\end{itemize}

\subsection*{Rust for Python}\label{rust-for-python}
\addcontentsline{toc}{subsection}{Rust for Python}

\begin{itemize}
\tightlist
\item
  \textbf{PyO3 Guide}: \url{https://pyo3.rs/latest/}
\item
  \textbf{Maturin Guide}: \url{https://www.maturin.rs/}
\end{itemize}

\section*{Community}\label{community}
\addcontentsline{toc}{section}{Community}

\markright{Community}

\subsection*{Contact}\label{contact}
\addcontentsline{toc}{subsection}{Contact}

\begin{itemize}
\tightlist
\item
  \textbf{Author}: Krishnamoorthy Sankaran
\item
  \textbf{Email}: krishnamoorthy.sankaran@sekrad.org
\end{itemize}

\subsection*{Contributing}\label{contributing}
\addcontentsline{toc}{subsection}{Contributing}

Contributions are welcome! Please feel free to submit a Pull Request on
GitHub.

\subsection*{Support}\label{support}
\addcontentsline{toc}{subsection}{Support}

For issues, questions, or feature requests: - Open an issue on
\href{https://github.com/sekarkrishna/additory/issues}{GitHub Issues} -
Email the author directly

\section*{Citation}\label{citation}
\addcontentsline{toc}{section}{Citation}

\markright{Citation}

If you use additory in your research or project, please cite:

\begin{Shaded}
\begin{Highlighting}[]
\CommentTok{@software\{additory2026,}
\CommentTok{  author = \{Sankaran, Krishnamoorthy\},}
\CommentTok{  title = \{Additory: Elegant data operations for DataFrames\},}
\CommentTok{  year = \{2026\},}
\CommentTok{  url = \{https://github.com/sekarkrishna/additory\},}
\CommentTok{  version = \{0.1.3a10\}}
\CommentTok{\}}
\end{Highlighting}
\end{Shaded}

\section*{License}\label{license-1}
\addcontentsline{toc}{section}{License}

\markright{License}

Additory is released under the MIT License.

\begin{verbatim}
MIT License

Copyright (c) 2026 Krishnamoorthy Sankaran

Permission is hereby granted, free of charge, to any person obtaining a copy
of this software and associated documentation files (the "Software"), to deal
in the Software without restriction, including without limitation the rights
to use, copy, modify, merge, publish, distribute, sublicense, and/or sell
copies of the Software, and to permit persons to whom the Software is
furnished to do so, subject to the following conditions:

The above copyright notice and this permission notice shall be included in all
copies or substantial portions of the Software.

THE SOFTWARE IS PROVIDED "AS IS", WITHOUT WARRANTY OF ANY KIND, EXPRESS OR
IMPLIED, INCLUDING BUT NOT LIMITED TO THE WARRANTIES OF MERCHANTABILITY,
FITNESS FOR A PARTICULAR PURPOSE AND NONINFRINGEMENT. IN NO EVENT SHALL THE
AUTHORS OR COPYRIGHT HOLDERS BE LIABLE FOR ANY CLAIM, DAMAGES OR OTHER
LIABILITY, WHETHER IN AN ACTION OF CONTRACT, TORT OR OTHERWISE, ARISING FROM,
OUT OF OR IN CONNECTION WITH THE SOFTWARE OR THE USE OR OTHER DEALINGS IN THE
SOFTWARE.
\end{verbatim}

\section*{Acknowledgments}\label{acknowledgments}
\addcontentsline{toc}{section}{Acknowledgments}

\markright{Acknowledgments}

Built with: - \href{https://www.rust-lang.org/}{Rust} - Systems
programming language - \href{https://pyo3.rs/}{PyO3} - Rust bindings for
Python - \href{https://www.pola.rs/}{Polars} - Lightning-fast DataFrame
library - \href{https://github.com/PyO3/maturin}{Maturin} - Build and
publish Rust-based Python packages

Special thanks to the open-source community for their amazing tools and
libraries.

\begin{center}\rule{0.5\linewidth}{0.5pt}\end{center}

\textbf{Made with ❤️ for the data science community}


\backmatter


\end{document}
